% UPUTSTVO: https://bg.ac.rs/files/sr/studije/studije-uni/Uputstvo-izrada-master-rad.pdf

\documentclass[12pt]{report}
\usepackage{fontspec}
\usepackage[english,serbian]{babel}
\usepackage{microtype}
\usepackage{lmodern}
\usepackage{indentfirst} % pise u uputstvu da ne treba
\usepackage[ddmmyyyy]{datetime}
\renewcommand*\contentsname{Sadržaj}
\usepackage{url}
\usepackage{enumitem}
\renewcommand{\chaptername}{}
\usepackage{pgf}
\usepackage{amsmath}
\usepackage{commath}
\usepackage{import}
\usepackage{graphicx}
\setlength{\parskip}{1em}
\setlength{\baselineskip}{1.6em} % ovo mozda menjati
\renewcommand{\figurename}{Slika}
\usepackage{amsfonts}
% margine
\usepackage[
	a4paper,	
	margin=2.5cm		
	]{geometry}
	
% linkovi
\usepackage{color}   %May be necessary if you want to color links
\usepackage{hyperref}
\hypersetup{
    colorlinks=true, %set true if you want colored links
    linktoc=all,     %set to all if you want both sections and subsections linked
    linkcolor=black,  %choose some color if you want links to stand out
	citecolor=violet
}	

%tikz
\usepackage{tikz}
\usetikzlibrary{shapes,arrows,positioning} 

% indentacija liste
\setlist[itemize,1]{leftmargin=\dimexpr +1.5cm}

% podesavanje headera
\usepackage{fancyhdr}
\setlength{\headheight}{20pt}
\pagestyle{fancy}
\fancyhf{}
\lhead{\leftmark} % \fontsize{10}{12} \selectfont -- mozes da ubacis u \lhead{}\\
\rhead{\thepage}

% podesavanje headera za strane gde pocinju poglavlja (plain style)
\fancypagestyle{plain}{%
\fancyhf{}
\rhead{\thepage}
\setlength{\headheight}{20pt}}

% sledeca dva su da mi pokazuje broj pored jednacine samo ako postoji referenca
% na nju negde u tekstu!
\usepackage{mathtools}  
\mathtoolsset{showonlyrefs}

%\rfoot{Page \thepage}

% skidam broj sa sadrzaja

% nadji nacin samo da prikazes broj samo pored formula na koje postoji referenca! uradio valjda

\addto{\captionsserbian}{\renewcommand{\bibname}{Literatura}}
\bibliographystyle{plain}

%\newcommand\eqdef{\stackrel{\mathclap{\footnotesize\mbox{def}}}{=}}
\newcommand{\eqdef}{\ensuremath{\stackrel{\text{def}}{=}}}
\newcommand{\eqivdef}{\ensuremath{\stackrel{\text{def}}{\iff}}}

\usepackage{amsthm}
\newtheorem*{theorem*}{Teorema}

\DeclareMathOperator*{\argmax}{\arg\!\max}


\begin{document}

\pagenumbering{roman}
% sredi margine za naslovnu!
\begin{titlepage}
    \begin{center}
    
        \vspace*{1cm}
        \hrule
        \vspace*{1cm}
        \textbf{Primena dubokog Q učenja na automatsko igranje video igara}
        \vspace*{1cm}
        \hrule
        \vspace{1cm}
        
		\resizebox{.2\linewidth}{!}{\includegraphics{img/matf_logo}}
		
        \begin{center}
        	Matematički fakultet \\
        	Univerzitet u Beogradu
        \end{center}
        \vspace{1.5cm}
        Student:\hfill Mentor:\\
        \textbf{Nikola Milev}
		\hfill        
        \textbf{Mladen Nikolić}
        \vfill
		Beograd, 2018.
    \end{center}
\end{titlepage}

\newpage ~ \thispagestyle{empty}

\addtocontents{toc}{\protect\thispagestyle{empty}}
\tableofcontents{}

\newpage
\pagenumbering{arabic}
\chapter{Uvod}



U maju 1997. godine, Gari Kasparov, tadašnji svetski šampion u šahu, izgubio je meč protiv računarskog sistema pod nazivom Deep Blue. Skoro dvadeset godina kasnije, program pod nazivom AlphaGo pobedio je profesionalnog ljudskog igrača u igri go. Iako su obe igre strateške i igraju se na tabli, između šaha i igre go postoji ogromna razlika. Pravila igre go dosta su jednostavna u odnosu na šah ali je prostor koji opisuje poteze igre go više od $10^{100}$ puta veći od prostora koji opisuje poteze šaha. Programi koji igraju šah često se zasnivaju na korišćenju stabala pretrage i ovaj pristup jednostavno nije primenljiv na igru go. \par

Na čemu je onda zasnovan AlphaGo? U pitanju je učenje potkrepljivanjem (eng. reinforcement learning). Ovo je vrsta mašinskog učenja koja počiva na sistemu kazne i nagrade. Podrazumeva se da se sistem sastoji od agenta i okruženja u kom agent dela (vrši akcije) i dobija numeričku nagradu i informaciju o promeni stanja okruženja. Osnovni dijagram ove komunikacije može se videti na slici \ref{fig:rl_diag}. Poput dresiranja psa, nagradama se ohrabruje poželjno ponašanje dok se nepoželjno kažnjava. Cilj jeste ostvariti što veću dugoročnu nagradu. Međutim, agent mora sam kroz istraživanje da shvati kako da dostigne najveću nagradu tako što isprobava različite akcije. Takođe, preduzete akcije mogu da utiču i na nagradu koja se pojavljuje dugo nakon što je sama akcija preduzeta. Ovo zahteva da se uvede pojam dugoročne nagrade. Pojmovi istraživanja i dugoročne nagrade su ključni pri učenju potkrepljivanjem. 

\begin{figure}
	\centering
	\resizebox{.4\linewidth}{!}{\tikzset{
    %Define standard arrow tip
    >=stealth',
    %Define style for boxes
    okruzenje/.style={
           rectangle,
           rounded corners,
           draw=black, very thick,
           text width=6.5em,
           minimum height=5em,
           text centered},
    agent/.style={
           ellipse,
           draw=black, very thick,
           text width=4.5em,
           minimum height=3em,
           text centered},   
    % Define arrow style
    pil/.style={
           ->,
           thick}
}

\begin{tikzpicture}[node distance=1cm, auto,]

\node[okruzenje] (okruzenje) {OKRUŽENJE} ;
\node [agent, below left=7em of okruzenje] (agent) {AGENT} ;

\draw (agent) edge[->, out=0, in=-90] node[right=0.5em] {AKCIJA} (okruzenje);
\draw (okruzenje) edge[->, out=180, in=90] node[align=left,left=0.5em] {NAGRADA\\~\\STANJE} (agent) ;

\end{tikzpicture}
}
	\caption{Dijagram komunikacije agenta sa okruženjem pri učenju potkrepljivanjem}
	\label{fig:rl_diag}
\end{figure}

\par
Pri učenju potkrepljivanjem, najčešće se pretpostavlja da je skup svih mogućih stanja okruženja diskretan. Ovo dozvoljava primenu Markovljevih procesa odlučivanja i omogućuje jednostavan formalan opis problema koji se rešava i pristupa njegovog rešavanja. Formalno predstavljanje problema i rešenja dato je u poglavlju \ref{ch:rl}.
\par
Učenje potkrepljivanjem jedna je od tri vrste mašinskog učenja, pored nadgledanog i nenadgledanog učenja. Pri nadgledanom učenju, sistem dobija skup ulaznih i izlaznih podataka s ciljem da izvrši generalizaciju nad tim podacima i uspešno generiše izlazne podatke na osnovu do sada neviđenih ulaznih podataka. Pri učenju potkrepljivanjem, ne postoje unapred poznate akcije koje treba preduzeti već sistem na osnovu nagrade mora zaključi koji je optimalni sled akcija. Iako široko korišćeno, nadgledano učenje nije prikladno za učenje iz novih iskustava, kada izlazni podaci nisu dostupni.  Kod nenadgledanog učenja, često je neophodno pronaći neku strukturu u podacima nad kojima se uči bez ikakvog pređašnjeg znanja o njima. Iako učenje potkrepljivanjem liči i na nadgledano i na nenadgledano učenje, agent ne traži strukturu niti mu je unapred data informacija o optimalnom ponašanju već teži maksimizaciji nagrade koju dobija od okruženja. \par
\par
Učenje potkrepljivanjem ima mnogobrojne primene kao što su samostalna vožnja automobila i letelica, automatsko konfigurisanje algoritama, trgovina na berzi, igranje igara, itd. Ovaj vid mašinskog učenja pokazao se kao dobar i za igranje video igara.  U radu objavljenom 2015. godine u časopisu Nature, DeepMind predstavlja sistem koji uči da igra video igre sa konzole Atari 2600, neke čak i daleko bolje od ljudi \cite{dqn_dm}. U avgustu 2017. godine, OpenAI predstavlja agenta koji isključivo kroz igranje igre i bez pređašnjeg znanja o igri stiče nivo umeća dovoljan da pobedi i neke od najboljih ljudskih takmičara u video igri Dota 2.\footnote{https://blog.openai.com/dota-2/} 
\par 
U naučnom radu koji je objavila kompanija DeepMind u časopisu Nature predložen je novi algoritam, DQN (eng.~{\em deep Q - network}), koji koristi spoj učenja uslovljavanjem i duboke neuronske mreže i uspešno savladava razne igre za Atari 2600 konzolu. Sve informacije dostupne agentu jesu pikseli sa ekrana, trenutni rezultat u igri i signal za kraj igre. U sklopu ovog rada, ispitana je struktura algoritma DQN i data je implementacija čije su performanse testirane na manjoj skali od one date u radu, zbog ograničenih resursa. Takođe je eksperimentisano sa elementima samog algoritma i opisano je kako oni utiču na njegovo ponašanje. Dva elementa, ciljna mreža i memorija, ključni su za uspešan proces učenja. Dobijeni rezultati pokazuju njihovu važnost. 
\par 
U sklopu rada opisani su osnovni pojmovi mašinskog učenja (glava ~\ref{ch:ml}), zadržavajući se na neuronskim mrežama uopšte (glava ~\ref{ch:nn}) i na konvolutivnim neuronskim mrežama (deo ~\ref{sec:cnn}). Glava ~\ref{ch:rl} posvećena je učenju potkrepljivanjem dok je algoritam DQN u celosti opisan u glavi ~\ref{ch:dqn}. U glavi ~\ref{ch:implementacija} data je implementacija, njena evaluacija, kao i eksperimenti i njihovi rezultati.
\chapter{Mašinsko učenje}
\label{ch:ml}
Mašinsko učenje počelo je da stiče veliku popularnost devedesetih godina prošlog veka zahvaljujući potrebi i mogućnosti da se uči iz velike količine dostupnih podataka i uspešnosti ovog pristupa u tome. Za popularizaciju mašinskog učenja početkom 21. veka najzaslužnije su neuronske mreže, u toj meri da je pojam mašinskog učenja među laicima često poistovećen sa pojmom neuronskih mreža. Ovo, naravno, nije tačno; sem neuronskih mreža, postoje razne druge tehnike, kao što su metod potpornih vektora, linearni modeli, probabilistički grafovski modeli, itd. \par

Mašinsko učenje nastalo je iz čovekove želje da oponaša prirodne mehanizme učenja kod čoveka i životinja kao jedne od osnovnih svojstava inteligencije i korišćenja dobijenih rezultata u praktične svrhe. Termin mašinsko učenje prvi je upotrebio pionir veštačke inteligencije, Artur Semjuel \cite{samuel_ai}, koji je doprineo razvoju veštačke inteligencije istražujući igru dame (eng. checkers) i tražeći način da stvori računarski program koji na osnovu iskusva može da savlada ovu igru \cite{samuel}. 
\par

Mašinsko učenje može se definisati kao disciplina koja se bavi izgradnjom prilagodljivih računarskih sistema koji su sposobni da poboljšaju svoje performanse koristeći informacije iz iskustva \cite{janicic_vi}. No, u biti mašinskog učenja leži generalizacija, tj. indukcija. Dve vrste zaključivanja, indukcija\footnote{Indukcija -- zaključivanje od pojedinačnog ka opštem} i dedukcija\footnote{Dedukcija -- zaključivanje od opšteg ka konkretnom} imaju svoje odgovarajuće discipline u sklopu veštačke inteligencije: mašinsko učenje i automatsko rezonovanje. Kao što se indukcija i dedukcija razlikuju, i mašinsko učenje i automatsko rezonovanje imaju različite oblasti primene. Automatsko rezonovanje zasnovano je na matematičkoj logici i koristi se kada je problem relativno lako formulisati ali ga, često zbog velikog prostora mogućih rešenja, nije jednostavno rešiti. U ovoj oblasti, neophodno je dobiti apsolutno tačna rešenja, ne dopuštajući nikakav nivo greške.
S druge strane, mašinsko učenje pogodnije je kada problem nije moguće precizno formulisati i kada se očekuje neki novo greške. Čovek neke od ovih problema lako rešava a neke ne. 
Ukoliko je neophodno napraviti sistem koji prepoznaje životinje na slikama, kako definisati problem? Koji su tačno elementi oblika životinje? Kako ih prepoznati?  Metodama automatskog rezonovanja bilo bi nemoguće definisati ovaj problem i rešiti ga. Mašinsko učenje, s druge strane, pokazalo se kao dobar pristup. Ono što je još karakteristično za mašinsko učenje jeste da rešenje ne mora biti savršeno tačno, iako se tome teži, i nivo prihvatljivog odstupanja zavisi od problema i konteksta primene. \par

Ova oblast je kroz manje od 20 godina od popularizacije postala deo svakodnevice. U sklopu društvene mreže Fejsbuk (eng. Facebook) implementiran je sistem za prepoznavanje lica koji preporučuje profile osoba koje se nalaze na slikama. Razni veb servisi koriste metode mašinskog učenja radi stvaranja sistema za preporuke (artikala u prodavnicama, video sadržaja na platformama za njihovo gledanje, itd.) i sistema za detekciju prevara. Mnoge firme koje se bave trgovinom na berzi imaju sisteme koji automatski trguju deonicama. U medicini, jedna od primena mašinskog učenja jeste za uspostavljanje dijagnoze. Još neke primene su u marketingu, za procesiranje prirodnih jezika, bezbednost, itd.

%[ODNOS MASINSKOG UCENJA I OSTALIH OBLASTI MATEMATIKE?]

\section{Vrste mašinskog učenja}

Kada se govori o određenoj vrsti mašinskog učenja, podrazumevaju se vrste problema, kao i načini za njihovo rešavanje. Prema problemima koji se rešavaju, mašinsko učenje deli se na tri vrste: nadgledano učenje (eng. supervised learning), nenadgledano učenje (eng. unsupervised learning) i učenje potkrepljivanjem (eng. reinforcement learning). Iako se podela mnogih autora sastoji samo iz nadgledanog i nenadgledanog učenja, postoji razlika između učenja potkrepljivanjem i preostale dve vrste. U nastavku su dati opisi pristupa nadgledanog i nenadgledanog učenja. Učenju uslovljavanjem, kao centralnoj temi ovog rada, posvećeno je više pažnje u poglavlju~ \ref{ch:rl}. 

\subsection{Nadgledano mašinsko učenje}

Pri nadgledanom mašinskom učenju, date su vrednosti ulaza i izlaza koje im odgovaraju za određeni broj slučajeva. Sistem treba na osnovu već datih veza za pojedinačne parove da ustanovi kakva veza postoji između tih parova i izvrši generalizaciju, odnosno, ukoliko ulazne podatke označimo sa $x$ a izlazne sa $y$, sistem treba da odredi funkciju $f$ takvu da važi
\begin{center}
	$y \approx f(x)$
\end{center}
Pri uspešno rešenom problemu nadgledanog učenja, funkcija $f$ davaće tačna ili približno tačna rešenja i za podatke koji do sada nisu viđeni.
Ulazne vrednosti nazivaju se atributima (eng. features) a izlazne ciljnim promenljivima (eng. target variables). Ovim opisom nije određema dimenzionalnost ni za ulazne ni za izlazne promenljive, iako je dimenzija izlazne promenljive uglavnom 1. Funkcija $f$ naziva se modelom. \par
Skup svih mogućih funkcija odgovarajuće dimenzionalnosti bio bi previše veliki za pretragu i zbog toga se uvode pretpostavke o samom modelu. Pretpostavlja da je definisan skup svih dopustivih modela i da je potrebno naći najpogodniji element tog skupa. Najčešće je taj skup određen parametrima, tj. uzima se da funkcija zavisi od nekog parametra $w$ koji je u opštem slušaju višedimenzioni i tada se funkcija označava sa $f_w(x)$.

\par
Neophodno je uvesti funkciju greške modela (eng. loss function), odnosno funkciju koja opisuje koliko dati model dobro određuje izlaz za dati ulaz. Ova funkcija se najčešće označava sa $L$ i $L(y, f_w(x))$ predstavlja razliku između željene i dobijene vrednosti za pojedinačni par promenljivih. No, nijedan par vrednosti promenljivih nije dovoljan za opis kvaliteta modela već treba naći funkciju koja globalno ocenjuje odstupanje modela od stvarnih vrednosti. U praksi, podrazumeva se postojanje uzorka:
\begin{center}
	$D=\{(x_i, y_i)|i=1,...,N\}$
\end{center}
i uvodi sledeća funkcija:
\begin{center}
	$E(w, D) = \frac{1}{N}\sum_{i=1}^{N}L(y_i, f_w(x_i))$
\end{center}
koja se naziva prosečnom greškom. Uobičajeno, algoritmi nadgledanog mašinskog učenja zasnivaju se na minimizaciji prosečne greške. \par

Postoje dva osnovna tipa zadataka nadgledanog mašinskog učenja:
\begin{itemize}
	\item Klasifikacija 
	\item Regresija
\end{itemize}

Klasifikacija (eng.~{\em classification}) predstavlja zadatak mašinskog učenja gde je cilj predvideti klasu u kojoj se ciljna promenljiva nalazi. Neki od primera klasifikacije su svrstavanje slika na one koje sadrže ili ne sadrže lice, označavanje nepoželjne (eng.~{\em spam}) elektronske pošte i prepoznavanje objekata na slikama. 
Jednostavan primer klasifikacije može se videti na slici \ref{fig:bin_klas}, gde su trouglovima označeni podaci iznad prave $y=2x+1$ a krugovima podaci ispod date prave. Dakle, ta prava je klasifikacioni model sa parametrima $(w_0, w_1) = (1, 2)$.

\begin{figure}
	\centering
	\resizebox{.8\linewidth}{!}{%% Creator: Matplotlib, PGF backend
%%
%% To include the figure in your LaTeX document, write
%%   \input{<filename>.pgf}
%%
%% Make sure the required packages are loaded in your preamble
%%   \usepackage{pgf}
%%
%% Figures using additional raster images can only be included by \input if
%% they are in the same directory as the main LaTeX file. For loading figures
%% from other directories you can use the `import` package
%%   \usepackage{import}
%% and then include the figures with
%%   \import{<path to file>}{<filename>.pgf}
%%
%% Matplotlib used the following preamble
%%   \usepackage{fontspec}
%%   \setmainfont{DejaVu Serif}
%%   \setsansfont{DejaVu Sans}
%%   \setmonofont{DejaVu Sans Mono}
%%
\begingroup%
\makeatletter%
\begin{pgfpicture}%
\pgfpathrectangle{\pgfpointorigin}{\pgfqpoint{6.400000in}{4.800000in}}%
\pgfusepath{use as bounding box, clip}%
\begin{pgfscope}%
\pgfsetbuttcap%
\pgfsetmiterjoin%
\definecolor{currentfill}{rgb}{1.000000,1.000000,1.000000}%
\pgfsetfillcolor{currentfill}%
\pgfsetlinewidth{0.000000pt}%
\definecolor{currentstroke}{rgb}{1.000000,1.000000,1.000000}%
\pgfsetstrokecolor{currentstroke}%
\pgfsetdash{}{0pt}%
\pgfpathmoveto{\pgfqpoint{0.000000in}{0.000000in}}%
\pgfpathlineto{\pgfqpoint{6.400000in}{0.000000in}}%
\pgfpathlineto{\pgfqpoint{6.400000in}{4.800000in}}%
\pgfpathlineto{\pgfqpoint{0.000000in}{4.800000in}}%
\pgfpathclose%
\pgfusepath{fill}%
\end{pgfscope}%
\begin{pgfscope}%
\pgfsetbuttcap%
\pgfsetmiterjoin%
\definecolor{currentfill}{rgb}{1.000000,1.000000,1.000000}%
\pgfsetfillcolor{currentfill}%
\pgfsetlinewidth{0.000000pt}%
\definecolor{currentstroke}{rgb}{0.000000,0.000000,0.000000}%
\pgfsetstrokecolor{currentstroke}%
\pgfsetstrokeopacity{0.000000}%
\pgfsetdash{}{0pt}%
\pgfpathmoveto{\pgfqpoint{0.800000in}{0.528000in}}%
\pgfpathlineto{\pgfqpoint{5.760000in}{0.528000in}}%
\pgfpathlineto{\pgfqpoint{5.760000in}{4.224000in}}%
\pgfpathlineto{\pgfqpoint{0.800000in}{4.224000in}}%
\pgfpathclose%
\pgfusepath{fill}%
\end{pgfscope}%
\begin{pgfscope}%
\pgfpathrectangle{\pgfqpoint{0.800000in}{0.528000in}}{\pgfqpoint{4.960000in}{3.696000in}} %
\pgfusepath{clip}%
\pgfsetbuttcap%
\pgfsetroundjoin%
\definecolor{currentfill}{rgb}{1.000000,0.647059,0.000000}%
\pgfsetfillcolor{currentfill}%
\pgfsetlinewidth{1.003750pt}%
\definecolor{currentstroke}{rgb}{1.000000,0.647059,0.000000}%
\pgfsetstrokecolor{currentstroke}%
\pgfsetdash{}{0pt}%
\pgfpathmoveto{\pgfqpoint{0.800000in}{1.589805in}}%
\pgfpathlineto{\pgfqpoint{0.841667in}{1.673139in}}%
\pgfpathlineto{\pgfqpoint{0.758333in}{1.673139in}}%
\pgfpathclose%
\pgfusepath{stroke,fill}%
\end{pgfscope}%
\begin{pgfscope}%
\pgfpathrectangle{\pgfqpoint{0.800000in}{0.528000in}}{\pgfqpoint{4.960000in}{3.696000in}} %
\pgfusepath{clip}%
\pgfsetbuttcap%
\pgfsetroundjoin%
\definecolor{currentfill}{rgb}{1.000000,0.647059,0.000000}%
\pgfsetfillcolor{currentfill}%
\pgfsetlinewidth{1.003750pt}%
\definecolor{currentstroke}{rgb}{1.000000,0.647059,0.000000}%
\pgfsetstrokecolor{currentstroke}%
\pgfsetdash{}{0pt}%
\pgfpathmoveto{\pgfqpoint{0.850101in}{4.970687in}}%
\pgfpathlineto{\pgfqpoint{0.891768in}{5.054020in}}%
\pgfpathlineto{\pgfqpoint{0.808434in}{5.054020in}}%
\pgfpathclose%
\pgfusepath{stroke,fill}%
\end{pgfscope}%
\begin{pgfscope}%
\pgfpathrectangle{\pgfqpoint{0.800000in}{0.528000in}}{\pgfqpoint{4.960000in}{3.696000in}} %
\pgfusepath{clip}%
\pgfsetbuttcap%
\pgfsetroundjoin%
\definecolor{currentfill}{rgb}{1.000000,0.647059,0.000000}%
\pgfsetfillcolor{currentfill}%
\pgfsetlinewidth{1.003750pt}%
\definecolor{currentstroke}{rgb}{1.000000,0.647059,0.000000}%
\pgfsetstrokecolor{currentstroke}%
\pgfsetdash{}{0pt}%
\pgfpathmoveto{\pgfqpoint{0.900202in}{5.344999in}}%
\pgfpathlineto{\pgfqpoint{0.941869in}{5.428332in}}%
\pgfpathlineto{\pgfqpoint{0.858535in}{5.428332in}}%
\pgfpathclose%
\pgfusepath{stroke,fill}%
\end{pgfscope}%
\begin{pgfscope}%
\pgfpathrectangle{\pgfqpoint{0.800000in}{0.528000in}}{\pgfqpoint{4.960000in}{3.696000in}} %
\pgfusepath{clip}%
\pgfsetbuttcap%
\pgfsetroundjoin%
\definecolor{currentfill}{rgb}{1.000000,0.647059,0.000000}%
\pgfsetfillcolor{currentfill}%
\pgfsetlinewidth{1.003750pt}%
\definecolor{currentstroke}{rgb}{1.000000,0.647059,0.000000}%
\pgfsetstrokecolor{currentstroke}%
\pgfsetdash{}{0pt}%
\pgfpathmoveto{\pgfqpoint{0.950303in}{2.702496in}}%
\pgfpathlineto{\pgfqpoint{0.991970in}{2.785829in}}%
\pgfpathlineto{\pgfqpoint{0.908636in}{2.785829in}}%
\pgfpathclose%
\pgfusepath{stroke,fill}%
\end{pgfscope}%
\begin{pgfscope}%
\pgfpathrectangle{\pgfqpoint{0.800000in}{0.528000in}}{\pgfqpoint{4.960000in}{3.696000in}} %
\pgfusepath{clip}%
\pgfsetbuttcap%
\pgfsetroundjoin%
\definecolor{currentfill}{rgb}{1.000000,0.647059,0.000000}%
\pgfsetfillcolor{currentfill}%
\pgfsetlinewidth{1.003750pt}%
\definecolor{currentstroke}{rgb}{1.000000,0.647059,0.000000}%
\pgfsetstrokecolor{currentstroke}%
\pgfsetdash{}{0pt}%
\pgfpathmoveto{\pgfqpoint{1.000404in}{2.642758in}}%
\pgfpathlineto{\pgfqpoint{1.042071in}{2.726092in}}%
\pgfpathlineto{\pgfqpoint{0.958737in}{2.726092in}}%
\pgfpathclose%
\pgfusepath{stroke,fill}%
\end{pgfscope}%
\begin{pgfscope}%
\pgfpathrectangle{\pgfqpoint{0.800000in}{0.528000in}}{\pgfqpoint{4.960000in}{3.696000in}} %
\pgfusepath{clip}%
\pgfsetbuttcap%
\pgfsetroundjoin%
\definecolor{currentfill}{rgb}{1.000000,0.647059,0.000000}%
\pgfsetfillcolor{currentfill}%
\pgfsetlinewidth{1.003750pt}%
\definecolor{currentstroke}{rgb}{1.000000,0.647059,0.000000}%
\pgfsetstrokecolor{currentstroke}%
\pgfsetdash{}{0pt}%
\pgfpathmoveto{\pgfqpoint{1.050505in}{5.287553in}}%
\pgfpathlineto{\pgfqpoint{1.092172in}{5.370886in}}%
\pgfpathlineto{\pgfqpoint{1.008838in}{5.370886in}}%
\pgfpathclose%
\pgfusepath{stroke,fill}%
\end{pgfscope}%
\begin{pgfscope}%
\pgfpathrectangle{\pgfqpoint{0.800000in}{0.528000in}}{\pgfqpoint{4.960000in}{3.696000in}} %
\pgfusepath{clip}%
\pgfsetbuttcap%
\pgfsetroundjoin%
\definecolor{currentfill}{rgb}{1.000000,0.647059,0.000000}%
\pgfsetfillcolor{currentfill}%
\pgfsetlinewidth{1.003750pt}%
\definecolor{currentstroke}{rgb}{1.000000,0.647059,0.000000}%
\pgfsetstrokecolor{currentstroke}%
\pgfsetdash{}{0pt}%
\pgfpathmoveto{\pgfqpoint{1.100606in}{1.148267in}}%
\pgfpathlineto{\pgfqpoint{1.142273in}{1.231601in}}%
\pgfpathlineto{\pgfqpoint{1.058939in}{1.231601in}}%
\pgfpathclose%
\pgfusepath{stroke,fill}%
\end{pgfscope}%
\begin{pgfscope}%
\pgfpathrectangle{\pgfqpoint{0.800000in}{0.528000in}}{\pgfqpoint{4.960000in}{3.696000in}} %
\pgfusepath{clip}%
\pgfsetbuttcap%
\pgfsetroundjoin%
\definecolor{currentfill}{rgb}{1.000000,0.647059,0.000000}%
\pgfsetfillcolor{currentfill}%
\pgfsetlinewidth{1.003750pt}%
\definecolor{currentstroke}{rgb}{1.000000,0.647059,0.000000}%
\pgfsetstrokecolor{currentstroke}%
\pgfsetdash{}{0pt}%
\pgfpathmoveto{\pgfqpoint{1.150707in}{2.771214in}}%
\pgfpathlineto{\pgfqpoint{1.192374in}{2.854548in}}%
\pgfpathlineto{\pgfqpoint{1.109040in}{2.854548in}}%
\pgfpathclose%
\pgfusepath{stroke,fill}%
\end{pgfscope}%
\begin{pgfscope}%
\pgfpathrectangle{\pgfqpoint{0.800000in}{0.528000in}}{\pgfqpoint{4.960000in}{3.696000in}} %
\pgfusepath{clip}%
\pgfsetbuttcap%
\pgfsetroundjoin%
\definecolor{currentfill}{rgb}{1.000000,0.647059,0.000000}%
\pgfsetfillcolor{currentfill}%
\pgfsetlinewidth{1.003750pt}%
\definecolor{currentstroke}{rgb}{1.000000,0.647059,0.000000}%
\pgfsetstrokecolor{currentstroke}%
\pgfsetdash{}{0pt}%
\pgfpathmoveto{\pgfqpoint{1.200808in}{1.874288in}}%
\pgfpathlineto{\pgfqpoint{1.242475in}{1.957621in}}%
\pgfpathlineto{\pgfqpoint{1.159141in}{1.957621in}}%
\pgfpathclose%
\pgfusepath{stroke,fill}%
\end{pgfscope}%
\begin{pgfscope}%
\pgfpathrectangle{\pgfqpoint{0.800000in}{0.528000in}}{\pgfqpoint{4.960000in}{3.696000in}} %
\pgfusepath{clip}%
\pgfsetbuttcap%
\pgfsetroundjoin%
\definecolor{currentfill}{rgb}{1.000000,0.647059,0.000000}%
\pgfsetfillcolor{currentfill}%
\pgfsetlinewidth{1.003750pt}%
\definecolor{currentstroke}{rgb}{1.000000,0.647059,0.000000}%
\pgfsetstrokecolor{currentstroke}%
\pgfsetdash{}{0pt}%
\pgfpathmoveto{\pgfqpoint{1.250909in}{2.035079in}}%
\pgfpathlineto{\pgfqpoint{1.292576in}{2.118413in}}%
\pgfpathlineto{\pgfqpoint{1.209242in}{2.118413in}}%
\pgfpathclose%
\pgfusepath{stroke,fill}%
\end{pgfscope}%
\begin{pgfscope}%
\pgfpathrectangle{\pgfqpoint{0.800000in}{0.528000in}}{\pgfqpoint{4.960000in}{3.696000in}} %
\pgfusepath{clip}%
\pgfsetbuttcap%
\pgfsetroundjoin%
\definecolor{currentfill}{rgb}{1.000000,0.647059,0.000000}%
\pgfsetfillcolor{currentfill}%
\pgfsetlinewidth{1.003750pt}%
\definecolor{currentstroke}{rgb}{1.000000,0.647059,0.000000}%
\pgfsetstrokecolor{currentstroke}%
\pgfsetdash{}{0pt}%
\pgfpathmoveto{\pgfqpoint{1.301010in}{2.153348in}}%
\pgfpathlineto{\pgfqpoint{1.342677in}{2.236681in}}%
\pgfpathlineto{\pgfqpoint{1.259343in}{2.236681in}}%
\pgfpathclose%
\pgfusepath{stroke,fill}%
\end{pgfscope}%
\begin{pgfscope}%
\pgfpathrectangle{\pgfqpoint{0.800000in}{0.528000in}}{\pgfqpoint{4.960000in}{3.696000in}} %
\pgfusepath{clip}%
\pgfsetbuttcap%
\pgfsetroundjoin%
\definecolor{currentfill}{rgb}{1.000000,0.647059,0.000000}%
\pgfsetfillcolor{currentfill}%
\pgfsetlinewidth{1.003750pt}%
\definecolor{currentstroke}{rgb}{1.000000,0.647059,0.000000}%
\pgfsetstrokecolor{currentstroke}%
\pgfsetdash{}{0pt}%
\pgfpathmoveto{\pgfqpoint{1.351111in}{1.666477in}}%
\pgfpathlineto{\pgfqpoint{1.392778in}{1.749811in}}%
\pgfpathlineto{\pgfqpoint{1.309444in}{1.749811in}}%
\pgfpathclose%
\pgfusepath{stroke,fill}%
\end{pgfscope}%
\begin{pgfscope}%
\pgfpathrectangle{\pgfqpoint{0.800000in}{0.528000in}}{\pgfqpoint{4.960000in}{3.696000in}} %
\pgfusepath{clip}%
\pgfsetbuttcap%
\pgfsetroundjoin%
\definecolor{currentfill}{rgb}{1.000000,0.647059,0.000000}%
\pgfsetfillcolor{currentfill}%
\pgfsetlinewidth{1.003750pt}%
\definecolor{currentstroke}{rgb}{1.000000,0.647059,0.000000}%
\pgfsetstrokecolor{currentstroke}%
\pgfsetdash{}{0pt}%
\pgfpathmoveto{\pgfqpoint{1.401212in}{2.161171in}}%
\pgfpathlineto{\pgfqpoint{1.442879in}{2.244504in}}%
\pgfpathlineto{\pgfqpoint{1.359545in}{2.244504in}}%
\pgfpathclose%
\pgfusepath{stroke,fill}%
\end{pgfscope}%
\begin{pgfscope}%
\pgfpathrectangle{\pgfqpoint{0.800000in}{0.528000in}}{\pgfqpoint{4.960000in}{3.696000in}} %
\pgfusepath{clip}%
\pgfsetbuttcap%
\pgfsetroundjoin%
\definecolor{currentfill}{rgb}{1.000000,0.647059,0.000000}%
\pgfsetfillcolor{currentfill}%
\pgfsetlinewidth{1.003750pt}%
\definecolor{currentstroke}{rgb}{1.000000,0.647059,0.000000}%
\pgfsetstrokecolor{currentstroke}%
\pgfsetdash{}{0pt}%
\pgfpathmoveto{\pgfqpoint{1.451313in}{5.088047in}}%
\pgfpathlineto{\pgfqpoint{1.492980in}{5.171380in}}%
\pgfpathlineto{\pgfqpoint{1.409646in}{5.171380in}}%
\pgfpathclose%
\pgfusepath{stroke,fill}%
\end{pgfscope}%
\begin{pgfscope}%
\pgfpathrectangle{\pgfqpoint{0.800000in}{0.528000in}}{\pgfqpoint{4.960000in}{3.696000in}} %
\pgfusepath{clip}%
\pgfsetbuttcap%
\pgfsetroundjoin%
\definecolor{currentfill}{rgb}{1.000000,0.647059,0.000000}%
\pgfsetfillcolor{currentfill}%
\pgfsetlinewidth{1.003750pt}%
\definecolor{currentstroke}{rgb}{1.000000,0.647059,0.000000}%
\pgfsetstrokecolor{currentstroke}%
\pgfsetdash{}{0pt}%
\pgfpathmoveto{\pgfqpoint{1.501414in}{4.137191in}}%
\pgfpathlineto{\pgfqpoint{1.543081in}{4.220524in}}%
\pgfpathlineto{\pgfqpoint{1.459747in}{4.220524in}}%
\pgfpathclose%
\pgfusepath{stroke,fill}%
\end{pgfscope}%
\begin{pgfscope}%
\pgfpathrectangle{\pgfqpoint{0.800000in}{0.528000in}}{\pgfqpoint{4.960000in}{3.696000in}} %
\pgfusepath{clip}%
\pgfsetbuttcap%
\pgfsetroundjoin%
\definecolor{currentfill}{rgb}{1.000000,0.647059,0.000000}%
\pgfsetfillcolor{currentfill}%
\pgfsetlinewidth{1.003750pt}%
\definecolor{currentstroke}{rgb}{1.000000,0.647059,0.000000}%
\pgfsetstrokecolor{currentstroke}%
\pgfsetdash{}{0pt}%
\pgfpathmoveto{\pgfqpoint{1.551515in}{1.834261in}}%
\pgfpathlineto{\pgfqpoint{1.593182in}{1.917594in}}%
\pgfpathlineto{\pgfqpoint{1.509848in}{1.917594in}}%
\pgfpathclose%
\pgfusepath{stroke,fill}%
\end{pgfscope}%
\begin{pgfscope}%
\pgfpathrectangle{\pgfqpoint{0.800000in}{0.528000in}}{\pgfqpoint{4.960000in}{3.696000in}} %
\pgfusepath{clip}%
\pgfsetbuttcap%
\pgfsetroundjoin%
\definecolor{currentfill}{rgb}{1.000000,0.647059,0.000000}%
\pgfsetfillcolor{currentfill}%
\pgfsetlinewidth{1.003750pt}%
\definecolor{currentstroke}{rgb}{1.000000,0.647059,0.000000}%
\pgfsetstrokecolor{currentstroke}%
\pgfsetdash{}{0pt}%
\pgfpathmoveto{\pgfqpoint{1.601616in}{4.738554in}}%
\pgfpathlineto{\pgfqpoint{1.643283in}{4.821887in}}%
\pgfpathlineto{\pgfqpoint{1.559949in}{4.821887in}}%
\pgfpathclose%
\pgfusepath{stroke,fill}%
\end{pgfscope}%
\begin{pgfscope}%
\pgfpathrectangle{\pgfqpoint{0.800000in}{0.528000in}}{\pgfqpoint{4.960000in}{3.696000in}} %
\pgfusepath{clip}%
\pgfsetbuttcap%
\pgfsetroundjoin%
\definecolor{currentfill}{rgb}{1.000000,0.647059,0.000000}%
\pgfsetfillcolor{currentfill}%
\pgfsetlinewidth{1.003750pt}%
\definecolor{currentstroke}{rgb}{1.000000,0.647059,0.000000}%
\pgfsetstrokecolor{currentstroke}%
\pgfsetdash{}{0pt}%
\pgfpathmoveto{\pgfqpoint{1.651717in}{3.607466in}}%
\pgfpathlineto{\pgfqpoint{1.693384in}{3.690800in}}%
\pgfpathlineto{\pgfqpoint{1.610051in}{3.690800in}}%
\pgfpathclose%
\pgfusepath{stroke,fill}%
\end{pgfscope}%
\begin{pgfscope}%
\pgfpathrectangle{\pgfqpoint{0.800000in}{0.528000in}}{\pgfqpoint{4.960000in}{3.696000in}} %
\pgfusepath{clip}%
\pgfsetbuttcap%
\pgfsetroundjoin%
\definecolor{currentfill}{rgb}{1.000000,0.647059,0.000000}%
\pgfsetfillcolor{currentfill}%
\pgfsetlinewidth{1.003750pt}%
\definecolor{currentstroke}{rgb}{1.000000,0.647059,0.000000}%
\pgfsetstrokecolor{currentstroke}%
\pgfsetdash{}{0pt}%
\pgfpathmoveto{\pgfqpoint{1.701818in}{2.528560in}}%
\pgfpathlineto{\pgfqpoint{1.743485in}{2.611894in}}%
\pgfpathlineto{\pgfqpoint{1.660152in}{2.611894in}}%
\pgfpathclose%
\pgfusepath{stroke,fill}%
\end{pgfscope}%
\begin{pgfscope}%
\pgfpathrectangle{\pgfqpoint{0.800000in}{0.528000in}}{\pgfqpoint{4.960000in}{3.696000in}} %
\pgfusepath{clip}%
\pgfsetbuttcap%
\pgfsetroundjoin%
\definecolor{currentfill}{rgb}{1.000000,0.647059,0.000000}%
\pgfsetfillcolor{currentfill}%
\pgfsetlinewidth{1.003750pt}%
\definecolor{currentstroke}{rgb}{1.000000,0.647059,0.000000}%
\pgfsetstrokecolor{currentstroke}%
\pgfsetdash{}{0pt}%
\pgfpathmoveto{\pgfqpoint{1.751919in}{1.902015in}}%
\pgfpathlineto{\pgfqpoint{1.793586in}{1.985348in}}%
\pgfpathlineto{\pgfqpoint{1.710253in}{1.985348in}}%
\pgfpathclose%
\pgfusepath{stroke,fill}%
\end{pgfscope}%
\begin{pgfscope}%
\pgfpathrectangle{\pgfqpoint{0.800000in}{0.528000in}}{\pgfqpoint{4.960000in}{3.696000in}} %
\pgfusepath{clip}%
\pgfsetbuttcap%
\pgfsetroundjoin%
\definecolor{currentfill}{rgb}{1.000000,0.647059,0.000000}%
\pgfsetfillcolor{currentfill}%
\pgfsetlinewidth{1.003750pt}%
\definecolor{currentstroke}{rgb}{1.000000,0.647059,0.000000}%
\pgfsetstrokecolor{currentstroke}%
\pgfsetdash{}{0pt}%
\pgfpathmoveto{\pgfqpoint{1.802020in}{2.412656in}}%
\pgfpathlineto{\pgfqpoint{1.843687in}{2.495989in}}%
\pgfpathlineto{\pgfqpoint{1.760354in}{2.495989in}}%
\pgfpathclose%
\pgfusepath{stroke,fill}%
\end{pgfscope}%
\begin{pgfscope}%
\pgfpathrectangle{\pgfqpoint{0.800000in}{0.528000in}}{\pgfqpoint{4.960000in}{3.696000in}} %
\pgfusepath{clip}%
\pgfsetbuttcap%
\pgfsetroundjoin%
\definecolor{currentfill}{rgb}{1.000000,0.647059,0.000000}%
\pgfsetfillcolor{currentfill}%
\pgfsetlinewidth{1.003750pt}%
\definecolor{currentstroke}{rgb}{1.000000,0.647059,0.000000}%
\pgfsetstrokecolor{currentstroke}%
\pgfsetdash{}{0pt}%
\pgfpathmoveto{\pgfqpoint{1.852121in}{2.548203in}}%
\pgfpathlineto{\pgfqpoint{1.893788in}{2.631536in}}%
\pgfpathlineto{\pgfqpoint{1.810455in}{2.631536in}}%
\pgfpathclose%
\pgfusepath{stroke,fill}%
\end{pgfscope}%
\begin{pgfscope}%
\pgfpathrectangle{\pgfqpoint{0.800000in}{0.528000in}}{\pgfqpoint{4.960000in}{3.696000in}} %
\pgfusepath{clip}%
\pgfsetbuttcap%
\pgfsetroundjoin%
\definecolor{currentfill}{rgb}{1.000000,0.647059,0.000000}%
\pgfsetfillcolor{currentfill}%
\pgfsetlinewidth{1.003750pt}%
\definecolor{currentstroke}{rgb}{1.000000,0.647059,0.000000}%
\pgfsetstrokecolor{currentstroke}%
\pgfsetdash{}{0pt}%
\pgfpathmoveto{\pgfqpoint{1.902222in}{1.907558in}}%
\pgfpathlineto{\pgfqpoint{1.943889in}{1.990891in}}%
\pgfpathlineto{\pgfqpoint{1.860556in}{1.990891in}}%
\pgfpathclose%
\pgfusepath{stroke,fill}%
\end{pgfscope}%
\begin{pgfscope}%
\pgfpathrectangle{\pgfqpoint{0.800000in}{0.528000in}}{\pgfqpoint{4.960000in}{3.696000in}} %
\pgfusepath{clip}%
\pgfsetbuttcap%
\pgfsetroundjoin%
\definecolor{currentfill}{rgb}{1.000000,0.647059,0.000000}%
\pgfsetfillcolor{currentfill}%
\pgfsetlinewidth{1.003750pt}%
\definecolor{currentstroke}{rgb}{1.000000,0.647059,0.000000}%
\pgfsetstrokecolor{currentstroke}%
\pgfsetdash{}{0pt}%
\pgfpathmoveto{\pgfqpoint{1.952323in}{3.122247in}}%
\pgfpathlineto{\pgfqpoint{1.993990in}{3.205581in}}%
\pgfpathlineto{\pgfqpoint{1.910657in}{3.205581in}}%
\pgfpathclose%
\pgfusepath{stroke,fill}%
\end{pgfscope}%
\begin{pgfscope}%
\pgfpathrectangle{\pgfqpoint{0.800000in}{0.528000in}}{\pgfqpoint{4.960000in}{3.696000in}} %
\pgfusepath{clip}%
\pgfsetbuttcap%
\pgfsetroundjoin%
\definecolor{currentfill}{rgb}{1.000000,0.647059,0.000000}%
\pgfsetfillcolor{currentfill}%
\pgfsetlinewidth{1.003750pt}%
\definecolor{currentstroke}{rgb}{1.000000,0.647059,0.000000}%
\pgfsetstrokecolor{currentstroke}%
\pgfsetdash{}{0pt}%
\pgfpathmoveto{\pgfqpoint{2.002424in}{6.410703in}}%
\pgfpathlineto{\pgfqpoint{2.044091in}{6.494036in}}%
\pgfpathlineto{\pgfqpoint{1.960758in}{6.494036in}}%
\pgfpathclose%
\pgfusepath{stroke,fill}%
\end{pgfscope}%
\begin{pgfscope}%
\pgfpathrectangle{\pgfqpoint{0.800000in}{0.528000in}}{\pgfqpoint{4.960000in}{3.696000in}} %
\pgfusepath{clip}%
\pgfsetbuttcap%
\pgfsetroundjoin%
\definecolor{currentfill}{rgb}{1.000000,0.647059,0.000000}%
\pgfsetfillcolor{currentfill}%
\pgfsetlinewidth{1.003750pt}%
\definecolor{currentstroke}{rgb}{1.000000,0.647059,0.000000}%
\pgfsetstrokecolor{currentstroke}%
\pgfsetdash{}{0pt}%
\pgfpathmoveto{\pgfqpoint{2.052525in}{5.313310in}}%
\pgfpathlineto{\pgfqpoint{2.094192in}{5.396643in}}%
\pgfpathlineto{\pgfqpoint{2.010859in}{5.396643in}}%
\pgfpathclose%
\pgfusepath{stroke,fill}%
\end{pgfscope}%
\begin{pgfscope}%
\pgfpathrectangle{\pgfqpoint{0.800000in}{0.528000in}}{\pgfqpoint{4.960000in}{3.696000in}} %
\pgfusepath{clip}%
\pgfsetbuttcap%
\pgfsetroundjoin%
\definecolor{currentfill}{rgb}{1.000000,0.647059,0.000000}%
\pgfsetfillcolor{currentfill}%
\pgfsetlinewidth{1.003750pt}%
\definecolor{currentstroke}{rgb}{1.000000,0.647059,0.000000}%
\pgfsetstrokecolor{currentstroke}%
\pgfsetdash{}{0pt}%
\pgfpathmoveto{\pgfqpoint{2.102626in}{3.923920in}}%
\pgfpathlineto{\pgfqpoint{2.144293in}{4.007254in}}%
\pgfpathlineto{\pgfqpoint{2.060960in}{4.007254in}}%
\pgfpathclose%
\pgfusepath{stroke,fill}%
\end{pgfscope}%
\begin{pgfscope}%
\pgfpathrectangle{\pgfqpoint{0.800000in}{0.528000in}}{\pgfqpoint{4.960000in}{3.696000in}} %
\pgfusepath{clip}%
\pgfsetbuttcap%
\pgfsetroundjoin%
\definecolor{currentfill}{rgb}{1.000000,0.647059,0.000000}%
\pgfsetfillcolor{currentfill}%
\pgfsetlinewidth{1.003750pt}%
\definecolor{currentstroke}{rgb}{1.000000,0.647059,0.000000}%
\pgfsetstrokecolor{currentstroke}%
\pgfsetdash{}{0pt}%
\pgfpathmoveto{\pgfqpoint{2.152727in}{6.173705in}}%
\pgfpathlineto{\pgfqpoint{2.194394in}{6.257038in}}%
\pgfpathlineto{\pgfqpoint{2.111061in}{6.257038in}}%
\pgfpathclose%
\pgfusepath{stroke,fill}%
\end{pgfscope}%
\begin{pgfscope}%
\pgfpathrectangle{\pgfqpoint{0.800000in}{0.528000in}}{\pgfqpoint{4.960000in}{3.696000in}} %
\pgfusepath{clip}%
\pgfsetbuttcap%
\pgfsetroundjoin%
\definecolor{currentfill}{rgb}{1.000000,0.647059,0.000000}%
\pgfsetfillcolor{currentfill}%
\pgfsetlinewidth{1.003750pt}%
\definecolor{currentstroke}{rgb}{1.000000,0.647059,0.000000}%
\pgfsetstrokecolor{currentstroke}%
\pgfsetdash{}{0pt}%
\pgfpathmoveto{\pgfqpoint{2.202828in}{3.907119in}}%
\pgfpathlineto{\pgfqpoint{2.244495in}{3.990453in}}%
\pgfpathlineto{\pgfqpoint{2.161162in}{3.990453in}}%
\pgfpathclose%
\pgfusepath{stroke,fill}%
\end{pgfscope}%
\begin{pgfscope}%
\pgfpathrectangle{\pgfqpoint{0.800000in}{0.528000in}}{\pgfqpoint{4.960000in}{3.696000in}} %
\pgfusepath{clip}%
\pgfsetbuttcap%
\pgfsetroundjoin%
\definecolor{currentfill}{rgb}{1.000000,0.647059,0.000000}%
\pgfsetfillcolor{currentfill}%
\pgfsetlinewidth{1.003750pt}%
\definecolor{currentstroke}{rgb}{1.000000,0.647059,0.000000}%
\pgfsetstrokecolor{currentstroke}%
\pgfsetdash{}{0pt}%
\pgfpathmoveto{\pgfqpoint{2.252929in}{6.329591in}}%
\pgfpathlineto{\pgfqpoint{2.294596in}{6.412925in}}%
\pgfpathlineto{\pgfqpoint{2.211263in}{6.412925in}}%
\pgfpathclose%
\pgfusepath{stroke,fill}%
\end{pgfscope}%
\begin{pgfscope}%
\pgfpathrectangle{\pgfqpoint{0.800000in}{0.528000in}}{\pgfqpoint{4.960000in}{3.696000in}} %
\pgfusepath{clip}%
\pgfsetbuttcap%
\pgfsetroundjoin%
\definecolor{currentfill}{rgb}{1.000000,0.647059,0.000000}%
\pgfsetfillcolor{currentfill}%
\pgfsetlinewidth{1.003750pt}%
\definecolor{currentstroke}{rgb}{1.000000,0.647059,0.000000}%
\pgfsetstrokecolor{currentstroke}%
\pgfsetdash{}{0pt}%
\pgfpathmoveto{\pgfqpoint{2.303030in}{5.495913in}}%
\pgfpathlineto{\pgfqpoint{2.344697in}{5.579246in}}%
\pgfpathlineto{\pgfqpoint{2.261364in}{5.579246in}}%
\pgfpathclose%
\pgfusepath{stroke,fill}%
\end{pgfscope}%
\begin{pgfscope}%
\pgfpathrectangle{\pgfqpoint{0.800000in}{0.528000in}}{\pgfqpoint{4.960000in}{3.696000in}} %
\pgfusepath{clip}%
\pgfsetbuttcap%
\pgfsetroundjoin%
\definecolor{currentfill}{rgb}{1.000000,0.647059,0.000000}%
\pgfsetfillcolor{currentfill}%
\pgfsetlinewidth{1.003750pt}%
\definecolor{currentstroke}{rgb}{1.000000,0.647059,0.000000}%
\pgfsetstrokecolor{currentstroke}%
\pgfsetdash{}{0pt}%
\pgfpathmoveto{\pgfqpoint{2.353131in}{5.284425in}}%
\pgfpathlineto{\pgfqpoint{2.394798in}{5.367758in}}%
\pgfpathlineto{\pgfqpoint{2.311465in}{5.367758in}}%
\pgfpathclose%
\pgfusepath{stroke,fill}%
\end{pgfscope}%
\begin{pgfscope}%
\pgfpathrectangle{\pgfqpoint{0.800000in}{0.528000in}}{\pgfqpoint{4.960000in}{3.696000in}} %
\pgfusepath{clip}%
\pgfsetbuttcap%
\pgfsetroundjoin%
\definecolor{currentfill}{rgb}{1.000000,0.647059,0.000000}%
\pgfsetfillcolor{currentfill}%
\pgfsetlinewidth{1.003750pt}%
\definecolor{currentstroke}{rgb}{1.000000,0.647059,0.000000}%
\pgfsetstrokecolor{currentstroke}%
\pgfsetdash{}{0pt}%
\pgfpathmoveto{\pgfqpoint{2.403232in}{4.847724in}}%
\pgfpathlineto{\pgfqpoint{2.444899in}{4.931058in}}%
\pgfpathlineto{\pgfqpoint{2.361566in}{4.931058in}}%
\pgfpathclose%
\pgfusepath{stroke,fill}%
\end{pgfscope}%
\begin{pgfscope}%
\pgfpathrectangle{\pgfqpoint{0.800000in}{0.528000in}}{\pgfqpoint{4.960000in}{3.696000in}} %
\pgfusepath{clip}%
\pgfsetbuttcap%
\pgfsetroundjoin%
\definecolor{currentfill}{rgb}{1.000000,0.647059,0.000000}%
\pgfsetfillcolor{currentfill}%
\pgfsetlinewidth{1.003750pt}%
\definecolor{currentstroke}{rgb}{1.000000,0.647059,0.000000}%
\pgfsetstrokecolor{currentstroke}%
\pgfsetdash{}{0pt}%
\pgfpathmoveto{\pgfqpoint{2.453333in}{4.103277in}}%
\pgfpathlineto{\pgfqpoint{2.495000in}{4.186610in}}%
\pgfpathlineto{\pgfqpoint{2.411667in}{4.186610in}}%
\pgfpathclose%
\pgfusepath{stroke,fill}%
\end{pgfscope}%
\begin{pgfscope}%
\pgfpathrectangle{\pgfqpoint{0.800000in}{0.528000in}}{\pgfqpoint{4.960000in}{3.696000in}} %
\pgfusepath{clip}%
\pgfsetbuttcap%
\pgfsetroundjoin%
\definecolor{currentfill}{rgb}{1.000000,0.647059,0.000000}%
\pgfsetfillcolor{currentfill}%
\pgfsetlinewidth{1.003750pt}%
\definecolor{currentstroke}{rgb}{1.000000,0.647059,0.000000}%
\pgfsetstrokecolor{currentstroke}%
\pgfsetdash{}{0pt}%
\pgfpathmoveto{\pgfqpoint{2.503434in}{5.221209in}}%
\pgfpathlineto{\pgfqpoint{2.545101in}{5.304542in}}%
\pgfpathlineto{\pgfqpoint{2.461768in}{5.304542in}}%
\pgfpathclose%
\pgfusepath{stroke,fill}%
\end{pgfscope}%
\begin{pgfscope}%
\pgfpathrectangle{\pgfqpoint{0.800000in}{0.528000in}}{\pgfqpoint{4.960000in}{3.696000in}} %
\pgfusepath{clip}%
\pgfsetbuttcap%
\pgfsetroundjoin%
\definecolor{currentfill}{rgb}{1.000000,0.647059,0.000000}%
\pgfsetfillcolor{currentfill}%
\pgfsetlinewidth{1.003750pt}%
\definecolor{currentstroke}{rgb}{1.000000,0.647059,0.000000}%
\pgfsetstrokecolor{currentstroke}%
\pgfsetdash{}{0pt}%
\pgfpathmoveto{\pgfqpoint{2.553535in}{2.489146in}}%
\pgfpathlineto{\pgfqpoint{2.595202in}{2.572480in}}%
\pgfpathlineto{\pgfqpoint{2.511869in}{2.572480in}}%
\pgfpathclose%
\pgfusepath{stroke,fill}%
\end{pgfscope}%
\begin{pgfscope}%
\pgfpathrectangle{\pgfqpoint{0.800000in}{0.528000in}}{\pgfqpoint{4.960000in}{3.696000in}} %
\pgfusepath{clip}%
\pgfsetbuttcap%
\pgfsetroundjoin%
\definecolor{currentfill}{rgb}{1.000000,0.647059,0.000000}%
\pgfsetfillcolor{currentfill}%
\pgfsetlinewidth{1.003750pt}%
\definecolor{currentstroke}{rgb}{1.000000,0.647059,0.000000}%
\pgfsetstrokecolor{currentstroke}%
\pgfsetdash{}{0pt}%
\pgfpathmoveto{\pgfqpoint{2.603636in}{3.692620in}}%
\pgfpathlineto{\pgfqpoint{2.645303in}{3.775953in}}%
\pgfpathlineto{\pgfqpoint{2.561970in}{3.775953in}}%
\pgfpathclose%
\pgfusepath{stroke,fill}%
\end{pgfscope}%
\begin{pgfscope}%
\pgfpathrectangle{\pgfqpoint{0.800000in}{0.528000in}}{\pgfqpoint{4.960000in}{3.696000in}} %
\pgfusepath{clip}%
\pgfsetbuttcap%
\pgfsetroundjoin%
\definecolor{currentfill}{rgb}{1.000000,0.647059,0.000000}%
\pgfsetfillcolor{currentfill}%
\pgfsetlinewidth{1.003750pt}%
\definecolor{currentstroke}{rgb}{1.000000,0.647059,0.000000}%
\pgfsetstrokecolor{currentstroke}%
\pgfsetdash{}{0pt}%
\pgfpathmoveto{\pgfqpoint{2.653737in}{1.965013in}}%
\pgfpathlineto{\pgfqpoint{2.695404in}{2.048347in}}%
\pgfpathlineto{\pgfqpoint{2.612071in}{2.048347in}}%
\pgfpathclose%
\pgfusepath{stroke,fill}%
\end{pgfscope}%
\begin{pgfscope}%
\pgfpathrectangle{\pgfqpoint{0.800000in}{0.528000in}}{\pgfqpoint{4.960000in}{3.696000in}} %
\pgfusepath{clip}%
\pgfsetbuttcap%
\pgfsetroundjoin%
\definecolor{currentfill}{rgb}{1.000000,0.647059,0.000000}%
\pgfsetfillcolor{currentfill}%
\pgfsetlinewidth{1.003750pt}%
\definecolor{currentstroke}{rgb}{1.000000,0.647059,0.000000}%
\pgfsetstrokecolor{currentstroke}%
\pgfsetdash{}{0pt}%
\pgfpathmoveto{\pgfqpoint{2.703838in}{2.839286in}}%
\pgfpathlineto{\pgfqpoint{2.745505in}{2.922620in}}%
\pgfpathlineto{\pgfqpoint{2.662172in}{2.922620in}}%
\pgfpathclose%
\pgfusepath{stroke,fill}%
\end{pgfscope}%
\begin{pgfscope}%
\pgfpathrectangle{\pgfqpoint{0.800000in}{0.528000in}}{\pgfqpoint{4.960000in}{3.696000in}} %
\pgfusepath{clip}%
\pgfsetbuttcap%
\pgfsetroundjoin%
\definecolor{currentfill}{rgb}{1.000000,0.647059,0.000000}%
\pgfsetfillcolor{currentfill}%
\pgfsetlinewidth{1.003750pt}%
\definecolor{currentstroke}{rgb}{1.000000,0.647059,0.000000}%
\pgfsetstrokecolor{currentstroke}%
\pgfsetdash{}{0pt}%
\pgfpathmoveto{\pgfqpoint{2.753939in}{4.757648in}}%
\pgfpathlineto{\pgfqpoint{2.795606in}{4.840981in}}%
\pgfpathlineto{\pgfqpoint{2.712273in}{4.840981in}}%
\pgfpathclose%
\pgfusepath{stroke,fill}%
\end{pgfscope}%
\begin{pgfscope}%
\pgfpathrectangle{\pgfqpoint{0.800000in}{0.528000in}}{\pgfqpoint{4.960000in}{3.696000in}} %
\pgfusepath{clip}%
\pgfsetbuttcap%
\pgfsetroundjoin%
\definecolor{currentfill}{rgb}{1.000000,0.647059,0.000000}%
\pgfsetfillcolor{currentfill}%
\pgfsetlinewidth{1.003750pt}%
\definecolor{currentstroke}{rgb}{1.000000,0.647059,0.000000}%
\pgfsetstrokecolor{currentstroke}%
\pgfsetdash{}{0pt}%
\pgfpathmoveto{\pgfqpoint{2.804040in}{2.356106in}}%
\pgfpathlineto{\pgfqpoint{2.845707in}{2.439439in}}%
\pgfpathlineto{\pgfqpoint{2.762374in}{2.439439in}}%
\pgfpathclose%
\pgfusepath{stroke,fill}%
\end{pgfscope}%
\begin{pgfscope}%
\pgfpathrectangle{\pgfqpoint{0.800000in}{0.528000in}}{\pgfqpoint{4.960000in}{3.696000in}} %
\pgfusepath{clip}%
\pgfsetbuttcap%
\pgfsetroundjoin%
\definecolor{currentfill}{rgb}{1.000000,0.647059,0.000000}%
\pgfsetfillcolor{currentfill}%
\pgfsetlinewidth{1.003750pt}%
\definecolor{currentstroke}{rgb}{1.000000,0.647059,0.000000}%
\pgfsetstrokecolor{currentstroke}%
\pgfsetdash{}{0pt}%
\pgfpathmoveto{\pgfqpoint{2.854141in}{2.935956in}}%
\pgfpathlineto{\pgfqpoint{2.895808in}{3.019290in}}%
\pgfpathlineto{\pgfqpoint{2.812475in}{3.019290in}}%
\pgfpathclose%
\pgfusepath{stroke,fill}%
\end{pgfscope}%
\begin{pgfscope}%
\pgfpathrectangle{\pgfqpoint{0.800000in}{0.528000in}}{\pgfqpoint{4.960000in}{3.696000in}} %
\pgfusepath{clip}%
\pgfsetbuttcap%
\pgfsetroundjoin%
\definecolor{currentfill}{rgb}{1.000000,0.647059,0.000000}%
\pgfsetfillcolor{currentfill}%
\pgfsetlinewidth{1.003750pt}%
\definecolor{currentstroke}{rgb}{1.000000,0.647059,0.000000}%
\pgfsetstrokecolor{currentstroke}%
\pgfsetdash{}{0pt}%
\pgfpathmoveto{\pgfqpoint{2.904242in}{3.155769in}}%
\pgfpathlineto{\pgfqpoint{2.945909in}{3.239103in}}%
\pgfpathlineto{\pgfqpoint{2.862576in}{3.239103in}}%
\pgfpathclose%
\pgfusepath{stroke,fill}%
\end{pgfscope}%
\begin{pgfscope}%
\pgfpathrectangle{\pgfqpoint{0.800000in}{0.528000in}}{\pgfqpoint{4.960000in}{3.696000in}} %
\pgfusepath{clip}%
\pgfsetbuttcap%
\pgfsetroundjoin%
\definecolor{currentfill}{rgb}{1.000000,0.647059,0.000000}%
\pgfsetfillcolor{currentfill}%
\pgfsetlinewidth{1.003750pt}%
\definecolor{currentstroke}{rgb}{1.000000,0.647059,0.000000}%
\pgfsetstrokecolor{currentstroke}%
\pgfsetdash{}{0pt}%
\pgfpathmoveto{\pgfqpoint{2.954343in}{5.642572in}}%
\pgfpathlineto{\pgfqpoint{2.996010in}{5.725906in}}%
\pgfpathlineto{\pgfqpoint{2.912677in}{5.725906in}}%
\pgfpathclose%
\pgfusepath{stroke,fill}%
\end{pgfscope}%
\begin{pgfscope}%
\pgfpathrectangle{\pgfqpoint{0.800000in}{0.528000in}}{\pgfqpoint{4.960000in}{3.696000in}} %
\pgfusepath{clip}%
\pgfsetbuttcap%
\pgfsetroundjoin%
\definecolor{currentfill}{rgb}{1.000000,0.647059,0.000000}%
\pgfsetfillcolor{currentfill}%
\pgfsetlinewidth{1.003750pt}%
\definecolor{currentstroke}{rgb}{1.000000,0.647059,0.000000}%
\pgfsetstrokecolor{currentstroke}%
\pgfsetdash{}{0pt}%
\pgfpathmoveto{\pgfqpoint{3.004444in}{3.368145in}}%
\pgfpathlineto{\pgfqpoint{3.046111in}{3.451479in}}%
\pgfpathlineto{\pgfqpoint{2.962778in}{3.451479in}}%
\pgfpathclose%
\pgfusepath{stroke,fill}%
\end{pgfscope}%
\begin{pgfscope}%
\pgfpathrectangle{\pgfqpoint{0.800000in}{0.528000in}}{\pgfqpoint{4.960000in}{3.696000in}} %
\pgfusepath{clip}%
\pgfsetbuttcap%
\pgfsetroundjoin%
\definecolor{currentfill}{rgb}{1.000000,0.647059,0.000000}%
\pgfsetfillcolor{currentfill}%
\pgfsetlinewidth{1.003750pt}%
\definecolor{currentstroke}{rgb}{1.000000,0.647059,0.000000}%
\pgfsetstrokecolor{currentstroke}%
\pgfsetdash{}{0pt}%
\pgfpathmoveto{\pgfqpoint{3.054545in}{4.336455in}}%
\pgfpathlineto{\pgfqpoint{3.096212in}{4.419788in}}%
\pgfpathlineto{\pgfqpoint{3.012879in}{4.419788in}}%
\pgfpathclose%
\pgfusepath{stroke,fill}%
\end{pgfscope}%
\begin{pgfscope}%
\pgfpathrectangle{\pgfqpoint{0.800000in}{0.528000in}}{\pgfqpoint{4.960000in}{3.696000in}} %
\pgfusepath{clip}%
\pgfsetbuttcap%
\pgfsetroundjoin%
\definecolor{currentfill}{rgb}{1.000000,0.647059,0.000000}%
\pgfsetfillcolor{currentfill}%
\pgfsetlinewidth{1.003750pt}%
\definecolor{currentstroke}{rgb}{1.000000,0.647059,0.000000}%
\pgfsetstrokecolor{currentstroke}%
\pgfsetdash{}{0pt}%
\pgfpathmoveto{\pgfqpoint{3.104646in}{5.985662in}}%
\pgfpathlineto{\pgfqpoint{3.146313in}{6.068995in}}%
\pgfpathlineto{\pgfqpoint{3.062980in}{6.068995in}}%
\pgfpathclose%
\pgfusepath{stroke,fill}%
\end{pgfscope}%
\begin{pgfscope}%
\pgfpathrectangle{\pgfqpoint{0.800000in}{0.528000in}}{\pgfqpoint{4.960000in}{3.696000in}} %
\pgfusepath{clip}%
\pgfsetbuttcap%
\pgfsetroundjoin%
\definecolor{currentfill}{rgb}{1.000000,0.647059,0.000000}%
\pgfsetfillcolor{currentfill}%
\pgfsetlinewidth{1.003750pt}%
\definecolor{currentstroke}{rgb}{1.000000,0.647059,0.000000}%
\pgfsetstrokecolor{currentstroke}%
\pgfsetdash{}{0pt}%
\pgfpathmoveto{\pgfqpoint{3.154747in}{5.332747in}}%
\pgfpathlineto{\pgfqpoint{3.196414in}{5.416080in}}%
\pgfpathlineto{\pgfqpoint{3.113081in}{5.416080in}}%
\pgfpathclose%
\pgfusepath{stroke,fill}%
\end{pgfscope}%
\begin{pgfscope}%
\pgfpathrectangle{\pgfqpoint{0.800000in}{0.528000in}}{\pgfqpoint{4.960000in}{3.696000in}} %
\pgfusepath{clip}%
\pgfsetbuttcap%
\pgfsetroundjoin%
\definecolor{currentfill}{rgb}{1.000000,0.647059,0.000000}%
\pgfsetfillcolor{currentfill}%
\pgfsetlinewidth{1.003750pt}%
\definecolor{currentstroke}{rgb}{1.000000,0.647059,0.000000}%
\pgfsetstrokecolor{currentstroke}%
\pgfsetdash{}{0pt}%
\pgfpathmoveto{\pgfqpoint{3.204848in}{6.902353in}}%
\pgfpathlineto{\pgfqpoint{3.246515in}{6.985686in}}%
\pgfpathlineto{\pgfqpoint{3.163182in}{6.985686in}}%
\pgfpathclose%
\pgfusepath{stroke,fill}%
\end{pgfscope}%
\begin{pgfscope}%
\pgfpathrectangle{\pgfqpoint{0.800000in}{0.528000in}}{\pgfqpoint{4.960000in}{3.696000in}} %
\pgfusepath{clip}%
\pgfsetbuttcap%
\pgfsetroundjoin%
\definecolor{currentfill}{rgb}{1.000000,0.647059,0.000000}%
\pgfsetfillcolor{currentfill}%
\pgfsetlinewidth{1.003750pt}%
\definecolor{currentstroke}{rgb}{1.000000,0.647059,0.000000}%
\pgfsetstrokecolor{currentstroke}%
\pgfsetdash{}{0pt}%
\pgfpathmoveto{\pgfqpoint{3.254949in}{5.456591in}}%
\pgfpathlineto{\pgfqpoint{3.296616in}{5.539924in}}%
\pgfpathlineto{\pgfqpoint{3.213283in}{5.539924in}}%
\pgfpathclose%
\pgfusepath{stroke,fill}%
\end{pgfscope}%
\begin{pgfscope}%
\pgfpathrectangle{\pgfqpoint{0.800000in}{0.528000in}}{\pgfqpoint{4.960000in}{3.696000in}} %
\pgfusepath{clip}%
\pgfsetbuttcap%
\pgfsetroundjoin%
\definecolor{currentfill}{rgb}{1.000000,0.647059,0.000000}%
\pgfsetfillcolor{currentfill}%
\pgfsetlinewidth{1.003750pt}%
\definecolor{currentstroke}{rgb}{1.000000,0.647059,0.000000}%
\pgfsetstrokecolor{currentstroke}%
\pgfsetdash{}{0pt}%
\pgfpathmoveto{\pgfqpoint{3.305051in}{4.884177in}}%
\pgfpathlineto{\pgfqpoint{3.346717in}{4.967511in}}%
\pgfpathlineto{\pgfqpoint{3.263384in}{4.967511in}}%
\pgfpathclose%
\pgfusepath{stroke,fill}%
\end{pgfscope}%
\begin{pgfscope}%
\pgfpathrectangle{\pgfqpoint{0.800000in}{0.528000in}}{\pgfqpoint{4.960000in}{3.696000in}} %
\pgfusepath{clip}%
\pgfsetbuttcap%
\pgfsetroundjoin%
\definecolor{currentfill}{rgb}{1.000000,0.647059,0.000000}%
\pgfsetfillcolor{currentfill}%
\pgfsetlinewidth{1.003750pt}%
\definecolor{currentstroke}{rgb}{1.000000,0.647059,0.000000}%
\pgfsetstrokecolor{currentstroke}%
\pgfsetdash{}{0pt}%
\pgfpathmoveto{\pgfqpoint{3.355152in}{5.022157in}}%
\pgfpathlineto{\pgfqpoint{3.396818in}{5.105490in}}%
\pgfpathlineto{\pgfqpoint{3.313485in}{5.105490in}}%
\pgfpathclose%
\pgfusepath{stroke,fill}%
\end{pgfscope}%
\begin{pgfscope}%
\pgfpathrectangle{\pgfqpoint{0.800000in}{0.528000in}}{\pgfqpoint{4.960000in}{3.696000in}} %
\pgfusepath{clip}%
\pgfsetbuttcap%
\pgfsetroundjoin%
\definecolor{currentfill}{rgb}{1.000000,0.647059,0.000000}%
\pgfsetfillcolor{currentfill}%
\pgfsetlinewidth{1.003750pt}%
\definecolor{currentstroke}{rgb}{1.000000,0.647059,0.000000}%
\pgfsetstrokecolor{currentstroke}%
\pgfsetdash{}{0pt}%
\pgfpathmoveto{\pgfqpoint{3.405253in}{3.991922in}}%
\pgfpathlineto{\pgfqpoint{3.446919in}{4.075255in}}%
\pgfpathlineto{\pgfqpoint{3.363586in}{4.075255in}}%
\pgfpathclose%
\pgfusepath{stroke,fill}%
\end{pgfscope}%
\begin{pgfscope}%
\pgfpathrectangle{\pgfqpoint{0.800000in}{0.528000in}}{\pgfqpoint{4.960000in}{3.696000in}} %
\pgfusepath{clip}%
\pgfsetbuttcap%
\pgfsetroundjoin%
\definecolor{currentfill}{rgb}{1.000000,0.647059,0.000000}%
\pgfsetfillcolor{currentfill}%
\pgfsetlinewidth{1.003750pt}%
\definecolor{currentstroke}{rgb}{1.000000,0.647059,0.000000}%
\pgfsetstrokecolor{currentstroke}%
\pgfsetdash{}{0pt}%
\pgfpathmoveto{\pgfqpoint{3.455354in}{6.839616in}}%
\pgfpathlineto{\pgfqpoint{3.497020in}{6.922949in}}%
\pgfpathlineto{\pgfqpoint{3.413687in}{6.922949in}}%
\pgfpathclose%
\pgfusepath{stroke,fill}%
\end{pgfscope}%
\begin{pgfscope}%
\pgfpathrectangle{\pgfqpoint{0.800000in}{0.528000in}}{\pgfqpoint{4.960000in}{3.696000in}} %
\pgfusepath{clip}%
\pgfsetbuttcap%
\pgfsetroundjoin%
\definecolor{currentfill}{rgb}{1.000000,0.647059,0.000000}%
\pgfsetfillcolor{currentfill}%
\pgfsetlinewidth{1.003750pt}%
\definecolor{currentstroke}{rgb}{1.000000,0.647059,0.000000}%
\pgfsetstrokecolor{currentstroke}%
\pgfsetdash{}{0pt}%
\pgfpathmoveto{\pgfqpoint{3.505455in}{5.133087in}}%
\pgfpathlineto{\pgfqpoint{3.547121in}{5.216421in}}%
\pgfpathlineto{\pgfqpoint{3.463788in}{5.216421in}}%
\pgfpathclose%
\pgfusepath{stroke,fill}%
\end{pgfscope}%
\begin{pgfscope}%
\pgfpathrectangle{\pgfqpoint{0.800000in}{0.528000in}}{\pgfqpoint{4.960000in}{3.696000in}} %
\pgfusepath{clip}%
\pgfsetbuttcap%
\pgfsetroundjoin%
\definecolor{currentfill}{rgb}{1.000000,0.647059,0.000000}%
\pgfsetfillcolor{currentfill}%
\pgfsetlinewidth{1.003750pt}%
\definecolor{currentstroke}{rgb}{1.000000,0.647059,0.000000}%
\pgfsetstrokecolor{currentstroke}%
\pgfsetdash{}{0pt}%
\pgfpathmoveto{\pgfqpoint{3.555556in}{7.312538in}}%
\pgfpathlineto{\pgfqpoint{3.597222in}{7.395872in}}%
\pgfpathlineto{\pgfqpoint{3.513889in}{7.395872in}}%
\pgfpathclose%
\pgfusepath{stroke,fill}%
\end{pgfscope}%
\begin{pgfscope}%
\pgfpathrectangle{\pgfqpoint{0.800000in}{0.528000in}}{\pgfqpoint{4.960000in}{3.696000in}} %
\pgfusepath{clip}%
\pgfsetbuttcap%
\pgfsetroundjoin%
\definecolor{currentfill}{rgb}{1.000000,0.647059,0.000000}%
\pgfsetfillcolor{currentfill}%
\pgfsetlinewidth{1.003750pt}%
\definecolor{currentstroke}{rgb}{1.000000,0.647059,0.000000}%
\pgfsetstrokecolor{currentstroke}%
\pgfsetdash{}{0pt}%
\pgfpathmoveto{\pgfqpoint{3.605657in}{7.035221in}}%
\pgfpathlineto{\pgfqpoint{3.647323in}{7.118555in}}%
\pgfpathlineto{\pgfqpoint{3.563990in}{7.118555in}}%
\pgfpathclose%
\pgfusepath{stroke,fill}%
\end{pgfscope}%
\begin{pgfscope}%
\pgfpathrectangle{\pgfqpoint{0.800000in}{0.528000in}}{\pgfqpoint{4.960000in}{3.696000in}} %
\pgfusepath{clip}%
\pgfsetbuttcap%
\pgfsetroundjoin%
\definecolor{currentfill}{rgb}{1.000000,0.647059,0.000000}%
\pgfsetfillcolor{currentfill}%
\pgfsetlinewidth{1.003750pt}%
\definecolor{currentstroke}{rgb}{1.000000,0.647059,0.000000}%
\pgfsetstrokecolor{currentstroke}%
\pgfsetdash{}{0pt}%
\pgfpathmoveto{\pgfqpoint{3.655758in}{2.947918in}}%
\pgfpathlineto{\pgfqpoint{3.697424in}{3.031251in}}%
\pgfpathlineto{\pgfqpoint{3.614091in}{3.031251in}}%
\pgfpathclose%
\pgfusepath{stroke,fill}%
\end{pgfscope}%
\begin{pgfscope}%
\pgfpathrectangle{\pgfqpoint{0.800000in}{0.528000in}}{\pgfqpoint{4.960000in}{3.696000in}} %
\pgfusepath{clip}%
\pgfsetbuttcap%
\pgfsetroundjoin%
\definecolor{currentfill}{rgb}{1.000000,0.647059,0.000000}%
\pgfsetfillcolor{currentfill}%
\pgfsetlinewidth{1.003750pt}%
\definecolor{currentstroke}{rgb}{1.000000,0.647059,0.000000}%
\pgfsetstrokecolor{currentstroke}%
\pgfsetdash{}{0pt}%
\pgfpathmoveto{\pgfqpoint{3.705859in}{6.404804in}}%
\pgfpathlineto{\pgfqpoint{3.747525in}{6.488137in}}%
\pgfpathlineto{\pgfqpoint{3.664192in}{6.488137in}}%
\pgfpathclose%
\pgfusepath{stroke,fill}%
\end{pgfscope}%
\begin{pgfscope}%
\pgfpathrectangle{\pgfqpoint{0.800000in}{0.528000in}}{\pgfqpoint{4.960000in}{3.696000in}} %
\pgfusepath{clip}%
\pgfsetbuttcap%
\pgfsetroundjoin%
\definecolor{currentfill}{rgb}{1.000000,0.647059,0.000000}%
\pgfsetfillcolor{currentfill}%
\pgfsetlinewidth{1.003750pt}%
\definecolor{currentstroke}{rgb}{1.000000,0.647059,0.000000}%
\pgfsetstrokecolor{currentstroke}%
\pgfsetdash{}{0pt}%
\pgfpathmoveto{\pgfqpoint{3.755960in}{3.113258in}}%
\pgfpathlineto{\pgfqpoint{3.797626in}{3.196591in}}%
\pgfpathlineto{\pgfqpoint{3.714293in}{3.196591in}}%
\pgfpathclose%
\pgfusepath{stroke,fill}%
\end{pgfscope}%
\begin{pgfscope}%
\pgfpathrectangle{\pgfqpoint{0.800000in}{0.528000in}}{\pgfqpoint{4.960000in}{3.696000in}} %
\pgfusepath{clip}%
\pgfsetbuttcap%
\pgfsetroundjoin%
\definecolor{currentfill}{rgb}{1.000000,0.647059,0.000000}%
\pgfsetfillcolor{currentfill}%
\pgfsetlinewidth{1.003750pt}%
\definecolor{currentstroke}{rgb}{1.000000,0.647059,0.000000}%
\pgfsetstrokecolor{currentstroke}%
\pgfsetdash{}{0pt}%
\pgfpathmoveto{\pgfqpoint{3.806061in}{6.861077in}}%
\pgfpathlineto{\pgfqpoint{3.847727in}{6.944410in}}%
\pgfpathlineto{\pgfqpoint{3.764394in}{6.944410in}}%
\pgfpathclose%
\pgfusepath{stroke,fill}%
\end{pgfscope}%
\begin{pgfscope}%
\pgfpathrectangle{\pgfqpoint{0.800000in}{0.528000in}}{\pgfqpoint{4.960000in}{3.696000in}} %
\pgfusepath{clip}%
\pgfsetbuttcap%
\pgfsetroundjoin%
\definecolor{currentfill}{rgb}{1.000000,0.647059,0.000000}%
\pgfsetfillcolor{currentfill}%
\pgfsetlinewidth{1.003750pt}%
\definecolor{currentstroke}{rgb}{1.000000,0.647059,0.000000}%
\pgfsetstrokecolor{currentstroke}%
\pgfsetdash{}{0pt}%
\pgfpathmoveto{\pgfqpoint{3.856162in}{4.542422in}}%
\pgfpathlineto{\pgfqpoint{3.897828in}{4.625755in}}%
\pgfpathlineto{\pgfqpoint{3.814495in}{4.625755in}}%
\pgfpathclose%
\pgfusepath{stroke,fill}%
\end{pgfscope}%
\begin{pgfscope}%
\pgfpathrectangle{\pgfqpoint{0.800000in}{0.528000in}}{\pgfqpoint{4.960000in}{3.696000in}} %
\pgfusepath{clip}%
\pgfsetbuttcap%
\pgfsetroundjoin%
\definecolor{currentfill}{rgb}{1.000000,0.647059,0.000000}%
\pgfsetfillcolor{currentfill}%
\pgfsetlinewidth{1.003750pt}%
\definecolor{currentstroke}{rgb}{1.000000,0.647059,0.000000}%
\pgfsetstrokecolor{currentstroke}%
\pgfsetdash{}{0pt}%
\pgfpathmoveto{\pgfqpoint{3.906263in}{4.966856in}}%
\pgfpathlineto{\pgfqpoint{3.947929in}{5.050189in}}%
\pgfpathlineto{\pgfqpoint{3.864596in}{5.050189in}}%
\pgfpathclose%
\pgfusepath{stroke,fill}%
\end{pgfscope}%
\begin{pgfscope}%
\pgfpathrectangle{\pgfqpoint{0.800000in}{0.528000in}}{\pgfqpoint{4.960000in}{3.696000in}} %
\pgfusepath{clip}%
\pgfsetbuttcap%
\pgfsetroundjoin%
\definecolor{currentfill}{rgb}{1.000000,0.647059,0.000000}%
\pgfsetfillcolor{currentfill}%
\pgfsetlinewidth{1.003750pt}%
\definecolor{currentstroke}{rgb}{1.000000,0.647059,0.000000}%
\pgfsetstrokecolor{currentstroke}%
\pgfsetdash{}{0pt}%
\pgfpathmoveto{\pgfqpoint{3.956364in}{4.498762in}}%
\pgfpathlineto{\pgfqpoint{3.998030in}{4.582095in}}%
\pgfpathlineto{\pgfqpoint{3.914697in}{4.582095in}}%
\pgfpathclose%
\pgfusepath{stroke,fill}%
\end{pgfscope}%
\begin{pgfscope}%
\pgfpathrectangle{\pgfqpoint{0.800000in}{0.528000in}}{\pgfqpoint{4.960000in}{3.696000in}} %
\pgfusepath{clip}%
\pgfsetbuttcap%
\pgfsetroundjoin%
\definecolor{currentfill}{rgb}{1.000000,0.647059,0.000000}%
\pgfsetfillcolor{currentfill}%
\pgfsetlinewidth{1.003750pt}%
\definecolor{currentstroke}{rgb}{1.000000,0.647059,0.000000}%
\pgfsetstrokecolor{currentstroke}%
\pgfsetdash{}{0pt}%
\pgfpathmoveto{\pgfqpoint{4.006465in}{7.217320in}}%
\pgfpathlineto{\pgfqpoint{4.048131in}{7.300653in}}%
\pgfpathlineto{\pgfqpoint{3.964798in}{7.300653in}}%
\pgfpathclose%
\pgfusepath{stroke,fill}%
\end{pgfscope}%
\begin{pgfscope}%
\pgfpathrectangle{\pgfqpoint{0.800000in}{0.528000in}}{\pgfqpoint{4.960000in}{3.696000in}} %
\pgfusepath{clip}%
\pgfsetbuttcap%
\pgfsetroundjoin%
\definecolor{currentfill}{rgb}{1.000000,0.647059,0.000000}%
\pgfsetfillcolor{currentfill}%
\pgfsetlinewidth{1.003750pt}%
\definecolor{currentstroke}{rgb}{1.000000,0.647059,0.000000}%
\pgfsetstrokecolor{currentstroke}%
\pgfsetdash{}{0pt}%
\pgfpathmoveto{\pgfqpoint{4.056566in}{6.145431in}}%
\pgfpathlineto{\pgfqpoint{4.098232in}{6.228764in}}%
\pgfpathlineto{\pgfqpoint{4.014899in}{6.228764in}}%
\pgfpathclose%
\pgfusepath{stroke,fill}%
\end{pgfscope}%
\begin{pgfscope}%
\pgfpathrectangle{\pgfqpoint{0.800000in}{0.528000in}}{\pgfqpoint{4.960000in}{3.696000in}} %
\pgfusepath{clip}%
\pgfsetbuttcap%
\pgfsetroundjoin%
\definecolor{currentfill}{rgb}{1.000000,0.647059,0.000000}%
\pgfsetfillcolor{currentfill}%
\pgfsetlinewidth{1.003750pt}%
\definecolor{currentstroke}{rgb}{1.000000,0.647059,0.000000}%
\pgfsetstrokecolor{currentstroke}%
\pgfsetdash{}{0pt}%
\pgfpathmoveto{\pgfqpoint{4.106667in}{3.154392in}}%
\pgfpathlineto{\pgfqpoint{4.148333in}{3.237725in}}%
\pgfpathlineto{\pgfqpoint{4.065000in}{3.237725in}}%
\pgfpathclose%
\pgfusepath{stroke,fill}%
\end{pgfscope}%
\begin{pgfscope}%
\pgfpathrectangle{\pgfqpoint{0.800000in}{0.528000in}}{\pgfqpoint{4.960000in}{3.696000in}} %
\pgfusepath{clip}%
\pgfsetbuttcap%
\pgfsetroundjoin%
\definecolor{currentfill}{rgb}{1.000000,0.647059,0.000000}%
\pgfsetfillcolor{currentfill}%
\pgfsetlinewidth{1.003750pt}%
\definecolor{currentstroke}{rgb}{1.000000,0.647059,0.000000}%
\pgfsetstrokecolor{currentstroke}%
\pgfsetdash{}{0pt}%
\pgfpathmoveto{\pgfqpoint{4.156768in}{5.877524in}}%
\pgfpathlineto{\pgfqpoint{4.198434in}{5.960857in}}%
\pgfpathlineto{\pgfqpoint{4.115101in}{5.960857in}}%
\pgfpathclose%
\pgfusepath{stroke,fill}%
\end{pgfscope}%
\begin{pgfscope}%
\pgfpathrectangle{\pgfqpoint{0.800000in}{0.528000in}}{\pgfqpoint{4.960000in}{3.696000in}} %
\pgfusepath{clip}%
\pgfsetbuttcap%
\pgfsetroundjoin%
\definecolor{currentfill}{rgb}{1.000000,0.647059,0.000000}%
\pgfsetfillcolor{currentfill}%
\pgfsetlinewidth{1.003750pt}%
\definecolor{currentstroke}{rgb}{1.000000,0.647059,0.000000}%
\pgfsetstrokecolor{currentstroke}%
\pgfsetdash{}{0pt}%
\pgfpathmoveto{\pgfqpoint{4.206869in}{4.309612in}}%
\pgfpathlineto{\pgfqpoint{4.248535in}{4.392945in}}%
\pgfpathlineto{\pgfqpoint{4.165202in}{4.392945in}}%
\pgfpathclose%
\pgfusepath{stroke,fill}%
\end{pgfscope}%
\begin{pgfscope}%
\pgfpathrectangle{\pgfqpoint{0.800000in}{0.528000in}}{\pgfqpoint{4.960000in}{3.696000in}} %
\pgfusepath{clip}%
\pgfsetbuttcap%
\pgfsetroundjoin%
\definecolor{currentfill}{rgb}{1.000000,0.647059,0.000000}%
\pgfsetfillcolor{currentfill}%
\pgfsetlinewidth{1.003750pt}%
\definecolor{currentstroke}{rgb}{1.000000,0.647059,0.000000}%
\pgfsetstrokecolor{currentstroke}%
\pgfsetdash{}{0pt}%
\pgfpathmoveto{\pgfqpoint{4.256970in}{7.652998in}}%
\pgfpathlineto{\pgfqpoint{4.298636in}{7.736332in}}%
\pgfpathlineto{\pgfqpoint{4.215303in}{7.736332in}}%
\pgfpathclose%
\pgfusepath{stroke,fill}%
\end{pgfscope}%
\begin{pgfscope}%
\pgfpathrectangle{\pgfqpoint{0.800000in}{0.528000in}}{\pgfqpoint{4.960000in}{3.696000in}} %
\pgfusepath{clip}%
\pgfsetbuttcap%
\pgfsetroundjoin%
\definecolor{currentfill}{rgb}{1.000000,0.647059,0.000000}%
\pgfsetfillcolor{currentfill}%
\pgfsetlinewidth{1.003750pt}%
\definecolor{currentstroke}{rgb}{1.000000,0.647059,0.000000}%
\pgfsetstrokecolor{currentstroke}%
\pgfsetdash{}{0pt}%
\pgfpathmoveto{\pgfqpoint{4.307071in}{6.789426in}}%
\pgfpathlineto{\pgfqpoint{4.348737in}{6.872759in}}%
\pgfpathlineto{\pgfqpoint{4.265404in}{6.872759in}}%
\pgfpathclose%
\pgfusepath{stroke,fill}%
\end{pgfscope}%
\begin{pgfscope}%
\pgfpathrectangle{\pgfqpoint{0.800000in}{0.528000in}}{\pgfqpoint{4.960000in}{3.696000in}} %
\pgfusepath{clip}%
\pgfsetbuttcap%
\pgfsetroundjoin%
\definecolor{currentfill}{rgb}{1.000000,0.647059,0.000000}%
\pgfsetfillcolor{currentfill}%
\pgfsetlinewidth{1.003750pt}%
\definecolor{currentstroke}{rgb}{1.000000,0.647059,0.000000}%
\pgfsetstrokecolor{currentstroke}%
\pgfsetdash{}{0pt}%
\pgfpathmoveto{\pgfqpoint{4.357172in}{4.784806in}}%
\pgfpathlineto{\pgfqpoint{4.398838in}{4.868140in}}%
\pgfpathlineto{\pgfqpoint{4.315505in}{4.868140in}}%
\pgfpathclose%
\pgfusepath{stroke,fill}%
\end{pgfscope}%
\begin{pgfscope}%
\pgfpathrectangle{\pgfqpoint{0.800000in}{0.528000in}}{\pgfqpoint{4.960000in}{3.696000in}} %
\pgfusepath{clip}%
\pgfsetbuttcap%
\pgfsetroundjoin%
\definecolor{currentfill}{rgb}{1.000000,0.647059,0.000000}%
\pgfsetfillcolor{currentfill}%
\pgfsetlinewidth{1.003750pt}%
\definecolor{currentstroke}{rgb}{1.000000,0.647059,0.000000}%
\pgfsetstrokecolor{currentstroke}%
\pgfsetdash{}{0pt}%
\pgfpathmoveto{\pgfqpoint{4.407273in}{7.819725in}}%
\pgfpathlineto{\pgfqpoint{4.448939in}{7.903058in}}%
\pgfpathlineto{\pgfqpoint{4.365606in}{7.903058in}}%
\pgfpathclose%
\pgfusepath{stroke,fill}%
\end{pgfscope}%
\begin{pgfscope}%
\pgfpathrectangle{\pgfqpoint{0.800000in}{0.528000in}}{\pgfqpoint{4.960000in}{3.696000in}} %
\pgfusepath{clip}%
\pgfsetbuttcap%
\pgfsetroundjoin%
\definecolor{currentfill}{rgb}{1.000000,0.647059,0.000000}%
\pgfsetfillcolor{currentfill}%
\pgfsetlinewidth{1.003750pt}%
\definecolor{currentstroke}{rgb}{1.000000,0.647059,0.000000}%
\pgfsetstrokecolor{currentstroke}%
\pgfsetdash{}{0pt}%
\pgfpathmoveto{\pgfqpoint{4.457374in}{7.100948in}}%
\pgfpathlineto{\pgfqpoint{4.499040in}{7.184281in}}%
\pgfpathlineto{\pgfqpoint{4.415707in}{7.184281in}}%
\pgfpathclose%
\pgfusepath{stroke,fill}%
\end{pgfscope}%
\begin{pgfscope}%
\pgfpathrectangle{\pgfqpoint{0.800000in}{0.528000in}}{\pgfqpoint{4.960000in}{3.696000in}} %
\pgfusepath{clip}%
\pgfsetbuttcap%
\pgfsetroundjoin%
\definecolor{currentfill}{rgb}{1.000000,0.647059,0.000000}%
\pgfsetfillcolor{currentfill}%
\pgfsetlinewidth{1.003750pt}%
\definecolor{currentstroke}{rgb}{1.000000,0.647059,0.000000}%
\pgfsetstrokecolor{currentstroke}%
\pgfsetdash{}{0pt}%
\pgfpathmoveto{\pgfqpoint{4.507475in}{4.142787in}}%
\pgfpathlineto{\pgfqpoint{4.549141in}{4.226120in}}%
\pgfpathlineto{\pgfqpoint{4.465808in}{4.226120in}}%
\pgfpathclose%
\pgfusepath{stroke,fill}%
\end{pgfscope}%
\begin{pgfscope}%
\pgfpathrectangle{\pgfqpoint{0.800000in}{0.528000in}}{\pgfqpoint{4.960000in}{3.696000in}} %
\pgfusepath{clip}%
\pgfsetbuttcap%
\pgfsetroundjoin%
\definecolor{currentfill}{rgb}{1.000000,0.647059,0.000000}%
\pgfsetfillcolor{currentfill}%
\pgfsetlinewidth{1.003750pt}%
\definecolor{currentstroke}{rgb}{1.000000,0.647059,0.000000}%
\pgfsetstrokecolor{currentstroke}%
\pgfsetdash{}{0pt}%
\pgfpathmoveto{\pgfqpoint{4.557576in}{4.983101in}}%
\pgfpathlineto{\pgfqpoint{4.599242in}{5.066435in}}%
\pgfpathlineto{\pgfqpoint{4.515909in}{5.066435in}}%
\pgfpathclose%
\pgfusepath{stroke,fill}%
\end{pgfscope}%
\begin{pgfscope}%
\pgfpathrectangle{\pgfqpoint{0.800000in}{0.528000in}}{\pgfqpoint{4.960000in}{3.696000in}} %
\pgfusepath{clip}%
\pgfsetbuttcap%
\pgfsetroundjoin%
\definecolor{currentfill}{rgb}{1.000000,0.647059,0.000000}%
\pgfsetfillcolor{currentfill}%
\pgfsetlinewidth{1.003750pt}%
\definecolor{currentstroke}{rgb}{1.000000,0.647059,0.000000}%
\pgfsetstrokecolor{currentstroke}%
\pgfsetdash{}{0pt}%
\pgfpathmoveto{\pgfqpoint{4.607677in}{5.365968in}}%
\pgfpathlineto{\pgfqpoint{4.649343in}{5.449301in}}%
\pgfpathlineto{\pgfqpoint{4.566010in}{5.449301in}}%
\pgfpathclose%
\pgfusepath{stroke,fill}%
\end{pgfscope}%
\begin{pgfscope}%
\pgfpathrectangle{\pgfqpoint{0.800000in}{0.528000in}}{\pgfqpoint{4.960000in}{3.696000in}} %
\pgfusepath{clip}%
\pgfsetbuttcap%
\pgfsetroundjoin%
\definecolor{currentfill}{rgb}{1.000000,0.647059,0.000000}%
\pgfsetfillcolor{currentfill}%
\pgfsetlinewidth{1.003750pt}%
\definecolor{currentstroke}{rgb}{1.000000,0.647059,0.000000}%
\pgfsetstrokecolor{currentstroke}%
\pgfsetdash{}{0pt}%
\pgfpathmoveto{\pgfqpoint{4.657778in}{7.302378in}}%
\pgfpathlineto{\pgfqpoint{4.699444in}{7.385711in}}%
\pgfpathlineto{\pgfqpoint{4.616111in}{7.385711in}}%
\pgfpathclose%
\pgfusepath{stroke,fill}%
\end{pgfscope}%
\begin{pgfscope}%
\pgfpathrectangle{\pgfqpoint{0.800000in}{0.528000in}}{\pgfqpoint{4.960000in}{3.696000in}} %
\pgfusepath{clip}%
\pgfsetbuttcap%
\pgfsetroundjoin%
\definecolor{currentfill}{rgb}{1.000000,0.647059,0.000000}%
\pgfsetfillcolor{currentfill}%
\pgfsetlinewidth{1.003750pt}%
\definecolor{currentstroke}{rgb}{1.000000,0.647059,0.000000}%
\pgfsetstrokecolor{currentstroke}%
\pgfsetdash{}{0pt}%
\pgfpathmoveto{\pgfqpoint{4.707879in}{6.913270in}}%
\pgfpathlineto{\pgfqpoint{4.749545in}{6.996603in}}%
\pgfpathlineto{\pgfqpoint{4.666212in}{6.996603in}}%
\pgfpathclose%
\pgfusepath{stroke,fill}%
\end{pgfscope}%
\begin{pgfscope}%
\pgfpathrectangle{\pgfqpoint{0.800000in}{0.528000in}}{\pgfqpoint{4.960000in}{3.696000in}} %
\pgfusepath{clip}%
\pgfsetbuttcap%
\pgfsetroundjoin%
\definecolor{currentfill}{rgb}{1.000000,0.647059,0.000000}%
\pgfsetfillcolor{currentfill}%
\pgfsetlinewidth{1.003750pt}%
\definecolor{currentstroke}{rgb}{1.000000,0.647059,0.000000}%
\pgfsetstrokecolor{currentstroke}%
\pgfsetdash{}{0pt}%
\pgfpathmoveto{\pgfqpoint{4.757980in}{3.965022in}}%
\pgfpathlineto{\pgfqpoint{4.799646in}{4.048355in}}%
\pgfpathlineto{\pgfqpoint{4.716313in}{4.048355in}}%
\pgfpathclose%
\pgfusepath{stroke,fill}%
\end{pgfscope}%
\begin{pgfscope}%
\pgfpathrectangle{\pgfqpoint{0.800000in}{0.528000in}}{\pgfqpoint{4.960000in}{3.696000in}} %
\pgfusepath{clip}%
\pgfsetbuttcap%
\pgfsetroundjoin%
\definecolor{currentfill}{rgb}{1.000000,0.647059,0.000000}%
\pgfsetfillcolor{currentfill}%
\pgfsetlinewidth{1.003750pt}%
\definecolor{currentstroke}{rgb}{1.000000,0.647059,0.000000}%
\pgfsetstrokecolor{currentstroke}%
\pgfsetdash{}{0pt}%
\pgfpathmoveto{\pgfqpoint{4.808081in}{4.176958in}}%
\pgfpathlineto{\pgfqpoint{4.849747in}{4.260292in}}%
\pgfpathlineto{\pgfqpoint{4.766414in}{4.260292in}}%
\pgfpathclose%
\pgfusepath{stroke,fill}%
\end{pgfscope}%
\begin{pgfscope}%
\pgfpathrectangle{\pgfqpoint{0.800000in}{0.528000in}}{\pgfqpoint{4.960000in}{3.696000in}} %
\pgfusepath{clip}%
\pgfsetbuttcap%
\pgfsetroundjoin%
\definecolor{currentfill}{rgb}{1.000000,0.647059,0.000000}%
\pgfsetfillcolor{currentfill}%
\pgfsetlinewidth{1.003750pt}%
\definecolor{currentstroke}{rgb}{1.000000,0.647059,0.000000}%
\pgfsetstrokecolor{currentstroke}%
\pgfsetdash{}{0pt}%
\pgfpathmoveto{\pgfqpoint{4.858182in}{5.430264in}}%
\pgfpathlineto{\pgfqpoint{4.899848in}{5.513598in}}%
\pgfpathlineto{\pgfqpoint{4.816515in}{5.513598in}}%
\pgfpathclose%
\pgfusepath{stroke,fill}%
\end{pgfscope}%
\begin{pgfscope}%
\pgfpathrectangle{\pgfqpoint{0.800000in}{0.528000in}}{\pgfqpoint{4.960000in}{3.696000in}} %
\pgfusepath{clip}%
\pgfsetbuttcap%
\pgfsetroundjoin%
\definecolor{currentfill}{rgb}{1.000000,0.647059,0.000000}%
\pgfsetfillcolor{currentfill}%
\pgfsetlinewidth{1.003750pt}%
\definecolor{currentstroke}{rgb}{1.000000,0.647059,0.000000}%
\pgfsetstrokecolor{currentstroke}%
\pgfsetdash{}{0pt}%
\pgfpathmoveto{\pgfqpoint{4.908283in}{8.398532in}}%
\pgfpathlineto{\pgfqpoint{4.949949in}{8.481865in}}%
\pgfpathlineto{\pgfqpoint{4.866616in}{8.481865in}}%
\pgfpathclose%
\pgfusepath{stroke,fill}%
\end{pgfscope}%
\begin{pgfscope}%
\pgfpathrectangle{\pgfqpoint{0.800000in}{0.528000in}}{\pgfqpoint{4.960000in}{3.696000in}} %
\pgfusepath{clip}%
\pgfsetbuttcap%
\pgfsetroundjoin%
\definecolor{currentfill}{rgb}{1.000000,0.647059,0.000000}%
\pgfsetfillcolor{currentfill}%
\pgfsetlinewidth{1.003750pt}%
\definecolor{currentstroke}{rgb}{1.000000,0.647059,0.000000}%
\pgfsetstrokecolor{currentstroke}%
\pgfsetdash{}{0pt}%
\pgfpathmoveto{\pgfqpoint{4.958384in}{5.011813in}}%
\pgfpathlineto{\pgfqpoint{5.000051in}{5.095146in}}%
\pgfpathlineto{\pgfqpoint{4.916717in}{5.095146in}}%
\pgfpathclose%
\pgfusepath{stroke,fill}%
\end{pgfscope}%
\begin{pgfscope}%
\pgfpathrectangle{\pgfqpoint{0.800000in}{0.528000in}}{\pgfqpoint{4.960000in}{3.696000in}} %
\pgfusepath{clip}%
\pgfsetbuttcap%
\pgfsetroundjoin%
\definecolor{currentfill}{rgb}{1.000000,0.647059,0.000000}%
\pgfsetfillcolor{currentfill}%
\pgfsetlinewidth{1.003750pt}%
\definecolor{currentstroke}{rgb}{1.000000,0.647059,0.000000}%
\pgfsetstrokecolor{currentstroke}%
\pgfsetdash{}{0pt}%
\pgfpathmoveto{\pgfqpoint{5.008485in}{6.003844in}}%
\pgfpathlineto{\pgfqpoint{5.050152in}{6.087177in}}%
\pgfpathlineto{\pgfqpoint{4.966818in}{6.087177in}}%
\pgfpathclose%
\pgfusepath{stroke,fill}%
\end{pgfscope}%
\begin{pgfscope}%
\pgfpathrectangle{\pgfqpoint{0.800000in}{0.528000in}}{\pgfqpoint{4.960000in}{3.696000in}} %
\pgfusepath{clip}%
\pgfsetbuttcap%
\pgfsetroundjoin%
\definecolor{currentfill}{rgb}{1.000000,0.647059,0.000000}%
\pgfsetfillcolor{currentfill}%
\pgfsetlinewidth{1.003750pt}%
\definecolor{currentstroke}{rgb}{1.000000,0.647059,0.000000}%
\pgfsetstrokecolor{currentstroke}%
\pgfsetdash{}{0pt}%
\pgfpathmoveto{\pgfqpoint{5.058586in}{7.089155in}}%
\pgfpathlineto{\pgfqpoint{5.100253in}{7.172489in}}%
\pgfpathlineto{\pgfqpoint{5.016919in}{7.172489in}}%
\pgfpathclose%
\pgfusepath{stroke,fill}%
\end{pgfscope}%
\begin{pgfscope}%
\pgfpathrectangle{\pgfqpoint{0.800000in}{0.528000in}}{\pgfqpoint{4.960000in}{3.696000in}} %
\pgfusepath{clip}%
\pgfsetbuttcap%
\pgfsetroundjoin%
\definecolor{currentfill}{rgb}{1.000000,0.647059,0.000000}%
\pgfsetfillcolor{currentfill}%
\pgfsetlinewidth{1.003750pt}%
\definecolor{currentstroke}{rgb}{1.000000,0.647059,0.000000}%
\pgfsetstrokecolor{currentstroke}%
\pgfsetdash{}{0pt}%
\pgfpathmoveto{\pgfqpoint{5.108687in}{5.526080in}}%
\pgfpathlineto{\pgfqpoint{5.150354in}{5.609414in}}%
\pgfpathlineto{\pgfqpoint{5.067020in}{5.609414in}}%
\pgfpathclose%
\pgfusepath{stroke,fill}%
\end{pgfscope}%
\begin{pgfscope}%
\pgfpathrectangle{\pgfqpoint{0.800000in}{0.528000in}}{\pgfqpoint{4.960000in}{3.696000in}} %
\pgfusepath{clip}%
\pgfsetbuttcap%
\pgfsetroundjoin%
\definecolor{currentfill}{rgb}{1.000000,0.647059,0.000000}%
\pgfsetfillcolor{currentfill}%
\pgfsetlinewidth{1.003750pt}%
\definecolor{currentstroke}{rgb}{1.000000,0.647059,0.000000}%
\pgfsetstrokecolor{currentstroke}%
\pgfsetdash{}{0pt}%
\pgfpathmoveto{\pgfqpoint{5.158788in}{7.730654in}}%
\pgfpathlineto{\pgfqpoint{5.200455in}{7.813988in}}%
\pgfpathlineto{\pgfqpoint{5.117121in}{7.813988in}}%
\pgfpathclose%
\pgfusepath{stroke,fill}%
\end{pgfscope}%
\begin{pgfscope}%
\pgfpathrectangle{\pgfqpoint{0.800000in}{0.528000in}}{\pgfqpoint{4.960000in}{3.696000in}} %
\pgfusepath{clip}%
\pgfsetbuttcap%
\pgfsetroundjoin%
\definecolor{currentfill}{rgb}{1.000000,0.647059,0.000000}%
\pgfsetfillcolor{currentfill}%
\pgfsetlinewidth{1.003750pt}%
\definecolor{currentstroke}{rgb}{1.000000,0.647059,0.000000}%
\pgfsetstrokecolor{currentstroke}%
\pgfsetdash{}{0pt}%
\pgfpathmoveto{\pgfqpoint{5.208889in}{5.441664in}}%
\pgfpathlineto{\pgfqpoint{5.250556in}{5.524997in}}%
\pgfpathlineto{\pgfqpoint{5.167222in}{5.524997in}}%
\pgfpathclose%
\pgfusepath{stroke,fill}%
\end{pgfscope}%
\begin{pgfscope}%
\pgfpathrectangle{\pgfqpoint{0.800000in}{0.528000in}}{\pgfqpoint{4.960000in}{3.696000in}} %
\pgfusepath{clip}%
\pgfsetbuttcap%
\pgfsetroundjoin%
\definecolor{currentfill}{rgb}{1.000000,0.647059,0.000000}%
\pgfsetfillcolor{currentfill}%
\pgfsetlinewidth{1.003750pt}%
\definecolor{currentstroke}{rgb}{1.000000,0.647059,0.000000}%
\pgfsetstrokecolor{currentstroke}%
\pgfsetdash{}{0pt}%
\pgfpathmoveto{\pgfqpoint{5.258990in}{7.641787in}}%
\pgfpathlineto{\pgfqpoint{5.300657in}{7.725120in}}%
\pgfpathlineto{\pgfqpoint{5.217323in}{7.725120in}}%
\pgfpathclose%
\pgfusepath{stroke,fill}%
\end{pgfscope}%
\begin{pgfscope}%
\pgfpathrectangle{\pgfqpoint{0.800000in}{0.528000in}}{\pgfqpoint{4.960000in}{3.696000in}} %
\pgfusepath{clip}%
\pgfsetbuttcap%
\pgfsetroundjoin%
\definecolor{currentfill}{rgb}{1.000000,0.647059,0.000000}%
\pgfsetfillcolor{currentfill}%
\pgfsetlinewidth{1.003750pt}%
\definecolor{currentstroke}{rgb}{1.000000,0.647059,0.000000}%
\pgfsetstrokecolor{currentstroke}%
\pgfsetdash{}{0pt}%
\pgfpathmoveto{\pgfqpoint{5.309091in}{4.198569in}}%
\pgfpathlineto{\pgfqpoint{5.350758in}{4.281902in}}%
\pgfpathlineto{\pgfqpoint{5.267424in}{4.281902in}}%
\pgfpathclose%
\pgfusepath{stroke,fill}%
\end{pgfscope}%
\begin{pgfscope}%
\pgfpathrectangle{\pgfqpoint{0.800000in}{0.528000in}}{\pgfqpoint{4.960000in}{3.696000in}} %
\pgfusepath{clip}%
\pgfsetbuttcap%
\pgfsetroundjoin%
\definecolor{currentfill}{rgb}{1.000000,0.647059,0.000000}%
\pgfsetfillcolor{currentfill}%
\pgfsetlinewidth{1.003750pt}%
\definecolor{currentstroke}{rgb}{1.000000,0.647059,0.000000}%
\pgfsetstrokecolor{currentstroke}%
\pgfsetdash{}{0pt}%
\pgfpathmoveto{\pgfqpoint{5.359192in}{6.715156in}}%
\pgfpathlineto{\pgfqpoint{5.400859in}{6.798489in}}%
\pgfpathlineto{\pgfqpoint{5.317525in}{6.798489in}}%
\pgfpathclose%
\pgfusepath{stroke,fill}%
\end{pgfscope}%
\begin{pgfscope}%
\pgfpathrectangle{\pgfqpoint{0.800000in}{0.528000in}}{\pgfqpoint{4.960000in}{3.696000in}} %
\pgfusepath{clip}%
\pgfsetbuttcap%
\pgfsetroundjoin%
\definecolor{currentfill}{rgb}{1.000000,0.647059,0.000000}%
\pgfsetfillcolor{currentfill}%
\pgfsetlinewidth{1.003750pt}%
\definecolor{currentstroke}{rgb}{1.000000,0.647059,0.000000}%
\pgfsetstrokecolor{currentstroke}%
\pgfsetdash{}{0pt}%
\pgfpathmoveto{\pgfqpoint{5.409293in}{4.423664in}}%
\pgfpathlineto{\pgfqpoint{5.450960in}{4.506998in}}%
\pgfpathlineto{\pgfqpoint{5.367626in}{4.506998in}}%
\pgfpathclose%
\pgfusepath{stroke,fill}%
\end{pgfscope}%
\begin{pgfscope}%
\pgfpathrectangle{\pgfqpoint{0.800000in}{0.528000in}}{\pgfqpoint{4.960000in}{3.696000in}} %
\pgfusepath{clip}%
\pgfsetbuttcap%
\pgfsetroundjoin%
\definecolor{currentfill}{rgb}{1.000000,0.647059,0.000000}%
\pgfsetfillcolor{currentfill}%
\pgfsetlinewidth{1.003750pt}%
\definecolor{currentstroke}{rgb}{1.000000,0.647059,0.000000}%
\pgfsetstrokecolor{currentstroke}%
\pgfsetdash{}{0pt}%
\pgfpathmoveto{\pgfqpoint{5.459394in}{4.922161in}}%
\pgfpathlineto{\pgfqpoint{5.501061in}{5.005494in}}%
\pgfpathlineto{\pgfqpoint{5.417727in}{5.005494in}}%
\pgfpathclose%
\pgfusepath{stroke,fill}%
\end{pgfscope}%
\begin{pgfscope}%
\pgfpathrectangle{\pgfqpoint{0.800000in}{0.528000in}}{\pgfqpoint{4.960000in}{3.696000in}} %
\pgfusepath{clip}%
\pgfsetbuttcap%
\pgfsetroundjoin%
\definecolor{currentfill}{rgb}{1.000000,0.647059,0.000000}%
\pgfsetfillcolor{currentfill}%
\pgfsetlinewidth{1.003750pt}%
\definecolor{currentstroke}{rgb}{1.000000,0.647059,0.000000}%
\pgfsetstrokecolor{currentstroke}%
\pgfsetdash{}{0pt}%
\pgfpathmoveto{\pgfqpoint{5.509495in}{8.781688in}}%
\pgfpathlineto{\pgfqpoint{5.551162in}{8.865022in}}%
\pgfpathlineto{\pgfqpoint{5.467828in}{8.865022in}}%
\pgfpathclose%
\pgfusepath{stroke,fill}%
\end{pgfscope}%
\begin{pgfscope}%
\pgfpathrectangle{\pgfqpoint{0.800000in}{0.528000in}}{\pgfqpoint{4.960000in}{3.696000in}} %
\pgfusepath{clip}%
\pgfsetbuttcap%
\pgfsetroundjoin%
\definecolor{currentfill}{rgb}{1.000000,0.647059,0.000000}%
\pgfsetfillcolor{currentfill}%
\pgfsetlinewidth{1.003750pt}%
\definecolor{currentstroke}{rgb}{1.000000,0.647059,0.000000}%
\pgfsetstrokecolor{currentstroke}%
\pgfsetdash{}{0pt}%
\pgfpathmoveto{\pgfqpoint{5.559596in}{5.492128in}}%
\pgfpathlineto{\pgfqpoint{5.601263in}{5.575461in}}%
\pgfpathlineto{\pgfqpoint{5.517929in}{5.575461in}}%
\pgfpathclose%
\pgfusepath{stroke,fill}%
\end{pgfscope}%
\begin{pgfscope}%
\pgfpathrectangle{\pgfqpoint{0.800000in}{0.528000in}}{\pgfqpoint{4.960000in}{3.696000in}} %
\pgfusepath{clip}%
\pgfsetbuttcap%
\pgfsetroundjoin%
\definecolor{currentfill}{rgb}{1.000000,0.647059,0.000000}%
\pgfsetfillcolor{currentfill}%
\pgfsetlinewidth{1.003750pt}%
\definecolor{currentstroke}{rgb}{1.000000,0.647059,0.000000}%
\pgfsetstrokecolor{currentstroke}%
\pgfsetdash{}{0pt}%
\pgfpathmoveto{\pgfqpoint{5.609697in}{5.211182in}}%
\pgfpathlineto{\pgfqpoint{5.651364in}{5.294515in}}%
\pgfpathlineto{\pgfqpoint{5.568030in}{5.294515in}}%
\pgfpathclose%
\pgfusepath{stroke,fill}%
\end{pgfscope}%
\begin{pgfscope}%
\pgfpathrectangle{\pgfqpoint{0.800000in}{0.528000in}}{\pgfqpoint{4.960000in}{3.696000in}} %
\pgfusepath{clip}%
\pgfsetbuttcap%
\pgfsetroundjoin%
\definecolor{currentfill}{rgb}{1.000000,0.647059,0.000000}%
\pgfsetfillcolor{currentfill}%
\pgfsetlinewidth{1.003750pt}%
\definecolor{currentstroke}{rgb}{1.000000,0.647059,0.000000}%
\pgfsetstrokecolor{currentstroke}%
\pgfsetdash{}{0pt}%
\pgfpathmoveto{\pgfqpoint{5.659798in}{6.999419in}}%
\pgfpathlineto{\pgfqpoint{5.701465in}{7.082752in}}%
\pgfpathlineto{\pgfqpoint{5.618131in}{7.082752in}}%
\pgfpathclose%
\pgfusepath{stroke,fill}%
\end{pgfscope}%
\begin{pgfscope}%
\pgfpathrectangle{\pgfqpoint{0.800000in}{0.528000in}}{\pgfqpoint{4.960000in}{3.696000in}} %
\pgfusepath{clip}%
\pgfsetbuttcap%
\pgfsetroundjoin%
\definecolor{currentfill}{rgb}{1.000000,0.647059,0.000000}%
\pgfsetfillcolor{currentfill}%
\pgfsetlinewidth{1.003750pt}%
\definecolor{currentstroke}{rgb}{1.000000,0.647059,0.000000}%
\pgfsetstrokecolor{currentstroke}%
\pgfsetdash{}{0pt}%
\pgfpathmoveto{\pgfqpoint{5.709899in}{6.584428in}}%
\pgfpathlineto{\pgfqpoint{5.751566in}{6.667761in}}%
\pgfpathlineto{\pgfqpoint{5.668232in}{6.667761in}}%
\pgfpathclose%
\pgfusepath{stroke,fill}%
\end{pgfscope}%
\begin{pgfscope}%
\pgfpathrectangle{\pgfqpoint{0.800000in}{0.528000in}}{\pgfqpoint{4.960000in}{3.696000in}} %
\pgfusepath{clip}%
\pgfsetbuttcap%
\pgfsetroundjoin%
\definecolor{currentfill}{rgb}{1.000000,0.647059,0.000000}%
\pgfsetfillcolor{currentfill}%
\pgfsetlinewidth{1.003750pt}%
\definecolor{currentstroke}{rgb}{1.000000,0.647059,0.000000}%
\pgfsetstrokecolor{currentstroke}%
\pgfsetdash{}{0pt}%
\pgfpathmoveto{\pgfqpoint{5.760000in}{7.412436in}}%
\pgfpathlineto{\pgfqpoint{5.801667in}{7.495770in}}%
\pgfpathlineto{\pgfqpoint{5.718333in}{7.495770in}}%
\pgfpathclose%
\pgfusepath{stroke,fill}%
\end{pgfscope}%
\begin{pgfscope}%
\pgfpathrectangle{\pgfqpoint{0.800000in}{0.528000in}}{\pgfqpoint{4.960000in}{3.696000in}} %
\pgfusepath{clip}%
\pgfsetbuttcap%
\pgfsetroundjoin%
\definecolor{currentfill}{rgb}{0.000000,0.500000,0.000000}%
\pgfsetfillcolor{currentfill}%
\pgfsetlinewidth{1.003750pt}%
\definecolor{currentstroke}{rgb}{0.000000,0.500000,0.000000}%
\pgfsetstrokecolor{currentstroke}%
\pgfsetdash{}{0pt}%
\pgfpathmoveto{\pgfqpoint{0.800000in}{-1.187955in}}%
\pgfpathcurveto{\pgfqpoint{0.811050in}{-1.187955in}}{\pgfqpoint{0.821649in}{-1.183564in}}{\pgfqpoint{0.829463in}{-1.175751in}}%
\pgfpathcurveto{\pgfqpoint{0.837276in}{-1.167937in}}{\pgfqpoint{0.841667in}{-1.157338in}}{\pgfqpoint{0.841667in}{-1.146288in}}%
\pgfpathcurveto{\pgfqpoint{0.841667in}{-1.135238in}}{\pgfqpoint{0.837276in}{-1.124639in}}{\pgfqpoint{0.829463in}{-1.116825in}}%
\pgfpathcurveto{\pgfqpoint{0.821649in}{-1.109012in}}{\pgfqpoint{0.811050in}{-1.104621in}}{\pgfqpoint{0.800000in}{-1.104621in}}%
\pgfpathcurveto{\pgfqpoint{0.788950in}{-1.104621in}}{\pgfqpoint{0.778351in}{-1.109012in}}{\pgfqpoint{0.770537in}{-1.116825in}}%
\pgfpathcurveto{\pgfqpoint{0.762724in}{-1.124639in}}{\pgfqpoint{0.758333in}{-1.135238in}}{\pgfqpoint{0.758333in}{-1.146288in}}%
\pgfpathcurveto{\pgfqpoint{0.758333in}{-1.157338in}}{\pgfqpoint{0.762724in}{-1.167937in}}{\pgfqpoint{0.770537in}{-1.175751in}}%
\pgfpathcurveto{\pgfqpoint{0.778351in}{-1.183564in}}{\pgfqpoint{0.788950in}{-1.187955in}}{\pgfqpoint{0.800000in}{-1.187955in}}%
\pgfpathclose%
\pgfusepath{stroke,fill}%
\end{pgfscope}%
\begin{pgfscope}%
\pgfpathrectangle{\pgfqpoint{0.800000in}{0.528000in}}{\pgfqpoint{4.960000in}{3.696000in}} %
\pgfusepath{clip}%
\pgfsetbuttcap%
\pgfsetroundjoin%
\definecolor{currentfill}{rgb}{0.000000,0.500000,0.000000}%
\pgfsetfillcolor{currentfill}%
\pgfsetlinewidth{1.003750pt}%
\definecolor{currentstroke}{rgb}{0.000000,0.500000,0.000000}%
\pgfsetstrokecolor{currentstroke}%
\pgfsetdash{}{0pt}%
\pgfpathmoveto{\pgfqpoint{0.850101in}{-1.652354in}}%
\pgfpathcurveto{\pgfqpoint{0.861151in}{-1.652354in}}{\pgfqpoint{0.871750in}{-1.647964in}}{\pgfqpoint{0.879564in}{-1.640150in}}%
\pgfpathcurveto{\pgfqpoint{0.887377in}{-1.632337in}}{\pgfqpoint{0.891768in}{-1.621738in}}{\pgfqpoint{0.891768in}{-1.610687in}}%
\pgfpathcurveto{\pgfqpoint{0.891768in}{-1.599637in}}{\pgfqpoint{0.887377in}{-1.589038in}}{\pgfqpoint{0.879564in}{-1.581225in}}%
\pgfpathcurveto{\pgfqpoint{0.871750in}{-1.573411in}}{\pgfqpoint{0.861151in}{-1.569021in}}{\pgfqpoint{0.850101in}{-1.569021in}}%
\pgfpathcurveto{\pgfqpoint{0.839051in}{-1.569021in}}{\pgfqpoint{0.828452in}{-1.573411in}}{\pgfqpoint{0.820638in}{-1.581225in}}%
\pgfpathcurveto{\pgfqpoint{0.812825in}{-1.589038in}}{\pgfqpoint{0.808434in}{-1.599637in}}{\pgfqpoint{0.808434in}{-1.610687in}}%
\pgfpathcurveto{\pgfqpoint{0.808434in}{-1.621738in}}{\pgfqpoint{0.812825in}{-1.632337in}}{\pgfqpoint{0.820638in}{-1.640150in}}%
\pgfpathcurveto{\pgfqpoint{0.828452in}{-1.647964in}}{\pgfqpoint{0.839051in}{-1.652354in}}{\pgfqpoint{0.850101in}{-1.652354in}}%
\pgfpathclose%
\pgfusepath{stroke,fill}%
\end{pgfscope}%
\begin{pgfscope}%
\pgfpathrectangle{\pgfqpoint{0.800000in}{0.528000in}}{\pgfqpoint{4.960000in}{3.696000in}} %
\pgfusepath{clip}%
\pgfsetbuttcap%
\pgfsetroundjoin%
\definecolor{currentfill}{rgb}{0.000000,0.500000,0.000000}%
\pgfsetfillcolor{currentfill}%
\pgfsetlinewidth{1.003750pt}%
\definecolor{currentstroke}{rgb}{0.000000,0.500000,0.000000}%
\pgfsetstrokecolor{currentstroke}%
\pgfsetdash{}{0pt}%
\pgfpathmoveto{\pgfqpoint{0.900202in}{-0.741881in}}%
\pgfpathcurveto{\pgfqpoint{0.911252in}{-0.741881in}}{\pgfqpoint{0.921851in}{-0.737491in}}{\pgfqpoint{0.929665in}{-0.729677in}}%
\pgfpathcurveto{\pgfqpoint{0.937478in}{-0.721864in}}{\pgfqpoint{0.941869in}{-0.711265in}}{\pgfqpoint{0.941869in}{-0.700214in}}%
\pgfpathcurveto{\pgfqpoint{0.941869in}{-0.689164in}}{\pgfqpoint{0.937478in}{-0.678565in}}{\pgfqpoint{0.929665in}{-0.670752in}}%
\pgfpathcurveto{\pgfqpoint{0.921851in}{-0.662938in}}{\pgfqpoint{0.911252in}{-0.658548in}}{\pgfqpoint{0.900202in}{-0.658548in}}%
\pgfpathcurveto{\pgfqpoint{0.889152in}{-0.658548in}}{\pgfqpoint{0.878553in}{-0.662938in}}{\pgfqpoint{0.870739in}{-0.670752in}}%
\pgfpathcurveto{\pgfqpoint{0.862926in}{-0.678565in}}{\pgfqpoint{0.858535in}{-0.689164in}}{\pgfqpoint{0.858535in}{-0.700214in}}%
\pgfpathcurveto{\pgfqpoint{0.858535in}{-0.711265in}}{\pgfqpoint{0.862926in}{-0.721864in}}{\pgfqpoint{0.870739in}{-0.729677in}}%
\pgfpathcurveto{\pgfqpoint{0.878553in}{-0.737491in}}{\pgfqpoint{0.889152in}{-0.741881in}}{\pgfqpoint{0.900202in}{-0.741881in}}%
\pgfpathclose%
\pgfusepath{stroke,fill}%
\end{pgfscope}%
\begin{pgfscope}%
\pgfpathrectangle{\pgfqpoint{0.800000in}{0.528000in}}{\pgfqpoint{4.960000in}{3.696000in}} %
\pgfusepath{clip}%
\pgfsetbuttcap%
\pgfsetroundjoin%
\definecolor{currentfill}{rgb}{0.000000,0.500000,0.000000}%
\pgfsetfillcolor{currentfill}%
\pgfsetlinewidth{1.003750pt}%
\definecolor{currentstroke}{rgb}{0.000000,0.500000,0.000000}%
\pgfsetstrokecolor{currentstroke}%
\pgfsetdash{}{0pt}%
\pgfpathmoveto{\pgfqpoint{0.950303in}{-2.603040in}}%
\pgfpathcurveto{\pgfqpoint{0.961353in}{-2.603040in}}{\pgfqpoint{0.971952in}{-2.598649in}}{\pgfqpoint{0.979766in}{-2.590836in}}%
\pgfpathcurveto{\pgfqpoint{0.987579in}{-2.583022in}}{\pgfqpoint{0.991970in}{-2.572423in}}{\pgfqpoint{0.991970in}{-2.561373in}}%
\pgfpathcurveto{\pgfqpoint{0.991970in}{-2.550323in}}{\pgfqpoint{0.987579in}{-2.539724in}}{\pgfqpoint{0.979766in}{-2.531910in}}%
\pgfpathcurveto{\pgfqpoint{0.971952in}{-2.524096in}}{\pgfqpoint{0.961353in}{-2.519706in}}{\pgfqpoint{0.950303in}{-2.519706in}}%
\pgfpathcurveto{\pgfqpoint{0.939253in}{-2.519706in}}{\pgfqpoint{0.928654in}{-2.524096in}}{\pgfqpoint{0.920840in}{-2.531910in}}%
\pgfpathcurveto{\pgfqpoint{0.913027in}{-2.539724in}}{\pgfqpoint{0.908636in}{-2.550323in}}{\pgfqpoint{0.908636in}{-2.561373in}}%
\pgfpathcurveto{\pgfqpoint{0.908636in}{-2.572423in}}{\pgfqpoint{0.913027in}{-2.583022in}}{\pgfqpoint{0.920840in}{-2.590836in}}%
\pgfpathcurveto{\pgfqpoint{0.928654in}{-2.598649in}}{\pgfqpoint{0.939253in}{-2.603040in}}{\pgfqpoint{0.950303in}{-2.603040in}}%
\pgfpathclose%
\pgfusepath{stroke,fill}%
\end{pgfscope}%
\begin{pgfscope}%
\pgfpathrectangle{\pgfqpoint{0.800000in}{0.528000in}}{\pgfqpoint{4.960000in}{3.696000in}} %
\pgfusepath{clip}%
\pgfsetbuttcap%
\pgfsetroundjoin%
\definecolor{currentfill}{rgb}{0.000000,0.500000,0.000000}%
\pgfsetfillcolor{currentfill}%
\pgfsetlinewidth{1.003750pt}%
\definecolor{currentstroke}{rgb}{0.000000,0.500000,0.000000}%
\pgfsetstrokecolor{currentstroke}%
\pgfsetdash{}{0pt}%
\pgfpathmoveto{\pgfqpoint{1.000404in}{-1.063397in}}%
\pgfpathcurveto{\pgfqpoint{1.011454in}{-1.063397in}}{\pgfqpoint{1.022053in}{-1.059007in}}{\pgfqpoint{1.029867in}{-1.051193in}}%
\pgfpathcurveto{\pgfqpoint{1.037680in}{-1.043380in}}{\pgfqpoint{1.042071in}{-1.032781in}}{\pgfqpoint{1.042071in}{-1.021731in}}%
\pgfpathcurveto{\pgfqpoint{1.042071in}{-1.010681in}}{\pgfqpoint{1.037680in}{-1.000081in}}{\pgfqpoint{1.029867in}{-0.992268in}}%
\pgfpathcurveto{\pgfqpoint{1.022053in}{-0.984454in}}{\pgfqpoint{1.011454in}{-0.980064in}}{\pgfqpoint{1.000404in}{-0.980064in}}%
\pgfpathcurveto{\pgfqpoint{0.989354in}{-0.980064in}}{\pgfqpoint{0.978755in}{-0.984454in}}{\pgfqpoint{0.970941in}{-0.992268in}}%
\pgfpathcurveto{\pgfqpoint{0.963128in}{-1.000081in}}{\pgfqpoint{0.958737in}{-1.010681in}}{\pgfqpoint{0.958737in}{-1.021731in}}%
\pgfpathcurveto{\pgfqpoint{0.958737in}{-1.032781in}}{\pgfqpoint{0.963128in}{-1.043380in}}{\pgfqpoint{0.970941in}{-1.051193in}}%
\pgfpathcurveto{\pgfqpoint{0.978755in}{-1.059007in}}{\pgfqpoint{0.989354in}{-1.063397in}}{\pgfqpoint{1.000404in}{-1.063397in}}%
\pgfpathclose%
\pgfusepath{stroke,fill}%
\end{pgfscope}%
\begin{pgfscope}%
\pgfpathrectangle{\pgfqpoint{0.800000in}{0.528000in}}{\pgfqpoint{4.960000in}{3.696000in}} %
\pgfusepath{clip}%
\pgfsetbuttcap%
\pgfsetroundjoin%
\definecolor{currentfill}{rgb}{0.000000,0.500000,0.000000}%
\pgfsetfillcolor{currentfill}%
\pgfsetlinewidth{1.003750pt}%
\definecolor{currentstroke}{rgb}{0.000000,0.500000,0.000000}%
\pgfsetstrokecolor{currentstroke}%
\pgfsetdash{}{0pt}%
\pgfpathmoveto{\pgfqpoint{1.050505in}{-2.106921in}}%
\pgfpathcurveto{\pgfqpoint{1.061555in}{-2.106921in}}{\pgfqpoint{1.072154in}{-2.102531in}}{\pgfqpoint{1.079968in}{-2.094718in}}%
\pgfpathcurveto{\pgfqpoint{1.087781in}{-2.086904in}}{\pgfqpoint{1.092172in}{-2.076305in}}{\pgfqpoint{1.092172in}{-2.065255in}}%
\pgfpathcurveto{\pgfqpoint{1.092172in}{-2.054205in}}{\pgfqpoint{1.087781in}{-2.043606in}}{\pgfqpoint{1.079968in}{-2.035792in}}%
\pgfpathcurveto{\pgfqpoint{1.072154in}{-2.027978in}}{\pgfqpoint{1.061555in}{-2.023588in}}{\pgfqpoint{1.050505in}{-2.023588in}}%
\pgfpathcurveto{\pgfqpoint{1.039455in}{-2.023588in}}{\pgfqpoint{1.028856in}{-2.027978in}}{\pgfqpoint{1.021042in}{-2.035792in}}%
\pgfpathcurveto{\pgfqpoint{1.013229in}{-2.043606in}}{\pgfqpoint{1.008838in}{-2.054205in}}{\pgfqpoint{1.008838in}{-2.065255in}}%
\pgfpathcurveto{\pgfqpoint{1.008838in}{-2.076305in}}{\pgfqpoint{1.013229in}{-2.086904in}}{\pgfqpoint{1.021042in}{-2.094718in}}%
\pgfpathcurveto{\pgfqpoint{1.028856in}{-2.102531in}}{\pgfqpoint{1.039455in}{-2.106921in}}{\pgfqpoint{1.050505in}{-2.106921in}}%
\pgfpathclose%
\pgfusepath{stroke,fill}%
\end{pgfscope}%
\begin{pgfscope}%
\pgfpathrectangle{\pgfqpoint{0.800000in}{0.528000in}}{\pgfqpoint{4.960000in}{3.696000in}} %
\pgfusepath{clip}%
\pgfsetbuttcap%
\pgfsetroundjoin%
\definecolor{currentfill}{rgb}{0.000000,0.500000,0.000000}%
\pgfsetfillcolor{currentfill}%
\pgfsetlinewidth{1.003750pt}%
\definecolor{currentstroke}{rgb}{0.000000,0.500000,0.000000}%
\pgfsetstrokecolor{currentstroke}%
\pgfsetdash{}{0pt}%
\pgfpathmoveto{\pgfqpoint{1.100606in}{-0.291055in}}%
\pgfpathcurveto{\pgfqpoint{1.111656in}{-0.291055in}}{\pgfqpoint{1.122255in}{-0.286665in}}{\pgfqpoint{1.130069in}{-0.278851in}}%
\pgfpathcurveto{\pgfqpoint{1.137882in}{-0.271038in}}{\pgfqpoint{1.142273in}{-0.260439in}}{\pgfqpoint{1.142273in}{-0.249389in}}%
\pgfpathcurveto{\pgfqpoint{1.142273in}{-0.238338in}}{\pgfqpoint{1.137882in}{-0.227739in}}{\pgfqpoint{1.130069in}{-0.219926in}}%
\pgfpathcurveto{\pgfqpoint{1.122255in}{-0.212112in}}{\pgfqpoint{1.111656in}{-0.207722in}}{\pgfqpoint{1.100606in}{-0.207722in}}%
\pgfpathcurveto{\pgfqpoint{1.089556in}{-0.207722in}}{\pgfqpoint{1.078957in}{-0.212112in}}{\pgfqpoint{1.071143in}{-0.219926in}}%
\pgfpathcurveto{\pgfqpoint{1.063330in}{-0.227739in}}{\pgfqpoint{1.058939in}{-0.238338in}}{\pgfqpoint{1.058939in}{-0.249389in}}%
\pgfpathcurveto{\pgfqpoint{1.058939in}{-0.260439in}}{\pgfqpoint{1.063330in}{-0.271038in}}{\pgfqpoint{1.071143in}{-0.278851in}}%
\pgfpathcurveto{\pgfqpoint{1.078957in}{-0.286665in}}{\pgfqpoint{1.089556in}{-0.291055in}}{\pgfqpoint{1.100606in}{-0.291055in}}%
\pgfpathclose%
\pgfusepath{stroke,fill}%
\end{pgfscope}%
\begin{pgfscope}%
\pgfpathrectangle{\pgfqpoint{0.800000in}{0.528000in}}{\pgfqpoint{4.960000in}{3.696000in}} %
\pgfusepath{clip}%
\pgfsetbuttcap%
\pgfsetroundjoin%
\definecolor{currentfill}{rgb}{0.000000,0.500000,0.000000}%
\pgfsetfillcolor{currentfill}%
\pgfsetlinewidth{1.003750pt}%
\definecolor{currentstroke}{rgb}{0.000000,0.500000,0.000000}%
\pgfsetstrokecolor{currentstroke}%
\pgfsetdash{}{0pt}%
\pgfpathmoveto{\pgfqpoint{1.150707in}{-1.030032in}}%
\pgfpathcurveto{\pgfqpoint{1.161757in}{-1.030032in}}{\pgfqpoint{1.172356in}{-1.025642in}}{\pgfqpoint{1.180170in}{-1.017828in}}%
\pgfpathcurveto{\pgfqpoint{1.187983in}{-1.010015in}}{\pgfqpoint{1.192374in}{-0.999416in}}{\pgfqpoint{1.192374in}{-0.988366in}}%
\pgfpathcurveto{\pgfqpoint{1.192374in}{-0.977316in}}{\pgfqpoint{1.187983in}{-0.966716in}}{\pgfqpoint{1.180170in}{-0.958903in}}%
\pgfpathcurveto{\pgfqpoint{1.172356in}{-0.951089in}}{\pgfqpoint{1.161757in}{-0.946699in}}{\pgfqpoint{1.150707in}{-0.946699in}}%
\pgfpathcurveto{\pgfqpoint{1.139657in}{-0.946699in}}{\pgfqpoint{1.129058in}{-0.951089in}}{\pgfqpoint{1.121244in}{-0.958903in}}%
\pgfpathcurveto{\pgfqpoint{1.113431in}{-0.966716in}}{\pgfqpoint{1.109040in}{-0.977316in}}{\pgfqpoint{1.109040in}{-0.988366in}}%
\pgfpathcurveto{\pgfqpoint{1.109040in}{-0.999416in}}{\pgfqpoint{1.113431in}{-1.010015in}}{\pgfqpoint{1.121244in}{-1.017828in}}%
\pgfpathcurveto{\pgfqpoint{1.129058in}{-1.025642in}}{\pgfqpoint{1.139657in}{-1.030032in}}{\pgfqpoint{1.150707in}{-1.030032in}}%
\pgfpathclose%
\pgfusepath{stroke,fill}%
\end{pgfscope}%
\begin{pgfscope}%
\pgfpathrectangle{\pgfqpoint{0.800000in}{0.528000in}}{\pgfqpoint{4.960000in}{3.696000in}} %
\pgfusepath{clip}%
\pgfsetbuttcap%
\pgfsetroundjoin%
\definecolor{currentfill}{rgb}{0.000000,0.500000,0.000000}%
\pgfsetfillcolor{currentfill}%
\pgfsetlinewidth{1.003750pt}%
\definecolor{currentstroke}{rgb}{0.000000,0.500000,0.000000}%
\pgfsetstrokecolor{currentstroke}%
\pgfsetdash{}{0pt}%
\pgfpathmoveto{\pgfqpoint{1.200808in}{-0.041152in}}%
\pgfpathcurveto{\pgfqpoint{1.211858in}{-0.041152in}}{\pgfqpoint{1.222457in}{-0.036761in}}{\pgfqpoint{1.230271in}{-0.028948in}}%
\pgfpathcurveto{\pgfqpoint{1.238084in}{-0.021134in}}{\pgfqpoint{1.242475in}{-0.010535in}}{\pgfqpoint{1.242475in}{0.000515in}}%
\pgfpathcurveto{\pgfqpoint{1.242475in}{0.011565in}}{\pgfqpoint{1.238084in}{0.022164in}}{\pgfqpoint{1.230271in}{0.029978in}}%
\pgfpathcurveto{\pgfqpoint{1.222457in}{0.037791in}}{\pgfqpoint{1.211858in}{0.042182in}}{\pgfqpoint{1.200808in}{0.042182in}}%
\pgfpathcurveto{\pgfqpoint{1.189758in}{0.042182in}}{\pgfqpoint{1.179159in}{0.037791in}}{\pgfqpoint{1.171345in}{0.029978in}}%
\pgfpathcurveto{\pgfqpoint{1.163532in}{0.022164in}}{\pgfqpoint{1.159141in}{0.011565in}}{\pgfqpoint{1.159141in}{0.000515in}}%
\pgfpathcurveto{\pgfqpoint{1.159141in}{-0.010535in}}{\pgfqpoint{1.163532in}{-0.021134in}}{\pgfqpoint{1.171345in}{-0.028948in}}%
\pgfpathcurveto{\pgfqpoint{1.179159in}{-0.036761in}}{\pgfqpoint{1.189758in}{-0.041152in}}{\pgfqpoint{1.200808in}{-0.041152in}}%
\pgfpathclose%
\pgfusepath{stroke,fill}%
\end{pgfscope}%
\begin{pgfscope}%
\pgfpathrectangle{\pgfqpoint{0.800000in}{0.528000in}}{\pgfqpoint{4.960000in}{3.696000in}} %
\pgfusepath{clip}%
\pgfsetbuttcap%
\pgfsetroundjoin%
\definecolor{currentfill}{rgb}{0.000000,0.500000,0.000000}%
\pgfsetfillcolor{currentfill}%
\pgfsetlinewidth{1.003750pt}%
\definecolor{currentstroke}{rgb}{0.000000,0.500000,0.000000}%
\pgfsetstrokecolor{currentstroke}%
\pgfsetdash{}{0pt}%
\pgfpathmoveto{\pgfqpoint{1.250909in}{-1.392512in}}%
\pgfpathcurveto{\pgfqpoint{1.261959in}{-1.392512in}}{\pgfqpoint{1.272558in}{-1.388122in}}{\pgfqpoint{1.280372in}{-1.380308in}}%
\pgfpathcurveto{\pgfqpoint{1.288185in}{-1.372495in}}{\pgfqpoint{1.292576in}{-1.361895in}}{\pgfqpoint{1.292576in}{-1.350845in}}%
\pgfpathcurveto{\pgfqpoint{1.292576in}{-1.339795in}}{\pgfqpoint{1.288185in}{-1.329196in}}{\pgfqpoint{1.280372in}{-1.321383in}}%
\pgfpathcurveto{\pgfqpoint{1.272558in}{-1.313569in}}{\pgfqpoint{1.261959in}{-1.309179in}}{\pgfqpoint{1.250909in}{-1.309179in}}%
\pgfpathcurveto{\pgfqpoint{1.239859in}{-1.309179in}}{\pgfqpoint{1.229260in}{-1.313569in}}{\pgfqpoint{1.221446in}{-1.321383in}}%
\pgfpathcurveto{\pgfqpoint{1.213633in}{-1.329196in}}{\pgfqpoint{1.209242in}{-1.339795in}}{\pgfqpoint{1.209242in}{-1.350845in}}%
\pgfpathcurveto{\pgfqpoint{1.209242in}{-1.361895in}}{\pgfqpoint{1.213633in}{-1.372495in}}{\pgfqpoint{1.221446in}{-1.380308in}}%
\pgfpathcurveto{\pgfqpoint{1.229260in}{-1.388122in}}{\pgfqpoint{1.239859in}{-1.392512in}}{\pgfqpoint{1.250909in}{-1.392512in}}%
\pgfpathclose%
\pgfusepath{stroke,fill}%
\end{pgfscope}%
\begin{pgfscope}%
\pgfpathrectangle{\pgfqpoint{0.800000in}{0.528000in}}{\pgfqpoint{4.960000in}{3.696000in}} %
\pgfusepath{clip}%
\pgfsetbuttcap%
\pgfsetroundjoin%
\definecolor{currentfill}{rgb}{0.000000,0.500000,0.000000}%
\pgfsetfillcolor{currentfill}%
\pgfsetlinewidth{1.003750pt}%
\definecolor{currentstroke}{rgb}{0.000000,0.500000,0.000000}%
\pgfsetstrokecolor{currentstroke}%
\pgfsetdash{}{0pt}%
\pgfpathmoveto{\pgfqpoint{1.301010in}{-2.646803in}}%
\pgfpathcurveto{\pgfqpoint{1.312060in}{-2.646803in}}{\pgfqpoint{1.322659in}{-2.642413in}}{\pgfqpoint{1.330473in}{-2.634599in}}%
\pgfpathcurveto{\pgfqpoint{1.338287in}{-2.626786in}}{\pgfqpoint{1.342677in}{-2.616187in}}{\pgfqpoint{1.342677in}{-2.605137in}}%
\pgfpathcurveto{\pgfqpoint{1.342677in}{-2.594086in}}{\pgfqpoint{1.338287in}{-2.583487in}}{\pgfqpoint{1.330473in}{-2.575674in}}%
\pgfpathcurveto{\pgfqpoint{1.322659in}{-2.567860in}}{\pgfqpoint{1.312060in}{-2.563470in}}{\pgfqpoint{1.301010in}{-2.563470in}}%
\pgfpathcurveto{\pgfqpoint{1.289960in}{-2.563470in}}{\pgfqpoint{1.279361in}{-2.567860in}}{\pgfqpoint{1.271547in}{-2.575674in}}%
\pgfpathcurveto{\pgfqpoint{1.263734in}{-2.583487in}}{\pgfqpoint{1.259343in}{-2.594086in}}{\pgfqpoint{1.259343in}{-2.605137in}}%
\pgfpathcurveto{\pgfqpoint{1.259343in}{-2.616187in}}{\pgfqpoint{1.263734in}{-2.626786in}}{\pgfqpoint{1.271547in}{-2.634599in}}%
\pgfpathcurveto{\pgfqpoint{1.279361in}{-2.642413in}}{\pgfqpoint{1.289960in}{-2.646803in}}{\pgfqpoint{1.301010in}{-2.646803in}}%
\pgfpathclose%
\pgfusepath{stroke,fill}%
\end{pgfscope}%
\begin{pgfscope}%
\pgfpathrectangle{\pgfqpoint{0.800000in}{0.528000in}}{\pgfqpoint{4.960000in}{3.696000in}} %
\pgfusepath{clip}%
\pgfsetbuttcap%
\pgfsetroundjoin%
\definecolor{currentfill}{rgb}{0.000000,0.500000,0.000000}%
\pgfsetfillcolor{currentfill}%
\pgfsetlinewidth{1.003750pt}%
\definecolor{currentstroke}{rgb}{0.000000,0.500000,0.000000}%
\pgfsetstrokecolor{currentstroke}%
\pgfsetdash{}{0pt}%
\pgfpathmoveto{\pgfqpoint{1.351111in}{0.239201in}}%
\pgfpathcurveto{\pgfqpoint{1.362161in}{0.239201in}}{\pgfqpoint{1.372760in}{0.243592in}}{\pgfqpoint{1.380574in}{0.251405in}}%
\pgfpathcurveto{\pgfqpoint{1.388388in}{0.259219in}}{\pgfqpoint{1.392778in}{0.269818in}}{\pgfqpoint{1.392778in}{0.280868in}}%
\pgfpathcurveto{\pgfqpoint{1.392778in}{0.291918in}}{\pgfqpoint{1.388388in}{0.302517in}}{\pgfqpoint{1.380574in}{0.310331in}}%
\pgfpathcurveto{\pgfqpoint{1.372760in}{0.318144in}}{\pgfqpoint{1.362161in}{0.322535in}}{\pgfqpoint{1.351111in}{0.322535in}}%
\pgfpathcurveto{\pgfqpoint{1.340061in}{0.322535in}}{\pgfqpoint{1.329462in}{0.318144in}}{\pgfqpoint{1.321648in}{0.310331in}}%
\pgfpathcurveto{\pgfqpoint{1.313835in}{0.302517in}}{\pgfqpoint{1.309444in}{0.291918in}}{\pgfqpoint{1.309444in}{0.280868in}}%
\pgfpathcurveto{\pgfqpoint{1.309444in}{0.269818in}}{\pgfqpoint{1.313835in}{0.259219in}}{\pgfqpoint{1.321648in}{0.251405in}}%
\pgfpathcurveto{\pgfqpoint{1.329462in}{0.243592in}}{\pgfqpoint{1.340061in}{0.239201in}}{\pgfqpoint{1.351111in}{0.239201in}}%
\pgfpathclose%
\pgfusepath{stroke,fill}%
\end{pgfscope}%
\begin{pgfscope}%
\pgfpathrectangle{\pgfqpoint{0.800000in}{0.528000in}}{\pgfqpoint{4.960000in}{3.696000in}} %
\pgfusepath{clip}%
\pgfsetbuttcap%
\pgfsetroundjoin%
\definecolor{currentfill}{rgb}{0.000000,0.500000,0.000000}%
\pgfsetfillcolor{currentfill}%
\pgfsetlinewidth{1.003750pt}%
\definecolor{currentstroke}{rgb}{0.000000,0.500000,0.000000}%
\pgfsetstrokecolor{currentstroke}%
\pgfsetdash{}{0pt}%
\pgfpathmoveto{\pgfqpoint{1.401212in}{-1.100664in}}%
\pgfpathcurveto{\pgfqpoint{1.412262in}{-1.100664in}}{\pgfqpoint{1.422861in}{-1.096273in}}{\pgfqpoint{1.430675in}{-1.088460in}}%
\pgfpathcurveto{\pgfqpoint{1.438489in}{-1.080646in}}{\pgfqpoint{1.442879in}{-1.070047in}}{\pgfqpoint{1.442879in}{-1.058997in}}%
\pgfpathcurveto{\pgfqpoint{1.442879in}{-1.047947in}}{\pgfqpoint{1.438489in}{-1.037348in}}{\pgfqpoint{1.430675in}{-1.029534in}}%
\pgfpathcurveto{\pgfqpoint{1.422861in}{-1.021720in}}{\pgfqpoint{1.412262in}{-1.017330in}}{\pgfqpoint{1.401212in}{-1.017330in}}%
\pgfpathcurveto{\pgfqpoint{1.390162in}{-1.017330in}}{\pgfqpoint{1.379563in}{-1.021720in}}{\pgfqpoint{1.371749in}{-1.029534in}}%
\pgfpathcurveto{\pgfqpoint{1.363936in}{-1.037348in}}{\pgfqpoint{1.359545in}{-1.047947in}}{\pgfqpoint{1.359545in}{-1.058997in}}%
\pgfpathcurveto{\pgfqpoint{1.359545in}{-1.070047in}}{\pgfqpoint{1.363936in}{-1.080646in}}{\pgfqpoint{1.371749in}{-1.088460in}}%
\pgfpathcurveto{\pgfqpoint{1.379563in}{-1.096273in}}{\pgfqpoint{1.390162in}{-1.100664in}}{\pgfqpoint{1.401212in}{-1.100664in}}%
\pgfpathclose%
\pgfusepath{stroke,fill}%
\end{pgfscope}%
\begin{pgfscope}%
\pgfpathrectangle{\pgfqpoint{0.800000in}{0.528000in}}{\pgfqpoint{4.960000in}{3.696000in}} %
\pgfusepath{clip}%
\pgfsetbuttcap%
\pgfsetroundjoin%
\definecolor{currentfill}{rgb}{0.000000,0.500000,0.000000}%
\pgfsetfillcolor{currentfill}%
\pgfsetlinewidth{1.003750pt}%
\definecolor{currentstroke}{rgb}{0.000000,0.500000,0.000000}%
\pgfsetstrokecolor{currentstroke}%
\pgfsetdash{}{0pt}%
\pgfpathmoveto{\pgfqpoint{1.451313in}{-0.157967in}}%
\pgfpathcurveto{\pgfqpoint{1.462363in}{-0.157967in}}{\pgfqpoint{1.472962in}{-0.153577in}}{\pgfqpoint{1.480776in}{-0.145764in}}%
\pgfpathcurveto{\pgfqpoint{1.488590in}{-0.137950in}}{\pgfqpoint{1.492980in}{-0.127351in}}{\pgfqpoint{1.492980in}{-0.116301in}}%
\pgfpathcurveto{\pgfqpoint{1.492980in}{-0.105251in}}{\pgfqpoint{1.488590in}{-0.094652in}}{\pgfqpoint{1.480776in}{-0.086838in}}%
\pgfpathcurveto{\pgfqpoint{1.472962in}{-0.079024in}}{\pgfqpoint{1.462363in}{-0.074634in}}{\pgfqpoint{1.451313in}{-0.074634in}}%
\pgfpathcurveto{\pgfqpoint{1.440263in}{-0.074634in}}{\pgfqpoint{1.429664in}{-0.079024in}}{\pgfqpoint{1.421850in}{-0.086838in}}%
\pgfpathcurveto{\pgfqpoint{1.414037in}{-0.094652in}}{\pgfqpoint{1.409646in}{-0.105251in}}{\pgfqpoint{1.409646in}{-0.116301in}}%
\pgfpathcurveto{\pgfqpoint{1.409646in}{-0.127351in}}{\pgfqpoint{1.414037in}{-0.137950in}}{\pgfqpoint{1.421850in}{-0.145764in}}%
\pgfpathcurveto{\pgfqpoint{1.429664in}{-0.153577in}}{\pgfqpoint{1.440263in}{-0.157967in}}{\pgfqpoint{1.451313in}{-0.157967in}}%
\pgfpathclose%
\pgfusepath{stroke,fill}%
\end{pgfscope}%
\begin{pgfscope}%
\pgfpathrectangle{\pgfqpoint{0.800000in}{0.528000in}}{\pgfqpoint{4.960000in}{3.696000in}} %
\pgfusepath{clip}%
\pgfsetbuttcap%
\pgfsetroundjoin%
\definecolor{currentfill}{rgb}{0.000000,0.500000,0.000000}%
\pgfsetfillcolor{currentfill}%
\pgfsetlinewidth{1.003750pt}%
\definecolor{currentstroke}{rgb}{0.000000,0.500000,0.000000}%
\pgfsetstrokecolor{currentstroke}%
\pgfsetdash{}{0pt}%
\pgfpathmoveto{\pgfqpoint{1.501414in}{-1.938403in}}%
\pgfpathcurveto{\pgfqpoint{1.512464in}{-1.938403in}}{\pgfqpoint{1.523063in}{-1.934013in}}{\pgfqpoint{1.530877in}{-1.926199in}}%
\pgfpathcurveto{\pgfqpoint{1.538691in}{-1.918385in}}{\pgfqpoint{1.543081in}{-1.907786in}}{\pgfqpoint{1.543081in}{-1.896736in}}%
\pgfpathcurveto{\pgfqpoint{1.543081in}{-1.885686in}}{\pgfqpoint{1.538691in}{-1.875087in}}{\pgfqpoint{1.530877in}{-1.867274in}}%
\pgfpathcurveto{\pgfqpoint{1.523063in}{-1.859460in}}{\pgfqpoint{1.512464in}{-1.855070in}}{\pgfqpoint{1.501414in}{-1.855070in}}%
\pgfpathcurveto{\pgfqpoint{1.490364in}{-1.855070in}}{\pgfqpoint{1.479765in}{-1.859460in}}{\pgfqpoint{1.471951in}{-1.867274in}}%
\pgfpathcurveto{\pgfqpoint{1.464138in}{-1.875087in}}{\pgfqpoint{1.459747in}{-1.885686in}}{\pgfqpoint{1.459747in}{-1.896736in}}%
\pgfpathcurveto{\pgfqpoint{1.459747in}{-1.907786in}}{\pgfqpoint{1.464138in}{-1.918385in}}{\pgfqpoint{1.471951in}{-1.926199in}}%
\pgfpathcurveto{\pgfqpoint{1.479765in}{-1.934013in}}{\pgfqpoint{1.490364in}{-1.938403in}}{\pgfqpoint{1.501414in}{-1.938403in}}%
\pgfpathclose%
\pgfusepath{stroke,fill}%
\end{pgfscope}%
\begin{pgfscope}%
\pgfpathrectangle{\pgfqpoint{0.800000in}{0.528000in}}{\pgfqpoint{4.960000in}{3.696000in}} %
\pgfusepath{clip}%
\pgfsetbuttcap%
\pgfsetroundjoin%
\definecolor{currentfill}{rgb}{0.000000,0.500000,0.000000}%
\pgfsetfillcolor{currentfill}%
\pgfsetlinewidth{1.003750pt}%
\definecolor{currentstroke}{rgb}{0.000000,0.500000,0.000000}%
\pgfsetstrokecolor{currentstroke}%
\pgfsetdash{}{0pt}%
\pgfpathmoveto{\pgfqpoint{1.551515in}{-2.297595in}}%
\pgfpathcurveto{\pgfqpoint{1.562565in}{-2.297595in}}{\pgfqpoint{1.573164in}{-2.293205in}}{\pgfqpoint{1.580978in}{-2.285391in}}%
\pgfpathcurveto{\pgfqpoint{1.588792in}{-2.277577in}}{\pgfqpoint{1.593182in}{-2.266978in}}{\pgfqpoint{1.593182in}{-2.255928in}}%
\pgfpathcurveto{\pgfqpoint{1.593182in}{-2.244878in}}{\pgfqpoint{1.588792in}{-2.234279in}}{\pgfqpoint{1.580978in}{-2.226466in}}%
\pgfpathcurveto{\pgfqpoint{1.573164in}{-2.218652in}}{\pgfqpoint{1.562565in}{-2.214262in}}{\pgfqpoint{1.551515in}{-2.214262in}}%
\pgfpathcurveto{\pgfqpoint{1.540465in}{-2.214262in}}{\pgfqpoint{1.529866in}{-2.218652in}}{\pgfqpoint{1.522052in}{-2.226466in}}%
\pgfpathcurveto{\pgfqpoint{1.514239in}{-2.234279in}}{\pgfqpoint{1.509848in}{-2.244878in}}{\pgfqpoint{1.509848in}{-2.255928in}}%
\pgfpathcurveto{\pgfqpoint{1.509848in}{-2.266978in}}{\pgfqpoint{1.514239in}{-2.277577in}}{\pgfqpoint{1.522052in}{-2.285391in}}%
\pgfpathcurveto{\pgfqpoint{1.529866in}{-2.293205in}}{\pgfqpoint{1.540465in}{-2.297595in}}{\pgfqpoint{1.551515in}{-2.297595in}}%
\pgfpathclose%
\pgfusepath{stroke,fill}%
\end{pgfscope}%
\begin{pgfscope}%
\pgfpathrectangle{\pgfqpoint{0.800000in}{0.528000in}}{\pgfqpoint{4.960000in}{3.696000in}} %
\pgfusepath{clip}%
\pgfsetbuttcap%
\pgfsetroundjoin%
\definecolor{currentfill}{rgb}{0.000000,0.500000,0.000000}%
\pgfsetfillcolor{currentfill}%
\pgfsetlinewidth{1.003750pt}%
\definecolor{currentstroke}{rgb}{0.000000,0.500000,0.000000}%
\pgfsetstrokecolor{currentstroke}%
\pgfsetdash{}{0pt}%
\pgfpathmoveto{\pgfqpoint{1.601616in}{-3.654665in}}%
\pgfpathcurveto{\pgfqpoint{1.612666in}{-3.654665in}}{\pgfqpoint{1.623265in}{-3.650275in}}{\pgfqpoint{1.631079in}{-3.642461in}}%
\pgfpathcurveto{\pgfqpoint{1.638893in}{-3.634648in}}{\pgfqpoint{1.643283in}{-3.624049in}}{\pgfqpoint{1.643283in}{-3.612998in}}%
\pgfpathcurveto{\pgfqpoint{1.643283in}{-3.601948in}}{\pgfqpoint{1.638893in}{-3.591349in}}{\pgfqpoint{1.631079in}{-3.583536in}}%
\pgfpathcurveto{\pgfqpoint{1.623265in}{-3.575722in}}{\pgfqpoint{1.612666in}{-3.571332in}}{\pgfqpoint{1.601616in}{-3.571332in}}%
\pgfpathcurveto{\pgfqpoint{1.590566in}{-3.571332in}}{\pgfqpoint{1.579967in}{-3.575722in}}{\pgfqpoint{1.572153in}{-3.583536in}}%
\pgfpathcurveto{\pgfqpoint{1.564340in}{-3.591349in}}{\pgfqpoint{1.559949in}{-3.601948in}}{\pgfqpoint{1.559949in}{-3.612998in}}%
\pgfpathcurveto{\pgfqpoint{1.559949in}{-3.624049in}}{\pgfqpoint{1.564340in}{-3.634648in}}{\pgfqpoint{1.572153in}{-3.642461in}}%
\pgfpathcurveto{\pgfqpoint{1.579967in}{-3.650275in}}{\pgfqpoint{1.590566in}{-3.654665in}}{\pgfqpoint{1.601616in}{-3.654665in}}%
\pgfpathclose%
\pgfusepath{stroke,fill}%
\end{pgfscope}%
\begin{pgfscope}%
\pgfpathrectangle{\pgfqpoint{0.800000in}{0.528000in}}{\pgfqpoint{4.960000in}{3.696000in}} %
\pgfusepath{clip}%
\pgfsetbuttcap%
\pgfsetroundjoin%
\definecolor{currentfill}{rgb}{0.000000,0.500000,0.000000}%
\pgfsetfillcolor{currentfill}%
\pgfsetlinewidth{1.003750pt}%
\definecolor{currentstroke}{rgb}{0.000000,0.500000,0.000000}%
\pgfsetstrokecolor{currentstroke}%
\pgfsetdash{}{0pt}%
\pgfpathmoveto{\pgfqpoint{1.651717in}{0.080954in}}%
\pgfpathcurveto{\pgfqpoint{1.662767in}{0.080954in}}{\pgfqpoint{1.673366in}{0.085344in}}{\pgfqpoint{1.681180in}{0.093158in}}%
\pgfpathcurveto{\pgfqpoint{1.688994in}{0.100971in}}{\pgfqpoint{1.693384in}{0.111570in}}{\pgfqpoint{1.693384in}{0.122620in}}%
\pgfpathcurveto{\pgfqpoint{1.693384in}{0.133671in}}{\pgfqpoint{1.688994in}{0.144270in}}{\pgfqpoint{1.681180in}{0.152083in}}%
\pgfpathcurveto{\pgfqpoint{1.673366in}{0.159897in}}{\pgfqpoint{1.662767in}{0.164287in}}{\pgfqpoint{1.651717in}{0.164287in}}%
\pgfpathcurveto{\pgfqpoint{1.640667in}{0.164287in}}{\pgfqpoint{1.630068in}{0.159897in}}{\pgfqpoint{1.622254in}{0.152083in}}%
\pgfpathcurveto{\pgfqpoint{1.614441in}{0.144270in}}{\pgfqpoint{1.610051in}{0.133671in}}{\pgfqpoint{1.610051in}{0.122620in}}%
\pgfpathcurveto{\pgfqpoint{1.610051in}{0.111570in}}{\pgfqpoint{1.614441in}{0.100971in}}{\pgfqpoint{1.622254in}{0.093158in}}%
\pgfpathcurveto{\pgfqpoint{1.630068in}{0.085344in}}{\pgfqpoint{1.640667in}{0.080954in}}{\pgfqpoint{1.651717in}{0.080954in}}%
\pgfpathclose%
\pgfusepath{stroke,fill}%
\end{pgfscope}%
\begin{pgfscope}%
\pgfpathrectangle{\pgfqpoint{0.800000in}{0.528000in}}{\pgfqpoint{4.960000in}{3.696000in}} %
\pgfusepath{clip}%
\pgfsetbuttcap%
\pgfsetroundjoin%
\definecolor{currentfill}{rgb}{0.000000,0.500000,0.000000}%
\pgfsetfillcolor{currentfill}%
\pgfsetlinewidth{1.003750pt}%
\definecolor{currentstroke}{rgb}{0.000000,0.500000,0.000000}%
\pgfsetstrokecolor{currentstroke}%
\pgfsetdash{}{0pt}%
\pgfpathmoveto{\pgfqpoint{1.701818in}{0.105126in}}%
\pgfpathcurveto{\pgfqpoint{1.712868in}{0.105126in}}{\pgfqpoint{1.723467in}{0.109517in}}{\pgfqpoint{1.731281in}{0.117330in}}%
\pgfpathcurveto{\pgfqpoint{1.739095in}{0.125144in}}{\pgfqpoint{1.743485in}{0.135743in}}{\pgfqpoint{1.743485in}{0.146793in}}%
\pgfpathcurveto{\pgfqpoint{1.743485in}{0.157843in}}{\pgfqpoint{1.739095in}{0.168442in}}{\pgfqpoint{1.731281in}{0.176256in}}%
\pgfpathcurveto{\pgfqpoint{1.723467in}{0.184069in}}{\pgfqpoint{1.712868in}{0.188460in}}{\pgfqpoint{1.701818in}{0.188460in}}%
\pgfpathcurveto{\pgfqpoint{1.690768in}{0.188460in}}{\pgfqpoint{1.680169in}{0.184069in}}{\pgfqpoint{1.672355in}{0.176256in}}%
\pgfpathcurveto{\pgfqpoint{1.664542in}{0.168442in}}{\pgfqpoint{1.660152in}{0.157843in}}{\pgfqpoint{1.660152in}{0.146793in}}%
\pgfpathcurveto{\pgfqpoint{1.660152in}{0.135743in}}{\pgfqpoint{1.664542in}{0.125144in}}{\pgfqpoint{1.672355in}{0.117330in}}%
\pgfpathcurveto{\pgfqpoint{1.680169in}{0.109517in}}{\pgfqpoint{1.690768in}{0.105126in}}{\pgfqpoint{1.701818in}{0.105126in}}%
\pgfpathclose%
\pgfusepath{stroke,fill}%
\end{pgfscope}%
\begin{pgfscope}%
\pgfpathrectangle{\pgfqpoint{0.800000in}{0.528000in}}{\pgfqpoint{4.960000in}{3.696000in}} %
\pgfusepath{clip}%
\pgfsetbuttcap%
\pgfsetroundjoin%
\definecolor{currentfill}{rgb}{0.000000,0.500000,0.000000}%
\pgfsetfillcolor{currentfill}%
\pgfsetlinewidth{1.003750pt}%
\definecolor{currentstroke}{rgb}{0.000000,0.500000,0.000000}%
\pgfsetstrokecolor{currentstroke}%
\pgfsetdash{}{0pt}%
\pgfpathmoveto{\pgfqpoint{1.751919in}{-3.143259in}}%
\pgfpathcurveto{\pgfqpoint{1.762969in}{-3.143259in}}{\pgfqpoint{1.773568in}{-3.138869in}}{\pgfqpoint{1.781382in}{-3.131055in}}%
\pgfpathcurveto{\pgfqpoint{1.789196in}{-3.123242in}}{\pgfqpoint{1.793586in}{-3.112643in}}{\pgfqpoint{1.793586in}{-3.101592in}}%
\pgfpathcurveto{\pgfqpoint{1.793586in}{-3.090542in}}{\pgfqpoint{1.789196in}{-3.079943in}}{\pgfqpoint{1.781382in}{-3.072130in}}%
\pgfpathcurveto{\pgfqpoint{1.773568in}{-3.064316in}}{\pgfqpoint{1.762969in}{-3.059926in}}{\pgfqpoint{1.751919in}{-3.059926in}}%
\pgfpathcurveto{\pgfqpoint{1.740869in}{-3.059926in}}{\pgfqpoint{1.730270in}{-3.064316in}}{\pgfqpoint{1.722456in}{-3.072130in}}%
\pgfpathcurveto{\pgfqpoint{1.714643in}{-3.079943in}}{\pgfqpoint{1.710253in}{-3.090542in}}{\pgfqpoint{1.710253in}{-3.101592in}}%
\pgfpathcurveto{\pgfqpoint{1.710253in}{-3.112643in}}{\pgfqpoint{1.714643in}{-3.123242in}}{\pgfqpoint{1.722456in}{-3.131055in}}%
\pgfpathcurveto{\pgfqpoint{1.730270in}{-3.138869in}}{\pgfqpoint{1.740869in}{-3.143259in}}{\pgfqpoint{1.751919in}{-3.143259in}}%
\pgfpathclose%
\pgfusepath{stroke,fill}%
\end{pgfscope}%
\begin{pgfscope}%
\pgfpathrectangle{\pgfqpoint{0.800000in}{0.528000in}}{\pgfqpoint{4.960000in}{3.696000in}} %
\pgfusepath{clip}%
\pgfsetbuttcap%
\pgfsetroundjoin%
\definecolor{currentfill}{rgb}{0.000000,0.500000,0.000000}%
\pgfsetfillcolor{currentfill}%
\pgfsetlinewidth{1.003750pt}%
\definecolor{currentstroke}{rgb}{0.000000,0.500000,0.000000}%
\pgfsetstrokecolor{currentstroke}%
\pgfsetdash{}{0pt}%
\pgfpathmoveto{\pgfqpoint{1.802020in}{0.631821in}}%
\pgfpathcurveto{\pgfqpoint{1.813070in}{0.631821in}}{\pgfqpoint{1.823669in}{0.636211in}}{\pgfqpoint{1.831483in}{0.644025in}}%
\pgfpathcurveto{\pgfqpoint{1.839297in}{0.651838in}}{\pgfqpoint{1.843687in}{0.662437in}}{\pgfqpoint{1.843687in}{0.673487in}}%
\pgfpathcurveto{\pgfqpoint{1.843687in}{0.684537in}}{\pgfqpoint{1.839297in}{0.695136in}}{\pgfqpoint{1.831483in}{0.702950in}}%
\pgfpathcurveto{\pgfqpoint{1.823669in}{0.710764in}}{\pgfqpoint{1.813070in}{0.715154in}}{\pgfqpoint{1.802020in}{0.715154in}}%
\pgfpathcurveto{\pgfqpoint{1.790970in}{0.715154in}}{\pgfqpoint{1.780371in}{0.710764in}}{\pgfqpoint{1.772557in}{0.702950in}}%
\pgfpathcurveto{\pgfqpoint{1.764744in}{0.695136in}}{\pgfqpoint{1.760354in}{0.684537in}}{\pgfqpoint{1.760354in}{0.673487in}}%
\pgfpathcurveto{\pgfqpoint{1.760354in}{0.662437in}}{\pgfqpoint{1.764744in}{0.651838in}}{\pgfqpoint{1.772557in}{0.644025in}}%
\pgfpathcurveto{\pgfqpoint{1.780371in}{0.636211in}}{\pgfqpoint{1.790970in}{0.631821in}}{\pgfqpoint{1.802020in}{0.631821in}}%
\pgfpathclose%
\pgfusepath{stroke,fill}%
\end{pgfscope}%
\begin{pgfscope}%
\pgfpathrectangle{\pgfqpoint{0.800000in}{0.528000in}}{\pgfqpoint{4.960000in}{3.696000in}} %
\pgfusepath{clip}%
\pgfsetbuttcap%
\pgfsetroundjoin%
\definecolor{currentfill}{rgb}{0.000000,0.500000,0.000000}%
\pgfsetfillcolor{currentfill}%
\pgfsetlinewidth{1.003750pt}%
\definecolor{currentstroke}{rgb}{0.000000,0.500000,0.000000}%
\pgfsetstrokecolor{currentstroke}%
\pgfsetdash{}{0pt}%
\pgfpathmoveto{\pgfqpoint{1.852121in}{-1.499330in}}%
\pgfpathcurveto{\pgfqpoint{1.863171in}{-1.499330in}}{\pgfqpoint{1.873770in}{-1.494940in}}{\pgfqpoint{1.881584in}{-1.487127in}}%
\pgfpathcurveto{\pgfqpoint{1.889398in}{-1.479313in}}{\pgfqpoint{1.893788in}{-1.468714in}}{\pgfqpoint{1.893788in}{-1.457664in}}%
\pgfpathcurveto{\pgfqpoint{1.893788in}{-1.446614in}}{\pgfqpoint{1.889398in}{-1.436015in}}{\pgfqpoint{1.881584in}{-1.428201in}}%
\pgfpathcurveto{\pgfqpoint{1.873770in}{-1.420387in}}{\pgfqpoint{1.863171in}{-1.415997in}}{\pgfqpoint{1.852121in}{-1.415997in}}%
\pgfpathcurveto{\pgfqpoint{1.841071in}{-1.415997in}}{\pgfqpoint{1.830472in}{-1.420387in}}{\pgfqpoint{1.822658in}{-1.428201in}}%
\pgfpathcurveto{\pgfqpoint{1.814845in}{-1.436015in}}{\pgfqpoint{1.810455in}{-1.446614in}}{\pgfqpoint{1.810455in}{-1.457664in}}%
\pgfpathcurveto{\pgfqpoint{1.810455in}{-1.468714in}}{\pgfqpoint{1.814845in}{-1.479313in}}{\pgfqpoint{1.822658in}{-1.487127in}}%
\pgfpathcurveto{\pgfqpoint{1.830472in}{-1.494940in}}{\pgfqpoint{1.841071in}{-1.499330in}}{\pgfqpoint{1.852121in}{-1.499330in}}%
\pgfpathclose%
\pgfusepath{stroke,fill}%
\end{pgfscope}%
\begin{pgfscope}%
\pgfpathrectangle{\pgfqpoint{0.800000in}{0.528000in}}{\pgfqpoint{4.960000in}{3.696000in}} %
\pgfusepath{clip}%
\pgfsetbuttcap%
\pgfsetroundjoin%
\definecolor{currentfill}{rgb}{0.000000,0.500000,0.000000}%
\pgfsetfillcolor{currentfill}%
\pgfsetlinewidth{1.003750pt}%
\definecolor{currentstroke}{rgb}{0.000000,0.500000,0.000000}%
\pgfsetstrokecolor{currentstroke}%
\pgfsetdash{}{0pt}%
\pgfpathmoveto{\pgfqpoint{1.902222in}{-2.161634in}}%
\pgfpathcurveto{\pgfqpoint{1.913272in}{-2.161634in}}{\pgfqpoint{1.923871in}{-2.157244in}}{\pgfqpoint{1.931685in}{-2.149430in}}%
\pgfpathcurveto{\pgfqpoint{1.939499in}{-2.141617in}}{\pgfqpoint{1.943889in}{-2.131018in}}{\pgfqpoint{1.943889in}{-2.119968in}}%
\pgfpathcurveto{\pgfqpoint{1.943889in}{-2.108918in}}{\pgfqpoint{1.939499in}{-2.098319in}}{\pgfqpoint{1.931685in}{-2.090505in}}%
\pgfpathcurveto{\pgfqpoint{1.923871in}{-2.082691in}}{\pgfqpoint{1.913272in}{-2.078301in}}{\pgfqpoint{1.902222in}{-2.078301in}}%
\pgfpathcurveto{\pgfqpoint{1.891172in}{-2.078301in}}{\pgfqpoint{1.880573in}{-2.082691in}}{\pgfqpoint{1.872759in}{-2.090505in}}%
\pgfpathcurveto{\pgfqpoint{1.864946in}{-2.098319in}}{\pgfqpoint{1.860556in}{-2.108918in}}{\pgfqpoint{1.860556in}{-2.119968in}}%
\pgfpathcurveto{\pgfqpoint{1.860556in}{-2.131018in}}{\pgfqpoint{1.864946in}{-2.141617in}}{\pgfqpoint{1.872759in}{-2.149430in}}%
\pgfpathcurveto{\pgfqpoint{1.880573in}{-2.157244in}}{\pgfqpoint{1.891172in}{-2.161634in}}{\pgfqpoint{1.902222in}{-2.161634in}}%
\pgfpathclose%
\pgfusepath{stroke,fill}%
\end{pgfscope}%
\begin{pgfscope}%
\pgfpathrectangle{\pgfqpoint{0.800000in}{0.528000in}}{\pgfqpoint{4.960000in}{3.696000in}} %
\pgfusepath{clip}%
\pgfsetbuttcap%
\pgfsetroundjoin%
\definecolor{currentfill}{rgb}{0.000000,0.500000,0.000000}%
\pgfsetfillcolor{currentfill}%
\pgfsetlinewidth{1.003750pt}%
\definecolor{currentstroke}{rgb}{0.000000,0.500000,0.000000}%
\pgfsetstrokecolor{currentstroke}%
\pgfsetdash{}{0pt}%
\pgfpathmoveto{\pgfqpoint{1.952323in}{-2.331141in}}%
\pgfpathcurveto{\pgfqpoint{1.963373in}{-2.331141in}}{\pgfqpoint{1.973972in}{-2.326751in}}{\pgfqpoint{1.981786in}{-2.318937in}}%
\pgfpathcurveto{\pgfqpoint{1.989600in}{-2.311124in}}{\pgfqpoint{1.993990in}{-2.300525in}}{\pgfqpoint{1.993990in}{-2.289475in}}%
\pgfpathcurveto{\pgfqpoint{1.993990in}{-2.278425in}}{\pgfqpoint{1.989600in}{-2.267825in}}{\pgfqpoint{1.981786in}{-2.260012in}}%
\pgfpathcurveto{\pgfqpoint{1.973972in}{-2.252198in}}{\pgfqpoint{1.963373in}{-2.247808in}}{\pgfqpoint{1.952323in}{-2.247808in}}%
\pgfpathcurveto{\pgfqpoint{1.941273in}{-2.247808in}}{\pgfqpoint{1.930674in}{-2.252198in}}{\pgfqpoint{1.922860in}{-2.260012in}}%
\pgfpathcurveto{\pgfqpoint{1.915047in}{-2.267825in}}{\pgfqpoint{1.910657in}{-2.278425in}}{\pgfqpoint{1.910657in}{-2.289475in}}%
\pgfpathcurveto{\pgfqpoint{1.910657in}{-2.300525in}}{\pgfqpoint{1.915047in}{-2.311124in}}{\pgfqpoint{1.922860in}{-2.318937in}}%
\pgfpathcurveto{\pgfqpoint{1.930674in}{-2.326751in}}{\pgfqpoint{1.941273in}{-2.331141in}}{\pgfqpoint{1.952323in}{-2.331141in}}%
\pgfpathclose%
\pgfusepath{stroke,fill}%
\end{pgfscope}%
\begin{pgfscope}%
\pgfpathrectangle{\pgfqpoint{0.800000in}{0.528000in}}{\pgfqpoint{4.960000in}{3.696000in}} %
\pgfusepath{clip}%
\pgfsetbuttcap%
\pgfsetroundjoin%
\definecolor{currentfill}{rgb}{0.000000,0.500000,0.000000}%
\pgfsetfillcolor{currentfill}%
\pgfsetlinewidth{1.003750pt}%
\definecolor{currentstroke}{rgb}{0.000000,0.500000,0.000000}%
\pgfsetstrokecolor{currentstroke}%
\pgfsetdash{}{0pt}%
\pgfpathmoveto{\pgfqpoint{2.002424in}{-1.311050in}}%
\pgfpathcurveto{\pgfqpoint{2.013474in}{-1.311050in}}{\pgfqpoint{2.024073in}{-1.306660in}}{\pgfqpoint{2.031887in}{-1.298846in}}%
\pgfpathcurveto{\pgfqpoint{2.039701in}{-1.291033in}}{\pgfqpoint{2.044091in}{-1.280433in}}{\pgfqpoint{2.044091in}{-1.269383in}}%
\pgfpathcurveto{\pgfqpoint{2.044091in}{-1.258333in}}{\pgfqpoint{2.039701in}{-1.247734in}}{\pgfqpoint{2.031887in}{-1.239921in}}%
\pgfpathcurveto{\pgfqpoint{2.024073in}{-1.232107in}}{\pgfqpoint{2.013474in}{-1.227717in}}{\pgfqpoint{2.002424in}{-1.227717in}}%
\pgfpathcurveto{\pgfqpoint{1.991374in}{-1.227717in}}{\pgfqpoint{1.980775in}{-1.232107in}}{\pgfqpoint{1.972961in}{-1.239921in}}%
\pgfpathcurveto{\pgfqpoint{1.965148in}{-1.247734in}}{\pgfqpoint{1.960758in}{-1.258333in}}{\pgfqpoint{1.960758in}{-1.269383in}}%
\pgfpathcurveto{\pgfqpoint{1.960758in}{-1.280433in}}{\pgfqpoint{1.965148in}{-1.291033in}}{\pgfqpoint{1.972961in}{-1.298846in}}%
\pgfpathcurveto{\pgfqpoint{1.980775in}{-1.306660in}}{\pgfqpoint{1.991374in}{-1.311050in}}{\pgfqpoint{2.002424in}{-1.311050in}}%
\pgfpathclose%
\pgfusepath{stroke,fill}%
\end{pgfscope}%
\begin{pgfscope}%
\pgfpathrectangle{\pgfqpoint{0.800000in}{0.528000in}}{\pgfqpoint{4.960000in}{3.696000in}} %
\pgfusepath{clip}%
\pgfsetbuttcap%
\pgfsetroundjoin%
\definecolor{currentfill}{rgb}{0.000000,0.500000,0.000000}%
\pgfsetfillcolor{currentfill}%
\pgfsetlinewidth{1.003750pt}%
\definecolor{currentstroke}{rgb}{0.000000,0.500000,0.000000}%
\pgfsetstrokecolor{currentstroke}%
\pgfsetdash{}{0pt}%
\pgfpathmoveto{\pgfqpoint{2.052525in}{-0.684279in}}%
\pgfpathcurveto{\pgfqpoint{2.063575in}{-0.684279in}}{\pgfqpoint{2.074174in}{-0.679889in}}{\pgfqpoint{2.081988in}{-0.672076in}}%
\pgfpathcurveto{\pgfqpoint{2.089802in}{-0.664262in}}{\pgfqpoint{2.094192in}{-0.653663in}}{\pgfqpoint{2.094192in}{-0.642613in}}%
\pgfpathcurveto{\pgfqpoint{2.094192in}{-0.631563in}}{\pgfqpoint{2.089802in}{-0.620964in}}{\pgfqpoint{2.081988in}{-0.613150in}}%
\pgfpathcurveto{\pgfqpoint{2.074174in}{-0.605336in}}{\pgfqpoint{2.063575in}{-0.600946in}}{\pgfqpoint{2.052525in}{-0.600946in}}%
\pgfpathcurveto{\pgfqpoint{2.041475in}{-0.600946in}}{\pgfqpoint{2.030876in}{-0.605336in}}{\pgfqpoint{2.023062in}{-0.613150in}}%
\pgfpathcurveto{\pgfqpoint{2.015249in}{-0.620964in}}{\pgfqpoint{2.010859in}{-0.631563in}}{\pgfqpoint{2.010859in}{-0.642613in}}%
\pgfpathcurveto{\pgfqpoint{2.010859in}{-0.653663in}}{\pgfqpoint{2.015249in}{-0.664262in}}{\pgfqpoint{2.023062in}{-0.672076in}}%
\pgfpathcurveto{\pgfqpoint{2.030876in}{-0.679889in}}{\pgfqpoint{2.041475in}{-0.684279in}}{\pgfqpoint{2.052525in}{-0.684279in}}%
\pgfpathclose%
\pgfusepath{stroke,fill}%
\end{pgfscope}%
\begin{pgfscope}%
\pgfpathrectangle{\pgfqpoint{0.800000in}{0.528000in}}{\pgfqpoint{4.960000in}{3.696000in}} %
\pgfusepath{clip}%
\pgfsetbuttcap%
\pgfsetroundjoin%
\definecolor{currentfill}{rgb}{0.000000,0.500000,0.000000}%
\pgfsetfillcolor{currentfill}%
\pgfsetlinewidth{1.003750pt}%
\definecolor{currentstroke}{rgb}{0.000000,0.500000,0.000000}%
\pgfsetstrokecolor{currentstroke}%
\pgfsetdash{}{0pt}%
\pgfpathmoveto{\pgfqpoint{2.102626in}{0.814808in}}%
\pgfpathcurveto{\pgfqpoint{2.113676in}{0.814808in}}{\pgfqpoint{2.124275in}{0.819199in}}{\pgfqpoint{2.132089in}{0.827012in}}%
\pgfpathcurveto{\pgfqpoint{2.139903in}{0.834826in}}{\pgfqpoint{2.144293in}{0.845425in}}{\pgfqpoint{2.144293in}{0.856475in}}%
\pgfpathcurveto{\pgfqpoint{2.144293in}{0.867525in}}{\pgfqpoint{2.139903in}{0.878124in}}{\pgfqpoint{2.132089in}{0.885938in}}%
\pgfpathcurveto{\pgfqpoint{2.124275in}{0.893751in}}{\pgfqpoint{2.113676in}{0.898142in}}{\pgfqpoint{2.102626in}{0.898142in}}%
\pgfpathcurveto{\pgfqpoint{2.091576in}{0.898142in}}{\pgfqpoint{2.080977in}{0.893751in}}{\pgfqpoint{2.073163in}{0.885938in}}%
\pgfpathcurveto{\pgfqpoint{2.065350in}{0.878124in}}{\pgfqpoint{2.060960in}{0.867525in}}{\pgfqpoint{2.060960in}{0.856475in}}%
\pgfpathcurveto{\pgfqpoint{2.060960in}{0.845425in}}{\pgfqpoint{2.065350in}{0.834826in}}{\pgfqpoint{2.073163in}{0.827012in}}%
\pgfpathcurveto{\pgfqpoint{2.080977in}{0.819199in}}{\pgfqpoint{2.091576in}{0.814808in}}{\pgfqpoint{2.102626in}{0.814808in}}%
\pgfpathclose%
\pgfusepath{stroke,fill}%
\end{pgfscope}%
\begin{pgfscope}%
\pgfpathrectangle{\pgfqpoint{0.800000in}{0.528000in}}{\pgfqpoint{4.960000in}{3.696000in}} %
\pgfusepath{clip}%
\pgfsetbuttcap%
\pgfsetroundjoin%
\definecolor{currentfill}{rgb}{0.000000,0.500000,0.000000}%
\pgfsetfillcolor{currentfill}%
\pgfsetlinewidth{1.003750pt}%
\definecolor{currentstroke}{rgb}{0.000000,0.500000,0.000000}%
\pgfsetstrokecolor{currentstroke}%
\pgfsetdash{}{0pt}%
\pgfpathmoveto{\pgfqpoint{2.152727in}{-1.218463in}}%
\pgfpathcurveto{\pgfqpoint{2.163777in}{-1.218463in}}{\pgfqpoint{2.174376in}{-1.214073in}}{\pgfqpoint{2.182190in}{-1.206259in}}%
\pgfpathcurveto{\pgfqpoint{2.190004in}{-1.198445in}}{\pgfqpoint{2.194394in}{-1.187846in}}{\pgfqpoint{2.194394in}{-1.176796in}}%
\pgfpathcurveto{\pgfqpoint{2.194394in}{-1.165746in}}{\pgfqpoint{2.190004in}{-1.155147in}}{\pgfqpoint{2.182190in}{-1.147333in}}%
\pgfpathcurveto{\pgfqpoint{2.174376in}{-1.139520in}}{\pgfqpoint{2.163777in}{-1.135130in}}{\pgfqpoint{2.152727in}{-1.135130in}}%
\pgfpathcurveto{\pgfqpoint{2.141677in}{-1.135130in}}{\pgfqpoint{2.131078in}{-1.139520in}}{\pgfqpoint{2.123264in}{-1.147333in}}%
\pgfpathcurveto{\pgfqpoint{2.115451in}{-1.155147in}}{\pgfqpoint{2.111061in}{-1.165746in}}{\pgfqpoint{2.111061in}{-1.176796in}}%
\pgfpathcurveto{\pgfqpoint{2.111061in}{-1.187846in}}{\pgfqpoint{2.115451in}{-1.198445in}}{\pgfqpoint{2.123264in}{-1.206259in}}%
\pgfpathcurveto{\pgfqpoint{2.131078in}{-1.214073in}}{\pgfqpoint{2.141677in}{-1.218463in}}{\pgfqpoint{2.152727in}{-1.218463in}}%
\pgfpathclose%
\pgfusepath{stroke,fill}%
\end{pgfscope}%
\begin{pgfscope}%
\pgfpathrectangle{\pgfqpoint{0.800000in}{0.528000in}}{\pgfqpoint{4.960000in}{3.696000in}} %
\pgfusepath{clip}%
\pgfsetbuttcap%
\pgfsetroundjoin%
\definecolor{currentfill}{rgb}{0.000000,0.500000,0.000000}%
\pgfsetfillcolor{currentfill}%
\pgfsetlinewidth{1.003750pt}%
\definecolor{currentstroke}{rgb}{0.000000,0.500000,0.000000}%
\pgfsetstrokecolor{currentstroke}%
\pgfsetdash{}{0pt}%
\pgfpathmoveto{\pgfqpoint{2.202828in}{0.536603in}}%
\pgfpathcurveto{\pgfqpoint{2.213878in}{0.536603in}}{\pgfqpoint{2.224477in}{0.540993in}}{\pgfqpoint{2.232291in}{0.548807in}}%
\pgfpathcurveto{\pgfqpoint{2.240105in}{0.556620in}}{\pgfqpoint{2.244495in}{0.567219in}}{\pgfqpoint{2.244495in}{0.578269in}}%
\pgfpathcurveto{\pgfqpoint{2.244495in}{0.589320in}}{\pgfqpoint{2.240105in}{0.599919in}}{\pgfqpoint{2.232291in}{0.607732in}}%
\pgfpathcurveto{\pgfqpoint{2.224477in}{0.615546in}}{\pgfqpoint{2.213878in}{0.619936in}}{\pgfqpoint{2.202828in}{0.619936in}}%
\pgfpathcurveto{\pgfqpoint{2.191778in}{0.619936in}}{\pgfqpoint{2.181179in}{0.615546in}}{\pgfqpoint{2.173366in}{0.607732in}}%
\pgfpathcurveto{\pgfqpoint{2.165552in}{0.599919in}}{\pgfqpoint{2.161162in}{0.589320in}}{\pgfqpoint{2.161162in}{0.578269in}}%
\pgfpathcurveto{\pgfqpoint{2.161162in}{0.567219in}}{\pgfqpoint{2.165552in}{0.556620in}}{\pgfqpoint{2.173366in}{0.548807in}}%
\pgfpathcurveto{\pgfqpoint{2.181179in}{0.540993in}}{\pgfqpoint{2.191778in}{0.536603in}}{\pgfqpoint{2.202828in}{0.536603in}}%
\pgfpathclose%
\pgfusepath{stroke,fill}%
\end{pgfscope}%
\begin{pgfscope}%
\pgfpathrectangle{\pgfqpoint{0.800000in}{0.528000in}}{\pgfqpoint{4.960000in}{3.696000in}} %
\pgfusepath{clip}%
\pgfsetbuttcap%
\pgfsetroundjoin%
\definecolor{currentfill}{rgb}{0.000000,0.500000,0.000000}%
\pgfsetfillcolor{currentfill}%
\pgfsetlinewidth{1.003750pt}%
\definecolor{currentstroke}{rgb}{0.000000,0.500000,0.000000}%
\pgfsetstrokecolor{currentstroke}%
\pgfsetdash{}{0pt}%
\pgfpathmoveto{\pgfqpoint{2.252929in}{-1.034681in}}%
\pgfpathcurveto{\pgfqpoint{2.263979in}{-1.034681in}}{\pgfqpoint{2.274578in}{-1.030291in}}{\pgfqpoint{2.282392in}{-1.022477in}}%
\pgfpathcurveto{\pgfqpoint{2.290206in}{-1.014663in}}{\pgfqpoint{2.294596in}{-1.004064in}}{\pgfqpoint{2.294596in}{-0.993014in}}%
\pgfpathcurveto{\pgfqpoint{2.294596in}{-0.981964in}}{\pgfqpoint{2.290206in}{-0.971365in}}{\pgfqpoint{2.282392in}{-0.963551in}}%
\pgfpathcurveto{\pgfqpoint{2.274578in}{-0.955738in}}{\pgfqpoint{2.263979in}{-0.951348in}}{\pgfqpoint{2.252929in}{-0.951348in}}%
\pgfpathcurveto{\pgfqpoint{2.241879in}{-0.951348in}}{\pgfqpoint{2.231280in}{-0.955738in}}{\pgfqpoint{2.223467in}{-0.963551in}}%
\pgfpathcurveto{\pgfqpoint{2.215653in}{-0.971365in}}{\pgfqpoint{2.211263in}{-0.981964in}}{\pgfqpoint{2.211263in}{-0.993014in}}%
\pgfpathcurveto{\pgfqpoint{2.211263in}{-1.004064in}}{\pgfqpoint{2.215653in}{-1.014663in}}{\pgfqpoint{2.223467in}{-1.022477in}}%
\pgfpathcurveto{\pgfqpoint{2.231280in}{-1.030291in}}{\pgfqpoint{2.241879in}{-1.034681in}}{\pgfqpoint{2.252929in}{-1.034681in}}%
\pgfpathclose%
\pgfusepath{stroke,fill}%
\end{pgfscope}%
\begin{pgfscope}%
\pgfpathrectangle{\pgfqpoint{0.800000in}{0.528000in}}{\pgfqpoint{4.960000in}{3.696000in}} %
\pgfusepath{clip}%
\pgfsetbuttcap%
\pgfsetroundjoin%
\definecolor{currentfill}{rgb}{0.000000,0.500000,0.000000}%
\pgfsetfillcolor{currentfill}%
\pgfsetlinewidth{1.003750pt}%
\definecolor{currentstroke}{rgb}{0.000000,0.500000,0.000000}%
\pgfsetstrokecolor{currentstroke}%
\pgfsetdash{}{0pt}%
\pgfpathmoveto{\pgfqpoint{2.303030in}{0.535905in}}%
\pgfpathcurveto{\pgfqpoint{2.314080in}{0.535905in}}{\pgfqpoint{2.324679in}{0.540295in}}{\pgfqpoint{2.332493in}{0.548109in}}%
\pgfpathcurveto{\pgfqpoint{2.340307in}{0.555922in}}{\pgfqpoint{2.344697in}{0.566521in}}{\pgfqpoint{2.344697in}{0.577571in}}%
\pgfpathcurveto{\pgfqpoint{2.344697in}{0.588621in}}{\pgfqpoint{2.340307in}{0.599220in}}{\pgfqpoint{2.332493in}{0.607034in}}%
\pgfpathcurveto{\pgfqpoint{2.324679in}{0.614848in}}{\pgfqpoint{2.314080in}{0.619238in}}{\pgfqpoint{2.303030in}{0.619238in}}%
\pgfpathcurveto{\pgfqpoint{2.291980in}{0.619238in}}{\pgfqpoint{2.281381in}{0.614848in}}{\pgfqpoint{2.273568in}{0.607034in}}%
\pgfpathcurveto{\pgfqpoint{2.265754in}{0.599220in}}{\pgfqpoint{2.261364in}{0.588621in}}{\pgfqpoint{2.261364in}{0.577571in}}%
\pgfpathcurveto{\pgfqpoint{2.261364in}{0.566521in}}{\pgfqpoint{2.265754in}{0.555922in}}{\pgfqpoint{2.273568in}{0.548109in}}%
\pgfpathcurveto{\pgfqpoint{2.281381in}{0.540295in}}{\pgfqpoint{2.291980in}{0.535905in}}{\pgfqpoint{2.303030in}{0.535905in}}%
\pgfpathclose%
\pgfusepath{stroke,fill}%
\end{pgfscope}%
\begin{pgfscope}%
\pgfpathrectangle{\pgfqpoint{0.800000in}{0.528000in}}{\pgfqpoint{4.960000in}{3.696000in}} %
\pgfusepath{clip}%
\pgfsetbuttcap%
\pgfsetroundjoin%
\definecolor{currentfill}{rgb}{0.000000,0.500000,0.000000}%
\pgfsetfillcolor{currentfill}%
\pgfsetlinewidth{1.003750pt}%
\definecolor{currentstroke}{rgb}{0.000000,0.500000,0.000000}%
\pgfsetstrokecolor{currentstroke}%
\pgfsetdash{}{0pt}%
\pgfpathmoveto{\pgfqpoint{2.353131in}{-2.529487in}}%
\pgfpathcurveto{\pgfqpoint{2.364181in}{-2.529487in}}{\pgfqpoint{2.374780in}{-2.525097in}}{\pgfqpoint{2.382594in}{-2.517283in}}%
\pgfpathcurveto{\pgfqpoint{2.390408in}{-2.509470in}}{\pgfqpoint{2.394798in}{-2.498871in}}{\pgfqpoint{2.394798in}{-2.487821in}}%
\pgfpathcurveto{\pgfqpoint{2.394798in}{-2.476771in}}{\pgfqpoint{2.390408in}{-2.466171in}}{\pgfqpoint{2.382594in}{-2.458358in}}%
\pgfpathcurveto{\pgfqpoint{2.374780in}{-2.450544in}}{\pgfqpoint{2.364181in}{-2.446154in}}{\pgfqpoint{2.353131in}{-2.446154in}}%
\pgfpathcurveto{\pgfqpoint{2.342081in}{-2.446154in}}{\pgfqpoint{2.331482in}{-2.450544in}}{\pgfqpoint{2.323669in}{-2.458358in}}%
\pgfpathcurveto{\pgfqpoint{2.315855in}{-2.466171in}}{\pgfqpoint{2.311465in}{-2.476771in}}{\pgfqpoint{2.311465in}{-2.487821in}}%
\pgfpathcurveto{\pgfqpoint{2.311465in}{-2.498871in}}{\pgfqpoint{2.315855in}{-2.509470in}}{\pgfqpoint{2.323669in}{-2.517283in}}%
\pgfpathcurveto{\pgfqpoint{2.331482in}{-2.525097in}}{\pgfqpoint{2.342081in}{-2.529487in}}{\pgfqpoint{2.353131in}{-2.529487in}}%
\pgfpathclose%
\pgfusepath{stroke,fill}%
\end{pgfscope}%
\begin{pgfscope}%
\pgfpathrectangle{\pgfqpoint{0.800000in}{0.528000in}}{\pgfqpoint{4.960000in}{3.696000in}} %
\pgfusepath{clip}%
\pgfsetbuttcap%
\pgfsetroundjoin%
\definecolor{currentfill}{rgb}{0.000000,0.500000,0.000000}%
\pgfsetfillcolor{currentfill}%
\pgfsetlinewidth{1.003750pt}%
\definecolor{currentstroke}{rgb}{0.000000,0.500000,0.000000}%
\pgfsetstrokecolor{currentstroke}%
\pgfsetdash{}{0pt}%
\pgfpathmoveto{\pgfqpoint{2.403232in}{-0.708868in}}%
\pgfpathcurveto{\pgfqpoint{2.414282in}{-0.708868in}}{\pgfqpoint{2.424881in}{-0.704478in}}{\pgfqpoint{2.432695in}{-0.696665in}}%
\pgfpathcurveto{\pgfqpoint{2.440509in}{-0.688851in}}{\pgfqpoint{2.444899in}{-0.678252in}}{\pgfqpoint{2.444899in}{-0.667202in}}%
\pgfpathcurveto{\pgfqpoint{2.444899in}{-0.656152in}}{\pgfqpoint{2.440509in}{-0.645553in}}{\pgfqpoint{2.432695in}{-0.637739in}}%
\pgfpathcurveto{\pgfqpoint{2.424881in}{-0.629925in}}{\pgfqpoint{2.414282in}{-0.625535in}}{\pgfqpoint{2.403232in}{-0.625535in}}%
\pgfpathcurveto{\pgfqpoint{2.392182in}{-0.625535in}}{\pgfqpoint{2.381583in}{-0.629925in}}{\pgfqpoint{2.373770in}{-0.637739in}}%
\pgfpathcurveto{\pgfqpoint{2.365956in}{-0.645553in}}{\pgfqpoint{2.361566in}{-0.656152in}}{\pgfqpoint{2.361566in}{-0.667202in}}%
\pgfpathcurveto{\pgfqpoint{2.361566in}{-0.678252in}}{\pgfqpoint{2.365956in}{-0.688851in}}{\pgfqpoint{2.373770in}{-0.696665in}}%
\pgfpathcurveto{\pgfqpoint{2.381583in}{-0.704478in}}{\pgfqpoint{2.392182in}{-0.708868in}}{\pgfqpoint{2.403232in}{-0.708868in}}%
\pgfpathclose%
\pgfusepath{stroke,fill}%
\end{pgfscope}%
\begin{pgfscope}%
\pgfpathrectangle{\pgfqpoint{0.800000in}{0.528000in}}{\pgfqpoint{4.960000in}{3.696000in}} %
\pgfusepath{clip}%
\pgfsetbuttcap%
\pgfsetroundjoin%
\definecolor{currentfill}{rgb}{0.000000,0.500000,0.000000}%
\pgfsetfillcolor{currentfill}%
\pgfsetlinewidth{1.003750pt}%
\definecolor{currentstroke}{rgb}{0.000000,0.500000,0.000000}%
\pgfsetstrokecolor{currentstroke}%
\pgfsetdash{}{0pt}%
\pgfpathmoveto{\pgfqpoint{2.453333in}{-0.228484in}}%
\pgfpathcurveto{\pgfqpoint{2.464383in}{-0.228484in}}{\pgfqpoint{2.474982in}{-0.224094in}}{\pgfqpoint{2.482796in}{-0.216280in}}%
\pgfpathcurveto{\pgfqpoint{2.490610in}{-0.208466in}}{\pgfqpoint{2.495000in}{-0.197867in}}{\pgfqpoint{2.495000in}{-0.186817in}}%
\pgfpathcurveto{\pgfqpoint{2.495000in}{-0.175767in}}{\pgfqpoint{2.490610in}{-0.165168in}}{\pgfqpoint{2.482796in}{-0.157354in}}%
\pgfpathcurveto{\pgfqpoint{2.474982in}{-0.149541in}}{\pgfqpoint{2.464383in}{-0.145151in}}{\pgfqpoint{2.453333in}{-0.145151in}}%
\pgfpathcurveto{\pgfqpoint{2.442283in}{-0.145151in}}{\pgfqpoint{2.431684in}{-0.149541in}}{\pgfqpoint{2.423871in}{-0.157354in}}%
\pgfpathcurveto{\pgfqpoint{2.416057in}{-0.165168in}}{\pgfqpoint{2.411667in}{-0.175767in}}{\pgfqpoint{2.411667in}{-0.186817in}}%
\pgfpathcurveto{\pgfqpoint{2.411667in}{-0.197867in}}{\pgfqpoint{2.416057in}{-0.208466in}}{\pgfqpoint{2.423871in}{-0.216280in}}%
\pgfpathcurveto{\pgfqpoint{2.431684in}{-0.224094in}}{\pgfqpoint{2.442283in}{-0.228484in}}{\pgfqpoint{2.453333in}{-0.228484in}}%
\pgfpathclose%
\pgfusepath{stroke,fill}%
\end{pgfscope}%
\begin{pgfscope}%
\pgfpathrectangle{\pgfqpoint{0.800000in}{0.528000in}}{\pgfqpoint{4.960000in}{3.696000in}} %
\pgfusepath{clip}%
\pgfsetbuttcap%
\pgfsetroundjoin%
\definecolor{currentfill}{rgb}{0.000000,0.500000,0.000000}%
\pgfsetfillcolor{currentfill}%
\pgfsetlinewidth{1.003750pt}%
\definecolor{currentstroke}{rgb}{0.000000,0.500000,0.000000}%
\pgfsetstrokecolor{currentstroke}%
\pgfsetdash{}{0pt}%
\pgfpathmoveto{\pgfqpoint{2.503434in}{-0.829136in}}%
\pgfpathcurveto{\pgfqpoint{2.514484in}{-0.829136in}}{\pgfqpoint{2.525084in}{-0.824746in}}{\pgfqpoint{2.532897in}{-0.816932in}}%
\pgfpathcurveto{\pgfqpoint{2.540711in}{-0.809119in}}{\pgfqpoint{2.545101in}{-0.798520in}}{\pgfqpoint{2.545101in}{-0.787469in}}%
\pgfpathcurveto{\pgfqpoint{2.545101in}{-0.776419in}}{\pgfqpoint{2.540711in}{-0.765820in}}{\pgfqpoint{2.532897in}{-0.758007in}}%
\pgfpathcurveto{\pgfqpoint{2.525084in}{-0.750193in}}{\pgfqpoint{2.514484in}{-0.745803in}}{\pgfqpoint{2.503434in}{-0.745803in}}%
\pgfpathcurveto{\pgfqpoint{2.492384in}{-0.745803in}}{\pgfqpoint{2.481785in}{-0.750193in}}{\pgfqpoint{2.473972in}{-0.758007in}}%
\pgfpathcurveto{\pgfqpoint{2.466158in}{-0.765820in}}{\pgfqpoint{2.461768in}{-0.776419in}}{\pgfqpoint{2.461768in}{-0.787469in}}%
\pgfpathcurveto{\pgfqpoint{2.461768in}{-0.798520in}}{\pgfqpoint{2.466158in}{-0.809119in}}{\pgfqpoint{2.473972in}{-0.816932in}}%
\pgfpathcurveto{\pgfqpoint{2.481785in}{-0.824746in}}{\pgfqpoint{2.492384in}{-0.829136in}}{\pgfqpoint{2.503434in}{-0.829136in}}%
\pgfpathclose%
\pgfusepath{stroke,fill}%
\end{pgfscope}%
\begin{pgfscope}%
\pgfpathrectangle{\pgfqpoint{0.800000in}{0.528000in}}{\pgfqpoint{4.960000in}{3.696000in}} %
\pgfusepath{clip}%
\pgfsetbuttcap%
\pgfsetroundjoin%
\definecolor{currentfill}{rgb}{0.000000,0.500000,0.000000}%
\pgfsetfillcolor{currentfill}%
\pgfsetlinewidth{1.003750pt}%
\definecolor{currentstroke}{rgb}{0.000000,0.500000,0.000000}%
\pgfsetstrokecolor{currentstroke}%
\pgfsetdash{}{0pt}%
\pgfpathmoveto{\pgfqpoint{2.553535in}{-0.687790in}}%
\pgfpathcurveto{\pgfqpoint{2.564585in}{-0.687790in}}{\pgfqpoint{2.575185in}{-0.683400in}}{\pgfqpoint{2.582998in}{-0.675586in}}%
\pgfpathcurveto{\pgfqpoint{2.590812in}{-0.667772in}}{\pgfqpoint{2.595202in}{-0.657173in}}{\pgfqpoint{2.595202in}{-0.646123in}}%
\pgfpathcurveto{\pgfqpoint{2.595202in}{-0.635073in}}{\pgfqpoint{2.590812in}{-0.624474in}}{\pgfqpoint{2.582998in}{-0.616660in}}%
\pgfpathcurveto{\pgfqpoint{2.575185in}{-0.608847in}}{\pgfqpoint{2.564585in}{-0.604457in}}{\pgfqpoint{2.553535in}{-0.604457in}}%
\pgfpathcurveto{\pgfqpoint{2.542485in}{-0.604457in}}{\pgfqpoint{2.531886in}{-0.608847in}}{\pgfqpoint{2.524073in}{-0.616660in}}%
\pgfpathcurveto{\pgfqpoint{2.516259in}{-0.624474in}}{\pgfqpoint{2.511869in}{-0.635073in}}{\pgfqpoint{2.511869in}{-0.646123in}}%
\pgfpathcurveto{\pgfqpoint{2.511869in}{-0.657173in}}{\pgfqpoint{2.516259in}{-0.667772in}}{\pgfqpoint{2.524073in}{-0.675586in}}%
\pgfpathcurveto{\pgfqpoint{2.531886in}{-0.683400in}}{\pgfqpoint{2.542485in}{-0.687790in}}{\pgfqpoint{2.553535in}{-0.687790in}}%
\pgfpathclose%
\pgfusepath{stroke,fill}%
\end{pgfscope}%
\begin{pgfscope}%
\pgfpathrectangle{\pgfqpoint{0.800000in}{0.528000in}}{\pgfqpoint{4.960000in}{3.696000in}} %
\pgfusepath{clip}%
\pgfsetbuttcap%
\pgfsetroundjoin%
\definecolor{currentfill}{rgb}{0.000000,0.500000,0.000000}%
\pgfsetfillcolor{currentfill}%
\pgfsetlinewidth{1.003750pt}%
\definecolor{currentstroke}{rgb}{0.000000,0.500000,0.000000}%
\pgfsetstrokecolor{currentstroke}%
\pgfsetdash{}{0pt}%
\pgfpathmoveto{\pgfqpoint{2.603636in}{1.149831in}}%
\pgfpathcurveto{\pgfqpoint{2.614686in}{1.149831in}}{\pgfqpoint{2.625286in}{1.154221in}}{\pgfqpoint{2.633099in}{1.162034in}}%
\pgfpathcurveto{\pgfqpoint{2.640913in}{1.169848in}}{\pgfqpoint{2.645303in}{1.180447in}}{\pgfqpoint{2.645303in}{1.191497in}}%
\pgfpathcurveto{\pgfqpoint{2.645303in}{1.202547in}}{\pgfqpoint{2.640913in}{1.213146in}}{\pgfqpoint{2.633099in}{1.220960in}}%
\pgfpathcurveto{\pgfqpoint{2.625286in}{1.228774in}}{\pgfqpoint{2.614686in}{1.233164in}}{\pgfqpoint{2.603636in}{1.233164in}}%
\pgfpathcurveto{\pgfqpoint{2.592586in}{1.233164in}}{\pgfqpoint{2.581987in}{1.228774in}}{\pgfqpoint{2.574174in}{1.220960in}}%
\pgfpathcurveto{\pgfqpoint{2.566360in}{1.213146in}}{\pgfqpoint{2.561970in}{1.202547in}}{\pgfqpoint{2.561970in}{1.191497in}}%
\pgfpathcurveto{\pgfqpoint{2.561970in}{1.180447in}}{\pgfqpoint{2.566360in}{1.169848in}}{\pgfqpoint{2.574174in}{1.162034in}}%
\pgfpathcurveto{\pgfqpoint{2.581987in}{1.154221in}}{\pgfqpoint{2.592586in}{1.149831in}}{\pgfqpoint{2.603636in}{1.149831in}}%
\pgfpathclose%
\pgfusepath{stroke,fill}%
\end{pgfscope}%
\begin{pgfscope}%
\pgfpathrectangle{\pgfqpoint{0.800000in}{0.528000in}}{\pgfqpoint{4.960000in}{3.696000in}} %
\pgfusepath{clip}%
\pgfsetbuttcap%
\pgfsetroundjoin%
\definecolor{currentfill}{rgb}{0.000000,0.500000,0.000000}%
\pgfsetfillcolor{currentfill}%
\pgfsetlinewidth{1.003750pt}%
\definecolor{currentstroke}{rgb}{0.000000,0.500000,0.000000}%
\pgfsetstrokecolor{currentstroke}%
\pgfsetdash{}{0pt}%
\pgfpathmoveto{\pgfqpoint{2.653737in}{-1.444467in}}%
\pgfpathcurveto{\pgfqpoint{2.664788in}{-1.444467in}}{\pgfqpoint{2.675387in}{-1.440077in}}{\pgfqpoint{2.683200in}{-1.432263in}}%
\pgfpathcurveto{\pgfqpoint{2.691014in}{-1.424450in}}{\pgfqpoint{2.695404in}{-1.413851in}}{\pgfqpoint{2.695404in}{-1.402801in}}%
\pgfpathcurveto{\pgfqpoint{2.695404in}{-1.391751in}}{\pgfqpoint{2.691014in}{-1.381152in}}{\pgfqpoint{2.683200in}{-1.373338in}}%
\pgfpathcurveto{\pgfqpoint{2.675387in}{-1.365524in}}{\pgfqpoint{2.664788in}{-1.361134in}}{\pgfqpoint{2.653737in}{-1.361134in}}%
\pgfpathcurveto{\pgfqpoint{2.642687in}{-1.361134in}}{\pgfqpoint{2.632088in}{-1.365524in}}{\pgfqpoint{2.624275in}{-1.373338in}}%
\pgfpathcurveto{\pgfqpoint{2.616461in}{-1.381152in}}{\pgfqpoint{2.612071in}{-1.391751in}}{\pgfqpoint{2.612071in}{-1.402801in}}%
\pgfpathcurveto{\pgfqpoint{2.612071in}{-1.413851in}}{\pgfqpoint{2.616461in}{-1.424450in}}{\pgfqpoint{2.624275in}{-1.432263in}}%
\pgfpathcurveto{\pgfqpoint{2.632088in}{-1.440077in}}{\pgfqpoint{2.642687in}{-1.444467in}}{\pgfqpoint{2.653737in}{-1.444467in}}%
\pgfpathclose%
\pgfusepath{stroke,fill}%
\end{pgfscope}%
\begin{pgfscope}%
\pgfpathrectangle{\pgfqpoint{0.800000in}{0.528000in}}{\pgfqpoint{4.960000in}{3.696000in}} %
\pgfusepath{clip}%
\pgfsetbuttcap%
\pgfsetroundjoin%
\definecolor{currentfill}{rgb}{0.000000,0.500000,0.000000}%
\pgfsetfillcolor{currentfill}%
\pgfsetlinewidth{1.003750pt}%
\definecolor{currentstroke}{rgb}{0.000000,0.500000,0.000000}%
\pgfsetstrokecolor{currentstroke}%
\pgfsetdash{}{0pt}%
\pgfpathmoveto{\pgfqpoint{2.703838in}{-2.981878in}}%
\pgfpathcurveto{\pgfqpoint{2.714889in}{-2.981878in}}{\pgfqpoint{2.725488in}{-2.977487in}}{\pgfqpoint{2.733301in}{-2.969674in}}%
\pgfpathcurveto{\pgfqpoint{2.741115in}{-2.961860in}}{\pgfqpoint{2.745505in}{-2.951261in}}{\pgfqpoint{2.745505in}{-2.940211in}}%
\pgfpathcurveto{\pgfqpoint{2.745505in}{-2.929161in}}{\pgfqpoint{2.741115in}{-2.918562in}}{\pgfqpoint{2.733301in}{-2.910748in}}%
\pgfpathcurveto{\pgfqpoint{2.725488in}{-2.902935in}}{\pgfqpoint{2.714889in}{-2.898544in}}{\pgfqpoint{2.703838in}{-2.898544in}}%
\pgfpathcurveto{\pgfqpoint{2.692788in}{-2.898544in}}{\pgfqpoint{2.682189in}{-2.902935in}}{\pgfqpoint{2.674376in}{-2.910748in}}%
\pgfpathcurveto{\pgfqpoint{2.666562in}{-2.918562in}}{\pgfqpoint{2.662172in}{-2.929161in}}{\pgfqpoint{2.662172in}{-2.940211in}}%
\pgfpathcurveto{\pgfqpoint{2.662172in}{-2.951261in}}{\pgfqpoint{2.666562in}{-2.961860in}}{\pgfqpoint{2.674376in}{-2.969674in}}%
\pgfpathcurveto{\pgfqpoint{2.682189in}{-2.977487in}}{\pgfqpoint{2.692788in}{-2.981878in}}{\pgfqpoint{2.703838in}{-2.981878in}}%
\pgfpathclose%
\pgfusepath{stroke,fill}%
\end{pgfscope}%
\begin{pgfscope}%
\pgfpathrectangle{\pgfqpoint{0.800000in}{0.528000in}}{\pgfqpoint{4.960000in}{3.696000in}} %
\pgfusepath{clip}%
\pgfsetbuttcap%
\pgfsetroundjoin%
\definecolor{currentfill}{rgb}{0.000000,0.500000,0.000000}%
\pgfsetfillcolor{currentfill}%
\pgfsetlinewidth{1.003750pt}%
\definecolor{currentstroke}{rgb}{0.000000,0.500000,0.000000}%
\pgfsetstrokecolor{currentstroke}%
\pgfsetdash{}{0pt}%
\pgfpathmoveto{\pgfqpoint{2.753939in}{-0.686957in}}%
\pgfpathcurveto{\pgfqpoint{2.764990in}{-0.686957in}}{\pgfqpoint{2.775589in}{-0.682567in}}{\pgfqpoint{2.783402in}{-0.674753in}}%
\pgfpathcurveto{\pgfqpoint{2.791216in}{-0.666939in}}{\pgfqpoint{2.795606in}{-0.656340in}}{\pgfqpoint{2.795606in}{-0.645290in}}%
\pgfpathcurveto{\pgfqpoint{2.795606in}{-0.634240in}}{\pgfqpoint{2.791216in}{-0.623641in}}{\pgfqpoint{2.783402in}{-0.615827in}}%
\pgfpathcurveto{\pgfqpoint{2.775589in}{-0.608014in}}{\pgfqpoint{2.764990in}{-0.603624in}}{\pgfqpoint{2.753939in}{-0.603624in}}%
\pgfpathcurveto{\pgfqpoint{2.742889in}{-0.603624in}}{\pgfqpoint{2.732290in}{-0.608014in}}{\pgfqpoint{2.724477in}{-0.615827in}}%
\pgfpathcurveto{\pgfqpoint{2.716663in}{-0.623641in}}{\pgfqpoint{2.712273in}{-0.634240in}}{\pgfqpoint{2.712273in}{-0.645290in}}%
\pgfpathcurveto{\pgfqpoint{2.712273in}{-0.656340in}}{\pgfqpoint{2.716663in}{-0.666939in}}{\pgfqpoint{2.724477in}{-0.674753in}}%
\pgfpathcurveto{\pgfqpoint{2.732290in}{-0.682567in}}{\pgfqpoint{2.742889in}{-0.686957in}}{\pgfqpoint{2.753939in}{-0.686957in}}%
\pgfpathclose%
\pgfusepath{stroke,fill}%
\end{pgfscope}%
\begin{pgfscope}%
\pgfpathrectangle{\pgfqpoint{0.800000in}{0.528000in}}{\pgfqpoint{4.960000in}{3.696000in}} %
\pgfusepath{clip}%
\pgfsetbuttcap%
\pgfsetroundjoin%
\definecolor{currentfill}{rgb}{0.000000,0.500000,0.000000}%
\pgfsetfillcolor{currentfill}%
\pgfsetlinewidth{1.003750pt}%
\definecolor{currentstroke}{rgb}{0.000000,0.500000,0.000000}%
\pgfsetstrokecolor{currentstroke}%
\pgfsetdash{}{0pt}%
\pgfpathmoveto{\pgfqpoint{2.804040in}{1.378506in}}%
\pgfpathcurveto{\pgfqpoint{2.815091in}{1.378506in}}{\pgfqpoint{2.825690in}{1.382896in}}{\pgfqpoint{2.833503in}{1.390710in}}%
\pgfpathcurveto{\pgfqpoint{2.841317in}{1.398523in}}{\pgfqpoint{2.845707in}{1.409122in}}{\pgfqpoint{2.845707in}{1.420172in}}%
\pgfpathcurveto{\pgfqpoint{2.845707in}{1.431223in}}{\pgfqpoint{2.841317in}{1.441822in}}{\pgfqpoint{2.833503in}{1.449635in}}%
\pgfpathcurveto{\pgfqpoint{2.825690in}{1.457449in}}{\pgfqpoint{2.815091in}{1.461839in}}{\pgfqpoint{2.804040in}{1.461839in}}%
\pgfpathcurveto{\pgfqpoint{2.792990in}{1.461839in}}{\pgfqpoint{2.782391in}{1.457449in}}{\pgfqpoint{2.774578in}{1.449635in}}%
\pgfpathcurveto{\pgfqpoint{2.766764in}{1.441822in}}{\pgfqpoint{2.762374in}{1.431223in}}{\pgfqpoint{2.762374in}{1.420172in}}%
\pgfpathcurveto{\pgfqpoint{2.762374in}{1.409122in}}{\pgfqpoint{2.766764in}{1.398523in}}{\pgfqpoint{2.774578in}{1.390710in}}%
\pgfpathcurveto{\pgfqpoint{2.782391in}{1.382896in}}{\pgfqpoint{2.792990in}{1.378506in}}{\pgfqpoint{2.804040in}{1.378506in}}%
\pgfpathclose%
\pgfusepath{stroke,fill}%
\end{pgfscope}%
\begin{pgfscope}%
\pgfpathrectangle{\pgfqpoint{0.800000in}{0.528000in}}{\pgfqpoint{4.960000in}{3.696000in}} %
\pgfusepath{clip}%
\pgfsetbuttcap%
\pgfsetroundjoin%
\definecolor{currentfill}{rgb}{0.000000,0.500000,0.000000}%
\pgfsetfillcolor{currentfill}%
\pgfsetlinewidth{1.003750pt}%
\definecolor{currentstroke}{rgb}{0.000000,0.500000,0.000000}%
\pgfsetstrokecolor{currentstroke}%
\pgfsetdash{}{0pt}%
\pgfpathmoveto{\pgfqpoint{2.854141in}{-0.991782in}}%
\pgfpathcurveto{\pgfqpoint{2.865192in}{-0.991782in}}{\pgfqpoint{2.875791in}{-0.987392in}}{\pgfqpoint{2.883604in}{-0.979578in}}%
\pgfpathcurveto{\pgfqpoint{2.891418in}{-0.971765in}}{\pgfqpoint{2.895808in}{-0.961166in}}{\pgfqpoint{2.895808in}{-0.950115in}}%
\pgfpathcurveto{\pgfqpoint{2.895808in}{-0.939065in}}{\pgfqpoint{2.891418in}{-0.928466in}}{\pgfqpoint{2.883604in}{-0.920653in}}%
\pgfpathcurveto{\pgfqpoint{2.875791in}{-0.912839in}}{\pgfqpoint{2.865192in}{-0.908449in}}{\pgfqpoint{2.854141in}{-0.908449in}}%
\pgfpathcurveto{\pgfqpoint{2.843091in}{-0.908449in}}{\pgfqpoint{2.832492in}{-0.912839in}}{\pgfqpoint{2.824679in}{-0.920653in}}%
\pgfpathcurveto{\pgfqpoint{2.816865in}{-0.928466in}}{\pgfqpoint{2.812475in}{-0.939065in}}{\pgfqpoint{2.812475in}{-0.950115in}}%
\pgfpathcurveto{\pgfqpoint{2.812475in}{-0.961166in}}{\pgfqpoint{2.816865in}{-0.971765in}}{\pgfqpoint{2.824679in}{-0.979578in}}%
\pgfpathcurveto{\pgfqpoint{2.832492in}{-0.987392in}}{\pgfqpoint{2.843091in}{-0.991782in}}{\pgfqpoint{2.854141in}{-0.991782in}}%
\pgfpathclose%
\pgfusepath{stroke,fill}%
\end{pgfscope}%
\begin{pgfscope}%
\pgfpathrectangle{\pgfqpoint{0.800000in}{0.528000in}}{\pgfqpoint{4.960000in}{3.696000in}} %
\pgfusepath{clip}%
\pgfsetbuttcap%
\pgfsetroundjoin%
\definecolor{currentfill}{rgb}{0.000000,0.500000,0.000000}%
\pgfsetfillcolor{currentfill}%
\pgfsetlinewidth{1.003750pt}%
\definecolor{currentstroke}{rgb}{0.000000,0.500000,0.000000}%
\pgfsetstrokecolor{currentstroke}%
\pgfsetdash{}{0pt}%
\pgfpathmoveto{\pgfqpoint{2.904242in}{-1.786507in}}%
\pgfpathcurveto{\pgfqpoint{2.915293in}{-1.786507in}}{\pgfqpoint{2.925892in}{-1.782117in}}{\pgfqpoint{2.933705in}{-1.774303in}}%
\pgfpathcurveto{\pgfqpoint{2.941519in}{-1.766489in}}{\pgfqpoint{2.945909in}{-1.755890in}}{\pgfqpoint{2.945909in}{-1.744840in}}%
\pgfpathcurveto{\pgfqpoint{2.945909in}{-1.733790in}}{\pgfqpoint{2.941519in}{-1.723191in}}{\pgfqpoint{2.933705in}{-1.715377in}}%
\pgfpathcurveto{\pgfqpoint{2.925892in}{-1.707564in}}{\pgfqpoint{2.915293in}{-1.703174in}}{\pgfqpoint{2.904242in}{-1.703174in}}%
\pgfpathcurveto{\pgfqpoint{2.893192in}{-1.703174in}}{\pgfqpoint{2.882593in}{-1.707564in}}{\pgfqpoint{2.874780in}{-1.715377in}}%
\pgfpathcurveto{\pgfqpoint{2.866966in}{-1.723191in}}{\pgfqpoint{2.862576in}{-1.733790in}}{\pgfqpoint{2.862576in}{-1.744840in}}%
\pgfpathcurveto{\pgfqpoint{2.862576in}{-1.755890in}}{\pgfqpoint{2.866966in}{-1.766489in}}{\pgfqpoint{2.874780in}{-1.774303in}}%
\pgfpathcurveto{\pgfqpoint{2.882593in}{-1.782117in}}{\pgfqpoint{2.893192in}{-1.786507in}}{\pgfqpoint{2.904242in}{-1.786507in}}%
\pgfpathclose%
\pgfusepath{stroke,fill}%
\end{pgfscope}%
\begin{pgfscope}%
\pgfpathrectangle{\pgfqpoint{0.800000in}{0.528000in}}{\pgfqpoint{4.960000in}{3.696000in}} %
\pgfusepath{clip}%
\pgfsetbuttcap%
\pgfsetroundjoin%
\definecolor{currentfill}{rgb}{0.000000,0.500000,0.000000}%
\pgfsetfillcolor{currentfill}%
\pgfsetlinewidth{1.003750pt}%
\definecolor{currentstroke}{rgb}{0.000000,0.500000,0.000000}%
\pgfsetstrokecolor{currentstroke}%
\pgfsetdash{}{0pt}%
\pgfpathmoveto{\pgfqpoint{2.954343in}{1.508471in}}%
\pgfpathcurveto{\pgfqpoint{2.965394in}{1.508471in}}{\pgfqpoint{2.975993in}{1.512862in}}{\pgfqpoint{2.983806in}{1.520675in}}%
\pgfpathcurveto{\pgfqpoint{2.991620in}{1.528489in}}{\pgfqpoint{2.996010in}{1.539088in}}{\pgfqpoint{2.996010in}{1.550138in}}%
\pgfpathcurveto{\pgfqpoint{2.996010in}{1.561188in}}{\pgfqpoint{2.991620in}{1.571787in}}{\pgfqpoint{2.983806in}{1.579601in}}%
\pgfpathcurveto{\pgfqpoint{2.975993in}{1.587414in}}{\pgfqpoint{2.965394in}{1.591805in}}{\pgfqpoint{2.954343in}{1.591805in}}%
\pgfpathcurveto{\pgfqpoint{2.943293in}{1.591805in}}{\pgfqpoint{2.932694in}{1.587414in}}{\pgfqpoint{2.924881in}{1.579601in}}%
\pgfpathcurveto{\pgfqpoint{2.917067in}{1.571787in}}{\pgfqpoint{2.912677in}{1.561188in}}{\pgfqpoint{2.912677in}{1.550138in}}%
\pgfpathcurveto{\pgfqpoint{2.912677in}{1.539088in}}{\pgfqpoint{2.917067in}{1.528489in}}{\pgfqpoint{2.924881in}{1.520675in}}%
\pgfpathcurveto{\pgfqpoint{2.932694in}{1.512862in}}{\pgfqpoint{2.943293in}{1.508471in}}{\pgfqpoint{2.954343in}{1.508471in}}%
\pgfpathclose%
\pgfusepath{stroke,fill}%
\end{pgfscope}%
\begin{pgfscope}%
\pgfpathrectangle{\pgfqpoint{0.800000in}{0.528000in}}{\pgfqpoint{4.960000in}{3.696000in}} %
\pgfusepath{clip}%
\pgfsetbuttcap%
\pgfsetroundjoin%
\definecolor{currentfill}{rgb}{0.000000,0.500000,0.000000}%
\pgfsetfillcolor{currentfill}%
\pgfsetlinewidth{1.003750pt}%
\definecolor{currentstroke}{rgb}{0.000000,0.500000,0.000000}%
\pgfsetstrokecolor{currentstroke}%
\pgfsetdash{}{0pt}%
\pgfpathmoveto{\pgfqpoint{3.004444in}{1.486842in}}%
\pgfpathcurveto{\pgfqpoint{3.015495in}{1.486842in}}{\pgfqpoint{3.026094in}{1.491232in}}{\pgfqpoint{3.033907in}{1.499046in}}%
\pgfpathcurveto{\pgfqpoint{3.041721in}{1.506860in}}{\pgfqpoint{3.046111in}{1.517459in}}{\pgfqpoint{3.046111in}{1.528509in}}%
\pgfpathcurveto{\pgfqpoint{3.046111in}{1.539559in}}{\pgfqpoint{3.041721in}{1.550158in}}{\pgfqpoint{3.033907in}{1.557972in}}%
\pgfpathcurveto{\pgfqpoint{3.026094in}{1.565785in}}{\pgfqpoint{3.015495in}{1.570175in}}{\pgfqpoint{3.004444in}{1.570175in}}%
\pgfpathcurveto{\pgfqpoint{2.993394in}{1.570175in}}{\pgfqpoint{2.982795in}{1.565785in}}{\pgfqpoint{2.974982in}{1.557972in}}%
\pgfpathcurveto{\pgfqpoint{2.967168in}{1.550158in}}{\pgfqpoint{2.962778in}{1.539559in}}{\pgfqpoint{2.962778in}{1.528509in}}%
\pgfpathcurveto{\pgfqpoint{2.962778in}{1.517459in}}{\pgfqpoint{2.967168in}{1.506860in}}{\pgfqpoint{2.974982in}{1.499046in}}%
\pgfpathcurveto{\pgfqpoint{2.982795in}{1.491232in}}{\pgfqpoint{2.993394in}{1.486842in}}{\pgfqpoint{3.004444in}{1.486842in}}%
\pgfpathclose%
\pgfusepath{stroke,fill}%
\end{pgfscope}%
\begin{pgfscope}%
\pgfpathrectangle{\pgfqpoint{0.800000in}{0.528000in}}{\pgfqpoint{4.960000in}{3.696000in}} %
\pgfusepath{clip}%
\pgfsetbuttcap%
\pgfsetroundjoin%
\definecolor{currentfill}{rgb}{0.000000,0.500000,0.000000}%
\pgfsetfillcolor{currentfill}%
\pgfsetlinewidth{1.003750pt}%
\definecolor{currentstroke}{rgb}{0.000000,0.500000,0.000000}%
\pgfsetstrokecolor{currentstroke}%
\pgfsetdash{}{0pt}%
\pgfpathmoveto{\pgfqpoint{3.054545in}{-2.846312in}}%
\pgfpathcurveto{\pgfqpoint{3.065596in}{-2.846312in}}{\pgfqpoint{3.076195in}{-2.841922in}}{\pgfqpoint{3.084008in}{-2.834109in}}%
\pgfpathcurveto{\pgfqpoint{3.091822in}{-2.826295in}}{\pgfqpoint{3.096212in}{-2.815696in}}{\pgfqpoint{3.096212in}{-2.804646in}}%
\pgfpathcurveto{\pgfqpoint{3.096212in}{-2.793596in}}{\pgfqpoint{3.091822in}{-2.782997in}}{\pgfqpoint{3.084008in}{-2.775183in}}%
\pgfpathcurveto{\pgfqpoint{3.076195in}{-2.767369in}}{\pgfqpoint{3.065596in}{-2.762979in}}{\pgfqpoint{3.054545in}{-2.762979in}}%
\pgfpathcurveto{\pgfqpoint{3.043495in}{-2.762979in}}{\pgfqpoint{3.032896in}{-2.767369in}}{\pgfqpoint{3.025083in}{-2.775183in}}%
\pgfpathcurveto{\pgfqpoint{3.017269in}{-2.782997in}}{\pgfqpoint{3.012879in}{-2.793596in}}{\pgfqpoint{3.012879in}{-2.804646in}}%
\pgfpathcurveto{\pgfqpoint{3.012879in}{-2.815696in}}{\pgfqpoint{3.017269in}{-2.826295in}}{\pgfqpoint{3.025083in}{-2.834109in}}%
\pgfpathcurveto{\pgfqpoint{3.032896in}{-2.841922in}}{\pgfqpoint{3.043495in}{-2.846312in}}{\pgfqpoint{3.054545in}{-2.846312in}}%
\pgfpathclose%
\pgfusepath{stroke,fill}%
\end{pgfscope}%
\begin{pgfscope}%
\pgfpathrectangle{\pgfqpoint{0.800000in}{0.528000in}}{\pgfqpoint{4.960000in}{3.696000in}} %
\pgfusepath{clip}%
\pgfsetbuttcap%
\pgfsetroundjoin%
\definecolor{currentfill}{rgb}{0.000000,0.500000,0.000000}%
\pgfsetfillcolor{currentfill}%
\pgfsetlinewidth{1.003750pt}%
\definecolor{currentstroke}{rgb}{0.000000,0.500000,0.000000}%
\pgfsetstrokecolor{currentstroke}%
\pgfsetdash{}{0pt}%
\pgfpathmoveto{\pgfqpoint{3.104646in}{1.595799in}}%
\pgfpathcurveto{\pgfqpoint{3.115697in}{1.595799in}}{\pgfqpoint{3.126296in}{1.600189in}}{\pgfqpoint{3.134109in}{1.608002in}}%
\pgfpathcurveto{\pgfqpoint{3.141923in}{1.615816in}}{\pgfqpoint{3.146313in}{1.626415in}}{\pgfqpoint{3.146313in}{1.637465in}}%
\pgfpathcurveto{\pgfqpoint{3.146313in}{1.648515in}}{\pgfqpoint{3.141923in}{1.659114in}}{\pgfqpoint{3.134109in}{1.666928in}}%
\pgfpathcurveto{\pgfqpoint{3.126296in}{1.674742in}}{\pgfqpoint{3.115697in}{1.679132in}}{\pgfqpoint{3.104646in}{1.679132in}}%
\pgfpathcurveto{\pgfqpoint{3.093596in}{1.679132in}}{\pgfqpoint{3.082997in}{1.674742in}}{\pgfqpoint{3.075184in}{1.666928in}}%
\pgfpathcurveto{\pgfqpoint{3.067370in}{1.659114in}}{\pgfqpoint{3.062980in}{1.648515in}}{\pgfqpoint{3.062980in}{1.637465in}}%
\pgfpathcurveto{\pgfqpoint{3.062980in}{1.626415in}}{\pgfqpoint{3.067370in}{1.615816in}}{\pgfqpoint{3.075184in}{1.608002in}}%
\pgfpathcurveto{\pgfqpoint{3.082997in}{1.600189in}}{\pgfqpoint{3.093596in}{1.595799in}}{\pgfqpoint{3.104646in}{1.595799in}}%
\pgfpathclose%
\pgfusepath{stroke,fill}%
\end{pgfscope}%
\begin{pgfscope}%
\pgfpathrectangle{\pgfqpoint{0.800000in}{0.528000in}}{\pgfqpoint{4.960000in}{3.696000in}} %
\pgfusepath{clip}%
\pgfsetbuttcap%
\pgfsetroundjoin%
\definecolor{currentfill}{rgb}{0.000000,0.500000,0.000000}%
\pgfsetfillcolor{currentfill}%
\pgfsetlinewidth{1.003750pt}%
\definecolor{currentstroke}{rgb}{0.000000,0.500000,0.000000}%
\pgfsetstrokecolor{currentstroke}%
\pgfsetdash{}{0pt}%
\pgfpathmoveto{\pgfqpoint{3.154747in}{-0.943956in}}%
\pgfpathcurveto{\pgfqpoint{3.165798in}{-0.943956in}}{\pgfqpoint{3.176397in}{-0.939565in}}{\pgfqpoint{3.184210in}{-0.931752in}}%
\pgfpathcurveto{\pgfqpoint{3.192024in}{-0.923938in}}{\pgfqpoint{3.196414in}{-0.913339in}}{\pgfqpoint{3.196414in}{-0.902289in}}%
\pgfpathcurveto{\pgfqpoint{3.196414in}{-0.891239in}}{\pgfqpoint{3.192024in}{-0.880640in}}{\pgfqpoint{3.184210in}{-0.872826in}}%
\pgfpathcurveto{\pgfqpoint{3.176397in}{-0.865012in}}{\pgfqpoint{3.165798in}{-0.860622in}}{\pgfqpoint{3.154747in}{-0.860622in}}%
\pgfpathcurveto{\pgfqpoint{3.143697in}{-0.860622in}}{\pgfqpoint{3.133098in}{-0.865012in}}{\pgfqpoint{3.125285in}{-0.872826in}}%
\pgfpathcurveto{\pgfqpoint{3.117471in}{-0.880640in}}{\pgfqpoint{3.113081in}{-0.891239in}}{\pgfqpoint{3.113081in}{-0.902289in}}%
\pgfpathcurveto{\pgfqpoint{3.113081in}{-0.913339in}}{\pgfqpoint{3.117471in}{-0.923938in}}{\pgfqpoint{3.125285in}{-0.931752in}}%
\pgfpathcurveto{\pgfqpoint{3.133098in}{-0.939565in}}{\pgfqpoint{3.143697in}{-0.943956in}}{\pgfqpoint{3.154747in}{-0.943956in}}%
\pgfpathclose%
\pgfusepath{stroke,fill}%
\end{pgfscope}%
\begin{pgfscope}%
\pgfpathrectangle{\pgfqpoint{0.800000in}{0.528000in}}{\pgfqpoint{4.960000in}{3.696000in}} %
\pgfusepath{clip}%
\pgfsetbuttcap%
\pgfsetroundjoin%
\definecolor{currentfill}{rgb}{0.000000,0.500000,0.000000}%
\pgfsetfillcolor{currentfill}%
\pgfsetlinewidth{1.003750pt}%
\definecolor{currentstroke}{rgb}{0.000000,0.500000,0.000000}%
\pgfsetstrokecolor{currentstroke}%
\pgfsetdash{}{0pt}%
\pgfpathmoveto{\pgfqpoint{3.204848in}{-1.671124in}}%
\pgfpathcurveto{\pgfqpoint{3.215899in}{-1.671124in}}{\pgfqpoint{3.226498in}{-1.666734in}}{\pgfqpoint{3.234311in}{-1.658920in}}%
\pgfpathcurveto{\pgfqpoint{3.242125in}{-1.651107in}}{\pgfqpoint{3.246515in}{-1.640508in}}{\pgfqpoint{3.246515in}{-1.629458in}}%
\pgfpathcurveto{\pgfqpoint{3.246515in}{-1.618408in}}{\pgfqpoint{3.242125in}{-1.607809in}}{\pgfqpoint{3.234311in}{-1.599995in}}%
\pgfpathcurveto{\pgfqpoint{3.226498in}{-1.592181in}}{\pgfqpoint{3.215899in}{-1.587791in}}{\pgfqpoint{3.204848in}{-1.587791in}}%
\pgfpathcurveto{\pgfqpoint{3.193798in}{-1.587791in}}{\pgfqpoint{3.183199in}{-1.592181in}}{\pgfqpoint{3.175386in}{-1.599995in}}%
\pgfpathcurveto{\pgfqpoint{3.167572in}{-1.607809in}}{\pgfqpoint{3.163182in}{-1.618408in}}{\pgfqpoint{3.163182in}{-1.629458in}}%
\pgfpathcurveto{\pgfqpoint{3.163182in}{-1.640508in}}{\pgfqpoint{3.167572in}{-1.651107in}}{\pgfqpoint{3.175386in}{-1.658920in}}%
\pgfpathcurveto{\pgfqpoint{3.183199in}{-1.666734in}}{\pgfqpoint{3.193798in}{-1.671124in}}{\pgfqpoint{3.204848in}{-1.671124in}}%
\pgfpathclose%
\pgfusepath{stroke,fill}%
\end{pgfscope}%
\begin{pgfscope}%
\pgfpathrectangle{\pgfqpoint{0.800000in}{0.528000in}}{\pgfqpoint{4.960000in}{3.696000in}} %
\pgfusepath{clip}%
\pgfsetbuttcap%
\pgfsetroundjoin%
\definecolor{currentfill}{rgb}{0.000000,0.500000,0.000000}%
\pgfsetfillcolor{currentfill}%
\pgfsetlinewidth{1.003750pt}%
\definecolor{currentstroke}{rgb}{0.000000,0.500000,0.000000}%
\pgfsetstrokecolor{currentstroke}%
\pgfsetdash{}{0pt}%
\pgfpathmoveto{\pgfqpoint{3.254949in}{1.103529in}}%
\pgfpathcurveto{\pgfqpoint{3.266000in}{1.103529in}}{\pgfqpoint{3.276599in}{1.107920in}}{\pgfqpoint{3.284412in}{1.115733in}}%
\pgfpathcurveto{\pgfqpoint{3.292226in}{1.123547in}}{\pgfqpoint{3.296616in}{1.134146in}}{\pgfqpoint{3.296616in}{1.145196in}}%
\pgfpathcurveto{\pgfqpoint{3.296616in}{1.156246in}}{\pgfqpoint{3.292226in}{1.166845in}}{\pgfqpoint{3.284412in}{1.174659in}}%
\pgfpathcurveto{\pgfqpoint{3.276599in}{1.182472in}}{\pgfqpoint{3.266000in}{1.186863in}}{\pgfqpoint{3.254949in}{1.186863in}}%
\pgfpathcurveto{\pgfqpoint{3.243899in}{1.186863in}}{\pgfqpoint{3.233300in}{1.182472in}}{\pgfqpoint{3.225487in}{1.174659in}}%
\pgfpathcurveto{\pgfqpoint{3.217673in}{1.166845in}}{\pgfqpoint{3.213283in}{1.156246in}}{\pgfqpoint{3.213283in}{1.145196in}}%
\pgfpathcurveto{\pgfqpoint{3.213283in}{1.134146in}}{\pgfqpoint{3.217673in}{1.123547in}}{\pgfqpoint{3.225487in}{1.115733in}}%
\pgfpathcurveto{\pgfqpoint{3.233300in}{1.107920in}}{\pgfqpoint{3.243899in}{1.103529in}}{\pgfqpoint{3.254949in}{1.103529in}}%
\pgfpathclose%
\pgfusepath{stroke,fill}%
\end{pgfscope}%
\begin{pgfscope}%
\pgfpathrectangle{\pgfqpoint{0.800000in}{0.528000in}}{\pgfqpoint{4.960000in}{3.696000in}} %
\pgfusepath{clip}%
\pgfsetbuttcap%
\pgfsetroundjoin%
\definecolor{currentfill}{rgb}{0.000000,0.500000,0.000000}%
\pgfsetfillcolor{currentfill}%
\pgfsetlinewidth{1.003750pt}%
\definecolor{currentstroke}{rgb}{0.000000,0.500000,0.000000}%
\pgfsetstrokecolor{currentstroke}%
\pgfsetdash{}{0pt}%
\pgfpathmoveto{\pgfqpoint{3.305051in}{-0.956379in}}%
\pgfpathcurveto{\pgfqpoint{3.316101in}{-0.956379in}}{\pgfqpoint{3.326700in}{-0.951989in}}{\pgfqpoint{3.334513in}{-0.944175in}}%
\pgfpathcurveto{\pgfqpoint{3.342327in}{-0.936361in}}{\pgfqpoint{3.346717in}{-0.925762in}}{\pgfqpoint{3.346717in}{-0.914712in}}%
\pgfpathcurveto{\pgfqpoint{3.346717in}{-0.903662in}}{\pgfqpoint{3.342327in}{-0.893063in}}{\pgfqpoint{3.334513in}{-0.885249in}}%
\pgfpathcurveto{\pgfqpoint{3.326700in}{-0.877436in}}{\pgfqpoint{3.316101in}{-0.873046in}}{\pgfqpoint{3.305051in}{-0.873046in}}%
\pgfpathcurveto{\pgfqpoint{3.294000in}{-0.873046in}}{\pgfqpoint{3.283401in}{-0.877436in}}{\pgfqpoint{3.275588in}{-0.885249in}}%
\pgfpathcurveto{\pgfqpoint{3.267774in}{-0.893063in}}{\pgfqpoint{3.263384in}{-0.903662in}}{\pgfqpoint{3.263384in}{-0.914712in}}%
\pgfpathcurveto{\pgfqpoint{3.263384in}{-0.925762in}}{\pgfqpoint{3.267774in}{-0.936361in}}{\pgfqpoint{3.275588in}{-0.944175in}}%
\pgfpathcurveto{\pgfqpoint{3.283401in}{-0.951989in}}{\pgfqpoint{3.294000in}{-0.956379in}}{\pgfqpoint{3.305051in}{-0.956379in}}%
\pgfpathclose%
\pgfusepath{stroke,fill}%
\end{pgfscope}%
\begin{pgfscope}%
\pgfpathrectangle{\pgfqpoint{0.800000in}{0.528000in}}{\pgfqpoint{4.960000in}{3.696000in}} %
\pgfusepath{clip}%
\pgfsetbuttcap%
\pgfsetroundjoin%
\definecolor{currentfill}{rgb}{0.000000,0.500000,0.000000}%
\pgfsetfillcolor{currentfill}%
\pgfsetlinewidth{1.003750pt}%
\definecolor{currentstroke}{rgb}{0.000000,0.500000,0.000000}%
\pgfsetstrokecolor{currentstroke}%
\pgfsetdash{}{0pt}%
\pgfpathmoveto{\pgfqpoint{3.355152in}{0.766928in}}%
\pgfpathcurveto{\pgfqpoint{3.366202in}{0.766928in}}{\pgfqpoint{3.376801in}{0.771319in}}{\pgfqpoint{3.384614in}{0.779132in}}%
\pgfpathcurveto{\pgfqpoint{3.392428in}{0.786946in}}{\pgfqpoint{3.396818in}{0.797545in}}{\pgfqpoint{3.396818in}{0.808595in}}%
\pgfpathcurveto{\pgfqpoint{3.396818in}{0.819645in}}{\pgfqpoint{3.392428in}{0.830244in}}{\pgfqpoint{3.384614in}{0.838058in}}%
\pgfpathcurveto{\pgfqpoint{3.376801in}{0.845872in}}{\pgfqpoint{3.366202in}{0.850262in}}{\pgfqpoint{3.355152in}{0.850262in}}%
\pgfpathcurveto{\pgfqpoint{3.344101in}{0.850262in}}{\pgfqpoint{3.333502in}{0.845872in}}{\pgfqpoint{3.325689in}{0.838058in}}%
\pgfpathcurveto{\pgfqpoint{3.317875in}{0.830244in}}{\pgfqpoint{3.313485in}{0.819645in}}{\pgfqpoint{3.313485in}{0.808595in}}%
\pgfpathcurveto{\pgfqpoint{3.313485in}{0.797545in}}{\pgfqpoint{3.317875in}{0.786946in}}{\pgfqpoint{3.325689in}{0.779132in}}%
\pgfpathcurveto{\pgfqpoint{3.333502in}{0.771319in}}{\pgfqpoint{3.344101in}{0.766928in}}{\pgfqpoint{3.355152in}{0.766928in}}%
\pgfpathclose%
\pgfusepath{stroke,fill}%
\end{pgfscope}%
\begin{pgfscope}%
\pgfpathrectangle{\pgfqpoint{0.800000in}{0.528000in}}{\pgfqpoint{4.960000in}{3.696000in}} %
\pgfusepath{clip}%
\pgfsetbuttcap%
\pgfsetroundjoin%
\definecolor{currentfill}{rgb}{0.000000,0.500000,0.000000}%
\pgfsetfillcolor{currentfill}%
\pgfsetlinewidth{1.003750pt}%
\definecolor{currentstroke}{rgb}{0.000000,0.500000,0.000000}%
\pgfsetstrokecolor{currentstroke}%
\pgfsetdash{}{0pt}%
\pgfpathmoveto{\pgfqpoint{3.405253in}{1.733081in}}%
\pgfpathcurveto{\pgfqpoint{3.416303in}{1.733081in}}{\pgfqpoint{3.426902in}{1.737472in}}{\pgfqpoint{3.434715in}{1.745285in}}%
\pgfpathcurveto{\pgfqpoint{3.442529in}{1.753099in}}{\pgfqpoint{3.446919in}{1.763698in}}{\pgfqpoint{3.446919in}{1.774748in}}%
\pgfpathcurveto{\pgfqpoint{3.446919in}{1.785798in}}{\pgfqpoint{3.442529in}{1.796397in}}{\pgfqpoint{3.434715in}{1.804211in}}%
\pgfpathcurveto{\pgfqpoint{3.426902in}{1.812024in}}{\pgfqpoint{3.416303in}{1.816415in}}{\pgfqpoint{3.405253in}{1.816415in}}%
\pgfpathcurveto{\pgfqpoint{3.394202in}{1.816415in}}{\pgfqpoint{3.383603in}{1.812024in}}{\pgfqpoint{3.375790in}{1.804211in}}%
\pgfpathcurveto{\pgfqpoint{3.367976in}{1.796397in}}{\pgfqpoint{3.363586in}{1.785798in}}{\pgfqpoint{3.363586in}{1.774748in}}%
\pgfpathcurveto{\pgfqpoint{3.363586in}{1.763698in}}{\pgfqpoint{3.367976in}{1.753099in}}{\pgfqpoint{3.375790in}{1.745285in}}%
\pgfpathcurveto{\pgfqpoint{3.383603in}{1.737472in}}{\pgfqpoint{3.394202in}{1.733081in}}{\pgfqpoint{3.405253in}{1.733081in}}%
\pgfpathclose%
\pgfusepath{stroke,fill}%
\end{pgfscope}%
\begin{pgfscope}%
\pgfpathrectangle{\pgfqpoint{0.800000in}{0.528000in}}{\pgfqpoint{4.960000in}{3.696000in}} %
\pgfusepath{clip}%
\pgfsetbuttcap%
\pgfsetroundjoin%
\definecolor{currentfill}{rgb}{0.000000,0.500000,0.000000}%
\pgfsetfillcolor{currentfill}%
\pgfsetlinewidth{1.003750pt}%
\definecolor{currentstroke}{rgb}{0.000000,0.500000,0.000000}%
\pgfsetstrokecolor{currentstroke}%
\pgfsetdash{}{0pt}%
\pgfpathmoveto{\pgfqpoint{3.455354in}{-0.976987in}}%
\pgfpathcurveto{\pgfqpoint{3.466404in}{-0.976987in}}{\pgfqpoint{3.477003in}{-0.972597in}}{\pgfqpoint{3.484816in}{-0.964783in}}%
\pgfpathcurveto{\pgfqpoint{3.492630in}{-0.956969in}}{\pgfqpoint{3.497020in}{-0.946370in}}{\pgfqpoint{3.497020in}{-0.935320in}}%
\pgfpathcurveto{\pgfqpoint{3.497020in}{-0.924270in}}{\pgfqpoint{3.492630in}{-0.913671in}}{\pgfqpoint{3.484816in}{-0.905858in}}%
\pgfpathcurveto{\pgfqpoint{3.477003in}{-0.898044in}}{\pgfqpoint{3.466404in}{-0.893654in}}{\pgfqpoint{3.455354in}{-0.893654in}}%
\pgfpathcurveto{\pgfqpoint{3.444303in}{-0.893654in}}{\pgfqpoint{3.433704in}{-0.898044in}}{\pgfqpoint{3.425891in}{-0.905858in}}%
\pgfpathcurveto{\pgfqpoint{3.418077in}{-0.913671in}}{\pgfqpoint{3.413687in}{-0.924270in}}{\pgfqpoint{3.413687in}{-0.935320in}}%
\pgfpathcurveto{\pgfqpoint{3.413687in}{-0.946370in}}{\pgfqpoint{3.418077in}{-0.956969in}}{\pgfqpoint{3.425891in}{-0.964783in}}%
\pgfpathcurveto{\pgfqpoint{3.433704in}{-0.972597in}}{\pgfqpoint{3.444303in}{-0.976987in}}{\pgfqpoint{3.455354in}{-0.976987in}}%
\pgfpathclose%
\pgfusepath{stroke,fill}%
\end{pgfscope}%
\begin{pgfscope}%
\pgfpathrectangle{\pgfqpoint{0.800000in}{0.528000in}}{\pgfqpoint{4.960000in}{3.696000in}} %
\pgfusepath{clip}%
\pgfsetbuttcap%
\pgfsetroundjoin%
\definecolor{currentfill}{rgb}{0.000000,0.500000,0.000000}%
\pgfsetfillcolor{currentfill}%
\pgfsetlinewidth{1.003750pt}%
\definecolor{currentstroke}{rgb}{0.000000,0.500000,0.000000}%
\pgfsetstrokecolor{currentstroke}%
\pgfsetdash{}{0pt}%
\pgfpathmoveto{\pgfqpoint{3.505455in}{1.249754in}}%
\pgfpathcurveto{\pgfqpoint{3.516505in}{1.249754in}}{\pgfqpoint{3.527104in}{1.254144in}}{\pgfqpoint{3.534917in}{1.261958in}}%
\pgfpathcurveto{\pgfqpoint{3.542731in}{1.269771in}}{\pgfqpoint{3.547121in}{1.280370in}}{\pgfqpoint{3.547121in}{1.291420in}}%
\pgfpathcurveto{\pgfqpoint{3.547121in}{1.302471in}}{\pgfqpoint{3.542731in}{1.313070in}}{\pgfqpoint{3.534917in}{1.320883in}}%
\pgfpathcurveto{\pgfqpoint{3.527104in}{1.328697in}}{\pgfqpoint{3.516505in}{1.333087in}}{\pgfqpoint{3.505455in}{1.333087in}}%
\pgfpathcurveto{\pgfqpoint{3.494404in}{1.333087in}}{\pgfqpoint{3.483805in}{1.328697in}}{\pgfqpoint{3.475992in}{1.320883in}}%
\pgfpathcurveto{\pgfqpoint{3.468178in}{1.313070in}}{\pgfqpoint{3.463788in}{1.302471in}}{\pgfqpoint{3.463788in}{1.291420in}}%
\pgfpathcurveto{\pgfqpoint{3.463788in}{1.280370in}}{\pgfqpoint{3.468178in}{1.269771in}}{\pgfqpoint{3.475992in}{1.261958in}}%
\pgfpathcurveto{\pgfqpoint{3.483805in}{1.254144in}}{\pgfqpoint{3.494404in}{1.249754in}}{\pgfqpoint{3.505455in}{1.249754in}}%
\pgfpathclose%
\pgfusepath{stroke,fill}%
\end{pgfscope}%
\begin{pgfscope}%
\pgfpathrectangle{\pgfqpoint{0.800000in}{0.528000in}}{\pgfqpoint{4.960000in}{3.696000in}} %
\pgfusepath{clip}%
\pgfsetbuttcap%
\pgfsetroundjoin%
\definecolor{currentfill}{rgb}{0.000000,0.500000,0.000000}%
\pgfsetfillcolor{currentfill}%
\pgfsetlinewidth{1.003750pt}%
\definecolor{currentstroke}{rgb}{0.000000,0.500000,0.000000}%
\pgfsetstrokecolor{currentstroke}%
\pgfsetdash{}{0pt}%
\pgfpathmoveto{\pgfqpoint{3.555556in}{2.326571in}}%
\pgfpathcurveto{\pgfqpoint{3.566606in}{2.326571in}}{\pgfqpoint{3.577205in}{2.330962in}}{\pgfqpoint{3.585018in}{2.338775in}}%
\pgfpathcurveto{\pgfqpoint{3.592832in}{2.346589in}}{\pgfqpoint{3.597222in}{2.357188in}}{\pgfqpoint{3.597222in}{2.368238in}}%
\pgfpathcurveto{\pgfqpoint{3.597222in}{2.379288in}}{\pgfqpoint{3.592832in}{2.389887in}}{\pgfqpoint{3.585018in}{2.397701in}}%
\pgfpathcurveto{\pgfqpoint{3.577205in}{2.405514in}}{\pgfqpoint{3.566606in}{2.409905in}}{\pgfqpoint{3.555556in}{2.409905in}}%
\pgfpathcurveto{\pgfqpoint{3.544505in}{2.409905in}}{\pgfqpoint{3.533906in}{2.405514in}}{\pgfqpoint{3.526093in}{2.397701in}}%
\pgfpathcurveto{\pgfqpoint{3.518279in}{2.389887in}}{\pgfqpoint{3.513889in}{2.379288in}}{\pgfqpoint{3.513889in}{2.368238in}}%
\pgfpathcurveto{\pgfqpoint{3.513889in}{2.357188in}}{\pgfqpoint{3.518279in}{2.346589in}}{\pgfqpoint{3.526093in}{2.338775in}}%
\pgfpathcurveto{\pgfqpoint{3.533906in}{2.330962in}}{\pgfqpoint{3.544505in}{2.326571in}}{\pgfqpoint{3.555556in}{2.326571in}}%
\pgfpathclose%
\pgfusepath{stroke,fill}%
\end{pgfscope}%
\begin{pgfscope}%
\pgfpathrectangle{\pgfqpoint{0.800000in}{0.528000in}}{\pgfqpoint{4.960000in}{3.696000in}} %
\pgfusepath{clip}%
\pgfsetbuttcap%
\pgfsetroundjoin%
\definecolor{currentfill}{rgb}{0.000000,0.500000,0.000000}%
\pgfsetfillcolor{currentfill}%
\pgfsetlinewidth{1.003750pt}%
\definecolor{currentstroke}{rgb}{0.000000,0.500000,0.000000}%
\pgfsetstrokecolor{currentstroke}%
\pgfsetdash{}{0pt}%
\pgfpathmoveto{\pgfqpoint{3.605657in}{1.189003in}}%
\pgfpathcurveto{\pgfqpoint{3.616707in}{1.189003in}}{\pgfqpoint{3.627306in}{1.193393in}}{\pgfqpoint{3.635119in}{1.201207in}}%
\pgfpathcurveto{\pgfqpoint{3.642933in}{1.209021in}}{\pgfqpoint{3.647323in}{1.219620in}}{\pgfqpoint{3.647323in}{1.230670in}}%
\pgfpathcurveto{\pgfqpoint{3.647323in}{1.241720in}}{\pgfqpoint{3.642933in}{1.252319in}}{\pgfqpoint{3.635119in}{1.260133in}}%
\pgfpathcurveto{\pgfqpoint{3.627306in}{1.267946in}}{\pgfqpoint{3.616707in}{1.272336in}}{\pgfqpoint{3.605657in}{1.272336in}}%
\pgfpathcurveto{\pgfqpoint{3.594606in}{1.272336in}}{\pgfqpoint{3.584007in}{1.267946in}}{\pgfqpoint{3.576194in}{1.260133in}}%
\pgfpathcurveto{\pgfqpoint{3.568380in}{1.252319in}}{\pgfqpoint{3.563990in}{1.241720in}}{\pgfqpoint{3.563990in}{1.230670in}}%
\pgfpathcurveto{\pgfqpoint{3.563990in}{1.219620in}}{\pgfqpoint{3.568380in}{1.209021in}}{\pgfqpoint{3.576194in}{1.201207in}}%
\pgfpathcurveto{\pgfqpoint{3.584007in}{1.193393in}}{\pgfqpoint{3.594606in}{1.189003in}}{\pgfqpoint{3.605657in}{1.189003in}}%
\pgfpathclose%
\pgfusepath{stroke,fill}%
\end{pgfscope}%
\begin{pgfscope}%
\pgfpathrectangle{\pgfqpoint{0.800000in}{0.528000in}}{\pgfqpoint{4.960000in}{3.696000in}} %
\pgfusepath{clip}%
\pgfsetbuttcap%
\pgfsetroundjoin%
\definecolor{currentfill}{rgb}{0.000000,0.500000,0.000000}%
\pgfsetfillcolor{currentfill}%
\pgfsetlinewidth{1.003750pt}%
\definecolor{currentstroke}{rgb}{0.000000,0.500000,0.000000}%
\pgfsetstrokecolor{currentstroke}%
\pgfsetdash{}{0pt}%
\pgfpathmoveto{\pgfqpoint{3.655758in}{0.537667in}}%
\pgfpathcurveto{\pgfqpoint{3.666808in}{0.537667in}}{\pgfqpoint{3.677407in}{0.542057in}}{\pgfqpoint{3.685220in}{0.549871in}}%
\pgfpathcurveto{\pgfqpoint{3.693034in}{0.557684in}}{\pgfqpoint{3.697424in}{0.568283in}}{\pgfqpoint{3.697424in}{0.579333in}}%
\pgfpathcurveto{\pgfqpoint{3.697424in}{0.590384in}}{\pgfqpoint{3.693034in}{0.600983in}}{\pgfqpoint{3.685220in}{0.608796in}}%
\pgfpathcurveto{\pgfqpoint{3.677407in}{0.616610in}}{\pgfqpoint{3.666808in}{0.621000in}}{\pgfqpoint{3.655758in}{0.621000in}}%
\pgfpathcurveto{\pgfqpoint{3.644707in}{0.621000in}}{\pgfqpoint{3.634108in}{0.616610in}}{\pgfqpoint{3.626295in}{0.608796in}}%
\pgfpathcurveto{\pgfqpoint{3.618481in}{0.600983in}}{\pgfqpoint{3.614091in}{0.590384in}}{\pgfqpoint{3.614091in}{0.579333in}}%
\pgfpathcurveto{\pgfqpoint{3.614091in}{0.568283in}}{\pgfqpoint{3.618481in}{0.557684in}}{\pgfqpoint{3.626295in}{0.549871in}}%
\pgfpathcurveto{\pgfqpoint{3.634108in}{0.542057in}}{\pgfqpoint{3.644707in}{0.537667in}}{\pgfqpoint{3.655758in}{0.537667in}}%
\pgfpathclose%
\pgfusepath{stroke,fill}%
\end{pgfscope}%
\begin{pgfscope}%
\pgfpathrectangle{\pgfqpoint{0.800000in}{0.528000in}}{\pgfqpoint{4.960000in}{3.696000in}} %
\pgfusepath{clip}%
\pgfsetbuttcap%
\pgfsetroundjoin%
\definecolor{currentfill}{rgb}{0.000000,0.500000,0.000000}%
\pgfsetfillcolor{currentfill}%
\pgfsetlinewidth{1.003750pt}%
\definecolor{currentstroke}{rgb}{0.000000,0.500000,0.000000}%
\pgfsetstrokecolor{currentstroke}%
\pgfsetdash{}{0pt}%
\pgfpathmoveto{\pgfqpoint{3.705859in}{2.017167in}}%
\pgfpathcurveto{\pgfqpoint{3.716909in}{2.017167in}}{\pgfqpoint{3.727508in}{2.021557in}}{\pgfqpoint{3.735321in}{2.029371in}}%
\pgfpathcurveto{\pgfqpoint{3.743135in}{2.037184in}}{\pgfqpoint{3.747525in}{2.047783in}}{\pgfqpoint{3.747525in}{2.058833in}}%
\pgfpathcurveto{\pgfqpoint{3.747525in}{2.069883in}}{\pgfqpoint{3.743135in}{2.080483in}}{\pgfqpoint{3.735321in}{2.088296in}}%
\pgfpathcurveto{\pgfqpoint{3.727508in}{2.096110in}}{\pgfqpoint{3.716909in}{2.100500in}}{\pgfqpoint{3.705859in}{2.100500in}}%
\pgfpathcurveto{\pgfqpoint{3.694808in}{2.100500in}}{\pgfqpoint{3.684209in}{2.096110in}}{\pgfqpoint{3.676396in}{2.088296in}}%
\pgfpathcurveto{\pgfqpoint{3.668582in}{2.080483in}}{\pgfqpoint{3.664192in}{2.069883in}}{\pgfqpoint{3.664192in}{2.058833in}}%
\pgfpathcurveto{\pgfqpoint{3.664192in}{2.047783in}}{\pgfqpoint{3.668582in}{2.037184in}}{\pgfqpoint{3.676396in}{2.029371in}}%
\pgfpathcurveto{\pgfqpoint{3.684209in}{2.021557in}}{\pgfqpoint{3.694808in}{2.017167in}}{\pgfqpoint{3.705859in}{2.017167in}}%
\pgfpathclose%
\pgfusepath{stroke,fill}%
\end{pgfscope}%
\begin{pgfscope}%
\pgfpathrectangle{\pgfqpoint{0.800000in}{0.528000in}}{\pgfqpoint{4.960000in}{3.696000in}} %
\pgfusepath{clip}%
\pgfsetbuttcap%
\pgfsetroundjoin%
\definecolor{currentfill}{rgb}{0.000000,0.500000,0.000000}%
\pgfsetfillcolor{currentfill}%
\pgfsetlinewidth{1.003750pt}%
\definecolor{currentstroke}{rgb}{0.000000,0.500000,0.000000}%
\pgfsetstrokecolor{currentstroke}%
\pgfsetdash{}{0pt}%
\pgfpathmoveto{\pgfqpoint{3.755960in}{-0.614694in}}%
\pgfpathcurveto{\pgfqpoint{3.767010in}{-0.614694in}}{\pgfqpoint{3.777609in}{-0.610304in}}{\pgfqpoint{3.785422in}{-0.602490in}}%
\pgfpathcurveto{\pgfqpoint{3.793236in}{-0.594677in}}{\pgfqpoint{3.797626in}{-0.584078in}}{\pgfqpoint{3.797626in}{-0.573027in}}%
\pgfpathcurveto{\pgfqpoint{3.797626in}{-0.561977in}}{\pgfqpoint{3.793236in}{-0.551378in}}{\pgfqpoint{3.785422in}{-0.543565in}}%
\pgfpathcurveto{\pgfqpoint{3.777609in}{-0.535751in}}{\pgfqpoint{3.767010in}{-0.531361in}}{\pgfqpoint{3.755960in}{-0.531361in}}%
\pgfpathcurveto{\pgfqpoint{3.744909in}{-0.531361in}}{\pgfqpoint{3.734310in}{-0.535751in}}{\pgfqpoint{3.726497in}{-0.543565in}}%
\pgfpathcurveto{\pgfqpoint{3.718683in}{-0.551378in}}{\pgfqpoint{3.714293in}{-0.561977in}}{\pgfqpoint{3.714293in}{-0.573027in}}%
\pgfpathcurveto{\pgfqpoint{3.714293in}{-0.584078in}}{\pgfqpoint{3.718683in}{-0.594677in}}{\pgfqpoint{3.726497in}{-0.602490in}}%
\pgfpathcurveto{\pgfqpoint{3.734310in}{-0.610304in}}{\pgfqpoint{3.744909in}{-0.614694in}}{\pgfqpoint{3.755960in}{-0.614694in}}%
\pgfpathclose%
\pgfusepath{stroke,fill}%
\end{pgfscope}%
\begin{pgfscope}%
\pgfpathrectangle{\pgfqpoint{0.800000in}{0.528000in}}{\pgfqpoint{4.960000in}{3.696000in}} %
\pgfusepath{clip}%
\pgfsetbuttcap%
\pgfsetroundjoin%
\definecolor{currentfill}{rgb}{0.000000,0.500000,0.000000}%
\pgfsetfillcolor{currentfill}%
\pgfsetlinewidth{1.003750pt}%
\definecolor{currentstroke}{rgb}{0.000000,0.500000,0.000000}%
\pgfsetstrokecolor{currentstroke}%
\pgfsetdash{}{0pt}%
\pgfpathmoveto{\pgfqpoint{3.806061in}{0.616924in}}%
\pgfpathcurveto{\pgfqpoint{3.817111in}{0.616924in}}{\pgfqpoint{3.827710in}{0.621314in}}{\pgfqpoint{3.835523in}{0.629128in}}%
\pgfpathcurveto{\pgfqpoint{3.843337in}{0.636941in}}{\pgfqpoint{3.847727in}{0.647540in}}{\pgfqpoint{3.847727in}{0.658590in}}%
\pgfpathcurveto{\pgfqpoint{3.847727in}{0.669641in}}{\pgfqpoint{3.843337in}{0.680240in}}{\pgfqpoint{3.835523in}{0.688053in}}%
\pgfpathcurveto{\pgfqpoint{3.827710in}{0.695867in}}{\pgfqpoint{3.817111in}{0.700257in}}{\pgfqpoint{3.806061in}{0.700257in}}%
\pgfpathcurveto{\pgfqpoint{3.795010in}{0.700257in}}{\pgfqpoint{3.784411in}{0.695867in}}{\pgfqpoint{3.776598in}{0.688053in}}%
\pgfpathcurveto{\pgfqpoint{3.768784in}{0.680240in}}{\pgfqpoint{3.764394in}{0.669641in}}{\pgfqpoint{3.764394in}{0.658590in}}%
\pgfpathcurveto{\pgfqpoint{3.764394in}{0.647540in}}{\pgfqpoint{3.768784in}{0.636941in}}{\pgfqpoint{3.776598in}{0.629128in}}%
\pgfpathcurveto{\pgfqpoint{3.784411in}{0.621314in}}{\pgfqpoint{3.795010in}{0.616924in}}{\pgfqpoint{3.806061in}{0.616924in}}%
\pgfpathclose%
\pgfusepath{stroke,fill}%
\end{pgfscope}%
\begin{pgfscope}%
\pgfpathrectangle{\pgfqpoint{0.800000in}{0.528000in}}{\pgfqpoint{4.960000in}{3.696000in}} %
\pgfusepath{clip}%
\pgfsetbuttcap%
\pgfsetroundjoin%
\definecolor{currentfill}{rgb}{0.000000,0.500000,0.000000}%
\pgfsetfillcolor{currentfill}%
\pgfsetlinewidth{1.003750pt}%
\definecolor{currentstroke}{rgb}{0.000000,0.500000,0.000000}%
\pgfsetstrokecolor{currentstroke}%
\pgfsetdash{}{0pt}%
\pgfpathmoveto{\pgfqpoint{3.856162in}{-0.479322in}}%
\pgfpathcurveto{\pgfqpoint{3.867212in}{-0.479322in}}{\pgfqpoint{3.877811in}{-0.474932in}}{\pgfqpoint{3.885624in}{-0.467119in}}%
\pgfpathcurveto{\pgfqpoint{3.893438in}{-0.459305in}}{\pgfqpoint{3.897828in}{-0.448706in}}{\pgfqpoint{3.897828in}{-0.437656in}}%
\pgfpathcurveto{\pgfqpoint{3.897828in}{-0.426606in}}{\pgfqpoint{3.893438in}{-0.416007in}}{\pgfqpoint{3.885624in}{-0.408193in}}%
\pgfpathcurveto{\pgfqpoint{3.877811in}{-0.400379in}}{\pgfqpoint{3.867212in}{-0.395989in}}{\pgfqpoint{3.856162in}{-0.395989in}}%
\pgfpathcurveto{\pgfqpoint{3.845111in}{-0.395989in}}{\pgfqpoint{3.834512in}{-0.400379in}}{\pgfqpoint{3.826699in}{-0.408193in}}%
\pgfpathcurveto{\pgfqpoint{3.818885in}{-0.416007in}}{\pgfqpoint{3.814495in}{-0.426606in}}{\pgfqpoint{3.814495in}{-0.437656in}}%
\pgfpathcurveto{\pgfqpoint{3.814495in}{-0.448706in}}{\pgfqpoint{3.818885in}{-0.459305in}}{\pgfqpoint{3.826699in}{-0.467119in}}%
\pgfpathcurveto{\pgfqpoint{3.834512in}{-0.474932in}}{\pgfqpoint{3.845111in}{-0.479322in}}{\pgfqpoint{3.856162in}{-0.479322in}}%
\pgfpathclose%
\pgfusepath{stroke,fill}%
\end{pgfscope}%
\begin{pgfscope}%
\pgfpathrectangle{\pgfqpoint{0.800000in}{0.528000in}}{\pgfqpoint{4.960000in}{3.696000in}} %
\pgfusepath{clip}%
\pgfsetbuttcap%
\pgfsetroundjoin%
\definecolor{currentfill}{rgb}{0.000000,0.500000,0.000000}%
\pgfsetfillcolor{currentfill}%
\pgfsetlinewidth{1.003750pt}%
\definecolor{currentstroke}{rgb}{0.000000,0.500000,0.000000}%
\pgfsetstrokecolor{currentstroke}%
\pgfsetdash{}{0pt}%
\pgfpathmoveto{\pgfqpoint{3.906263in}{1.193995in}}%
\pgfpathcurveto{\pgfqpoint{3.917313in}{1.193995in}}{\pgfqpoint{3.927912in}{1.198385in}}{\pgfqpoint{3.935725in}{1.206199in}}%
\pgfpathcurveto{\pgfqpoint{3.943539in}{1.214013in}}{\pgfqpoint{3.947929in}{1.224612in}}{\pgfqpoint{3.947929in}{1.235662in}}%
\pgfpathcurveto{\pgfqpoint{3.947929in}{1.246712in}}{\pgfqpoint{3.943539in}{1.257311in}}{\pgfqpoint{3.935725in}{1.265125in}}%
\pgfpathcurveto{\pgfqpoint{3.927912in}{1.272938in}}{\pgfqpoint{3.917313in}{1.277329in}}{\pgfqpoint{3.906263in}{1.277329in}}%
\pgfpathcurveto{\pgfqpoint{3.895212in}{1.277329in}}{\pgfqpoint{3.884613in}{1.272938in}}{\pgfqpoint{3.876800in}{1.265125in}}%
\pgfpathcurveto{\pgfqpoint{3.868986in}{1.257311in}}{\pgfqpoint{3.864596in}{1.246712in}}{\pgfqpoint{3.864596in}{1.235662in}}%
\pgfpathcurveto{\pgfqpoint{3.864596in}{1.224612in}}{\pgfqpoint{3.868986in}{1.214013in}}{\pgfqpoint{3.876800in}{1.206199in}}%
\pgfpathcurveto{\pgfqpoint{3.884613in}{1.198385in}}{\pgfqpoint{3.895212in}{1.193995in}}{\pgfqpoint{3.906263in}{1.193995in}}%
\pgfpathclose%
\pgfusepath{stroke,fill}%
\end{pgfscope}%
\begin{pgfscope}%
\pgfpathrectangle{\pgfqpoint{0.800000in}{0.528000in}}{\pgfqpoint{4.960000in}{3.696000in}} %
\pgfusepath{clip}%
\pgfsetbuttcap%
\pgfsetroundjoin%
\definecolor{currentfill}{rgb}{0.000000,0.500000,0.000000}%
\pgfsetfillcolor{currentfill}%
\pgfsetlinewidth{1.003750pt}%
\definecolor{currentstroke}{rgb}{0.000000,0.500000,0.000000}%
\pgfsetstrokecolor{currentstroke}%
\pgfsetdash{}{0pt}%
\pgfpathmoveto{\pgfqpoint{3.956364in}{-1.543579in}}%
\pgfpathcurveto{\pgfqpoint{3.967414in}{-1.543579in}}{\pgfqpoint{3.978013in}{-1.539189in}}{\pgfqpoint{3.985826in}{-1.531375in}}%
\pgfpathcurveto{\pgfqpoint{3.993640in}{-1.523562in}}{\pgfqpoint{3.998030in}{-1.512963in}}{\pgfqpoint{3.998030in}{-1.501912in}}%
\pgfpathcurveto{\pgfqpoint{3.998030in}{-1.490862in}}{\pgfqpoint{3.993640in}{-1.480263in}}{\pgfqpoint{3.985826in}{-1.472450in}}%
\pgfpathcurveto{\pgfqpoint{3.978013in}{-1.464636in}}{\pgfqpoint{3.967414in}{-1.460246in}}{\pgfqpoint{3.956364in}{-1.460246in}}%
\pgfpathcurveto{\pgfqpoint{3.945314in}{-1.460246in}}{\pgfqpoint{3.934714in}{-1.464636in}}{\pgfqpoint{3.926901in}{-1.472450in}}%
\pgfpathcurveto{\pgfqpoint{3.919087in}{-1.480263in}}{\pgfqpoint{3.914697in}{-1.490862in}}{\pgfqpoint{3.914697in}{-1.501912in}}%
\pgfpathcurveto{\pgfqpoint{3.914697in}{-1.512963in}}{\pgfqpoint{3.919087in}{-1.523562in}}{\pgfqpoint{3.926901in}{-1.531375in}}%
\pgfpathcurveto{\pgfqpoint{3.934714in}{-1.539189in}}{\pgfqpoint{3.945314in}{-1.543579in}}{\pgfqpoint{3.956364in}{-1.543579in}}%
\pgfpathclose%
\pgfusepath{stroke,fill}%
\end{pgfscope}%
\begin{pgfscope}%
\pgfpathrectangle{\pgfqpoint{0.800000in}{0.528000in}}{\pgfqpoint{4.960000in}{3.696000in}} %
\pgfusepath{clip}%
\pgfsetbuttcap%
\pgfsetroundjoin%
\definecolor{currentfill}{rgb}{0.000000,0.500000,0.000000}%
\pgfsetfillcolor{currentfill}%
\pgfsetlinewidth{1.003750pt}%
\definecolor{currentstroke}{rgb}{0.000000,0.500000,0.000000}%
\pgfsetstrokecolor{currentstroke}%
\pgfsetdash{}{0pt}%
\pgfpathmoveto{\pgfqpoint{4.006465in}{0.024514in}}%
\pgfpathcurveto{\pgfqpoint{4.017515in}{0.024514in}}{\pgfqpoint{4.028114in}{0.028904in}}{\pgfqpoint{4.035927in}{0.036718in}}%
\pgfpathcurveto{\pgfqpoint{4.043741in}{0.044532in}}{\pgfqpoint{4.048131in}{0.055131in}}{\pgfqpoint{4.048131in}{0.066181in}}%
\pgfpathcurveto{\pgfqpoint{4.048131in}{0.077231in}}{\pgfqpoint{4.043741in}{0.087830in}}{\pgfqpoint{4.035927in}{0.095643in}}%
\pgfpathcurveto{\pgfqpoint{4.028114in}{0.103457in}}{\pgfqpoint{4.017515in}{0.107847in}}{\pgfqpoint{4.006465in}{0.107847in}}%
\pgfpathcurveto{\pgfqpoint{3.995415in}{0.107847in}}{\pgfqpoint{3.984815in}{0.103457in}}{\pgfqpoint{3.977002in}{0.095643in}}%
\pgfpathcurveto{\pgfqpoint{3.969188in}{0.087830in}}{\pgfqpoint{3.964798in}{0.077231in}}{\pgfqpoint{3.964798in}{0.066181in}}%
\pgfpathcurveto{\pgfqpoint{3.964798in}{0.055131in}}{\pgfqpoint{3.969188in}{0.044532in}}{\pgfqpoint{3.977002in}{0.036718in}}%
\pgfpathcurveto{\pgfqpoint{3.984815in}{0.028904in}}{\pgfqpoint{3.995415in}{0.024514in}}{\pgfqpoint{4.006465in}{0.024514in}}%
\pgfpathclose%
\pgfusepath{stroke,fill}%
\end{pgfscope}%
\begin{pgfscope}%
\pgfpathrectangle{\pgfqpoint{0.800000in}{0.528000in}}{\pgfqpoint{4.960000in}{3.696000in}} %
\pgfusepath{clip}%
\pgfsetbuttcap%
\pgfsetroundjoin%
\definecolor{currentfill}{rgb}{0.000000,0.500000,0.000000}%
\pgfsetfillcolor{currentfill}%
\pgfsetlinewidth{1.003750pt}%
\definecolor{currentstroke}{rgb}{0.000000,0.500000,0.000000}%
\pgfsetstrokecolor{currentstroke}%
\pgfsetdash{}{0pt}%
\pgfpathmoveto{\pgfqpoint{4.056566in}{0.434837in}}%
\pgfpathcurveto{\pgfqpoint{4.067616in}{0.434837in}}{\pgfqpoint{4.078215in}{0.439227in}}{\pgfqpoint{4.086028in}{0.447041in}}%
\pgfpathcurveto{\pgfqpoint{4.093842in}{0.454854in}}{\pgfqpoint{4.098232in}{0.465454in}}{\pgfqpoint{4.098232in}{0.476504in}}%
\pgfpathcurveto{\pgfqpoint{4.098232in}{0.487554in}}{\pgfqpoint{4.093842in}{0.498153in}}{\pgfqpoint{4.086028in}{0.505966in}}%
\pgfpathcurveto{\pgfqpoint{4.078215in}{0.513780in}}{\pgfqpoint{4.067616in}{0.518170in}}{\pgfqpoint{4.056566in}{0.518170in}}%
\pgfpathcurveto{\pgfqpoint{4.045516in}{0.518170in}}{\pgfqpoint{4.034916in}{0.513780in}}{\pgfqpoint{4.027103in}{0.505966in}}%
\pgfpathcurveto{\pgfqpoint{4.019289in}{0.498153in}}{\pgfqpoint{4.014899in}{0.487554in}}{\pgfqpoint{4.014899in}{0.476504in}}%
\pgfpathcurveto{\pgfqpoint{4.014899in}{0.465454in}}{\pgfqpoint{4.019289in}{0.454854in}}{\pgfqpoint{4.027103in}{0.447041in}}%
\pgfpathcurveto{\pgfqpoint{4.034916in}{0.439227in}}{\pgfqpoint{4.045516in}{0.434837in}}{\pgfqpoint{4.056566in}{0.434837in}}%
\pgfpathclose%
\pgfusepath{stroke,fill}%
\end{pgfscope}%
\begin{pgfscope}%
\pgfpathrectangle{\pgfqpoint{0.800000in}{0.528000in}}{\pgfqpoint{4.960000in}{3.696000in}} %
\pgfusepath{clip}%
\pgfsetbuttcap%
\pgfsetroundjoin%
\definecolor{currentfill}{rgb}{0.000000,0.500000,0.000000}%
\pgfsetfillcolor{currentfill}%
\pgfsetlinewidth{1.003750pt}%
\definecolor{currentstroke}{rgb}{0.000000,0.500000,0.000000}%
\pgfsetstrokecolor{currentstroke}%
\pgfsetdash{}{0pt}%
\pgfpathmoveto{\pgfqpoint{4.106667in}{2.072927in}}%
\pgfpathcurveto{\pgfqpoint{4.117717in}{2.072927in}}{\pgfqpoint{4.128316in}{2.077318in}}{\pgfqpoint{4.136129in}{2.085131in}}%
\pgfpathcurveto{\pgfqpoint{4.143943in}{2.092945in}}{\pgfqpoint{4.148333in}{2.103544in}}{\pgfqpoint{4.148333in}{2.114594in}}%
\pgfpathcurveto{\pgfqpoint{4.148333in}{2.125644in}}{\pgfqpoint{4.143943in}{2.136243in}}{\pgfqpoint{4.136129in}{2.144057in}}%
\pgfpathcurveto{\pgfqpoint{4.128316in}{2.151870in}}{\pgfqpoint{4.117717in}{2.156261in}}{\pgfqpoint{4.106667in}{2.156261in}}%
\pgfpathcurveto{\pgfqpoint{4.095617in}{2.156261in}}{\pgfqpoint{4.085018in}{2.151870in}}{\pgfqpoint{4.077204in}{2.144057in}}%
\pgfpathcurveto{\pgfqpoint{4.069390in}{2.136243in}}{\pgfqpoint{4.065000in}{2.125644in}}{\pgfqpoint{4.065000in}{2.114594in}}%
\pgfpathcurveto{\pgfqpoint{4.065000in}{2.103544in}}{\pgfqpoint{4.069390in}{2.092945in}}{\pgfqpoint{4.077204in}{2.085131in}}%
\pgfpathcurveto{\pgfqpoint{4.085018in}{2.077318in}}{\pgfqpoint{4.095617in}{2.072927in}}{\pgfqpoint{4.106667in}{2.072927in}}%
\pgfpathclose%
\pgfusepath{stroke,fill}%
\end{pgfscope}%
\begin{pgfscope}%
\pgfpathrectangle{\pgfqpoint{0.800000in}{0.528000in}}{\pgfqpoint{4.960000in}{3.696000in}} %
\pgfusepath{clip}%
\pgfsetbuttcap%
\pgfsetroundjoin%
\definecolor{currentfill}{rgb}{0.000000,0.500000,0.000000}%
\pgfsetfillcolor{currentfill}%
\pgfsetlinewidth{1.003750pt}%
\definecolor{currentstroke}{rgb}{0.000000,0.500000,0.000000}%
\pgfsetstrokecolor{currentstroke}%
\pgfsetdash{}{0pt}%
\pgfpathmoveto{\pgfqpoint{4.156768in}{2.471377in}}%
\pgfpathcurveto{\pgfqpoint{4.167818in}{2.471377in}}{\pgfqpoint{4.178417in}{2.475767in}}{\pgfqpoint{4.186230in}{2.483581in}}%
\pgfpathcurveto{\pgfqpoint{4.194044in}{2.491394in}}{\pgfqpoint{4.198434in}{2.501993in}}{\pgfqpoint{4.198434in}{2.513043in}}%
\pgfpathcurveto{\pgfqpoint{4.198434in}{2.524093in}}{\pgfqpoint{4.194044in}{2.534693in}}{\pgfqpoint{4.186230in}{2.542506in}}%
\pgfpathcurveto{\pgfqpoint{4.178417in}{2.550320in}}{\pgfqpoint{4.167818in}{2.554710in}}{\pgfqpoint{4.156768in}{2.554710in}}%
\pgfpathcurveto{\pgfqpoint{4.145718in}{2.554710in}}{\pgfqpoint{4.135119in}{2.550320in}}{\pgfqpoint{4.127305in}{2.542506in}}%
\pgfpathcurveto{\pgfqpoint{4.119491in}{2.534693in}}{\pgfqpoint{4.115101in}{2.524093in}}{\pgfqpoint{4.115101in}{2.513043in}}%
\pgfpathcurveto{\pgfqpoint{4.115101in}{2.501993in}}{\pgfqpoint{4.119491in}{2.491394in}}{\pgfqpoint{4.127305in}{2.483581in}}%
\pgfpathcurveto{\pgfqpoint{4.135119in}{2.475767in}}{\pgfqpoint{4.145718in}{2.471377in}}{\pgfqpoint{4.156768in}{2.471377in}}%
\pgfpathclose%
\pgfusepath{stroke,fill}%
\end{pgfscope}%
\begin{pgfscope}%
\pgfpathrectangle{\pgfqpoint{0.800000in}{0.528000in}}{\pgfqpoint{4.960000in}{3.696000in}} %
\pgfusepath{clip}%
\pgfsetbuttcap%
\pgfsetroundjoin%
\definecolor{currentfill}{rgb}{0.000000,0.500000,0.000000}%
\pgfsetfillcolor{currentfill}%
\pgfsetlinewidth{1.003750pt}%
\definecolor{currentstroke}{rgb}{0.000000,0.500000,0.000000}%
\pgfsetstrokecolor{currentstroke}%
\pgfsetdash{}{0pt}%
\pgfpathmoveto{\pgfqpoint{4.206869in}{-1.229223in}}%
\pgfpathcurveto{\pgfqpoint{4.217919in}{-1.229223in}}{\pgfqpoint{4.228518in}{-1.224833in}}{\pgfqpoint{4.236331in}{-1.217019in}}%
\pgfpathcurveto{\pgfqpoint{4.244145in}{-1.209206in}}{\pgfqpoint{4.248535in}{-1.198607in}}{\pgfqpoint{4.248535in}{-1.187557in}}%
\pgfpathcurveto{\pgfqpoint{4.248535in}{-1.176507in}}{\pgfqpoint{4.244145in}{-1.165908in}}{\pgfqpoint{4.236331in}{-1.158094in}}%
\pgfpathcurveto{\pgfqpoint{4.228518in}{-1.150280in}}{\pgfqpoint{4.217919in}{-1.145890in}}{\pgfqpoint{4.206869in}{-1.145890in}}%
\pgfpathcurveto{\pgfqpoint{4.195819in}{-1.145890in}}{\pgfqpoint{4.185220in}{-1.150280in}}{\pgfqpoint{4.177406in}{-1.158094in}}%
\pgfpathcurveto{\pgfqpoint{4.169592in}{-1.165908in}}{\pgfqpoint{4.165202in}{-1.176507in}}{\pgfqpoint{4.165202in}{-1.187557in}}%
\pgfpathcurveto{\pgfqpoint{4.165202in}{-1.198607in}}{\pgfqpoint{4.169592in}{-1.209206in}}{\pgfqpoint{4.177406in}{-1.217019in}}%
\pgfpathcurveto{\pgfqpoint{4.185220in}{-1.224833in}}{\pgfqpoint{4.195819in}{-1.229223in}}{\pgfqpoint{4.206869in}{-1.229223in}}%
\pgfpathclose%
\pgfusepath{stroke,fill}%
\end{pgfscope}%
\begin{pgfscope}%
\pgfpathrectangle{\pgfqpoint{0.800000in}{0.528000in}}{\pgfqpoint{4.960000in}{3.696000in}} %
\pgfusepath{clip}%
\pgfsetbuttcap%
\pgfsetroundjoin%
\definecolor{currentfill}{rgb}{0.000000,0.500000,0.000000}%
\pgfsetfillcolor{currentfill}%
\pgfsetlinewidth{1.003750pt}%
\definecolor{currentstroke}{rgb}{0.000000,0.500000,0.000000}%
\pgfsetstrokecolor{currentstroke}%
\pgfsetdash{}{0pt}%
\pgfpathmoveto{\pgfqpoint{4.256970in}{-1.608056in}}%
\pgfpathcurveto{\pgfqpoint{4.268020in}{-1.608056in}}{\pgfqpoint{4.278619in}{-1.603665in}}{\pgfqpoint{4.286432in}{-1.595852in}}%
\pgfpathcurveto{\pgfqpoint{4.294246in}{-1.588038in}}{\pgfqpoint{4.298636in}{-1.577439in}}{\pgfqpoint{4.298636in}{-1.566389in}}%
\pgfpathcurveto{\pgfqpoint{4.298636in}{-1.555339in}}{\pgfqpoint{4.294246in}{-1.544740in}}{\pgfqpoint{4.286432in}{-1.536926in}}%
\pgfpathcurveto{\pgfqpoint{4.278619in}{-1.529113in}}{\pgfqpoint{4.268020in}{-1.524722in}}{\pgfqpoint{4.256970in}{-1.524722in}}%
\pgfpathcurveto{\pgfqpoint{4.245920in}{-1.524722in}}{\pgfqpoint{4.235321in}{-1.529113in}}{\pgfqpoint{4.227507in}{-1.536926in}}%
\pgfpathcurveto{\pgfqpoint{4.219693in}{-1.544740in}}{\pgfqpoint{4.215303in}{-1.555339in}}{\pgfqpoint{4.215303in}{-1.566389in}}%
\pgfpathcurveto{\pgfqpoint{4.215303in}{-1.577439in}}{\pgfqpoint{4.219693in}{-1.588038in}}{\pgfqpoint{4.227507in}{-1.595852in}}%
\pgfpathcurveto{\pgfqpoint{4.235321in}{-1.603665in}}{\pgfqpoint{4.245920in}{-1.608056in}}{\pgfqpoint{4.256970in}{-1.608056in}}%
\pgfpathclose%
\pgfusepath{stroke,fill}%
\end{pgfscope}%
\begin{pgfscope}%
\pgfpathrectangle{\pgfqpoint{0.800000in}{0.528000in}}{\pgfqpoint{4.960000in}{3.696000in}} %
\pgfusepath{clip}%
\pgfsetbuttcap%
\pgfsetroundjoin%
\definecolor{currentfill}{rgb}{0.000000,0.500000,0.000000}%
\pgfsetfillcolor{currentfill}%
\pgfsetlinewidth{1.003750pt}%
\definecolor{currentstroke}{rgb}{0.000000,0.500000,0.000000}%
\pgfsetstrokecolor{currentstroke}%
\pgfsetdash{}{0pt}%
\pgfpathmoveto{\pgfqpoint{4.307071in}{-0.487112in}}%
\pgfpathcurveto{\pgfqpoint{4.318121in}{-0.487112in}}{\pgfqpoint{4.328720in}{-0.482722in}}{\pgfqpoint{4.336533in}{-0.474908in}}%
\pgfpathcurveto{\pgfqpoint{4.344347in}{-0.467095in}}{\pgfqpoint{4.348737in}{-0.456495in}}{\pgfqpoint{4.348737in}{-0.445445in}}%
\pgfpathcurveto{\pgfqpoint{4.348737in}{-0.434395in}}{\pgfqpoint{4.344347in}{-0.423796in}}{\pgfqpoint{4.336533in}{-0.415983in}}%
\pgfpathcurveto{\pgfqpoint{4.328720in}{-0.408169in}}{\pgfqpoint{4.318121in}{-0.403779in}}{\pgfqpoint{4.307071in}{-0.403779in}}%
\pgfpathcurveto{\pgfqpoint{4.296021in}{-0.403779in}}{\pgfqpoint{4.285422in}{-0.408169in}}{\pgfqpoint{4.277608in}{-0.415983in}}%
\pgfpathcurveto{\pgfqpoint{4.269794in}{-0.423796in}}{\pgfqpoint{4.265404in}{-0.434395in}}{\pgfqpoint{4.265404in}{-0.445445in}}%
\pgfpathcurveto{\pgfqpoint{4.265404in}{-0.456495in}}{\pgfqpoint{4.269794in}{-0.467095in}}{\pgfqpoint{4.277608in}{-0.474908in}}%
\pgfpathcurveto{\pgfqpoint{4.285422in}{-0.482722in}}{\pgfqpoint{4.296021in}{-0.487112in}}{\pgfqpoint{4.307071in}{-0.487112in}}%
\pgfpathclose%
\pgfusepath{stroke,fill}%
\end{pgfscope}%
\begin{pgfscope}%
\pgfpathrectangle{\pgfqpoint{0.800000in}{0.528000in}}{\pgfqpoint{4.960000in}{3.696000in}} %
\pgfusepath{clip}%
\pgfsetbuttcap%
\pgfsetroundjoin%
\definecolor{currentfill}{rgb}{0.000000,0.500000,0.000000}%
\pgfsetfillcolor{currentfill}%
\pgfsetlinewidth{1.003750pt}%
\definecolor{currentstroke}{rgb}{0.000000,0.500000,0.000000}%
\pgfsetstrokecolor{currentstroke}%
\pgfsetdash{}{0pt}%
\pgfpathmoveto{\pgfqpoint{4.357172in}{1.878283in}}%
\pgfpathcurveto{\pgfqpoint{4.368222in}{1.878283in}}{\pgfqpoint{4.378821in}{1.882673in}}{\pgfqpoint{4.386634in}{1.890487in}}%
\pgfpathcurveto{\pgfqpoint{4.394448in}{1.898301in}}{\pgfqpoint{4.398838in}{1.908900in}}{\pgfqpoint{4.398838in}{1.919950in}}%
\pgfpathcurveto{\pgfqpoint{4.398838in}{1.931000in}}{\pgfqpoint{4.394448in}{1.941599in}}{\pgfqpoint{4.386634in}{1.949413in}}%
\pgfpathcurveto{\pgfqpoint{4.378821in}{1.957226in}}{\pgfqpoint{4.368222in}{1.961616in}}{\pgfqpoint{4.357172in}{1.961616in}}%
\pgfpathcurveto{\pgfqpoint{4.346122in}{1.961616in}}{\pgfqpoint{4.335523in}{1.957226in}}{\pgfqpoint{4.327709in}{1.949413in}}%
\pgfpathcurveto{\pgfqpoint{4.319895in}{1.941599in}}{\pgfqpoint{4.315505in}{1.931000in}}{\pgfqpoint{4.315505in}{1.919950in}}%
\pgfpathcurveto{\pgfqpoint{4.315505in}{1.908900in}}{\pgfqpoint{4.319895in}{1.898301in}}{\pgfqpoint{4.327709in}{1.890487in}}%
\pgfpathcurveto{\pgfqpoint{4.335523in}{1.882673in}}{\pgfqpoint{4.346122in}{1.878283in}}{\pgfqpoint{4.357172in}{1.878283in}}%
\pgfpathclose%
\pgfusepath{stroke,fill}%
\end{pgfscope}%
\begin{pgfscope}%
\pgfpathrectangle{\pgfqpoint{0.800000in}{0.528000in}}{\pgfqpoint{4.960000in}{3.696000in}} %
\pgfusepath{clip}%
\pgfsetbuttcap%
\pgfsetroundjoin%
\definecolor{currentfill}{rgb}{0.000000,0.500000,0.000000}%
\pgfsetfillcolor{currentfill}%
\pgfsetlinewidth{1.003750pt}%
\definecolor{currentstroke}{rgb}{0.000000,0.500000,0.000000}%
\pgfsetstrokecolor{currentstroke}%
\pgfsetdash{}{0pt}%
\pgfpathmoveto{\pgfqpoint{4.407273in}{-1.599351in}}%
\pgfpathcurveto{\pgfqpoint{4.418323in}{-1.599351in}}{\pgfqpoint{4.428922in}{-1.594961in}}{\pgfqpoint{4.436736in}{-1.587147in}}%
\pgfpathcurveto{\pgfqpoint{4.444549in}{-1.579333in}}{\pgfqpoint{4.448939in}{-1.568734in}}{\pgfqpoint{4.448939in}{-1.557684in}}%
\pgfpathcurveto{\pgfqpoint{4.448939in}{-1.546634in}}{\pgfqpoint{4.444549in}{-1.536035in}}{\pgfqpoint{4.436736in}{-1.528222in}}%
\pgfpathcurveto{\pgfqpoint{4.428922in}{-1.520408in}}{\pgfqpoint{4.418323in}{-1.516018in}}{\pgfqpoint{4.407273in}{-1.516018in}}%
\pgfpathcurveto{\pgfqpoint{4.396223in}{-1.516018in}}{\pgfqpoint{4.385624in}{-1.520408in}}{\pgfqpoint{4.377810in}{-1.528222in}}%
\pgfpathcurveto{\pgfqpoint{4.369996in}{-1.536035in}}{\pgfqpoint{4.365606in}{-1.546634in}}{\pgfqpoint{4.365606in}{-1.557684in}}%
\pgfpathcurveto{\pgfqpoint{4.365606in}{-1.568734in}}{\pgfqpoint{4.369996in}{-1.579333in}}{\pgfqpoint{4.377810in}{-1.587147in}}%
\pgfpathcurveto{\pgfqpoint{4.385624in}{-1.594961in}}{\pgfqpoint{4.396223in}{-1.599351in}}{\pgfqpoint{4.407273in}{-1.599351in}}%
\pgfpathclose%
\pgfusepath{stroke,fill}%
\end{pgfscope}%
\begin{pgfscope}%
\pgfpathrectangle{\pgfqpoint{0.800000in}{0.528000in}}{\pgfqpoint{4.960000in}{3.696000in}} %
\pgfusepath{clip}%
\pgfsetbuttcap%
\pgfsetroundjoin%
\definecolor{currentfill}{rgb}{0.000000,0.500000,0.000000}%
\pgfsetfillcolor{currentfill}%
\pgfsetlinewidth{1.003750pt}%
\definecolor{currentstroke}{rgb}{0.000000,0.500000,0.000000}%
\pgfsetstrokecolor{currentstroke}%
\pgfsetdash{}{0pt}%
\pgfpathmoveto{\pgfqpoint{4.457374in}{-0.033806in}}%
\pgfpathcurveto{\pgfqpoint{4.468424in}{-0.033806in}}{\pgfqpoint{4.479023in}{-0.029416in}}{\pgfqpoint{4.486837in}{-0.021602in}}%
\pgfpathcurveto{\pgfqpoint{4.494650in}{-0.013789in}}{\pgfqpoint{4.499040in}{-0.003190in}}{\pgfqpoint{4.499040in}{0.007861in}}%
\pgfpathcurveto{\pgfqpoint{4.499040in}{0.018911in}}{\pgfqpoint{4.494650in}{0.029510in}}{\pgfqpoint{4.486837in}{0.037323in}}%
\pgfpathcurveto{\pgfqpoint{4.479023in}{0.045137in}}{\pgfqpoint{4.468424in}{0.049527in}}{\pgfqpoint{4.457374in}{0.049527in}}%
\pgfpathcurveto{\pgfqpoint{4.446324in}{0.049527in}}{\pgfqpoint{4.435725in}{0.045137in}}{\pgfqpoint{4.427911in}{0.037323in}}%
\pgfpathcurveto{\pgfqpoint{4.420097in}{0.029510in}}{\pgfqpoint{4.415707in}{0.018911in}}{\pgfqpoint{4.415707in}{0.007861in}}%
\pgfpathcurveto{\pgfqpoint{4.415707in}{-0.003190in}}{\pgfqpoint{4.420097in}{-0.013789in}}{\pgfqpoint{4.427911in}{-0.021602in}}%
\pgfpathcurveto{\pgfqpoint{4.435725in}{-0.029416in}}{\pgfqpoint{4.446324in}{-0.033806in}}{\pgfqpoint{4.457374in}{-0.033806in}}%
\pgfpathclose%
\pgfusepath{stroke,fill}%
\end{pgfscope}%
\begin{pgfscope}%
\pgfpathrectangle{\pgfqpoint{0.800000in}{0.528000in}}{\pgfqpoint{4.960000in}{3.696000in}} %
\pgfusepath{clip}%
\pgfsetbuttcap%
\pgfsetroundjoin%
\definecolor{currentfill}{rgb}{0.000000,0.500000,0.000000}%
\pgfsetfillcolor{currentfill}%
\pgfsetlinewidth{1.003750pt}%
\definecolor{currentstroke}{rgb}{0.000000,0.500000,0.000000}%
\pgfsetstrokecolor{currentstroke}%
\pgfsetdash{}{0pt}%
\pgfpathmoveto{\pgfqpoint{4.507475in}{1.635490in}}%
\pgfpathcurveto{\pgfqpoint{4.518525in}{1.635490in}}{\pgfqpoint{4.529124in}{1.639880in}}{\pgfqpoint{4.536938in}{1.647694in}}%
\pgfpathcurveto{\pgfqpoint{4.544751in}{1.655508in}}{\pgfqpoint{4.549141in}{1.666107in}}{\pgfqpoint{4.549141in}{1.677157in}}%
\pgfpathcurveto{\pgfqpoint{4.549141in}{1.688207in}}{\pgfqpoint{4.544751in}{1.698806in}}{\pgfqpoint{4.536938in}{1.706620in}}%
\pgfpathcurveto{\pgfqpoint{4.529124in}{1.714433in}}{\pgfqpoint{4.518525in}{1.718823in}}{\pgfqpoint{4.507475in}{1.718823in}}%
\pgfpathcurveto{\pgfqpoint{4.496425in}{1.718823in}}{\pgfqpoint{4.485826in}{1.714433in}}{\pgfqpoint{4.478012in}{1.706620in}}%
\pgfpathcurveto{\pgfqpoint{4.470198in}{1.698806in}}{\pgfqpoint{4.465808in}{1.688207in}}{\pgfqpoint{4.465808in}{1.677157in}}%
\pgfpathcurveto{\pgfqpoint{4.465808in}{1.666107in}}{\pgfqpoint{4.470198in}{1.655508in}}{\pgfqpoint{4.478012in}{1.647694in}}%
\pgfpathcurveto{\pgfqpoint{4.485826in}{1.639880in}}{\pgfqpoint{4.496425in}{1.635490in}}{\pgfqpoint{4.507475in}{1.635490in}}%
\pgfpathclose%
\pgfusepath{stroke,fill}%
\end{pgfscope}%
\begin{pgfscope}%
\pgfpathrectangle{\pgfqpoint{0.800000in}{0.528000in}}{\pgfqpoint{4.960000in}{3.696000in}} %
\pgfusepath{clip}%
\pgfsetbuttcap%
\pgfsetroundjoin%
\definecolor{currentfill}{rgb}{0.000000,0.500000,0.000000}%
\pgfsetfillcolor{currentfill}%
\pgfsetlinewidth{1.003750pt}%
\definecolor{currentstroke}{rgb}{0.000000,0.500000,0.000000}%
\pgfsetstrokecolor{currentstroke}%
\pgfsetdash{}{0pt}%
\pgfpathmoveto{\pgfqpoint{4.557576in}{-0.499582in}}%
\pgfpathcurveto{\pgfqpoint{4.568626in}{-0.499582in}}{\pgfqpoint{4.579225in}{-0.495192in}}{\pgfqpoint{4.587039in}{-0.487378in}}%
\pgfpathcurveto{\pgfqpoint{4.594852in}{-0.479564in}}{\pgfqpoint{4.599242in}{-0.468965in}}{\pgfqpoint{4.599242in}{-0.457915in}}%
\pgfpathcurveto{\pgfqpoint{4.599242in}{-0.446865in}}{\pgfqpoint{4.594852in}{-0.436266in}}{\pgfqpoint{4.587039in}{-0.428452in}}%
\pgfpathcurveto{\pgfqpoint{4.579225in}{-0.420639in}}{\pgfqpoint{4.568626in}{-0.416249in}}{\pgfqpoint{4.557576in}{-0.416249in}}%
\pgfpathcurveto{\pgfqpoint{4.546526in}{-0.416249in}}{\pgfqpoint{4.535927in}{-0.420639in}}{\pgfqpoint{4.528113in}{-0.428452in}}%
\pgfpathcurveto{\pgfqpoint{4.520299in}{-0.436266in}}{\pgfqpoint{4.515909in}{-0.446865in}}{\pgfqpoint{4.515909in}{-0.457915in}}%
\pgfpathcurveto{\pgfqpoint{4.515909in}{-0.468965in}}{\pgfqpoint{4.520299in}{-0.479564in}}{\pgfqpoint{4.528113in}{-0.487378in}}%
\pgfpathcurveto{\pgfqpoint{4.535927in}{-0.495192in}}{\pgfqpoint{4.546526in}{-0.499582in}}{\pgfqpoint{4.557576in}{-0.499582in}}%
\pgfpathclose%
\pgfusepath{stroke,fill}%
\end{pgfscope}%
\begin{pgfscope}%
\pgfpathrectangle{\pgfqpoint{0.800000in}{0.528000in}}{\pgfqpoint{4.960000in}{3.696000in}} %
\pgfusepath{clip}%
\pgfsetbuttcap%
\pgfsetroundjoin%
\definecolor{currentfill}{rgb}{0.000000,0.500000,0.000000}%
\pgfsetfillcolor{currentfill}%
\pgfsetlinewidth{1.003750pt}%
\definecolor{currentstroke}{rgb}{0.000000,0.500000,0.000000}%
\pgfsetstrokecolor{currentstroke}%
\pgfsetdash{}{0pt}%
\pgfpathmoveto{\pgfqpoint{4.607677in}{2.092573in}}%
\pgfpathcurveto{\pgfqpoint{4.618727in}{2.092573in}}{\pgfqpoint{4.629326in}{2.096963in}}{\pgfqpoint{4.637140in}{2.104777in}}%
\pgfpathcurveto{\pgfqpoint{4.644953in}{2.112590in}}{\pgfqpoint{4.649343in}{2.123189in}}{\pgfqpoint{4.649343in}{2.134239in}}%
\pgfpathcurveto{\pgfqpoint{4.649343in}{2.145290in}}{\pgfqpoint{4.644953in}{2.155889in}}{\pgfqpoint{4.637140in}{2.163702in}}%
\pgfpathcurveto{\pgfqpoint{4.629326in}{2.171516in}}{\pgfqpoint{4.618727in}{2.175906in}}{\pgfqpoint{4.607677in}{2.175906in}}%
\pgfpathcurveto{\pgfqpoint{4.596627in}{2.175906in}}{\pgfqpoint{4.586028in}{2.171516in}}{\pgfqpoint{4.578214in}{2.163702in}}%
\pgfpathcurveto{\pgfqpoint{4.570400in}{2.155889in}}{\pgfqpoint{4.566010in}{2.145290in}}{\pgfqpoint{4.566010in}{2.134239in}}%
\pgfpathcurveto{\pgfqpoint{4.566010in}{2.123189in}}{\pgfqpoint{4.570400in}{2.112590in}}{\pgfqpoint{4.578214in}{2.104777in}}%
\pgfpathcurveto{\pgfqpoint{4.586028in}{2.096963in}}{\pgfqpoint{4.596627in}{2.092573in}}{\pgfqpoint{4.607677in}{2.092573in}}%
\pgfpathclose%
\pgfusepath{stroke,fill}%
\end{pgfscope}%
\begin{pgfscope}%
\pgfpathrectangle{\pgfqpoint{0.800000in}{0.528000in}}{\pgfqpoint{4.960000in}{3.696000in}} %
\pgfusepath{clip}%
\pgfsetbuttcap%
\pgfsetroundjoin%
\definecolor{currentfill}{rgb}{0.000000,0.500000,0.000000}%
\pgfsetfillcolor{currentfill}%
\pgfsetlinewidth{1.003750pt}%
\definecolor{currentstroke}{rgb}{0.000000,0.500000,0.000000}%
\pgfsetstrokecolor{currentstroke}%
\pgfsetdash{}{0pt}%
\pgfpathmoveto{\pgfqpoint{4.657778in}{-0.139076in}}%
\pgfpathcurveto{\pgfqpoint{4.668828in}{-0.139076in}}{\pgfqpoint{4.679427in}{-0.134686in}}{\pgfqpoint{4.687241in}{-0.126873in}}%
\pgfpathcurveto{\pgfqpoint{4.695054in}{-0.119059in}}{\pgfqpoint{4.699444in}{-0.108460in}}{\pgfqpoint{4.699444in}{-0.097410in}}%
\pgfpathcurveto{\pgfqpoint{4.699444in}{-0.086360in}}{\pgfqpoint{4.695054in}{-0.075761in}}{\pgfqpoint{4.687241in}{-0.067947in}}%
\pgfpathcurveto{\pgfqpoint{4.679427in}{-0.060133in}}{\pgfqpoint{4.668828in}{-0.055743in}}{\pgfqpoint{4.657778in}{-0.055743in}}%
\pgfpathcurveto{\pgfqpoint{4.646728in}{-0.055743in}}{\pgfqpoint{4.636129in}{-0.060133in}}{\pgfqpoint{4.628315in}{-0.067947in}}%
\pgfpathcurveto{\pgfqpoint{4.620501in}{-0.075761in}}{\pgfqpoint{4.616111in}{-0.086360in}}{\pgfqpoint{4.616111in}{-0.097410in}}%
\pgfpathcurveto{\pgfqpoint{4.616111in}{-0.108460in}}{\pgfqpoint{4.620501in}{-0.119059in}}{\pgfqpoint{4.628315in}{-0.126873in}}%
\pgfpathcurveto{\pgfqpoint{4.636129in}{-0.134686in}}{\pgfqpoint{4.646728in}{-0.139076in}}{\pgfqpoint{4.657778in}{-0.139076in}}%
\pgfpathclose%
\pgfusepath{stroke,fill}%
\end{pgfscope}%
\begin{pgfscope}%
\pgfpathrectangle{\pgfqpoint{0.800000in}{0.528000in}}{\pgfqpoint{4.960000in}{3.696000in}} %
\pgfusepath{clip}%
\pgfsetbuttcap%
\pgfsetroundjoin%
\definecolor{currentfill}{rgb}{0.000000,0.500000,0.000000}%
\pgfsetfillcolor{currentfill}%
\pgfsetlinewidth{1.003750pt}%
\definecolor{currentstroke}{rgb}{0.000000,0.500000,0.000000}%
\pgfsetstrokecolor{currentstroke}%
\pgfsetdash{}{0pt}%
\pgfpathmoveto{\pgfqpoint{4.707879in}{-1.167236in}}%
\pgfpathcurveto{\pgfqpoint{4.718929in}{-1.167236in}}{\pgfqpoint{4.729528in}{-1.162845in}}{\pgfqpoint{4.737342in}{-1.155032in}}%
\pgfpathcurveto{\pgfqpoint{4.745155in}{-1.147218in}}{\pgfqpoint{4.749545in}{-1.136619in}}{\pgfqpoint{4.749545in}{-1.125569in}}%
\pgfpathcurveto{\pgfqpoint{4.749545in}{-1.114519in}}{\pgfqpoint{4.745155in}{-1.103920in}}{\pgfqpoint{4.737342in}{-1.096106in}}%
\pgfpathcurveto{\pgfqpoint{4.729528in}{-1.088293in}}{\pgfqpoint{4.718929in}{-1.083902in}}{\pgfqpoint{4.707879in}{-1.083902in}}%
\pgfpathcurveto{\pgfqpoint{4.696829in}{-1.083902in}}{\pgfqpoint{4.686230in}{-1.088293in}}{\pgfqpoint{4.678416in}{-1.096106in}}%
\pgfpathcurveto{\pgfqpoint{4.670602in}{-1.103920in}}{\pgfqpoint{4.666212in}{-1.114519in}}{\pgfqpoint{4.666212in}{-1.125569in}}%
\pgfpathcurveto{\pgfqpoint{4.666212in}{-1.136619in}}{\pgfqpoint{4.670602in}{-1.147218in}}{\pgfqpoint{4.678416in}{-1.155032in}}%
\pgfpathcurveto{\pgfqpoint{4.686230in}{-1.162845in}}{\pgfqpoint{4.696829in}{-1.167236in}}{\pgfqpoint{4.707879in}{-1.167236in}}%
\pgfpathclose%
\pgfusepath{stroke,fill}%
\end{pgfscope}%
\begin{pgfscope}%
\pgfpathrectangle{\pgfqpoint{0.800000in}{0.528000in}}{\pgfqpoint{4.960000in}{3.696000in}} %
\pgfusepath{clip}%
\pgfsetbuttcap%
\pgfsetroundjoin%
\definecolor{currentfill}{rgb}{0.000000,0.500000,0.000000}%
\pgfsetfillcolor{currentfill}%
\pgfsetlinewidth{1.003750pt}%
\definecolor{currentstroke}{rgb}{0.000000,0.500000,0.000000}%
\pgfsetstrokecolor{currentstroke}%
\pgfsetdash{}{0pt}%
\pgfpathmoveto{\pgfqpoint{4.757980in}{-0.816509in}}%
\pgfpathcurveto{\pgfqpoint{4.769030in}{-0.816509in}}{\pgfqpoint{4.779629in}{-0.812119in}}{\pgfqpoint{4.787443in}{-0.804305in}}%
\pgfpathcurveto{\pgfqpoint{4.795256in}{-0.796491in}}{\pgfqpoint{4.799646in}{-0.785892in}}{\pgfqpoint{4.799646in}{-0.774842in}}%
\pgfpathcurveto{\pgfqpoint{4.799646in}{-0.763792in}}{\pgfqpoint{4.795256in}{-0.753193in}}{\pgfqpoint{4.787443in}{-0.745379in}}%
\pgfpathcurveto{\pgfqpoint{4.779629in}{-0.737566in}}{\pgfqpoint{4.769030in}{-0.733176in}}{\pgfqpoint{4.757980in}{-0.733176in}}%
\pgfpathcurveto{\pgfqpoint{4.746930in}{-0.733176in}}{\pgfqpoint{4.736331in}{-0.737566in}}{\pgfqpoint{4.728517in}{-0.745379in}}%
\pgfpathcurveto{\pgfqpoint{4.720703in}{-0.753193in}}{\pgfqpoint{4.716313in}{-0.763792in}}{\pgfqpoint{4.716313in}{-0.774842in}}%
\pgfpathcurveto{\pgfqpoint{4.716313in}{-0.785892in}}{\pgfqpoint{4.720703in}{-0.796491in}}{\pgfqpoint{4.728517in}{-0.804305in}}%
\pgfpathcurveto{\pgfqpoint{4.736331in}{-0.812119in}}{\pgfqpoint{4.746930in}{-0.816509in}}{\pgfqpoint{4.757980in}{-0.816509in}}%
\pgfpathclose%
\pgfusepath{stroke,fill}%
\end{pgfscope}%
\begin{pgfscope}%
\pgfpathrectangle{\pgfqpoint{0.800000in}{0.528000in}}{\pgfqpoint{4.960000in}{3.696000in}} %
\pgfusepath{clip}%
\pgfsetbuttcap%
\pgfsetroundjoin%
\definecolor{currentfill}{rgb}{0.000000,0.500000,0.000000}%
\pgfsetfillcolor{currentfill}%
\pgfsetlinewidth{1.003750pt}%
\definecolor{currentstroke}{rgb}{0.000000,0.500000,0.000000}%
\pgfsetstrokecolor{currentstroke}%
\pgfsetdash{}{0pt}%
\pgfpathmoveto{\pgfqpoint{4.808081in}{2.887262in}}%
\pgfpathcurveto{\pgfqpoint{4.819131in}{2.887262in}}{\pgfqpoint{4.829730in}{2.891652in}}{\pgfqpoint{4.837544in}{2.899466in}}%
\pgfpathcurveto{\pgfqpoint{4.845357in}{2.907279in}}{\pgfqpoint{4.849747in}{2.917878in}}{\pgfqpoint{4.849747in}{2.928928in}}%
\pgfpathcurveto{\pgfqpoint{4.849747in}{2.939979in}}{\pgfqpoint{4.845357in}{2.950578in}}{\pgfqpoint{4.837544in}{2.958391in}}%
\pgfpathcurveto{\pgfqpoint{4.829730in}{2.966205in}}{\pgfqpoint{4.819131in}{2.970595in}}{\pgfqpoint{4.808081in}{2.970595in}}%
\pgfpathcurveto{\pgfqpoint{4.797031in}{2.970595in}}{\pgfqpoint{4.786432in}{2.966205in}}{\pgfqpoint{4.778618in}{2.958391in}}%
\pgfpathcurveto{\pgfqpoint{4.770804in}{2.950578in}}{\pgfqpoint{4.766414in}{2.939979in}}{\pgfqpoint{4.766414in}{2.928928in}}%
\pgfpathcurveto{\pgfqpoint{4.766414in}{2.917878in}}{\pgfqpoint{4.770804in}{2.907279in}}{\pgfqpoint{4.778618in}{2.899466in}}%
\pgfpathcurveto{\pgfqpoint{4.786432in}{2.891652in}}{\pgfqpoint{4.797031in}{2.887262in}}{\pgfqpoint{4.808081in}{2.887262in}}%
\pgfpathclose%
\pgfusepath{stroke,fill}%
\end{pgfscope}%
\begin{pgfscope}%
\pgfpathrectangle{\pgfqpoint{0.800000in}{0.528000in}}{\pgfqpoint{4.960000in}{3.696000in}} %
\pgfusepath{clip}%
\pgfsetbuttcap%
\pgfsetroundjoin%
\definecolor{currentfill}{rgb}{0.000000,0.500000,0.000000}%
\pgfsetfillcolor{currentfill}%
\pgfsetlinewidth{1.003750pt}%
\definecolor{currentstroke}{rgb}{0.000000,0.500000,0.000000}%
\pgfsetstrokecolor{currentstroke}%
\pgfsetdash{}{0pt}%
\pgfpathmoveto{\pgfqpoint{4.858182in}{1.526323in}}%
\pgfpathcurveto{\pgfqpoint{4.869232in}{1.526323in}}{\pgfqpoint{4.879831in}{1.530713in}}{\pgfqpoint{4.887645in}{1.538527in}}%
\pgfpathcurveto{\pgfqpoint{4.895458in}{1.546341in}}{\pgfqpoint{4.899848in}{1.556940in}}{\pgfqpoint{4.899848in}{1.567990in}}%
\pgfpathcurveto{\pgfqpoint{4.899848in}{1.579040in}}{\pgfqpoint{4.895458in}{1.589639in}}{\pgfqpoint{4.887645in}{1.597452in}}%
\pgfpathcurveto{\pgfqpoint{4.879831in}{1.605266in}}{\pgfqpoint{4.869232in}{1.609656in}}{\pgfqpoint{4.858182in}{1.609656in}}%
\pgfpathcurveto{\pgfqpoint{4.847132in}{1.609656in}}{\pgfqpoint{4.836533in}{1.605266in}}{\pgfqpoint{4.828719in}{1.597452in}}%
\pgfpathcurveto{\pgfqpoint{4.820905in}{1.589639in}}{\pgfqpoint{4.816515in}{1.579040in}}{\pgfqpoint{4.816515in}{1.567990in}}%
\pgfpathcurveto{\pgfqpoint{4.816515in}{1.556940in}}{\pgfqpoint{4.820905in}{1.546341in}}{\pgfqpoint{4.828719in}{1.538527in}}%
\pgfpathcurveto{\pgfqpoint{4.836533in}{1.530713in}}{\pgfqpoint{4.847132in}{1.526323in}}{\pgfqpoint{4.858182in}{1.526323in}}%
\pgfpathclose%
\pgfusepath{stroke,fill}%
\end{pgfscope}%
\begin{pgfscope}%
\pgfpathrectangle{\pgfqpoint{0.800000in}{0.528000in}}{\pgfqpoint{4.960000in}{3.696000in}} %
\pgfusepath{clip}%
\pgfsetbuttcap%
\pgfsetroundjoin%
\definecolor{currentfill}{rgb}{0.000000,0.500000,0.000000}%
\pgfsetfillcolor{currentfill}%
\pgfsetlinewidth{1.003750pt}%
\definecolor{currentstroke}{rgb}{0.000000,0.500000,0.000000}%
\pgfsetstrokecolor{currentstroke}%
\pgfsetdash{}{0pt}%
\pgfpathmoveto{\pgfqpoint{4.908283in}{-1.316729in}}%
\pgfpathcurveto{\pgfqpoint{4.919333in}{-1.316729in}}{\pgfqpoint{4.929932in}{-1.312339in}}{\pgfqpoint{4.937746in}{-1.304525in}}%
\pgfpathcurveto{\pgfqpoint{4.945559in}{-1.296711in}}{\pgfqpoint{4.949949in}{-1.286112in}}{\pgfqpoint{4.949949in}{-1.275062in}}%
\pgfpathcurveto{\pgfqpoint{4.949949in}{-1.264012in}}{\pgfqpoint{4.945559in}{-1.253413in}}{\pgfqpoint{4.937746in}{-1.245599in}}%
\pgfpathcurveto{\pgfqpoint{4.929932in}{-1.237786in}}{\pgfqpoint{4.919333in}{-1.233396in}}{\pgfqpoint{4.908283in}{-1.233396in}}%
\pgfpathcurveto{\pgfqpoint{4.897233in}{-1.233396in}}{\pgfqpoint{4.886634in}{-1.237786in}}{\pgfqpoint{4.878820in}{-1.245599in}}%
\pgfpathcurveto{\pgfqpoint{4.871006in}{-1.253413in}}{\pgfqpoint{4.866616in}{-1.264012in}}{\pgfqpoint{4.866616in}{-1.275062in}}%
\pgfpathcurveto{\pgfqpoint{4.866616in}{-1.286112in}}{\pgfqpoint{4.871006in}{-1.296711in}}{\pgfqpoint{4.878820in}{-1.304525in}}%
\pgfpathcurveto{\pgfqpoint{4.886634in}{-1.312339in}}{\pgfqpoint{4.897233in}{-1.316729in}}{\pgfqpoint{4.908283in}{-1.316729in}}%
\pgfpathclose%
\pgfusepath{stroke,fill}%
\end{pgfscope}%
\begin{pgfscope}%
\pgfpathrectangle{\pgfqpoint{0.800000in}{0.528000in}}{\pgfqpoint{4.960000in}{3.696000in}} %
\pgfusepath{clip}%
\pgfsetbuttcap%
\pgfsetroundjoin%
\definecolor{currentfill}{rgb}{0.000000,0.500000,0.000000}%
\pgfsetfillcolor{currentfill}%
\pgfsetlinewidth{1.003750pt}%
\definecolor{currentstroke}{rgb}{0.000000,0.500000,0.000000}%
\pgfsetstrokecolor{currentstroke}%
\pgfsetdash{}{0pt}%
\pgfpathmoveto{\pgfqpoint{4.958384in}{1.926007in}}%
\pgfpathcurveto{\pgfqpoint{4.969434in}{1.926007in}}{\pgfqpoint{4.980033in}{1.930397in}}{\pgfqpoint{4.987847in}{1.938211in}}%
\pgfpathcurveto{\pgfqpoint{4.995660in}{1.946024in}}{\pgfqpoint{5.000051in}{1.956623in}}{\pgfqpoint{5.000051in}{1.967674in}}%
\pgfpathcurveto{\pgfqpoint{5.000051in}{1.978724in}}{\pgfqpoint{4.995660in}{1.989323in}}{\pgfqpoint{4.987847in}{1.997136in}}%
\pgfpathcurveto{\pgfqpoint{4.980033in}{2.004950in}}{\pgfqpoint{4.969434in}{2.009340in}}{\pgfqpoint{4.958384in}{2.009340in}}%
\pgfpathcurveto{\pgfqpoint{4.947334in}{2.009340in}}{\pgfqpoint{4.936735in}{2.004950in}}{\pgfqpoint{4.928921in}{1.997136in}}%
\pgfpathcurveto{\pgfqpoint{4.921107in}{1.989323in}}{\pgfqpoint{4.916717in}{1.978724in}}{\pgfqpoint{4.916717in}{1.967674in}}%
\pgfpathcurveto{\pgfqpoint{4.916717in}{1.956623in}}{\pgfqpoint{4.921107in}{1.946024in}}{\pgfqpoint{4.928921in}{1.938211in}}%
\pgfpathcurveto{\pgfqpoint{4.936735in}{1.930397in}}{\pgfqpoint{4.947334in}{1.926007in}}{\pgfqpoint{4.958384in}{1.926007in}}%
\pgfpathclose%
\pgfusepath{stroke,fill}%
\end{pgfscope}%
\begin{pgfscope}%
\pgfpathrectangle{\pgfqpoint{0.800000in}{0.528000in}}{\pgfqpoint{4.960000in}{3.696000in}} %
\pgfusepath{clip}%
\pgfsetbuttcap%
\pgfsetroundjoin%
\definecolor{currentfill}{rgb}{0.000000,0.500000,0.000000}%
\pgfsetfillcolor{currentfill}%
\pgfsetlinewidth{1.003750pt}%
\definecolor{currentstroke}{rgb}{0.000000,0.500000,0.000000}%
\pgfsetstrokecolor{currentstroke}%
\pgfsetdash{}{0pt}%
\pgfpathmoveto{\pgfqpoint{5.008485in}{2.410626in}}%
\pgfpathcurveto{\pgfqpoint{5.019535in}{2.410626in}}{\pgfqpoint{5.030134in}{2.415016in}}{\pgfqpoint{5.037948in}{2.422830in}}%
\pgfpathcurveto{\pgfqpoint{5.045761in}{2.430643in}}{\pgfqpoint{5.050152in}{2.441242in}}{\pgfqpoint{5.050152in}{2.452293in}}%
\pgfpathcurveto{\pgfqpoint{5.050152in}{2.463343in}}{\pgfqpoint{5.045761in}{2.473942in}}{\pgfqpoint{5.037948in}{2.481755in}}%
\pgfpathcurveto{\pgfqpoint{5.030134in}{2.489569in}}{\pgfqpoint{5.019535in}{2.493959in}}{\pgfqpoint{5.008485in}{2.493959in}}%
\pgfpathcurveto{\pgfqpoint{4.997435in}{2.493959in}}{\pgfqpoint{4.986836in}{2.489569in}}{\pgfqpoint{4.979022in}{2.481755in}}%
\pgfpathcurveto{\pgfqpoint{4.971208in}{2.473942in}}{\pgfqpoint{4.966818in}{2.463343in}}{\pgfqpoint{4.966818in}{2.452293in}}%
\pgfpathcurveto{\pgfqpoint{4.966818in}{2.441242in}}{\pgfqpoint{4.971208in}{2.430643in}}{\pgfqpoint{4.979022in}{2.422830in}}%
\pgfpathcurveto{\pgfqpoint{4.986836in}{2.415016in}}{\pgfqpoint{4.997435in}{2.410626in}}{\pgfqpoint{5.008485in}{2.410626in}}%
\pgfpathclose%
\pgfusepath{stroke,fill}%
\end{pgfscope}%
\begin{pgfscope}%
\pgfpathrectangle{\pgfqpoint{0.800000in}{0.528000in}}{\pgfqpoint{4.960000in}{3.696000in}} %
\pgfusepath{clip}%
\pgfsetbuttcap%
\pgfsetroundjoin%
\definecolor{currentfill}{rgb}{0.000000,0.500000,0.000000}%
\pgfsetfillcolor{currentfill}%
\pgfsetlinewidth{1.003750pt}%
\definecolor{currentstroke}{rgb}{0.000000,0.500000,0.000000}%
\pgfsetstrokecolor{currentstroke}%
\pgfsetdash{}{0pt}%
\pgfpathmoveto{\pgfqpoint{5.058586in}{3.393339in}}%
\pgfpathcurveto{\pgfqpoint{5.069636in}{3.393339in}}{\pgfqpoint{5.080235in}{3.397729in}}{\pgfqpoint{5.088049in}{3.405543in}}%
\pgfpathcurveto{\pgfqpoint{5.095862in}{3.413357in}}{\pgfqpoint{5.100253in}{3.423956in}}{\pgfqpoint{5.100253in}{3.435006in}}%
\pgfpathcurveto{\pgfqpoint{5.100253in}{3.446056in}}{\pgfqpoint{5.095862in}{3.456655in}}{\pgfqpoint{5.088049in}{3.464468in}}%
\pgfpathcurveto{\pgfqpoint{5.080235in}{3.472282in}}{\pgfqpoint{5.069636in}{3.476672in}}{\pgfqpoint{5.058586in}{3.476672in}}%
\pgfpathcurveto{\pgfqpoint{5.047536in}{3.476672in}}{\pgfqpoint{5.036937in}{3.472282in}}{\pgfqpoint{5.029123in}{3.464468in}}%
\pgfpathcurveto{\pgfqpoint{5.021309in}{3.456655in}}{\pgfqpoint{5.016919in}{3.446056in}}{\pgfqpoint{5.016919in}{3.435006in}}%
\pgfpathcurveto{\pgfqpoint{5.016919in}{3.423956in}}{\pgfqpoint{5.021309in}{3.413357in}}{\pgfqpoint{5.029123in}{3.405543in}}%
\pgfpathcurveto{\pgfqpoint{5.036937in}{3.397729in}}{\pgfqpoint{5.047536in}{3.393339in}}{\pgfqpoint{5.058586in}{3.393339in}}%
\pgfpathclose%
\pgfusepath{stroke,fill}%
\end{pgfscope}%
\begin{pgfscope}%
\pgfpathrectangle{\pgfqpoint{0.800000in}{0.528000in}}{\pgfqpoint{4.960000in}{3.696000in}} %
\pgfusepath{clip}%
\pgfsetbuttcap%
\pgfsetroundjoin%
\definecolor{currentfill}{rgb}{0.000000,0.500000,0.000000}%
\pgfsetfillcolor{currentfill}%
\pgfsetlinewidth{1.003750pt}%
\definecolor{currentstroke}{rgb}{0.000000,0.500000,0.000000}%
\pgfsetstrokecolor{currentstroke}%
\pgfsetdash{}{0pt}%
\pgfpathmoveto{\pgfqpoint{5.108687in}{-0.259895in}}%
\pgfpathcurveto{\pgfqpoint{5.119737in}{-0.259895in}}{\pgfqpoint{5.130336in}{-0.255504in}}{\pgfqpoint{5.138150in}{-0.247691in}}%
\pgfpathcurveto{\pgfqpoint{5.145963in}{-0.239877in}}{\pgfqpoint{5.150354in}{-0.229278in}}{\pgfqpoint{5.150354in}{-0.218228in}}%
\pgfpathcurveto{\pgfqpoint{5.150354in}{-0.207178in}}{\pgfqpoint{5.145963in}{-0.196579in}}{\pgfqpoint{5.138150in}{-0.188765in}}%
\pgfpathcurveto{\pgfqpoint{5.130336in}{-0.180952in}}{\pgfqpoint{5.119737in}{-0.176561in}}{\pgfqpoint{5.108687in}{-0.176561in}}%
\pgfpathcurveto{\pgfqpoint{5.097637in}{-0.176561in}}{\pgfqpoint{5.087038in}{-0.180952in}}{\pgfqpoint{5.079224in}{-0.188765in}}%
\pgfpathcurveto{\pgfqpoint{5.071410in}{-0.196579in}}{\pgfqpoint{5.067020in}{-0.207178in}}{\pgfqpoint{5.067020in}{-0.218228in}}%
\pgfpathcurveto{\pgfqpoint{5.067020in}{-0.229278in}}{\pgfqpoint{5.071410in}{-0.239877in}}{\pgfqpoint{5.079224in}{-0.247691in}}%
\pgfpathcurveto{\pgfqpoint{5.087038in}{-0.255504in}}{\pgfqpoint{5.097637in}{-0.259895in}}{\pgfqpoint{5.108687in}{-0.259895in}}%
\pgfpathclose%
\pgfusepath{stroke,fill}%
\end{pgfscope}%
\begin{pgfscope}%
\pgfpathrectangle{\pgfqpoint{0.800000in}{0.528000in}}{\pgfqpoint{4.960000in}{3.696000in}} %
\pgfusepath{clip}%
\pgfsetbuttcap%
\pgfsetroundjoin%
\definecolor{currentfill}{rgb}{0.000000,0.500000,0.000000}%
\pgfsetfillcolor{currentfill}%
\pgfsetlinewidth{1.003750pt}%
\definecolor{currentstroke}{rgb}{0.000000,0.500000,0.000000}%
\pgfsetstrokecolor{currentstroke}%
\pgfsetdash{}{0pt}%
\pgfpathmoveto{\pgfqpoint{5.158788in}{-0.308708in}}%
\pgfpathcurveto{\pgfqpoint{5.169838in}{-0.308708in}}{\pgfqpoint{5.180437in}{-0.304318in}}{\pgfqpoint{5.188251in}{-0.296504in}}%
\pgfpathcurveto{\pgfqpoint{5.196064in}{-0.288690in}}{\pgfqpoint{5.200455in}{-0.278091in}}{\pgfqpoint{5.200455in}{-0.267041in}}%
\pgfpathcurveto{\pgfqpoint{5.200455in}{-0.255991in}}{\pgfqpoint{5.196064in}{-0.245392in}}{\pgfqpoint{5.188251in}{-0.237578in}}%
\pgfpathcurveto{\pgfqpoint{5.180437in}{-0.229765in}}{\pgfqpoint{5.169838in}{-0.225375in}}{\pgfqpoint{5.158788in}{-0.225375in}}%
\pgfpathcurveto{\pgfqpoint{5.147738in}{-0.225375in}}{\pgfqpoint{5.137139in}{-0.229765in}}{\pgfqpoint{5.129325in}{-0.237578in}}%
\pgfpathcurveto{\pgfqpoint{5.121511in}{-0.245392in}}{\pgfqpoint{5.117121in}{-0.255991in}}{\pgfqpoint{5.117121in}{-0.267041in}}%
\pgfpathcurveto{\pgfqpoint{5.117121in}{-0.278091in}}{\pgfqpoint{5.121511in}{-0.288690in}}{\pgfqpoint{5.129325in}{-0.296504in}}%
\pgfpathcurveto{\pgfqpoint{5.137139in}{-0.304318in}}{\pgfqpoint{5.147738in}{-0.308708in}}{\pgfqpoint{5.158788in}{-0.308708in}}%
\pgfpathclose%
\pgfusepath{stroke,fill}%
\end{pgfscope}%
\begin{pgfscope}%
\pgfpathrectangle{\pgfqpoint{0.800000in}{0.528000in}}{\pgfqpoint{4.960000in}{3.696000in}} %
\pgfusepath{clip}%
\pgfsetbuttcap%
\pgfsetroundjoin%
\definecolor{currentfill}{rgb}{0.000000,0.500000,0.000000}%
\pgfsetfillcolor{currentfill}%
\pgfsetlinewidth{1.003750pt}%
\definecolor{currentstroke}{rgb}{0.000000,0.500000,0.000000}%
\pgfsetstrokecolor{currentstroke}%
\pgfsetdash{}{0pt}%
\pgfpathmoveto{\pgfqpoint{5.208889in}{2.282065in}}%
\pgfpathcurveto{\pgfqpoint{5.219939in}{2.282065in}}{\pgfqpoint{5.230538in}{2.286455in}}{\pgfqpoint{5.238352in}{2.294269in}}%
\pgfpathcurveto{\pgfqpoint{5.246165in}{2.302082in}}{\pgfqpoint{5.250556in}{2.312681in}}{\pgfqpoint{5.250556in}{2.323731in}}%
\pgfpathcurveto{\pgfqpoint{5.250556in}{2.334781in}}{\pgfqpoint{5.246165in}{2.345381in}}{\pgfqpoint{5.238352in}{2.353194in}}%
\pgfpathcurveto{\pgfqpoint{5.230538in}{2.361008in}}{\pgfqpoint{5.219939in}{2.365398in}}{\pgfqpoint{5.208889in}{2.365398in}}%
\pgfpathcurveto{\pgfqpoint{5.197839in}{2.365398in}}{\pgfqpoint{5.187240in}{2.361008in}}{\pgfqpoint{5.179426in}{2.353194in}}%
\pgfpathcurveto{\pgfqpoint{5.171612in}{2.345381in}}{\pgfqpoint{5.167222in}{2.334781in}}{\pgfqpoint{5.167222in}{2.323731in}}%
\pgfpathcurveto{\pgfqpoint{5.167222in}{2.312681in}}{\pgfqpoint{5.171612in}{2.302082in}}{\pgfqpoint{5.179426in}{2.294269in}}%
\pgfpathcurveto{\pgfqpoint{5.187240in}{2.286455in}}{\pgfqpoint{5.197839in}{2.282065in}}{\pgfqpoint{5.208889in}{2.282065in}}%
\pgfpathclose%
\pgfusepath{stroke,fill}%
\end{pgfscope}%
\begin{pgfscope}%
\pgfpathrectangle{\pgfqpoint{0.800000in}{0.528000in}}{\pgfqpoint{4.960000in}{3.696000in}} %
\pgfusepath{clip}%
\pgfsetbuttcap%
\pgfsetroundjoin%
\definecolor{currentfill}{rgb}{0.000000,0.500000,0.000000}%
\pgfsetfillcolor{currentfill}%
\pgfsetlinewidth{1.003750pt}%
\definecolor{currentstroke}{rgb}{0.000000,0.500000,0.000000}%
\pgfsetstrokecolor{currentstroke}%
\pgfsetdash{}{0pt}%
\pgfpathmoveto{\pgfqpoint{5.258990in}{2.944829in}}%
\pgfpathcurveto{\pgfqpoint{5.270040in}{2.944829in}}{\pgfqpoint{5.280639in}{2.949219in}}{\pgfqpoint{5.288453in}{2.957033in}}%
\pgfpathcurveto{\pgfqpoint{5.296266in}{2.964847in}}{\pgfqpoint{5.300657in}{2.975446in}}{\pgfqpoint{5.300657in}{2.986496in}}%
\pgfpathcurveto{\pgfqpoint{5.300657in}{2.997546in}}{\pgfqpoint{5.296266in}{3.008145in}}{\pgfqpoint{5.288453in}{3.015959in}}%
\pgfpathcurveto{\pgfqpoint{5.280639in}{3.023772in}}{\pgfqpoint{5.270040in}{3.028162in}}{\pgfqpoint{5.258990in}{3.028162in}}%
\pgfpathcurveto{\pgfqpoint{5.247940in}{3.028162in}}{\pgfqpoint{5.237341in}{3.023772in}}{\pgfqpoint{5.229527in}{3.015959in}}%
\pgfpathcurveto{\pgfqpoint{5.221713in}{3.008145in}}{\pgfqpoint{5.217323in}{2.997546in}}{\pgfqpoint{5.217323in}{2.986496in}}%
\pgfpathcurveto{\pgfqpoint{5.217323in}{2.975446in}}{\pgfqpoint{5.221713in}{2.964847in}}{\pgfqpoint{5.229527in}{2.957033in}}%
\pgfpathcurveto{\pgfqpoint{5.237341in}{2.949219in}}{\pgfqpoint{5.247940in}{2.944829in}}{\pgfqpoint{5.258990in}{2.944829in}}%
\pgfpathclose%
\pgfusepath{stroke,fill}%
\end{pgfscope}%
\begin{pgfscope}%
\pgfpathrectangle{\pgfqpoint{0.800000in}{0.528000in}}{\pgfqpoint{4.960000in}{3.696000in}} %
\pgfusepath{clip}%
\pgfsetbuttcap%
\pgfsetroundjoin%
\definecolor{currentfill}{rgb}{0.000000,0.500000,0.000000}%
\pgfsetfillcolor{currentfill}%
\pgfsetlinewidth{1.003750pt}%
\definecolor{currentstroke}{rgb}{0.000000,0.500000,0.000000}%
\pgfsetstrokecolor{currentstroke}%
\pgfsetdash{}{0pt}%
\pgfpathmoveto{\pgfqpoint{5.309091in}{-0.657110in}}%
\pgfpathcurveto{\pgfqpoint{5.320141in}{-0.657110in}}{\pgfqpoint{5.330740in}{-0.652720in}}{\pgfqpoint{5.338554in}{-0.644906in}}%
\pgfpathcurveto{\pgfqpoint{5.346367in}{-0.637092in}}{\pgfqpoint{5.350758in}{-0.626493in}}{\pgfqpoint{5.350758in}{-0.615443in}}%
\pgfpathcurveto{\pgfqpoint{5.350758in}{-0.604393in}}{\pgfqpoint{5.346367in}{-0.593794in}}{\pgfqpoint{5.338554in}{-0.585980in}}%
\pgfpathcurveto{\pgfqpoint{5.330740in}{-0.578167in}}{\pgfqpoint{5.320141in}{-0.573777in}}{\pgfqpoint{5.309091in}{-0.573777in}}%
\pgfpathcurveto{\pgfqpoint{5.298041in}{-0.573777in}}{\pgfqpoint{5.287442in}{-0.578167in}}{\pgfqpoint{5.279628in}{-0.585980in}}%
\pgfpathcurveto{\pgfqpoint{5.271815in}{-0.593794in}}{\pgfqpoint{5.267424in}{-0.604393in}}{\pgfqpoint{5.267424in}{-0.615443in}}%
\pgfpathcurveto{\pgfqpoint{5.267424in}{-0.626493in}}{\pgfqpoint{5.271815in}{-0.637092in}}{\pgfqpoint{5.279628in}{-0.644906in}}%
\pgfpathcurveto{\pgfqpoint{5.287442in}{-0.652720in}}{\pgfqpoint{5.298041in}{-0.657110in}}{\pgfqpoint{5.309091in}{-0.657110in}}%
\pgfpathclose%
\pgfusepath{stroke,fill}%
\end{pgfscope}%
\begin{pgfscope}%
\pgfpathrectangle{\pgfqpoint{0.800000in}{0.528000in}}{\pgfqpoint{4.960000in}{3.696000in}} %
\pgfusepath{clip}%
\pgfsetbuttcap%
\pgfsetroundjoin%
\definecolor{currentfill}{rgb}{0.000000,0.500000,0.000000}%
\pgfsetfillcolor{currentfill}%
\pgfsetlinewidth{1.003750pt}%
\definecolor{currentstroke}{rgb}{0.000000,0.500000,0.000000}%
\pgfsetstrokecolor{currentstroke}%
\pgfsetdash{}{0pt}%
\pgfpathmoveto{\pgfqpoint{5.359192in}{2.088842in}}%
\pgfpathcurveto{\pgfqpoint{5.370242in}{2.088842in}}{\pgfqpoint{5.380841in}{2.093233in}}{\pgfqpoint{5.388655in}{2.101046in}}%
\pgfpathcurveto{\pgfqpoint{5.396468in}{2.108860in}}{\pgfqpoint{5.400859in}{2.119459in}}{\pgfqpoint{5.400859in}{2.130509in}}%
\pgfpathcurveto{\pgfqpoint{5.400859in}{2.141559in}}{\pgfqpoint{5.396468in}{2.152158in}}{\pgfqpoint{5.388655in}{2.159972in}}%
\pgfpathcurveto{\pgfqpoint{5.380841in}{2.167786in}}{\pgfqpoint{5.370242in}{2.172176in}}{\pgfqpoint{5.359192in}{2.172176in}}%
\pgfpathcurveto{\pgfqpoint{5.348142in}{2.172176in}}{\pgfqpoint{5.337543in}{2.167786in}}{\pgfqpoint{5.329729in}{2.159972in}}%
\pgfpathcurveto{\pgfqpoint{5.321916in}{2.152158in}}{\pgfqpoint{5.317525in}{2.141559in}}{\pgfqpoint{5.317525in}{2.130509in}}%
\pgfpathcurveto{\pgfqpoint{5.317525in}{2.119459in}}{\pgfqpoint{5.321916in}{2.108860in}}{\pgfqpoint{5.329729in}{2.101046in}}%
\pgfpathcurveto{\pgfqpoint{5.337543in}{2.093233in}}{\pgfqpoint{5.348142in}{2.088842in}}{\pgfqpoint{5.359192in}{2.088842in}}%
\pgfpathclose%
\pgfusepath{stroke,fill}%
\end{pgfscope}%
\begin{pgfscope}%
\pgfpathrectangle{\pgfqpoint{0.800000in}{0.528000in}}{\pgfqpoint{4.960000in}{3.696000in}} %
\pgfusepath{clip}%
\pgfsetbuttcap%
\pgfsetroundjoin%
\definecolor{currentfill}{rgb}{0.000000,0.500000,0.000000}%
\pgfsetfillcolor{currentfill}%
\pgfsetlinewidth{1.003750pt}%
\definecolor{currentstroke}{rgb}{0.000000,0.500000,0.000000}%
\pgfsetstrokecolor{currentstroke}%
\pgfsetdash{}{0pt}%
\pgfpathmoveto{\pgfqpoint{5.409293in}{1.291735in}}%
\pgfpathcurveto{\pgfqpoint{5.420343in}{1.291735in}}{\pgfqpoint{5.430942in}{1.296126in}}{\pgfqpoint{5.438756in}{1.303939in}}%
\pgfpathcurveto{\pgfqpoint{5.446569in}{1.311753in}}{\pgfqpoint{5.450960in}{1.322352in}}{\pgfqpoint{5.450960in}{1.333402in}}%
\pgfpathcurveto{\pgfqpoint{5.450960in}{1.344452in}}{\pgfqpoint{5.446569in}{1.355051in}}{\pgfqpoint{5.438756in}{1.362865in}}%
\pgfpathcurveto{\pgfqpoint{5.430942in}{1.370679in}}{\pgfqpoint{5.420343in}{1.375069in}}{\pgfqpoint{5.409293in}{1.375069in}}%
\pgfpathcurveto{\pgfqpoint{5.398243in}{1.375069in}}{\pgfqpoint{5.387644in}{1.370679in}}{\pgfqpoint{5.379830in}{1.362865in}}%
\pgfpathcurveto{\pgfqpoint{5.372017in}{1.355051in}}{\pgfqpoint{5.367626in}{1.344452in}}{\pgfqpoint{5.367626in}{1.333402in}}%
\pgfpathcurveto{\pgfqpoint{5.367626in}{1.322352in}}{\pgfqpoint{5.372017in}{1.311753in}}{\pgfqpoint{5.379830in}{1.303939in}}%
\pgfpathcurveto{\pgfqpoint{5.387644in}{1.296126in}}{\pgfqpoint{5.398243in}{1.291735in}}{\pgfqpoint{5.409293in}{1.291735in}}%
\pgfpathclose%
\pgfusepath{stroke,fill}%
\end{pgfscope}%
\begin{pgfscope}%
\pgfpathrectangle{\pgfqpoint{0.800000in}{0.528000in}}{\pgfqpoint{4.960000in}{3.696000in}} %
\pgfusepath{clip}%
\pgfsetbuttcap%
\pgfsetroundjoin%
\definecolor{currentfill}{rgb}{0.000000,0.500000,0.000000}%
\pgfsetfillcolor{currentfill}%
\pgfsetlinewidth{1.003750pt}%
\definecolor{currentstroke}{rgb}{0.000000,0.500000,0.000000}%
\pgfsetstrokecolor{currentstroke}%
\pgfsetdash{}{0pt}%
\pgfpathmoveto{\pgfqpoint{5.459394in}{2.160995in}}%
\pgfpathcurveto{\pgfqpoint{5.470444in}{2.160995in}}{\pgfqpoint{5.481043in}{2.165385in}}{\pgfqpoint{5.488857in}{2.173199in}}%
\pgfpathcurveto{\pgfqpoint{5.496670in}{2.181012in}}{\pgfqpoint{5.501061in}{2.191611in}}{\pgfqpoint{5.501061in}{2.202662in}}%
\pgfpathcurveto{\pgfqpoint{5.501061in}{2.213712in}}{\pgfqpoint{5.496670in}{2.224311in}}{\pgfqpoint{5.488857in}{2.232124in}}%
\pgfpathcurveto{\pgfqpoint{5.481043in}{2.239938in}}{\pgfqpoint{5.470444in}{2.244328in}}{\pgfqpoint{5.459394in}{2.244328in}}%
\pgfpathcurveto{\pgfqpoint{5.448344in}{2.244328in}}{\pgfqpoint{5.437745in}{2.239938in}}{\pgfqpoint{5.429931in}{2.232124in}}%
\pgfpathcurveto{\pgfqpoint{5.422118in}{2.224311in}}{\pgfqpoint{5.417727in}{2.213712in}}{\pgfqpoint{5.417727in}{2.202662in}}%
\pgfpathcurveto{\pgfqpoint{5.417727in}{2.191611in}}{\pgfqpoint{5.422118in}{2.181012in}}{\pgfqpoint{5.429931in}{2.173199in}}%
\pgfpathcurveto{\pgfqpoint{5.437745in}{2.165385in}}{\pgfqpoint{5.448344in}{2.160995in}}{\pgfqpoint{5.459394in}{2.160995in}}%
\pgfpathclose%
\pgfusepath{stroke,fill}%
\end{pgfscope}%
\begin{pgfscope}%
\pgfpathrectangle{\pgfqpoint{0.800000in}{0.528000in}}{\pgfqpoint{4.960000in}{3.696000in}} %
\pgfusepath{clip}%
\pgfsetbuttcap%
\pgfsetroundjoin%
\definecolor{currentfill}{rgb}{0.000000,0.500000,0.000000}%
\pgfsetfillcolor{currentfill}%
\pgfsetlinewidth{1.003750pt}%
\definecolor{currentstroke}{rgb}{0.000000,0.500000,0.000000}%
\pgfsetstrokecolor{currentstroke}%
\pgfsetdash{}{0pt}%
\pgfpathmoveto{\pgfqpoint{5.509495in}{2.546164in}}%
\pgfpathcurveto{\pgfqpoint{5.520545in}{2.546164in}}{\pgfqpoint{5.531144in}{2.550554in}}{\pgfqpoint{5.538958in}{2.558368in}}%
\pgfpathcurveto{\pgfqpoint{5.546771in}{2.566181in}}{\pgfqpoint{5.551162in}{2.576780in}}{\pgfqpoint{5.551162in}{2.587830in}}%
\pgfpathcurveto{\pgfqpoint{5.551162in}{2.598881in}}{\pgfqpoint{5.546771in}{2.609480in}}{\pgfqpoint{5.538958in}{2.617293in}}%
\pgfpathcurveto{\pgfqpoint{5.531144in}{2.625107in}}{\pgfqpoint{5.520545in}{2.629497in}}{\pgfqpoint{5.509495in}{2.629497in}}%
\pgfpathcurveto{\pgfqpoint{5.498445in}{2.629497in}}{\pgfqpoint{5.487846in}{2.625107in}}{\pgfqpoint{5.480032in}{2.617293in}}%
\pgfpathcurveto{\pgfqpoint{5.472219in}{2.609480in}}{\pgfqpoint{5.467828in}{2.598881in}}{\pgfqpoint{5.467828in}{2.587830in}}%
\pgfpathcurveto{\pgfqpoint{5.467828in}{2.576780in}}{\pgfqpoint{5.472219in}{2.566181in}}{\pgfqpoint{5.480032in}{2.558368in}}%
\pgfpathcurveto{\pgfqpoint{5.487846in}{2.550554in}}{\pgfqpoint{5.498445in}{2.546164in}}{\pgfqpoint{5.509495in}{2.546164in}}%
\pgfpathclose%
\pgfusepath{stroke,fill}%
\end{pgfscope}%
\begin{pgfscope}%
\pgfpathrectangle{\pgfqpoint{0.800000in}{0.528000in}}{\pgfqpoint{4.960000in}{3.696000in}} %
\pgfusepath{clip}%
\pgfsetbuttcap%
\pgfsetroundjoin%
\definecolor{currentfill}{rgb}{0.000000,0.500000,0.000000}%
\pgfsetfillcolor{currentfill}%
\pgfsetlinewidth{1.003750pt}%
\definecolor{currentstroke}{rgb}{0.000000,0.500000,0.000000}%
\pgfsetstrokecolor{currentstroke}%
\pgfsetdash{}{0pt}%
\pgfpathmoveto{\pgfqpoint{5.559596in}{2.201759in}}%
\pgfpathcurveto{\pgfqpoint{5.570646in}{2.201759in}}{\pgfqpoint{5.581245in}{2.206149in}}{\pgfqpoint{5.589059in}{2.213963in}}%
\pgfpathcurveto{\pgfqpoint{5.596872in}{2.221777in}}{\pgfqpoint{5.601263in}{2.232376in}}{\pgfqpoint{5.601263in}{2.243426in}}%
\pgfpathcurveto{\pgfqpoint{5.601263in}{2.254476in}}{\pgfqpoint{5.596872in}{2.265075in}}{\pgfqpoint{5.589059in}{2.272888in}}%
\pgfpathcurveto{\pgfqpoint{5.581245in}{2.280702in}}{\pgfqpoint{5.570646in}{2.285092in}}{\pgfqpoint{5.559596in}{2.285092in}}%
\pgfpathcurveto{\pgfqpoint{5.548546in}{2.285092in}}{\pgfqpoint{5.537947in}{2.280702in}}{\pgfqpoint{5.530133in}{2.272888in}}%
\pgfpathcurveto{\pgfqpoint{5.522320in}{2.265075in}}{\pgfqpoint{5.517929in}{2.254476in}}{\pgfqpoint{5.517929in}{2.243426in}}%
\pgfpathcurveto{\pgfqpoint{5.517929in}{2.232376in}}{\pgfqpoint{5.522320in}{2.221777in}}{\pgfqpoint{5.530133in}{2.213963in}}%
\pgfpathcurveto{\pgfqpoint{5.537947in}{2.206149in}}{\pgfqpoint{5.548546in}{2.201759in}}{\pgfqpoint{5.559596in}{2.201759in}}%
\pgfpathclose%
\pgfusepath{stroke,fill}%
\end{pgfscope}%
\begin{pgfscope}%
\pgfpathrectangle{\pgfqpoint{0.800000in}{0.528000in}}{\pgfqpoint{4.960000in}{3.696000in}} %
\pgfusepath{clip}%
\pgfsetbuttcap%
\pgfsetroundjoin%
\definecolor{currentfill}{rgb}{0.000000,0.500000,0.000000}%
\pgfsetfillcolor{currentfill}%
\pgfsetlinewidth{1.003750pt}%
\definecolor{currentstroke}{rgb}{0.000000,0.500000,0.000000}%
\pgfsetstrokecolor{currentstroke}%
\pgfsetdash{}{0pt}%
\pgfpathmoveto{\pgfqpoint{5.609697in}{0.757946in}}%
\pgfpathcurveto{\pgfqpoint{5.620747in}{0.757946in}}{\pgfqpoint{5.631346in}{0.762336in}}{\pgfqpoint{5.639160in}{0.770149in}}%
\pgfpathcurveto{\pgfqpoint{5.646973in}{0.777963in}}{\pgfqpoint{5.651364in}{0.788562in}}{\pgfqpoint{5.651364in}{0.799612in}}%
\pgfpathcurveto{\pgfqpoint{5.651364in}{0.810662in}}{\pgfqpoint{5.646973in}{0.821261in}}{\pgfqpoint{5.639160in}{0.829075in}}%
\pgfpathcurveto{\pgfqpoint{5.631346in}{0.836889in}}{\pgfqpoint{5.620747in}{0.841279in}}{\pgfqpoint{5.609697in}{0.841279in}}%
\pgfpathcurveto{\pgfqpoint{5.598647in}{0.841279in}}{\pgfqpoint{5.588048in}{0.836889in}}{\pgfqpoint{5.580234in}{0.829075in}}%
\pgfpathcurveto{\pgfqpoint{5.572421in}{0.821261in}}{\pgfqpoint{5.568030in}{0.810662in}}{\pgfqpoint{5.568030in}{0.799612in}}%
\pgfpathcurveto{\pgfqpoint{5.568030in}{0.788562in}}{\pgfqpoint{5.572421in}{0.777963in}}{\pgfqpoint{5.580234in}{0.770149in}}%
\pgfpathcurveto{\pgfqpoint{5.588048in}{0.762336in}}{\pgfqpoint{5.598647in}{0.757946in}}{\pgfqpoint{5.609697in}{0.757946in}}%
\pgfpathclose%
\pgfusepath{stroke,fill}%
\end{pgfscope}%
\begin{pgfscope}%
\pgfpathrectangle{\pgfqpoint{0.800000in}{0.528000in}}{\pgfqpoint{4.960000in}{3.696000in}} %
\pgfusepath{clip}%
\pgfsetbuttcap%
\pgfsetroundjoin%
\definecolor{currentfill}{rgb}{0.000000,0.500000,0.000000}%
\pgfsetfillcolor{currentfill}%
\pgfsetlinewidth{1.003750pt}%
\definecolor{currentstroke}{rgb}{0.000000,0.500000,0.000000}%
\pgfsetstrokecolor{currentstroke}%
\pgfsetdash{}{0pt}%
\pgfpathmoveto{\pgfqpoint{5.659798in}{0.124574in}}%
\pgfpathcurveto{\pgfqpoint{5.670848in}{0.124574in}}{\pgfqpoint{5.681447in}{0.128965in}}{\pgfqpoint{5.689261in}{0.136778in}}%
\pgfpathcurveto{\pgfqpoint{5.697074in}{0.144592in}}{\pgfqpoint{5.701465in}{0.155191in}}{\pgfqpoint{5.701465in}{0.166241in}}%
\pgfpathcurveto{\pgfqpoint{5.701465in}{0.177291in}}{\pgfqpoint{5.697074in}{0.187890in}}{\pgfqpoint{5.689261in}{0.195704in}}%
\pgfpathcurveto{\pgfqpoint{5.681447in}{0.203518in}}{\pgfqpoint{5.670848in}{0.207908in}}{\pgfqpoint{5.659798in}{0.207908in}}%
\pgfpathcurveto{\pgfqpoint{5.648748in}{0.207908in}}{\pgfqpoint{5.638149in}{0.203518in}}{\pgfqpoint{5.630335in}{0.195704in}}%
\pgfpathcurveto{\pgfqpoint{5.622522in}{0.187890in}}{\pgfqpoint{5.618131in}{0.177291in}}{\pgfqpoint{5.618131in}{0.166241in}}%
\pgfpathcurveto{\pgfqpoint{5.618131in}{0.155191in}}{\pgfqpoint{5.622522in}{0.144592in}}{\pgfqpoint{5.630335in}{0.136778in}}%
\pgfpathcurveto{\pgfqpoint{5.638149in}{0.128965in}}{\pgfqpoint{5.648748in}{0.124574in}}{\pgfqpoint{5.659798in}{0.124574in}}%
\pgfpathclose%
\pgfusepath{stroke,fill}%
\end{pgfscope}%
\begin{pgfscope}%
\pgfpathrectangle{\pgfqpoint{0.800000in}{0.528000in}}{\pgfqpoint{4.960000in}{3.696000in}} %
\pgfusepath{clip}%
\pgfsetbuttcap%
\pgfsetroundjoin%
\definecolor{currentfill}{rgb}{0.000000,0.500000,0.000000}%
\pgfsetfillcolor{currentfill}%
\pgfsetlinewidth{1.003750pt}%
\definecolor{currentstroke}{rgb}{0.000000,0.500000,0.000000}%
\pgfsetstrokecolor{currentstroke}%
\pgfsetdash{}{0pt}%
\pgfpathmoveto{\pgfqpoint{5.709899in}{0.257535in}}%
\pgfpathcurveto{\pgfqpoint{5.720949in}{0.257535in}}{\pgfqpoint{5.731548in}{0.261925in}}{\pgfqpoint{5.739362in}{0.269739in}}%
\pgfpathcurveto{\pgfqpoint{5.747175in}{0.277552in}}{\pgfqpoint{5.751566in}{0.288151in}}{\pgfqpoint{5.751566in}{0.299201in}}%
\pgfpathcurveto{\pgfqpoint{5.751566in}{0.310251in}}{\pgfqpoint{5.747175in}{0.320851in}}{\pgfqpoint{5.739362in}{0.328664in}}%
\pgfpathcurveto{\pgfqpoint{5.731548in}{0.336478in}}{\pgfqpoint{5.720949in}{0.340868in}}{\pgfqpoint{5.709899in}{0.340868in}}%
\pgfpathcurveto{\pgfqpoint{5.698849in}{0.340868in}}{\pgfqpoint{5.688250in}{0.336478in}}{\pgfqpoint{5.680436in}{0.328664in}}%
\pgfpathcurveto{\pgfqpoint{5.672623in}{0.320851in}}{\pgfqpoint{5.668232in}{0.310251in}}{\pgfqpoint{5.668232in}{0.299201in}}%
\pgfpathcurveto{\pgfqpoint{5.668232in}{0.288151in}}{\pgfqpoint{5.672623in}{0.277552in}}{\pgfqpoint{5.680436in}{0.269739in}}%
\pgfpathcurveto{\pgfqpoint{5.688250in}{0.261925in}}{\pgfqpoint{5.698849in}{0.257535in}}{\pgfqpoint{5.709899in}{0.257535in}}%
\pgfpathclose%
\pgfusepath{stroke,fill}%
\end{pgfscope}%
\begin{pgfscope}%
\pgfpathrectangle{\pgfqpoint{0.800000in}{0.528000in}}{\pgfqpoint{4.960000in}{3.696000in}} %
\pgfusepath{clip}%
\pgfsetbuttcap%
\pgfsetroundjoin%
\definecolor{currentfill}{rgb}{0.000000,0.500000,0.000000}%
\pgfsetfillcolor{currentfill}%
\pgfsetlinewidth{1.003750pt}%
\definecolor{currentstroke}{rgb}{0.000000,0.500000,0.000000}%
\pgfsetstrokecolor{currentstroke}%
\pgfsetdash{}{0pt}%
\pgfpathmoveto{\pgfqpoint{5.760000in}{-0.075474in}}%
\pgfpathcurveto{\pgfqpoint{5.771050in}{-0.075474in}}{\pgfqpoint{5.781649in}{-0.071084in}}{\pgfqpoint{5.789463in}{-0.063270in}}%
\pgfpathcurveto{\pgfqpoint{5.797276in}{-0.055456in}}{\pgfqpoint{5.801667in}{-0.044857in}}{\pgfqpoint{5.801667in}{-0.033807in}}%
\pgfpathcurveto{\pgfqpoint{5.801667in}{-0.022757in}}{\pgfqpoint{5.797276in}{-0.012158in}}{\pgfqpoint{5.789463in}{-0.004344in}}%
\pgfpathcurveto{\pgfqpoint{5.781649in}{0.003469in}}{\pgfqpoint{5.771050in}{0.007859in}}{\pgfqpoint{5.760000in}{0.007859in}}%
\pgfpathcurveto{\pgfqpoint{5.748950in}{0.007859in}}{\pgfqpoint{5.738351in}{0.003469in}}{\pgfqpoint{5.730537in}{-0.004344in}}%
\pgfpathcurveto{\pgfqpoint{5.722724in}{-0.012158in}}{\pgfqpoint{5.718333in}{-0.022757in}}{\pgfqpoint{5.718333in}{-0.033807in}}%
\pgfpathcurveto{\pgfqpoint{5.718333in}{-0.044857in}}{\pgfqpoint{5.722724in}{-0.055456in}}{\pgfqpoint{5.730537in}{-0.063270in}}%
\pgfpathcurveto{\pgfqpoint{5.738351in}{-0.071084in}}{\pgfqpoint{5.748950in}{-0.075474in}}{\pgfqpoint{5.760000in}{-0.075474in}}%
\pgfpathclose%
\pgfusepath{stroke,fill}%
\end{pgfscope}%
\begin{pgfscope}%
\pgfsetbuttcap%
\pgfsetroundjoin%
\definecolor{currentfill}{rgb}{0.000000,0.000000,0.000000}%
\pgfsetfillcolor{currentfill}%
\pgfsetlinewidth{0.803000pt}%
\definecolor{currentstroke}{rgb}{0.000000,0.000000,0.000000}%
\pgfsetstrokecolor{currentstroke}%
\pgfsetdash{}{0pt}%
\pgfsys@defobject{currentmarker}{\pgfqpoint{0.000000in}{-0.048611in}}{\pgfqpoint{0.000000in}{0.000000in}}{%
\pgfpathmoveto{\pgfqpoint{0.000000in}{0.000000in}}%
\pgfpathlineto{\pgfqpoint{0.000000in}{-0.048611in}}%
\pgfusepath{stroke,fill}%
}%
\begin{pgfscope}%
\pgfsys@transformshift{0.800000in}{0.528000in}%
\pgfsys@useobject{currentmarker}{}%
\end{pgfscope}%
\end{pgfscope}%
\begin{pgfscope}%
\pgftext[x=0.800000in,y=0.430778in,,top]{\sffamily\fontsize{10.000000}{12.000000}\selectfont 1.0}%
\end{pgfscope}%
\begin{pgfscope}%
\pgfsetbuttcap%
\pgfsetroundjoin%
\definecolor{currentfill}{rgb}{0.000000,0.000000,0.000000}%
\pgfsetfillcolor{currentfill}%
\pgfsetlinewidth{0.803000pt}%
\definecolor{currentstroke}{rgb}{0.000000,0.000000,0.000000}%
\pgfsetstrokecolor{currentstroke}%
\pgfsetdash{}{0pt}%
\pgfsys@defobject{currentmarker}{\pgfqpoint{0.000000in}{-0.048611in}}{\pgfqpoint{0.000000in}{0.000000in}}{%
\pgfpathmoveto{\pgfqpoint{0.000000in}{0.000000in}}%
\pgfpathlineto{\pgfqpoint{0.000000in}{-0.048611in}}%
\pgfusepath{stroke,fill}%
}%
\begin{pgfscope}%
\pgfsys@transformshift{1.420000in}{0.528000in}%
\pgfsys@useobject{currentmarker}{}%
\end{pgfscope}%
\end{pgfscope}%
\begin{pgfscope}%
\pgftext[x=1.420000in,y=0.430778in,,top]{\sffamily\fontsize{10.000000}{12.000000}\selectfont 1.5}%
\end{pgfscope}%
\begin{pgfscope}%
\pgfsetbuttcap%
\pgfsetroundjoin%
\definecolor{currentfill}{rgb}{0.000000,0.000000,0.000000}%
\pgfsetfillcolor{currentfill}%
\pgfsetlinewidth{0.803000pt}%
\definecolor{currentstroke}{rgb}{0.000000,0.000000,0.000000}%
\pgfsetstrokecolor{currentstroke}%
\pgfsetdash{}{0pt}%
\pgfsys@defobject{currentmarker}{\pgfqpoint{0.000000in}{-0.048611in}}{\pgfqpoint{0.000000in}{0.000000in}}{%
\pgfpathmoveto{\pgfqpoint{0.000000in}{0.000000in}}%
\pgfpathlineto{\pgfqpoint{0.000000in}{-0.048611in}}%
\pgfusepath{stroke,fill}%
}%
\begin{pgfscope}%
\pgfsys@transformshift{2.040000in}{0.528000in}%
\pgfsys@useobject{currentmarker}{}%
\end{pgfscope}%
\end{pgfscope}%
\begin{pgfscope}%
\pgftext[x=2.040000in,y=0.430778in,,top]{\sffamily\fontsize{10.000000}{12.000000}\selectfont 2.0}%
\end{pgfscope}%
\begin{pgfscope}%
\pgfsetbuttcap%
\pgfsetroundjoin%
\definecolor{currentfill}{rgb}{0.000000,0.000000,0.000000}%
\pgfsetfillcolor{currentfill}%
\pgfsetlinewidth{0.803000pt}%
\definecolor{currentstroke}{rgb}{0.000000,0.000000,0.000000}%
\pgfsetstrokecolor{currentstroke}%
\pgfsetdash{}{0pt}%
\pgfsys@defobject{currentmarker}{\pgfqpoint{0.000000in}{-0.048611in}}{\pgfqpoint{0.000000in}{0.000000in}}{%
\pgfpathmoveto{\pgfqpoint{0.000000in}{0.000000in}}%
\pgfpathlineto{\pgfqpoint{0.000000in}{-0.048611in}}%
\pgfusepath{stroke,fill}%
}%
\begin{pgfscope}%
\pgfsys@transformshift{2.660000in}{0.528000in}%
\pgfsys@useobject{currentmarker}{}%
\end{pgfscope}%
\end{pgfscope}%
\begin{pgfscope}%
\pgftext[x=2.660000in,y=0.430778in,,top]{\sffamily\fontsize{10.000000}{12.000000}\selectfont 2.5}%
\end{pgfscope}%
\begin{pgfscope}%
\pgfsetbuttcap%
\pgfsetroundjoin%
\definecolor{currentfill}{rgb}{0.000000,0.000000,0.000000}%
\pgfsetfillcolor{currentfill}%
\pgfsetlinewidth{0.803000pt}%
\definecolor{currentstroke}{rgb}{0.000000,0.000000,0.000000}%
\pgfsetstrokecolor{currentstroke}%
\pgfsetdash{}{0pt}%
\pgfsys@defobject{currentmarker}{\pgfqpoint{0.000000in}{-0.048611in}}{\pgfqpoint{0.000000in}{0.000000in}}{%
\pgfpathmoveto{\pgfqpoint{0.000000in}{0.000000in}}%
\pgfpathlineto{\pgfqpoint{0.000000in}{-0.048611in}}%
\pgfusepath{stroke,fill}%
}%
\begin{pgfscope}%
\pgfsys@transformshift{3.280000in}{0.528000in}%
\pgfsys@useobject{currentmarker}{}%
\end{pgfscope}%
\end{pgfscope}%
\begin{pgfscope}%
\pgftext[x=3.280000in,y=0.430778in,,top]{\sffamily\fontsize{10.000000}{12.000000}\selectfont 3.0}%
\end{pgfscope}%
\begin{pgfscope}%
\pgfsetbuttcap%
\pgfsetroundjoin%
\definecolor{currentfill}{rgb}{0.000000,0.000000,0.000000}%
\pgfsetfillcolor{currentfill}%
\pgfsetlinewidth{0.803000pt}%
\definecolor{currentstroke}{rgb}{0.000000,0.000000,0.000000}%
\pgfsetstrokecolor{currentstroke}%
\pgfsetdash{}{0pt}%
\pgfsys@defobject{currentmarker}{\pgfqpoint{0.000000in}{-0.048611in}}{\pgfqpoint{0.000000in}{0.000000in}}{%
\pgfpathmoveto{\pgfqpoint{0.000000in}{0.000000in}}%
\pgfpathlineto{\pgfqpoint{0.000000in}{-0.048611in}}%
\pgfusepath{stroke,fill}%
}%
\begin{pgfscope}%
\pgfsys@transformshift{3.900000in}{0.528000in}%
\pgfsys@useobject{currentmarker}{}%
\end{pgfscope}%
\end{pgfscope}%
\begin{pgfscope}%
\pgftext[x=3.900000in,y=0.430778in,,top]{\sffamily\fontsize{10.000000}{12.000000}\selectfont 3.5}%
\end{pgfscope}%
\begin{pgfscope}%
\pgfsetbuttcap%
\pgfsetroundjoin%
\definecolor{currentfill}{rgb}{0.000000,0.000000,0.000000}%
\pgfsetfillcolor{currentfill}%
\pgfsetlinewidth{0.803000pt}%
\definecolor{currentstroke}{rgb}{0.000000,0.000000,0.000000}%
\pgfsetstrokecolor{currentstroke}%
\pgfsetdash{}{0pt}%
\pgfsys@defobject{currentmarker}{\pgfqpoint{0.000000in}{-0.048611in}}{\pgfqpoint{0.000000in}{0.000000in}}{%
\pgfpathmoveto{\pgfqpoint{0.000000in}{0.000000in}}%
\pgfpathlineto{\pgfqpoint{0.000000in}{-0.048611in}}%
\pgfusepath{stroke,fill}%
}%
\begin{pgfscope}%
\pgfsys@transformshift{4.520000in}{0.528000in}%
\pgfsys@useobject{currentmarker}{}%
\end{pgfscope}%
\end{pgfscope}%
\begin{pgfscope}%
\pgftext[x=4.520000in,y=0.430778in,,top]{\sffamily\fontsize{10.000000}{12.000000}\selectfont 4.0}%
\end{pgfscope}%
\begin{pgfscope}%
\pgfsetbuttcap%
\pgfsetroundjoin%
\definecolor{currentfill}{rgb}{0.000000,0.000000,0.000000}%
\pgfsetfillcolor{currentfill}%
\pgfsetlinewidth{0.803000pt}%
\definecolor{currentstroke}{rgb}{0.000000,0.000000,0.000000}%
\pgfsetstrokecolor{currentstroke}%
\pgfsetdash{}{0pt}%
\pgfsys@defobject{currentmarker}{\pgfqpoint{0.000000in}{-0.048611in}}{\pgfqpoint{0.000000in}{0.000000in}}{%
\pgfpathmoveto{\pgfqpoint{0.000000in}{0.000000in}}%
\pgfpathlineto{\pgfqpoint{0.000000in}{-0.048611in}}%
\pgfusepath{stroke,fill}%
}%
\begin{pgfscope}%
\pgfsys@transformshift{5.140000in}{0.528000in}%
\pgfsys@useobject{currentmarker}{}%
\end{pgfscope}%
\end{pgfscope}%
\begin{pgfscope}%
\pgftext[x=5.140000in,y=0.430778in,,top]{\sffamily\fontsize{10.000000}{12.000000}\selectfont 4.5}%
\end{pgfscope}%
\begin{pgfscope}%
\pgfsetbuttcap%
\pgfsetroundjoin%
\definecolor{currentfill}{rgb}{0.000000,0.000000,0.000000}%
\pgfsetfillcolor{currentfill}%
\pgfsetlinewidth{0.803000pt}%
\definecolor{currentstroke}{rgb}{0.000000,0.000000,0.000000}%
\pgfsetstrokecolor{currentstroke}%
\pgfsetdash{}{0pt}%
\pgfsys@defobject{currentmarker}{\pgfqpoint{0.000000in}{-0.048611in}}{\pgfqpoint{0.000000in}{0.000000in}}{%
\pgfpathmoveto{\pgfqpoint{0.000000in}{0.000000in}}%
\pgfpathlineto{\pgfqpoint{0.000000in}{-0.048611in}}%
\pgfusepath{stroke,fill}%
}%
\begin{pgfscope}%
\pgfsys@transformshift{5.760000in}{0.528000in}%
\pgfsys@useobject{currentmarker}{}%
\end{pgfscope}%
\end{pgfscope}%
\begin{pgfscope}%
\pgftext[x=5.760000in,y=0.430778in,,top]{\sffamily\fontsize{10.000000}{12.000000}\selectfont 5.0}%
\end{pgfscope}%
\begin{pgfscope}%
\pgfsetbuttcap%
\pgfsetroundjoin%
\definecolor{currentfill}{rgb}{0.000000,0.000000,0.000000}%
\pgfsetfillcolor{currentfill}%
\pgfsetlinewidth{0.803000pt}%
\definecolor{currentstroke}{rgb}{0.000000,0.000000,0.000000}%
\pgfsetstrokecolor{currentstroke}%
\pgfsetdash{}{0pt}%
\pgfsys@defobject{currentmarker}{\pgfqpoint{-0.048611in}{0.000000in}}{\pgfqpoint{0.000000in}{0.000000in}}{%
\pgfpathmoveto{\pgfqpoint{0.000000in}{0.000000in}}%
\pgfpathlineto{\pgfqpoint{-0.048611in}{0.000000in}}%
\pgfusepath{stroke,fill}%
}%
\begin{pgfscope}%
\pgfsys@transformshift{0.800000in}{0.528000in}%
\pgfsys@useobject{currentmarker}{}%
\end{pgfscope}%
\end{pgfscope}%
\begin{pgfscope}%
\pgftext[x=0.614413in,y=0.475238in,left,base]{\sffamily\fontsize{10.000000}{12.000000}\selectfont 3}%
\end{pgfscope}%
\begin{pgfscope}%
\pgfsetbuttcap%
\pgfsetroundjoin%
\definecolor{currentfill}{rgb}{0.000000,0.000000,0.000000}%
\pgfsetfillcolor{currentfill}%
\pgfsetlinewidth{0.803000pt}%
\definecolor{currentstroke}{rgb}{0.000000,0.000000,0.000000}%
\pgfsetstrokecolor{currentstroke}%
\pgfsetdash{}{0pt}%
\pgfsys@defobject{currentmarker}{\pgfqpoint{-0.048611in}{0.000000in}}{\pgfqpoint{0.000000in}{0.000000in}}{%
\pgfpathmoveto{\pgfqpoint{0.000000in}{0.000000in}}%
\pgfpathlineto{\pgfqpoint{-0.048611in}{0.000000in}}%
\pgfusepath{stroke,fill}%
}%
\begin{pgfscope}%
\pgfsys@transformshift{0.800000in}{0.990000in}%
\pgfsys@useobject{currentmarker}{}%
\end{pgfscope}%
\end{pgfscope}%
\begin{pgfscope}%
\pgftext[x=0.614413in,y=0.937238in,left,base]{\sffamily\fontsize{10.000000}{12.000000}\selectfont 4}%
\end{pgfscope}%
\begin{pgfscope}%
\pgfsetbuttcap%
\pgfsetroundjoin%
\definecolor{currentfill}{rgb}{0.000000,0.000000,0.000000}%
\pgfsetfillcolor{currentfill}%
\pgfsetlinewidth{0.803000pt}%
\definecolor{currentstroke}{rgb}{0.000000,0.000000,0.000000}%
\pgfsetstrokecolor{currentstroke}%
\pgfsetdash{}{0pt}%
\pgfsys@defobject{currentmarker}{\pgfqpoint{-0.048611in}{0.000000in}}{\pgfqpoint{0.000000in}{0.000000in}}{%
\pgfpathmoveto{\pgfqpoint{0.000000in}{0.000000in}}%
\pgfpathlineto{\pgfqpoint{-0.048611in}{0.000000in}}%
\pgfusepath{stroke,fill}%
}%
\begin{pgfscope}%
\pgfsys@transformshift{0.800000in}{1.452000in}%
\pgfsys@useobject{currentmarker}{}%
\end{pgfscope}%
\end{pgfscope}%
\begin{pgfscope}%
\pgftext[x=0.614413in,y=1.399238in,left,base]{\sffamily\fontsize{10.000000}{12.000000}\selectfont 5}%
\end{pgfscope}%
\begin{pgfscope}%
\pgfsetbuttcap%
\pgfsetroundjoin%
\definecolor{currentfill}{rgb}{0.000000,0.000000,0.000000}%
\pgfsetfillcolor{currentfill}%
\pgfsetlinewidth{0.803000pt}%
\definecolor{currentstroke}{rgb}{0.000000,0.000000,0.000000}%
\pgfsetstrokecolor{currentstroke}%
\pgfsetdash{}{0pt}%
\pgfsys@defobject{currentmarker}{\pgfqpoint{-0.048611in}{0.000000in}}{\pgfqpoint{0.000000in}{0.000000in}}{%
\pgfpathmoveto{\pgfqpoint{0.000000in}{0.000000in}}%
\pgfpathlineto{\pgfqpoint{-0.048611in}{0.000000in}}%
\pgfusepath{stroke,fill}%
}%
\begin{pgfscope}%
\pgfsys@transformshift{0.800000in}{1.914000in}%
\pgfsys@useobject{currentmarker}{}%
\end{pgfscope}%
\end{pgfscope}%
\begin{pgfscope}%
\pgftext[x=0.614413in,y=1.861238in,left,base]{\sffamily\fontsize{10.000000}{12.000000}\selectfont 6}%
\end{pgfscope}%
\begin{pgfscope}%
\pgfsetbuttcap%
\pgfsetroundjoin%
\definecolor{currentfill}{rgb}{0.000000,0.000000,0.000000}%
\pgfsetfillcolor{currentfill}%
\pgfsetlinewidth{0.803000pt}%
\definecolor{currentstroke}{rgb}{0.000000,0.000000,0.000000}%
\pgfsetstrokecolor{currentstroke}%
\pgfsetdash{}{0pt}%
\pgfsys@defobject{currentmarker}{\pgfqpoint{-0.048611in}{0.000000in}}{\pgfqpoint{0.000000in}{0.000000in}}{%
\pgfpathmoveto{\pgfqpoint{0.000000in}{0.000000in}}%
\pgfpathlineto{\pgfqpoint{-0.048611in}{0.000000in}}%
\pgfusepath{stroke,fill}%
}%
\begin{pgfscope}%
\pgfsys@transformshift{0.800000in}{2.376000in}%
\pgfsys@useobject{currentmarker}{}%
\end{pgfscope}%
\end{pgfscope}%
\begin{pgfscope}%
\pgftext[x=0.614413in,y=2.323238in,left,base]{\sffamily\fontsize{10.000000}{12.000000}\selectfont 7}%
\end{pgfscope}%
\begin{pgfscope}%
\pgfsetbuttcap%
\pgfsetroundjoin%
\definecolor{currentfill}{rgb}{0.000000,0.000000,0.000000}%
\pgfsetfillcolor{currentfill}%
\pgfsetlinewidth{0.803000pt}%
\definecolor{currentstroke}{rgb}{0.000000,0.000000,0.000000}%
\pgfsetstrokecolor{currentstroke}%
\pgfsetdash{}{0pt}%
\pgfsys@defobject{currentmarker}{\pgfqpoint{-0.048611in}{0.000000in}}{\pgfqpoint{0.000000in}{0.000000in}}{%
\pgfpathmoveto{\pgfqpoint{0.000000in}{0.000000in}}%
\pgfpathlineto{\pgfqpoint{-0.048611in}{0.000000in}}%
\pgfusepath{stroke,fill}%
}%
\begin{pgfscope}%
\pgfsys@transformshift{0.800000in}{2.838000in}%
\pgfsys@useobject{currentmarker}{}%
\end{pgfscope}%
\end{pgfscope}%
\begin{pgfscope}%
\pgftext[x=0.614413in,y=2.785238in,left,base]{\sffamily\fontsize{10.000000}{12.000000}\selectfont 8}%
\end{pgfscope}%
\begin{pgfscope}%
\pgfsetbuttcap%
\pgfsetroundjoin%
\definecolor{currentfill}{rgb}{0.000000,0.000000,0.000000}%
\pgfsetfillcolor{currentfill}%
\pgfsetlinewidth{0.803000pt}%
\definecolor{currentstroke}{rgb}{0.000000,0.000000,0.000000}%
\pgfsetstrokecolor{currentstroke}%
\pgfsetdash{}{0pt}%
\pgfsys@defobject{currentmarker}{\pgfqpoint{-0.048611in}{0.000000in}}{\pgfqpoint{0.000000in}{0.000000in}}{%
\pgfpathmoveto{\pgfqpoint{0.000000in}{0.000000in}}%
\pgfpathlineto{\pgfqpoint{-0.048611in}{0.000000in}}%
\pgfusepath{stroke,fill}%
}%
\begin{pgfscope}%
\pgfsys@transformshift{0.800000in}{3.300000in}%
\pgfsys@useobject{currentmarker}{}%
\end{pgfscope}%
\end{pgfscope}%
\begin{pgfscope}%
\pgftext[x=0.614413in,y=3.247238in,left,base]{\sffamily\fontsize{10.000000}{12.000000}\selectfont 9}%
\end{pgfscope}%
\begin{pgfscope}%
\pgfsetbuttcap%
\pgfsetroundjoin%
\definecolor{currentfill}{rgb}{0.000000,0.000000,0.000000}%
\pgfsetfillcolor{currentfill}%
\pgfsetlinewidth{0.803000pt}%
\definecolor{currentstroke}{rgb}{0.000000,0.000000,0.000000}%
\pgfsetstrokecolor{currentstroke}%
\pgfsetdash{}{0pt}%
\pgfsys@defobject{currentmarker}{\pgfqpoint{-0.048611in}{0.000000in}}{\pgfqpoint{0.000000in}{0.000000in}}{%
\pgfpathmoveto{\pgfqpoint{0.000000in}{0.000000in}}%
\pgfpathlineto{\pgfqpoint{-0.048611in}{0.000000in}}%
\pgfusepath{stroke,fill}%
}%
\begin{pgfscope}%
\pgfsys@transformshift{0.800000in}{3.762000in}%
\pgfsys@useobject{currentmarker}{}%
\end{pgfscope}%
\end{pgfscope}%
\begin{pgfscope}%
\pgftext[x=0.526047in,y=3.709238in,left,base]{\sffamily\fontsize{10.000000}{12.000000}\selectfont 10}%
\end{pgfscope}%
\begin{pgfscope}%
\pgfsetbuttcap%
\pgfsetroundjoin%
\definecolor{currentfill}{rgb}{0.000000,0.000000,0.000000}%
\pgfsetfillcolor{currentfill}%
\pgfsetlinewidth{0.803000pt}%
\definecolor{currentstroke}{rgb}{0.000000,0.000000,0.000000}%
\pgfsetstrokecolor{currentstroke}%
\pgfsetdash{}{0pt}%
\pgfsys@defobject{currentmarker}{\pgfqpoint{-0.048611in}{0.000000in}}{\pgfqpoint{0.000000in}{0.000000in}}{%
\pgfpathmoveto{\pgfqpoint{0.000000in}{0.000000in}}%
\pgfpathlineto{\pgfqpoint{-0.048611in}{0.000000in}}%
\pgfusepath{stroke,fill}%
}%
\begin{pgfscope}%
\pgfsys@transformshift{0.800000in}{4.224000in}%
\pgfsys@useobject{currentmarker}{}%
\end{pgfscope}%
\end{pgfscope}%
\begin{pgfscope}%
\pgftext[x=0.526047in,y=4.171238in,left,base]{\sffamily\fontsize{10.000000}{12.000000}\selectfont 11}%
\end{pgfscope}%
\begin{pgfscope}%
\pgfpathrectangle{\pgfqpoint{0.800000in}{0.528000in}}{\pgfqpoint{4.960000in}{3.696000in}} %
\pgfusepath{clip}%
\pgfsetrectcap%
\pgfsetroundjoin%
\pgfsetlinewidth{2.007500pt}%
\definecolor{currentstroke}{rgb}{0.121569,0.466667,0.705882}%
\pgfsetstrokecolor{currentstroke}%
\pgfsetdash{}{0pt}%
\pgfpathmoveto{\pgfqpoint{0.800000in}{0.528000in}}%
\pgfpathlineto{\pgfqpoint{0.850101in}{0.565333in}}%
\pgfpathlineto{\pgfqpoint{0.900202in}{0.602667in}}%
\pgfpathlineto{\pgfqpoint{0.950303in}{0.640000in}}%
\pgfpathlineto{\pgfqpoint{1.000404in}{0.677333in}}%
\pgfpathlineto{\pgfqpoint{1.050505in}{0.714667in}}%
\pgfpathlineto{\pgfqpoint{1.100606in}{0.752000in}}%
\pgfpathlineto{\pgfqpoint{1.150707in}{0.789333in}}%
\pgfpathlineto{\pgfqpoint{1.200808in}{0.826667in}}%
\pgfpathlineto{\pgfqpoint{1.250909in}{0.864000in}}%
\pgfpathlineto{\pgfqpoint{1.301010in}{0.901333in}}%
\pgfpathlineto{\pgfqpoint{1.351111in}{0.938667in}}%
\pgfpathlineto{\pgfqpoint{1.401212in}{0.976000in}}%
\pgfpathlineto{\pgfqpoint{1.451313in}{1.013333in}}%
\pgfpathlineto{\pgfqpoint{1.501414in}{1.050667in}}%
\pgfpathlineto{\pgfqpoint{1.551515in}{1.088000in}}%
\pgfpathlineto{\pgfqpoint{1.601616in}{1.125333in}}%
\pgfpathlineto{\pgfqpoint{1.651717in}{1.162667in}}%
\pgfpathlineto{\pgfqpoint{1.701818in}{1.200000in}}%
\pgfpathlineto{\pgfqpoint{1.751919in}{1.237333in}}%
\pgfpathlineto{\pgfqpoint{1.802020in}{1.274667in}}%
\pgfpathlineto{\pgfqpoint{1.852121in}{1.312000in}}%
\pgfpathlineto{\pgfqpoint{1.902222in}{1.349333in}}%
\pgfpathlineto{\pgfqpoint{1.952323in}{1.386667in}}%
\pgfpathlineto{\pgfqpoint{2.002424in}{1.424000in}}%
\pgfpathlineto{\pgfqpoint{2.052525in}{1.461333in}}%
\pgfpathlineto{\pgfqpoint{2.102626in}{1.498667in}}%
\pgfpathlineto{\pgfqpoint{2.152727in}{1.536000in}}%
\pgfpathlineto{\pgfqpoint{2.202828in}{1.573333in}}%
\pgfpathlineto{\pgfqpoint{2.252929in}{1.610667in}}%
\pgfpathlineto{\pgfqpoint{2.303030in}{1.648000in}}%
\pgfpathlineto{\pgfqpoint{2.353131in}{1.685333in}}%
\pgfpathlineto{\pgfqpoint{2.403232in}{1.722667in}}%
\pgfpathlineto{\pgfqpoint{2.453333in}{1.760000in}}%
\pgfpathlineto{\pgfqpoint{2.503434in}{1.797333in}}%
\pgfpathlineto{\pgfqpoint{2.553535in}{1.834667in}}%
\pgfpathlineto{\pgfqpoint{2.603636in}{1.872000in}}%
\pgfpathlineto{\pgfqpoint{2.653737in}{1.909333in}}%
\pgfpathlineto{\pgfqpoint{2.703838in}{1.946667in}}%
\pgfpathlineto{\pgfqpoint{2.753939in}{1.984000in}}%
\pgfpathlineto{\pgfqpoint{2.804040in}{2.021333in}}%
\pgfpathlineto{\pgfqpoint{2.854141in}{2.058667in}}%
\pgfpathlineto{\pgfqpoint{2.904242in}{2.096000in}}%
\pgfpathlineto{\pgfqpoint{2.954343in}{2.133333in}}%
\pgfpathlineto{\pgfqpoint{3.004444in}{2.170667in}}%
\pgfpathlineto{\pgfqpoint{3.054545in}{2.208000in}}%
\pgfpathlineto{\pgfqpoint{3.104646in}{2.245333in}}%
\pgfpathlineto{\pgfqpoint{3.154747in}{2.282667in}}%
\pgfpathlineto{\pgfqpoint{3.204848in}{2.320000in}}%
\pgfpathlineto{\pgfqpoint{3.254949in}{2.357333in}}%
\pgfpathlineto{\pgfqpoint{3.305051in}{2.394667in}}%
\pgfpathlineto{\pgfqpoint{3.355152in}{2.432000in}}%
\pgfpathlineto{\pgfqpoint{3.405253in}{2.469333in}}%
\pgfpathlineto{\pgfqpoint{3.455354in}{2.506667in}}%
\pgfpathlineto{\pgfqpoint{3.505455in}{2.544000in}}%
\pgfpathlineto{\pgfqpoint{3.555556in}{2.581333in}}%
\pgfpathlineto{\pgfqpoint{3.605657in}{2.618667in}}%
\pgfpathlineto{\pgfqpoint{3.655758in}{2.656000in}}%
\pgfpathlineto{\pgfqpoint{3.705859in}{2.693333in}}%
\pgfpathlineto{\pgfqpoint{3.755960in}{2.730667in}}%
\pgfpathlineto{\pgfqpoint{3.806061in}{2.768000in}}%
\pgfpathlineto{\pgfqpoint{3.856162in}{2.805333in}}%
\pgfpathlineto{\pgfqpoint{3.906263in}{2.842667in}}%
\pgfpathlineto{\pgfqpoint{3.956364in}{2.880000in}}%
\pgfpathlineto{\pgfqpoint{4.006465in}{2.917333in}}%
\pgfpathlineto{\pgfqpoint{4.056566in}{2.954667in}}%
\pgfpathlineto{\pgfqpoint{4.106667in}{2.992000in}}%
\pgfpathlineto{\pgfqpoint{4.156768in}{3.029333in}}%
\pgfpathlineto{\pgfqpoint{4.206869in}{3.066667in}}%
\pgfpathlineto{\pgfqpoint{4.256970in}{3.104000in}}%
\pgfpathlineto{\pgfqpoint{4.307071in}{3.141333in}}%
\pgfpathlineto{\pgfqpoint{4.357172in}{3.178667in}}%
\pgfpathlineto{\pgfqpoint{4.407273in}{3.216000in}}%
\pgfpathlineto{\pgfqpoint{4.457374in}{3.253333in}}%
\pgfpathlineto{\pgfqpoint{4.507475in}{3.290667in}}%
\pgfpathlineto{\pgfqpoint{4.557576in}{3.328000in}}%
\pgfpathlineto{\pgfqpoint{4.607677in}{3.365333in}}%
\pgfpathlineto{\pgfqpoint{4.657778in}{3.402667in}}%
\pgfpathlineto{\pgfqpoint{4.707879in}{3.440000in}}%
\pgfpathlineto{\pgfqpoint{4.757980in}{3.477333in}}%
\pgfpathlineto{\pgfqpoint{4.808081in}{3.514667in}}%
\pgfpathlineto{\pgfqpoint{4.858182in}{3.552000in}}%
\pgfpathlineto{\pgfqpoint{4.908283in}{3.589333in}}%
\pgfpathlineto{\pgfqpoint{4.958384in}{3.626667in}}%
\pgfpathlineto{\pgfqpoint{5.008485in}{3.664000in}}%
\pgfpathlineto{\pgfqpoint{5.058586in}{3.701333in}}%
\pgfpathlineto{\pgfqpoint{5.108687in}{3.738667in}}%
\pgfpathlineto{\pgfqpoint{5.158788in}{3.776000in}}%
\pgfpathlineto{\pgfqpoint{5.208889in}{3.813333in}}%
\pgfpathlineto{\pgfqpoint{5.258990in}{3.850667in}}%
\pgfpathlineto{\pgfqpoint{5.309091in}{3.888000in}}%
\pgfpathlineto{\pgfqpoint{5.359192in}{3.925333in}}%
\pgfpathlineto{\pgfqpoint{5.409293in}{3.962667in}}%
\pgfpathlineto{\pgfqpoint{5.459394in}{4.000000in}}%
\pgfpathlineto{\pgfqpoint{5.509495in}{4.037333in}}%
\pgfpathlineto{\pgfqpoint{5.559596in}{4.074667in}}%
\pgfpathlineto{\pgfqpoint{5.609697in}{4.112000in}}%
\pgfpathlineto{\pgfqpoint{5.659798in}{4.149333in}}%
\pgfpathlineto{\pgfqpoint{5.709899in}{4.186667in}}%
\pgfpathlineto{\pgfqpoint{5.760000in}{4.224000in}}%
\pgfusepath{stroke}%
\end{pgfscope}%
\begin{pgfscope}%
\pgfsetrectcap%
\pgfsetmiterjoin%
\pgfsetlinewidth{0.803000pt}%
\definecolor{currentstroke}{rgb}{0.000000,0.000000,0.000000}%
\pgfsetstrokecolor{currentstroke}%
\pgfsetdash{}{0pt}%
\pgfpathmoveto{\pgfqpoint{0.800000in}{0.528000in}}%
\pgfpathlineto{\pgfqpoint{0.800000in}{4.224000in}}%
\pgfusepath{stroke}%
\end{pgfscope}%
\begin{pgfscope}%
\pgfsetrectcap%
\pgfsetmiterjoin%
\pgfsetlinewidth{0.803000pt}%
\definecolor{currentstroke}{rgb}{0.000000,0.000000,0.000000}%
\pgfsetstrokecolor{currentstroke}%
\pgfsetdash{}{0pt}%
\pgfpathmoveto{\pgfqpoint{5.760000in}{0.528000in}}%
\pgfpathlineto{\pgfqpoint{5.760000in}{4.224000in}}%
\pgfusepath{stroke}%
\end{pgfscope}%
\begin{pgfscope}%
\pgfsetrectcap%
\pgfsetmiterjoin%
\pgfsetlinewidth{0.803000pt}%
\definecolor{currentstroke}{rgb}{0.000000,0.000000,0.000000}%
\pgfsetstrokecolor{currentstroke}%
\pgfsetdash{}{0pt}%
\pgfpathmoveto{\pgfqpoint{0.800000in}{0.528000in}}%
\pgfpathlineto{\pgfqpoint{5.760000in}{0.528000in}}%
\pgfusepath{stroke}%
\end{pgfscope}%
\begin{pgfscope}%
\pgfsetrectcap%
\pgfsetmiterjoin%
\pgfsetlinewidth{0.803000pt}%
\definecolor{currentstroke}{rgb}{0.000000,0.000000,0.000000}%
\pgfsetstrokecolor{currentstroke}%
\pgfsetdash{}{0pt}%
\pgfpathmoveto{\pgfqpoint{0.800000in}{4.224000in}}%
\pgfpathlineto{\pgfqpoint{5.760000in}{4.224000in}}%
\pgfusepath{stroke}%
\end{pgfscope}%
\end{pgfpicture}%
\makeatother%
\endgroup%
}
	\caption{Binarna klasifikacija tačaka u skladu sa položajem u odnosu na pravu $2x+1$}
	\label{fig:bin_klas}
\end{figure}

\par
Regresija je zadatak predviđanja neprekidne ciljne promenljive. Na primer, cene nekretnina mogu se predvideti na osnovu površine, lokacije, populacije koja živi u komšiluku, itd. Često korišćena vrsta regresije jeste linearna regresija. U slučaju linearne regresije, podrazumeva se da je funkcija $f_w(x)$ linearna u odnosu na parametar $w$. Iako se ovo na prvi pogled čini kao prilično jako ograničenje, to nije nužno slučaj; kako za atribute ne postoji zahtev za linearnosti, oni pre pravljenja linearne kombinacije mogu biti proizvoljno transformisani. Primer linearne regresije jeste aproksimacija polinomom:
\begin{center}
	$f_w(x) = w_0 + \sum_{i=1}^{N}w_ix^i$
\end{center}

\subsection{Nenadgledano mašinsko učenje}

Nenadgledano učenje obuhvata skup problema (i njihovih rešenja) u kojima sistem prihvata ulazne podatke bez izlaznih. Ovo znači da sistem sam mora da zaključi kakve zakonitosti važe u podacima.  Kako nije moguće odrediti preciznost sistema, cilj je naći najbolji model u odnosu na neki kriterijum koji je unapred zadat.
Jedan primer nenadgledanog mašinskog učenja je klasterovanje: sistem grupiše neoznačene podatke u odnosu na  neki kriterijum. Svaka grupa (klaster) sastoji se iz podataka koji su međusobno slični i različiti od elemenata preostalih grupa u odnosu na taj kriterijum. Cilj algoritma je određivanja tog kriterijuma grupisanja. Jednostavan primer klasterovanja po numeričkim atributima $x$ i $y$ može se videti na slici \ref{fig:klaster}.
	
\begin{figure}
	\centering
	\resizebox{.8\linewidth}{!}{%% Creator: Matplotlib, PGF backend
%%
%% To include the figure in your LaTeX document, write
%%   \input{<filename>.pgf}
%%
%% Make sure the required packages are loaded in your preamble
%%   \usepackage{pgf}
%%
%% Figures using additional raster images can only be included by \input if
%% they are in the same directory as the main LaTeX file. For loading figures
%% from other directories you can use the `import` package
%%   \usepackage{import}
%% and then include the figures with
%%   \import{<path to file>}{<filename>.pgf}
%%
%% Matplotlib used the following preamble
%%   \usepackage{fontspec}
%%   \setmainfont{DejaVu Serif}
%%   \setsansfont{DejaVu Sans}
%%   \setmonofont{DejaVu Sans Mono}
%%
\begingroup%
\makeatletter%
\begin{pgfpicture}%
\pgfpathrectangle{\pgfpointorigin}{\pgfqpoint{6.400000in}{4.800000in}}%
\pgfusepath{use as bounding box, clip}%
\begin{pgfscope}%
\pgfsetbuttcap%
\pgfsetmiterjoin%
\definecolor{currentfill}{rgb}{1.000000,1.000000,1.000000}%
\pgfsetfillcolor{currentfill}%
\pgfsetlinewidth{0.000000pt}%
\definecolor{currentstroke}{rgb}{1.000000,1.000000,1.000000}%
\pgfsetstrokecolor{currentstroke}%
\pgfsetdash{}{0pt}%
\pgfpathmoveto{\pgfqpoint{0.000000in}{0.000000in}}%
\pgfpathlineto{\pgfqpoint{6.400000in}{0.000000in}}%
\pgfpathlineto{\pgfqpoint{6.400000in}{4.800000in}}%
\pgfpathlineto{\pgfqpoint{0.000000in}{4.800000in}}%
\pgfpathclose%
\pgfusepath{fill}%
\end{pgfscope}%
\begin{pgfscope}%
\pgfsetbuttcap%
\pgfsetmiterjoin%
\definecolor{currentfill}{rgb}{1.000000,1.000000,1.000000}%
\pgfsetfillcolor{currentfill}%
\pgfsetlinewidth{0.000000pt}%
\definecolor{currentstroke}{rgb}{0.000000,0.000000,0.000000}%
\pgfsetstrokecolor{currentstroke}%
\pgfsetstrokeopacity{0.000000}%
\pgfsetdash{}{0pt}%
\pgfpathmoveto{\pgfqpoint{1.115200in}{0.528000in}}%
\pgfpathlineto{\pgfqpoint{5.444800in}{0.528000in}}%
\pgfpathlineto{\pgfqpoint{5.444800in}{4.224000in}}%
\pgfpathlineto{\pgfqpoint{1.115200in}{4.224000in}}%
\pgfpathclose%
\pgfusepath{fill}%
\end{pgfscope}%
\begin{pgfscope}%
\pgfpathrectangle{\pgfqpoint{1.115200in}{0.528000in}}{\pgfqpoint{4.329600in}{3.696000in}} %
\pgfusepath{clip}%
\pgfsetbuttcap%
\pgfsetroundjoin%
\definecolor{currentfill}{rgb}{0.000000,0.501961,0.000000}%
\pgfsetfillcolor{currentfill}%
\pgfsetlinewidth{1.003750pt}%
\definecolor{currentstroke}{rgb}{0.000000,0.501961,0.000000}%
\pgfsetstrokecolor{currentstroke}%
\pgfsetdash{}{0pt}%
\pgfpathmoveto{\pgfqpoint{3.794951in}{2.912157in}}%
\pgfpathcurveto{\pgfqpoint{3.806001in}{2.912157in}}{\pgfqpoint{3.816600in}{2.916547in}}{\pgfqpoint{3.824414in}{2.924361in}}%
\pgfpathcurveto{\pgfqpoint{3.832228in}{2.932174in}}{\pgfqpoint{3.836618in}{2.942773in}}{\pgfqpoint{3.836618in}{2.953824in}}%
\pgfpathcurveto{\pgfqpoint{3.836618in}{2.964874in}}{\pgfqpoint{3.832228in}{2.975473in}}{\pgfqpoint{3.824414in}{2.983286in}}%
\pgfpathcurveto{\pgfqpoint{3.816600in}{2.991100in}}{\pgfqpoint{3.806001in}{2.995490in}}{\pgfqpoint{3.794951in}{2.995490in}}%
\pgfpathcurveto{\pgfqpoint{3.783901in}{2.995490in}}{\pgfqpoint{3.773302in}{2.991100in}}{\pgfqpoint{3.765489in}{2.983286in}}%
\pgfpathcurveto{\pgfqpoint{3.757675in}{2.975473in}}{\pgfqpoint{3.753285in}{2.964874in}}{\pgfqpoint{3.753285in}{2.953824in}}%
\pgfpathcurveto{\pgfqpoint{3.753285in}{2.942773in}}{\pgfqpoint{3.757675in}{2.932174in}}{\pgfqpoint{3.765489in}{2.924361in}}%
\pgfpathcurveto{\pgfqpoint{3.773302in}{2.916547in}}{\pgfqpoint{3.783901in}{2.912157in}}{\pgfqpoint{3.794951in}{2.912157in}}%
\pgfpathclose%
\pgfusepath{stroke,fill}%
\end{pgfscope}%
\begin{pgfscope}%
\pgfpathrectangle{\pgfqpoint{1.115200in}{0.528000in}}{\pgfqpoint{4.329600in}{3.696000in}} %
\pgfusepath{clip}%
\pgfsetbuttcap%
\pgfsetroundjoin%
\definecolor{currentfill}{rgb}{0.000000,0.501961,0.000000}%
\pgfsetfillcolor{currentfill}%
\pgfsetlinewidth{1.003750pt}%
\definecolor{currentstroke}{rgb}{0.000000,0.501961,0.000000}%
\pgfsetstrokecolor{currentstroke}%
\pgfsetdash{}{0pt}%
\pgfpathmoveto{\pgfqpoint{4.014766in}{3.578152in}}%
\pgfpathcurveto{\pgfqpoint{4.025816in}{3.578152in}}{\pgfqpoint{4.036415in}{3.582542in}}{\pgfqpoint{4.044228in}{3.590355in}}%
\pgfpathcurveto{\pgfqpoint{4.052042in}{3.598169in}}{\pgfqpoint{4.056432in}{3.608768in}}{\pgfqpoint{4.056432in}{3.619818in}}%
\pgfpathcurveto{\pgfqpoint{4.056432in}{3.630868in}}{\pgfqpoint{4.052042in}{3.641467in}}{\pgfqpoint{4.044228in}{3.649281in}}%
\pgfpathcurveto{\pgfqpoint{4.036415in}{3.657095in}}{\pgfqpoint{4.025816in}{3.661485in}}{\pgfqpoint{4.014766in}{3.661485in}}%
\pgfpathcurveto{\pgfqpoint{4.003715in}{3.661485in}}{\pgfqpoint{3.993116in}{3.657095in}}{\pgfqpoint{3.985303in}{3.649281in}}%
\pgfpathcurveto{\pgfqpoint{3.977489in}{3.641467in}}{\pgfqpoint{3.973099in}{3.630868in}}{\pgfqpoint{3.973099in}{3.619818in}}%
\pgfpathcurveto{\pgfqpoint{3.973099in}{3.608768in}}{\pgfqpoint{3.977489in}{3.598169in}}{\pgfqpoint{3.985303in}{3.590355in}}%
\pgfpathcurveto{\pgfqpoint{3.993116in}{3.582542in}}{\pgfqpoint{4.003715in}{3.578152in}}{\pgfqpoint{4.014766in}{3.578152in}}%
\pgfpathclose%
\pgfusepath{stroke,fill}%
\end{pgfscope}%
\begin{pgfscope}%
\pgfpathrectangle{\pgfqpoint{1.115200in}{0.528000in}}{\pgfqpoint{4.329600in}{3.696000in}} %
\pgfusepath{clip}%
\pgfsetbuttcap%
\pgfsetroundjoin%
\definecolor{currentfill}{rgb}{0.000000,0.501961,0.000000}%
\pgfsetfillcolor{currentfill}%
\pgfsetlinewidth{1.003750pt}%
\definecolor{currentstroke}{rgb}{0.000000,0.501961,0.000000}%
\pgfsetstrokecolor{currentstroke}%
\pgfsetdash{}{0pt}%
\pgfpathmoveto{\pgfqpoint{3.906804in}{3.141140in}}%
\pgfpathcurveto{\pgfqpoint{3.917854in}{3.141140in}}{\pgfqpoint{3.928453in}{3.145530in}}{\pgfqpoint{3.936267in}{3.153344in}}%
\pgfpathcurveto{\pgfqpoint{3.944080in}{3.161157in}}{\pgfqpoint{3.948470in}{3.171756in}}{\pgfqpoint{3.948470in}{3.182806in}}%
\pgfpathcurveto{\pgfqpoint{3.948470in}{3.193857in}}{\pgfqpoint{3.944080in}{3.204456in}}{\pgfqpoint{3.936267in}{3.212269in}}%
\pgfpathcurveto{\pgfqpoint{3.928453in}{3.220083in}}{\pgfqpoint{3.917854in}{3.224473in}}{\pgfqpoint{3.906804in}{3.224473in}}%
\pgfpathcurveto{\pgfqpoint{3.895754in}{3.224473in}}{\pgfqpoint{3.885155in}{3.220083in}}{\pgfqpoint{3.877341in}{3.212269in}}%
\pgfpathcurveto{\pgfqpoint{3.869527in}{3.204456in}}{\pgfqpoint{3.865137in}{3.193857in}}{\pgfqpoint{3.865137in}{3.182806in}}%
\pgfpathcurveto{\pgfqpoint{3.865137in}{3.171756in}}{\pgfqpoint{3.869527in}{3.161157in}}{\pgfqpoint{3.877341in}{3.153344in}}%
\pgfpathcurveto{\pgfqpoint{3.885155in}{3.145530in}}{\pgfqpoint{3.895754in}{3.141140in}}{\pgfqpoint{3.906804in}{3.141140in}}%
\pgfpathclose%
\pgfusepath{stroke,fill}%
\end{pgfscope}%
\begin{pgfscope}%
\pgfpathrectangle{\pgfqpoint{1.115200in}{0.528000in}}{\pgfqpoint{4.329600in}{3.696000in}} %
\pgfusepath{clip}%
\pgfsetbuttcap%
\pgfsetroundjoin%
\definecolor{currentfill}{rgb}{0.000000,0.501961,0.000000}%
\pgfsetfillcolor{currentfill}%
\pgfsetlinewidth{1.003750pt}%
\definecolor{currentstroke}{rgb}{0.000000,0.501961,0.000000}%
\pgfsetstrokecolor{currentstroke}%
\pgfsetdash{}{0pt}%
\pgfpathmoveto{\pgfqpoint{3.670912in}{3.787660in}}%
\pgfpathcurveto{\pgfqpoint{3.681962in}{3.787660in}}{\pgfqpoint{3.692561in}{3.792050in}}{\pgfqpoint{3.700375in}{3.799864in}}%
\pgfpathcurveto{\pgfqpoint{3.708188in}{3.807677in}}{\pgfqpoint{3.712578in}{3.818276in}}{\pgfqpoint{3.712578in}{3.829326in}}%
\pgfpathcurveto{\pgfqpoint{3.712578in}{3.840377in}}{\pgfqpoint{3.708188in}{3.850976in}}{\pgfqpoint{3.700375in}{3.858789in}}%
\pgfpathcurveto{\pgfqpoint{3.692561in}{3.866603in}}{\pgfqpoint{3.681962in}{3.870993in}}{\pgfqpoint{3.670912in}{3.870993in}}%
\pgfpathcurveto{\pgfqpoint{3.659862in}{3.870993in}}{\pgfqpoint{3.649263in}{3.866603in}}{\pgfqpoint{3.641449in}{3.858789in}}%
\pgfpathcurveto{\pgfqpoint{3.633635in}{3.850976in}}{\pgfqpoint{3.629245in}{3.840377in}}{\pgfqpoint{3.629245in}{3.829326in}}%
\pgfpathcurveto{\pgfqpoint{3.629245in}{3.818276in}}{\pgfqpoint{3.633635in}{3.807677in}}{\pgfqpoint{3.641449in}{3.799864in}}%
\pgfpathcurveto{\pgfqpoint{3.649263in}{3.792050in}}{\pgfqpoint{3.659862in}{3.787660in}}{\pgfqpoint{3.670912in}{3.787660in}}%
\pgfpathclose%
\pgfusepath{stroke,fill}%
\end{pgfscope}%
\begin{pgfscope}%
\pgfpathrectangle{\pgfqpoint{1.115200in}{0.528000in}}{\pgfqpoint{4.329600in}{3.696000in}} %
\pgfusepath{clip}%
\pgfsetbuttcap%
\pgfsetroundjoin%
\definecolor{currentfill}{rgb}{0.000000,0.501961,0.000000}%
\pgfsetfillcolor{currentfill}%
\pgfsetlinewidth{1.003750pt}%
\definecolor{currentstroke}{rgb}{0.000000,0.501961,0.000000}%
\pgfsetstrokecolor{currentstroke}%
\pgfsetdash{}{0pt}%
\pgfpathmoveto{\pgfqpoint{3.335592in}{3.359038in}}%
\pgfpathcurveto{\pgfqpoint{3.346642in}{3.359038in}}{\pgfqpoint{3.357241in}{3.363428in}}{\pgfqpoint{3.365055in}{3.371242in}}%
\pgfpathcurveto{\pgfqpoint{3.372868in}{3.379056in}}{\pgfqpoint{3.377259in}{3.389655in}}{\pgfqpoint{3.377259in}{3.400705in}}%
\pgfpathcurveto{\pgfqpoint{3.377259in}{3.411755in}}{\pgfqpoint{3.372868in}{3.422354in}}{\pgfqpoint{3.365055in}{3.430168in}}%
\pgfpathcurveto{\pgfqpoint{3.357241in}{3.437981in}}{\pgfqpoint{3.346642in}{3.442372in}}{\pgfqpoint{3.335592in}{3.442372in}}%
\pgfpathcurveto{\pgfqpoint{3.324542in}{3.442372in}}{\pgfqpoint{3.313943in}{3.437981in}}{\pgfqpoint{3.306129in}{3.430168in}}%
\pgfpathcurveto{\pgfqpoint{3.298316in}{3.422354in}}{\pgfqpoint{3.293925in}{3.411755in}}{\pgfqpoint{3.293925in}{3.400705in}}%
\pgfpathcurveto{\pgfqpoint{3.293925in}{3.389655in}}{\pgfqpoint{3.298316in}{3.379056in}}{\pgfqpoint{3.306129in}{3.371242in}}%
\pgfpathcurveto{\pgfqpoint{3.313943in}{3.363428in}}{\pgfqpoint{3.324542in}{3.359038in}}{\pgfqpoint{3.335592in}{3.359038in}}%
\pgfpathclose%
\pgfusepath{stroke,fill}%
\end{pgfscope}%
\begin{pgfscope}%
\pgfpathrectangle{\pgfqpoint{1.115200in}{0.528000in}}{\pgfqpoint{4.329600in}{3.696000in}} %
\pgfusepath{clip}%
\pgfsetbuttcap%
\pgfsetroundjoin%
\definecolor{currentfill}{rgb}{0.000000,0.501961,0.000000}%
\pgfsetfillcolor{currentfill}%
\pgfsetlinewidth{1.003750pt}%
\definecolor{currentstroke}{rgb}{0.000000,0.501961,0.000000}%
\pgfsetstrokecolor{currentstroke}%
\pgfsetdash{}{0pt}%
\pgfpathmoveto{\pgfqpoint{4.053679in}{3.160995in}}%
\pgfpathcurveto{\pgfqpoint{4.064729in}{3.160995in}}{\pgfqpoint{4.075328in}{3.165385in}}{\pgfqpoint{4.083142in}{3.173199in}}%
\pgfpathcurveto{\pgfqpoint{4.090955in}{3.181013in}}{\pgfqpoint{4.095345in}{3.191612in}}{\pgfqpoint{4.095345in}{3.202662in}}%
\pgfpathcurveto{\pgfqpoint{4.095345in}{3.213712in}}{\pgfqpoint{4.090955in}{3.224311in}}{\pgfqpoint{4.083142in}{3.232125in}}%
\pgfpathcurveto{\pgfqpoint{4.075328in}{3.239938in}}{\pgfqpoint{4.064729in}{3.244328in}}{\pgfqpoint{4.053679in}{3.244328in}}%
\pgfpathcurveto{\pgfqpoint{4.042629in}{3.244328in}}{\pgfqpoint{4.032030in}{3.239938in}}{\pgfqpoint{4.024216in}{3.232125in}}%
\pgfpathcurveto{\pgfqpoint{4.016402in}{3.224311in}}{\pgfqpoint{4.012012in}{3.213712in}}{\pgfqpoint{4.012012in}{3.202662in}}%
\pgfpathcurveto{\pgfqpoint{4.012012in}{3.191612in}}{\pgfqpoint{4.016402in}{3.181013in}}{\pgfqpoint{4.024216in}{3.173199in}}%
\pgfpathcurveto{\pgfqpoint{4.032030in}{3.165385in}}{\pgfqpoint{4.042629in}{3.160995in}}{\pgfqpoint{4.053679in}{3.160995in}}%
\pgfpathclose%
\pgfusepath{stroke,fill}%
\end{pgfscope}%
\begin{pgfscope}%
\pgfpathrectangle{\pgfqpoint{1.115200in}{0.528000in}}{\pgfqpoint{4.329600in}{3.696000in}} %
\pgfusepath{clip}%
\pgfsetbuttcap%
\pgfsetroundjoin%
\definecolor{currentfill}{rgb}{0.000000,0.501961,0.000000}%
\pgfsetfillcolor{currentfill}%
\pgfsetlinewidth{1.003750pt}%
\definecolor{currentstroke}{rgb}{0.000000,0.501961,0.000000}%
\pgfsetstrokecolor{currentstroke}%
\pgfsetdash{}{0pt}%
\pgfpathmoveto{\pgfqpoint{4.098836in}{3.493344in}}%
\pgfpathcurveto{\pgfqpoint{4.109886in}{3.493344in}}{\pgfqpoint{4.120485in}{3.497734in}}{\pgfqpoint{4.128299in}{3.505547in}}%
\pgfpathcurveto{\pgfqpoint{4.136112in}{3.513361in}}{\pgfqpoint{4.140503in}{3.523960in}}{\pgfqpoint{4.140503in}{3.535010in}}%
\pgfpathcurveto{\pgfqpoint{4.140503in}{3.546060in}}{\pgfqpoint{4.136112in}{3.556659in}}{\pgfqpoint{4.128299in}{3.564473in}}%
\pgfpathcurveto{\pgfqpoint{4.120485in}{3.572287in}}{\pgfqpoint{4.109886in}{3.576677in}}{\pgfqpoint{4.098836in}{3.576677in}}%
\pgfpathcurveto{\pgfqpoint{4.087786in}{3.576677in}}{\pgfqpoint{4.077187in}{3.572287in}}{\pgfqpoint{4.069373in}{3.564473in}}%
\pgfpathcurveto{\pgfqpoint{4.061559in}{3.556659in}}{\pgfqpoint{4.057169in}{3.546060in}}{\pgfqpoint{4.057169in}{3.535010in}}%
\pgfpathcurveto{\pgfqpoint{4.057169in}{3.523960in}}{\pgfqpoint{4.061559in}{3.513361in}}{\pgfqpoint{4.069373in}{3.505547in}}%
\pgfpathcurveto{\pgfqpoint{4.077187in}{3.497734in}}{\pgfqpoint{4.087786in}{3.493344in}}{\pgfqpoint{4.098836in}{3.493344in}}%
\pgfpathclose%
\pgfusepath{stroke,fill}%
\end{pgfscope}%
\begin{pgfscope}%
\pgfpathrectangle{\pgfqpoint{1.115200in}{0.528000in}}{\pgfqpoint{4.329600in}{3.696000in}} %
\pgfusepath{clip}%
\pgfsetbuttcap%
\pgfsetroundjoin%
\definecolor{currentfill}{rgb}{1.000000,0.000000,0.000000}%
\pgfsetfillcolor{currentfill}%
\pgfsetlinewidth{1.003750pt}%
\definecolor{currentstroke}{rgb}{1.000000,0.000000,0.000000}%
\pgfsetstrokecolor{currentstroke}%
\pgfsetdash{}{0pt}%
\pgfpathmoveto{\pgfqpoint{4.635205in}{1.486440in}}%
\pgfpathcurveto{\pgfqpoint{4.646256in}{1.486440in}}{\pgfqpoint{4.656855in}{1.490830in}}{\pgfqpoint{4.664668in}{1.498644in}}%
\pgfpathcurveto{\pgfqpoint{4.672482in}{1.506457in}}{\pgfqpoint{4.676872in}{1.517056in}}{\pgfqpoint{4.676872in}{1.528107in}}%
\pgfpathcurveto{\pgfqpoint{4.676872in}{1.539157in}}{\pgfqpoint{4.672482in}{1.549756in}}{\pgfqpoint{4.664668in}{1.557569in}}%
\pgfpathcurveto{\pgfqpoint{4.656855in}{1.565383in}}{\pgfqpoint{4.646256in}{1.569773in}}{\pgfqpoint{4.635205in}{1.569773in}}%
\pgfpathcurveto{\pgfqpoint{4.624155in}{1.569773in}}{\pgfqpoint{4.613556in}{1.565383in}}{\pgfqpoint{4.605743in}{1.557569in}}%
\pgfpathcurveto{\pgfqpoint{4.597929in}{1.549756in}}{\pgfqpoint{4.593539in}{1.539157in}}{\pgfqpoint{4.593539in}{1.528107in}}%
\pgfpathcurveto{\pgfqpoint{4.593539in}{1.517056in}}{\pgfqpoint{4.597929in}{1.506457in}}{\pgfqpoint{4.605743in}{1.498644in}}%
\pgfpathcurveto{\pgfqpoint{4.613556in}{1.490830in}}{\pgfqpoint{4.624155in}{1.486440in}}{\pgfqpoint{4.635205in}{1.486440in}}%
\pgfpathclose%
\pgfusepath{stroke,fill}%
\end{pgfscope}%
\begin{pgfscope}%
\pgfpathrectangle{\pgfqpoint{1.115200in}{0.528000in}}{\pgfqpoint{4.329600in}{3.696000in}} %
\pgfusepath{clip}%
\pgfsetbuttcap%
\pgfsetroundjoin%
\definecolor{currentfill}{rgb}{1.000000,0.000000,0.000000}%
\pgfsetfillcolor{currentfill}%
\pgfsetlinewidth{1.003750pt}%
\definecolor{currentstroke}{rgb}{1.000000,0.000000,0.000000}%
\pgfsetstrokecolor{currentstroke}%
\pgfsetdash{}{0pt}%
\pgfpathmoveto{\pgfqpoint{4.640915in}{2.354764in}}%
\pgfpathcurveto{\pgfqpoint{4.651965in}{2.354764in}}{\pgfqpoint{4.662564in}{2.359154in}}{\pgfqpoint{4.670378in}{2.366968in}}%
\pgfpathcurveto{\pgfqpoint{4.678192in}{2.374781in}}{\pgfqpoint{4.682582in}{2.385380in}}{\pgfqpoint{4.682582in}{2.396430in}}%
\pgfpathcurveto{\pgfqpoint{4.682582in}{2.407480in}}{\pgfqpoint{4.678192in}{2.418079in}}{\pgfqpoint{4.670378in}{2.425893in}}%
\pgfpathcurveto{\pgfqpoint{4.662564in}{2.433707in}}{\pgfqpoint{4.651965in}{2.438097in}}{\pgfqpoint{4.640915in}{2.438097in}}%
\pgfpathcurveto{\pgfqpoint{4.629865in}{2.438097in}}{\pgfqpoint{4.619266in}{2.433707in}}{\pgfqpoint{4.611452in}{2.425893in}}%
\pgfpathcurveto{\pgfqpoint{4.603639in}{2.418079in}}{\pgfqpoint{4.599248in}{2.407480in}}{\pgfqpoint{4.599248in}{2.396430in}}%
\pgfpathcurveto{\pgfqpoint{4.599248in}{2.385380in}}{\pgfqpoint{4.603639in}{2.374781in}}{\pgfqpoint{4.611452in}{2.366968in}}%
\pgfpathcurveto{\pgfqpoint{4.619266in}{2.359154in}}{\pgfqpoint{4.629865in}{2.354764in}}{\pgfqpoint{4.640915in}{2.354764in}}%
\pgfpathclose%
\pgfusepath{stroke,fill}%
\end{pgfscope}%
\begin{pgfscope}%
\pgfpathrectangle{\pgfqpoint{1.115200in}{0.528000in}}{\pgfqpoint{4.329600in}{3.696000in}} %
\pgfusepath{clip}%
\pgfsetbuttcap%
\pgfsetroundjoin%
\definecolor{currentfill}{rgb}{1.000000,0.000000,0.000000}%
\pgfsetfillcolor{currentfill}%
\pgfsetlinewidth{1.003750pt}%
\definecolor{currentstroke}{rgb}{1.000000,0.000000,0.000000}%
\pgfsetstrokecolor{currentstroke}%
\pgfsetdash{}{0pt}%
\pgfpathmoveto{\pgfqpoint{4.624117in}{2.084918in}}%
\pgfpathcurveto{\pgfqpoint{4.635167in}{2.084918in}}{\pgfqpoint{4.645766in}{2.089309in}}{\pgfqpoint{4.653580in}{2.097122in}}%
\pgfpathcurveto{\pgfqpoint{4.661394in}{2.104936in}}{\pgfqpoint{4.665784in}{2.115535in}}{\pgfqpoint{4.665784in}{2.126585in}}%
\pgfpathcurveto{\pgfqpoint{4.665784in}{2.137635in}}{\pgfqpoint{4.661394in}{2.148234in}}{\pgfqpoint{4.653580in}{2.156048in}}%
\pgfpathcurveto{\pgfqpoint{4.645766in}{2.163861in}}{\pgfqpoint{4.635167in}{2.168252in}}{\pgfqpoint{4.624117in}{2.168252in}}%
\pgfpathcurveto{\pgfqpoint{4.613067in}{2.168252in}}{\pgfqpoint{4.602468in}{2.163861in}}{\pgfqpoint{4.594654in}{2.156048in}}%
\pgfpathcurveto{\pgfqpoint{4.586841in}{2.148234in}}{\pgfqpoint{4.582451in}{2.137635in}}{\pgfqpoint{4.582451in}{2.126585in}}%
\pgfpathcurveto{\pgfqpoint{4.582451in}{2.115535in}}{\pgfqpoint{4.586841in}{2.104936in}}{\pgfqpoint{4.594654in}{2.097122in}}%
\pgfpathcurveto{\pgfqpoint{4.602468in}{2.089309in}}{\pgfqpoint{4.613067in}{2.084918in}}{\pgfqpoint{4.624117in}{2.084918in}}%
\pgfpathclose%
\pgfusepath{stroke,fill}%
\end{pgfscope}%
\begin{pgfscope}%
\pgfpathrectangle{\pgfqpoint{1.115200in}{0.528000in}}{\pgfqpoint{4.329600in}{3.696000in}} %
\pgfusepath{clip}%
\pgfsetbuttcap%
\pgfsetroundjoin%
\definecolor{currentfill}{rgb}{1.000000,0.000000,0.000000}%
\pgfsetfillcolor{currentfill}%
\pgfsetlinewidth{1.003750pt}%
\definecolor{currentstroke}{rgb}{1.000000,0.000000,0.000000}%
\pgfsetstrokecolor{currentstroke}%
\pgfsetdash{}{0pt}%
\pgfpathmoveto{\pgfqpoint{4.550855in}{1.801790in}}%
\pgfpathcurveto{\pgfqpoint{4.561906in}{1.801790in}}{\pgfqpoint{4.572505in}{1.806181in}}{\pgfqpoint{4.580318in}{1.813994in}}%
\pgfpathcurveto{\pgfqpoint{4.588132in}{1.821808in}}{\pgfqpoint{4.592522in}{1.832407in}}{\pgfqpoint{4.592522in}{1.843457in}}%
\pgfpathcurveto{\pgfqpoint{4.592522in}{1.854507in}}{\pgfqpoint{4.588132in}{1.865106in}}{\pgfqpoint{4.580318in}{1.872920in}}%
\pgfpathcurveto{\pgfqpoint{4.572505in}{1.880734in}}{\pgfqpoint{4.561906in}{1.885124in}}{\pgfqpoint{4.550855in}{1.885124in}}%
\pgfpathcurveto{\pgfqpoint{4.539805in}{1.885124in}}{\pgfqpoint{4.529206in}{1.880734in}}{\pgfqpoint{4.521393in}{1.872920in}}%
\pgfpathcurveto{\pgfqpoint{4.513579in}{1.865106in}}{\pgfqpoint{4.509189in}{1.854507in}}{\pgfqpoint{4.509189in}{1.843457in}}%
\pgfpathcurveto{\pgfqpoint{4.509189in}{1.832407in}}{\pgfqpoint{4.513579in}{1.821808in}}{\pgfqpoint{4.521393in}{1.813994in}}%
\pgfpathcurveto{\pgfqpoint{4.529206in}{1.806181in}}{\pgfqpoint{4.539805in}{1.801790in}}{\pgfqpoint{4.550855in}{1.801790in}}%
\pgfpathclose%
\pgfusepath{stroke,fill}%
\end{pgfscope}%
\begin{pgfscope}%
\pgfpathrectangle{\pgfqpoint{1.115200in}{0.528000in}}{\pgfqpoint{4.329600in}{3.696000in}} %
\pgfusepath{clip}%
\pgfsetbuttcap%
\pgfsetroundjoin%
\definecolor{currentfill}{rgb}{1.000000,0.000000,0.000000}%
\pgfsetfillcolor{currentfill}%
\pgfsetlinewidth{1.003750pt}%
\definecolor{currentstroke}{rgb}{1.000000,0.000000,0.000000}%
\pgfsetstrokecolor{currentstroke}%
\pgfsetdash{}{0pt}%
\pgfpathmoveto{\pgfqpoint{4.800923in}{2.028368in}}%
\pgfpathcurveto{\pgfqpoint{4.811973in}{2.028368in}}{\pgfqpoint{4.822572in}{2.032758in}}{\pgfqpoint{4.830385in}{2.040572in}}%
\pgfpathcurveto{\pgfqpoint{4.838199in}{2.048385in}}{\pgfqpoint{4.842589in}{2.058984in}}{\pgfqpoint{4.842589in}{2.070034in}}%
\pgfpathcurveto{\pgfqpoint{4.842589in}{2.081085in}}{\pgfqpoint{4.838199in}{2.091684in}}{\pgfqpoint{4.830385in}{2.099497in}}%
\pgfpathcurveto{\pgfqpoint{4.822572in}{2.107311in}}{\pgfqpoint{4.811973in}{2.111701in}}{\pgfqpoint{4.800923in}{2.111701in}}%
\pgfpathcurveto{\pgfqpoint{4.789872in}{2.111701in}}{\pgfqpoint{4.779273in}{2.107311in}}{\pgfqpoint{4.771460in}{2.099497in}}%
\pgfpathcurveto{\pgfqpoint{4.763646in}{2.091684in}}{\pgfqpoint{4.759256in}{2.081085in}}{\pgfqpoint{4.759256in}{2.070034in}}%
\pgfpathcurveto{\pgfqpoint{4.759256in}{2.058984in}}{\pgfqpoint{4.763646in}{2.048385in}}{\pgfqpoint{4.771460in}{2.040572in}}%
\pgfpathcurveto{\pgfqpoint{4.779273in}{2.032758in}}{\pgfqpoint{4.789872in}{2.028368in}}{\pgfqpoint{4.800923in}{2.028368in}}%
\pgfpathclose%
\pgfusepath{stroke,fill}%
\end{pgfscope}%
\begin{pgfscope}%
\pgfpathrectangle{\pgfqpoint{1.115200in}{0.528000in}}{\pgfqpoint{4.329600in}{3.696000in}} %
\pgfusepath{clip}%
\pgfsetbuttcap%
\pgfsetroundjoin%
\definecolor{currentfill}{rgb}{1.000000,0.000000,0.000000}%
\pgfsetfillcolor{currentfill}%
\pgfsetlinewidth{1.003750pt}%
\definecolor{currentstroke}{rgb}{1.000000,0.000000,0.000000}%
\pgfsetstrokecolor{currentstroke}%
\pgfsetdash{}{0pt}%
\pgfpathmoveto{\pgfqpoint{4.856816in}{1.673294in}}%
\pgfpathcurveto{\pgfqpoint{4.867866in}{1.673294in}}{\pgfqpoint{4.878465in}{1.677685in}}{\pgfqpoint{4.886278in}{1.685498in}}%
\pgfpathcurveto{\pgfqpoint{4.894092in}{1.693312in}}{\pgfqpoint{4.898482in}{1.703911in}}{\pgfqpoint{4.898482in}{1.714961in}}%
\pgfpathcurveto{\pgfqpoint{4.898482in}{1.726011in}}{\pgfqpoint{4.894092in}{1.736610in}}{\pgfqpoint{4.886278in}{1.744424in}}%
\pgfpathcurveto{\pgfqpoint{4.878465in}{1.752237in}}{\pgfqpoint{4.867866in}{1.756628in}}{\pgfqpoint{4.856816in}{1.756628in}}%
\pgfpathcurveto{\pgfqpoint{4.845765in}{1.756628in}}{\pgfqpoint{4.835166in}{1.752237in}}{\pgfqpoint{4.827353in}{1.744424in}}%
\pgfpathcurveto{\pgfqpoint{4.819539in}{1.736610in}}{\pgfqpoint{4.815149in}{1.726011in}}{\pgfqpoint{4.815149in}{1.714961in}}%
\pgfpathcurveto{\pgfqpoint{4.815149in}{1.703911in}}{\pgfqpoint{4.819539in}{1.693312in}}{\pgfqpoint{4.827353in}{1.685498in}}%
\pgfpathcurveto{\pgfqpoint{4.835166in}{1.677685in}}{\pgfqpoint{4.845765in}{1.673294in}}{\pgfqpoint{4.856816in}{1.673294in}}%
\pgfpathclose%
\pgfusepath{stroke,fill}%
\end{pgfscope}%
\begin{pgfscope}%
\pgfpathrectangle{\pgfqpoint{1.115200in}{0.528000in}}{\pgfqpoint{4.329600in}{3.696000in}} %
\pgfusepath{clip}%
\pgfsetbuttcap%
\pgfsetroundjoin%
\definecolor{currentfill}{rgb}{1.000000,0.000000,0.000000}%
\pgfsetfillcolor{currentfill}%
\pgfsetlinewidth{1.003750pt}%
\definecolor{currentstroke}{rgb}{1.000000,0.000000,0.000000}%
\pgfsetstrokecolor{currentstroke}%
\pgfsetdash{}{0pt}%
\pgfpathmoveto{\pgfqpoint{4.712403in}{1.389682in}}%
\pgfpathcurveto{\pgfqpoint{4.723453in}{1.389682in}}{\pgfqpoint{4.734053in}{1.394072in}}{\pgfqpoint{4.741866in}{1.401886in}}%
\pgfpathcurveto{\pgfqpoint{4.749680in}{1.409700in}}{\pgfqpoint{4.754070in}{1.420299in}}{\pgfqpoint{4.754070in}{1.431349in}}%
\pgfpathcurveto{\pgfqpoint{4.754070in}{1.442399in}}{\pgfqpoint{4.749680in}{1.452998in}}{\pgfqpoint{4.741866in}{1.460812in}}%
\pgfpathcurveto{\pgfqpoint{4.734053in}{1.468625in}}{\pgfqpoint{4.723453in}{1.473016in}}{\pgfqpoint{4.712403in}{1.473016in}}%
\pgfpathcurveto{\pgfqpoint{4.701353in}{1.473016in}}{\pgfqpoint{4.690754in}{1.468625in}}{\pgfqpoint{4.682941in}{1.460812in}}%
\pgfpathcurveto{\pgfqpoint{4.675127in}{1.452998in}}{\pgfqpoint{4.670737in}{1.442399in}}{\pgfqpoint{4.670737in}{1.431349in}}%
\pgfpathcurveto{\pgfqpoint{4.670737in}{1.420299in}}{\pgfqpoint{4.675127in}{1.409700in}}{\pgfqpoint{4.682941in}{1.401886in}}%
\pgfpathcurveto{\pgfqpoint{4.690754in}{1.394072in}}{\pgfqpoint{4.701353in}{1.389682in}}{\pgfqpoint{4.712403in}{1.389682in}}%
\pgfpathclose%
\pgfusepath{stroke,fill}%
\end{pgfscope}%
\begin{pgfscope}%
\pgfpathrectangle{\pgfqpoint{1.115200in}{0.528000in}}{\pgfqpoint{4.329600in}{3.696000in}} %
\pgfusepath{clip}%
\pgfsetbuttcap%
\pgfsetroundjoin%
\definecolor{currentfill}{rgb}{1.000000,0.000000,0.000000}%
\pgfsetfillcolor{currentfill}%
\pgfsetlinewidth{1.003750pt}%
\definecolor{currentstroke}{rgb}{1.000000,0.000000,0.000000}%
\pgfsetstrokecolor{currentstroke}%
\pgfsetdash{}{0pt}%
\pgfpathmoveto{\pgfqpoint{5.174229in}{2.130709in}}%
\pgfpathcurveto{\pgfqpoint{5.185279in}{2.130709in}}{\pgfqpoint{5.195878in}{2.135100in}}{\pgfqpoint{5.203691in}{2.142913in}}%
\pgfpathcurveto{\pgfqpoint{5.211505in}{2.150727in}}{\pgfqpoint{5.215895in}{2.161326in}}{\pgfqpoint{5.215895in}{2.172376in}}%
\pgfpathcurveto{\pgfqpoint{5.215895in}{2.183426in}}{\pgfqpoint{5.211505in}{2.194025in}}{\pgfqpoint{5.203691in}{2.201839in}}%
\pgfpathcurveto{\pgfqpoint{5.195878in}{2.209652in}}{\pgfqpoint{5.185279in}{2.214043in}}{\pgfqpoint{5.174229in}{2.214043in}}%
\pgfpathcurveto{\pgfqpoint{5.163178in}{2.214043in}}{\pgfqpoint{5.152579in}{2.209652in}}{\pgfqpoint{5.144766in}{2.201839in}}%
\pgfpathcurveto{\pgfqpoint{5.136952in}{2.194025in}}{\pgfqpoint{5.132562in}{2.183426in}}{\pgfqpoint{5.132562in}{2.172376in}}%
\pgfpathcurveto{\pgfqpoint{5.132562in}{2.161326in}}{\pgfqpoint{5.136952in}{2.150727in}}{\pgfqpoint{5.144766in}{2.142913in}}%
\pgfpathcurveto{\pgfqpoint{5.152579in}{2.135100in}}{\pgfqpoint{5.163178in}{2.130709in}}{\pgfqpoint{5.174229in}{2.130709in}}%
\pgfpathclose%
\pgfusepath{stroke,fill}%
\end{pgfscope}%
\begin{pgfscope}%
\pgfpathrectangle{\pgfqpoint{1.115200in}{0.528000in}}{\pgfqpoint{4.329600in}{3.696000in}} %
\pgfusepath{clip}%
\pgfsetbuttcap%
\pgfsetroundjoin%
\definecolor{currentfill}{rgb}{1.000000,0.000000,0.000000}%
\pgfsetfillcolor{currentfill}%
\pgfsetlinewidth{1.003750pt}%
\definecolor{currentstroke}{rgb}{1.000000,0.000000,0.000000}%
\pgfsetstrokecolor{currentstroke}%
\pgfsetdash{}{0pt}%
\pgfpathmoveto{\pgfqpoint{4.170268in}{2.157907in}}%
\pgfpathcurveto{\pgfqpoint{4.181319in}{2.157907in}}{\pgfqpoint{4.191918in}{2.162297in}}{\pgfqpoint{4.199731in}{2.170111in}}%
\pgfpathcurveto{\pgfqpoint{4.207545in}{2.177925in}}{\pgfqpoint{4.211935in}{2.188524in}}{\pgfqpoint{4.211935in}{2.199574in}}%
\pgfpathcurveto{\pgfqpoint{4.211935in}{2.210624in}}{\pgfqpoint{4.207545in}{2.221223in}}{\pgfqpoint{4.199731in}{2.229037in}}%
\pgfpathcurveto{\pgfqpoint{4.191918in}{2.236850in}}{\pgfqpoint{4.181319in}{2.241240in}}{\pgfqpoint{4.170268in}{2.241240in}}%
\pgfpathcurveto{\pgfqpoint{4.159218in}{2.241240in}}{\pgfqpoint{4.148619in}{2.236850in}}{\pgfqpoint{4.140806in}{2.229037in}}%
\pgfpathcurveto{\pgfqpoint{4.132992in}{2.221223in}}{\pgfqpoint{4.128602in}{2.210624in}}{\pgfqpoint{4.128602in}{2.199574in}}%
\pgfpathcurveto{\pgfqpoint{4.128602in}{2.188524in}}{\pgfqpoint{4.132992in}{2.177925in}}{\pgfqpoint{4.140806in}{2.170111in}}%
\pgfpathcurveto{\pgfqpoint{4.148619in}{2.162297in}}{\pgfqpoint{4.159218in}{2.157907in}}{\pgfqpoint{4.170268in}{2.157907in}}%
\pgfpathclose%
\pgfusepath{stroke,fill}%
\end{pgfscope}%
\begin{pgfscope}%
\pgfpathrectangle{\pgfqpoint{1.115200in}{0.528000in}}{\pgfqpoint{4.329600in}{3.696000in}} %
\pgfusepath{clip}%
\pgfsetbuttcap%
\pgfsetroundjoin%
\definecolor{currentfill}{rgb}{1.000000,0.000000,0.000000}%
\pgfsetfillcolor{currentfill}%
\pgfsetlinewidth{1.003750pt}%
\definecolor{currentstroke}{rgb}{1.000000,0.000000,0.000000}%
\pgfsetstrokecolor{currentstroke}%
\pgfsetdash{}{0pt}%
\pgfpathmoveto{\pgfqpoint{4.453121in}{1.895388in}}%
\pgfpathcurveto{\pgfqpoint{4.464172in}{1.895388in}}{\pgfqpoint{4.474771in}{1.899779in}}{\pgfqpoint{4.482584in}{1.907592in}}%
\pgfpathcurveto{\pgfqpoint{4.490398in}{1.915406in}}{\pgfqpoint{4.494788in}{1.926005in}}{\pgfqpoint{4.494788in}{1.937055in}}%
\pgfpathcurveto{\pgfqpoint{4.494788in}{1.948105in}}{\pgfqpoint{4.490398in}{1.958704in}}{\pgfqpoint{4.482584in}{1.966518in}}%
\pgfpathcurveto{\pgfqpoint{4.474771in}{1.974331in}}{\pgfqpoint{4.464172in}{1.978722in}}{\pgfqpoint{4.453121in}{1.978722in}}%
\pgfpathcurveto{\pgfqpoint{4.442071in}{1.978722in}}{\pgfqpoint{4.431472in}{1.974331in}}{\pgfqpoint{4.423659in}{1.966518in}}%
\pgfpathcurveto{\pgfqpoint{4.415845in}{1.958704in}}{\pgfqpoint{4.411455in}{1.948105in}}{\pgfqpoint{4.411455in}{1.937055in}}%
\pgfpathcurveto{\pgfqpoint{4.411455in}{1.926005in}}{\pgfqpoint{4.415845in}{1.915406in}}{\pgfqpoint{4.423659in}{1.907592in}}%
\pgfpathcurveto{\pgfqpoint{4.431472in}{1.899779in}}{\pgfqpoint{4.442071in}{1.895388in}}{\pgfqpoint{4.453121in}{1.895388in}}%
\pgfpathclose%
\pgfusepath{stroke,fill}%
\end{pgfscope}%
\begin{pgfscope}%
\pgfpathrectangle{\pgfqpoint{1.115200in}{0.528000in}}{\pgfqpoint{4.329600in}{3.696000in}} %
\pgfusepath{clip}%
\pgfsetbuttcap%
\pgfsetroundjoin%
\definecolor{currentfill}{rgb}{1.000000,0.000000,0.000000}%
\pgfsetfillcolor{currentfill}%
\pgfsetlinewidth{1.003750pt}%
\definecolor{currentstroke}{rgb}{1.000000,0.000000,0.000000}%
\pgfsetstrokecolor{currentstroke}%
\pgfsetdash{}{0pt}%
\pgfpathmoveto{\pgfqpoint{4.861250in}{1.852177in}}%
\pgfpathcurveto{\pgfqpoint{4.872300in}{1.852177in}}{\pgfqpoint{4.882899in}{1.856568in}}{\pgfqpoint{4.890713in}{1.864381in}}%
\pgfpathcurveto{\pgfqpoint{4.898527in}{1.872195in}}{\pgfqpoint{4.902917in}{1.882794in}}{\pgfqpoint{4.902917in}{1.893844in}}%
\pgfpathcurveto{\pgfqpoint{4.902917in}{1.904894in}}{\pgfqpoint{4.898527in}{1.915493in}}{\pgfqpoint{4.890713in}{1.923307in}}%
\pgfpathcurveto{\pgfqpoint{4.882899in}{1.931121in}}{\pgfqpoint{4.872300in}{1.935511in}}{\pgfqpoint{4.861250in}{1.935511in}}%
\pgfpathcurveto{\pgfqpoint{4.850200in}{1.935511in}}{\pgfqpoint{4.839601in}{1.931121in}}{\pgfqpoint{4.831788in}{1.923307in}}%
\pgfpathcurveto{\pgfqpoint{4.823974in}{1.915493in}}{\pgfqpoint{4.819584in}{1.904894in}}{\pgfqpoint{4.819584in}{1.893844in}}%
\pgfpathcurveto{\pgfqpoint{4.819584in}{1.882794in}}{\pgfqpoint{4.823974in}{1.872195in}}{\pgfqpoint{4.831788in}{1.864381in}}%
\pgfpathcurveto{\pgfqpoint{4.839601in}{1.856568in}}{\pgfqpoint{4.850200in}{1.852177in}}{\pgfqpoint{4.861250in}{1.852177in}}%
\pgfpathclose%
\pgfusepath{stroke,fill}%
\end{pgfscope}%
\begin{pgfscope}%
\pgfpathrectangle{\pgfqpoint{1.115200in}{0.528000in}}{\pgfqpoint{4.329600in}{3.696000in}} %
\pgfusepath{clip}%
\pgfsetbuttcap%
\pgfsetroundjoin%
\definecolor{currentfill}{rgb}{0.000000,0.000000,1.000000}%
\pgfsetfillcolor{currentfill}%
\pgfsetlinewidth{1.003750pt}%
\definecolor{currentstroke}{rgb}{0.000000,0.000000,1.000000}%
\pgfsetstrokecolor{currentstroke}%
\pgfsetdash{}{0pt}%
\pgfpathmoveto{\pgfqpoint{1.834260in}{1.780375in}}%
\pgfpathcurveto{\pgfqpoint{1.845310in}{1.780375in}}{\pgfqpoint{1.855909in}{1.784765in}}{\pgfqpoint{1.863723in}{1.792578in}}%
\pgfpathcurveto{\pgfqpoint{1.871537in}{1.800392in}}{\pgfqpoint{1.875927in}{1.810991in}}{\pgfqpoint{1.875927in}{1.822041in}}%
\pgfpathcurveto{\pgfqpoint{1.875927in}{1.833091in}}{\pgfqpoint{1.871537in}{1.843690in}}{\pgfqpoint{1.863723in}{1.851504in}}%
\pgfpathcurveto{\pgfqpoint{1.855909in}{1.859318in}}{\pgfqpoint{1.845310in}{1.863708in}}{\pgfqpoint{1.834260in}{1.863708in}}%
\pgfpathcurveto{\pgfqpoint{1.823210in}{1.863708in}}{\pgfqpoint{1.812611in}{1.859318in}}{\pgfqpoint{1.804798in}{1.851504in}}%
\pgfpathcurveto{\pgfqpoint{1.796984in}{1.843690in}}{\pgfqpoint{1.792594in}{1.833091in}}{\pgfqpoint{1.792594in}{1.822041in}}%
\pgfpathcurveto{\pgfqpoint{1.792594in}{1.810991in}}{\pgfqpoint{1.796984in}{1.800392in}}{\pgfqpoint{1.804798in}{1.792578in}}%
\pgfpathcurveto{\pgfqpoint{1.812611in}{1.784765in}}{\pgfqpoint{1.823210in}{1.780375in}}{\pgfqpoint{1.834260in}{1.780375in}}%
\pgfpathclose%
\pgfusepath{stroke,fill}%
\end{pgfscope}%
\begin{pgfscope}%
\pgfpathrectangle{\pgfqpoint{1.115200in}{0.528000in}}{\pgfqpoint{4.329600in}{3.696000in}} %
\pgfusepath{clip}%
\pgfsetbuttcap%
\pgfsetroundjoin%
\definecolor{currentfill}{rgb}{0.000000,0.000000,1.000000}%
\pgfsetfillcolor{currentfill}%
\pgfsetlinewidth{1.003750pt}%
\definecolor{currentstroke}{rgb}{0.000000,0.000000,1.000000}%
\pgfsetstrokecolor{currentstroke}%
\pgfsetdash{}{0pt}%
\pgfpathmoveto{\pgfqpoint{2.039030in}{1.448013in}}%
\pgfpathcurveto{\pgfqpoint{2.050080in}{1.448013in}}{\pgfqpoint{2.060679in}{1.452403in}}{\pgfqpoint{2.068493in}{1.460216in}}%
\pgfpathcurveto{\pgfqpoint{2.076306in}{1.468030in}}{\pgfqpoint{2.080697in}{1.478629in}}{\pgfqpoint{2.080697in}{1.489679in}}%
\pgfpathcurveto{\pgfqpoint{2.080697in}{1.500729in}}{\pgfqpoint{2.076306in}{1.511328in}}{\pgfqpoint{2.068493in}{1.519142in}}%
\pgfpathcurveto{\pgfqpoint{2.060679in}{1.526956in}}{\pgfqpoint{2.050080in}{1.531346in}}{\pgfqpoint{2.039030in}{1.531346in}}%
\pgfpathcurveto{\pgfqpoint{2.027980in}{1.531346in}}{\pgfqpoint{2.017381in}{1.526956in}}{\pgfqpoint{2.009567in}{1.519142in}}%
\pgfpathcurveto{\pgfqpoint{2.001753in}{1.511328in}}{\pgfqpoint{1.997363in}{1.500729in}}{\pgfqpoint{1.997363in}{1.489679in}}%
\pgfpathcurveto{\pgfqpoint{1.997363in}{1.478629in}}{\pgfqpoint{2.001753in}{1.468030in}}{\pgfqpoint{2.009567in}{1.460216in}}%
\pgfpathcurveto{\pgfqpoint{2.017381in}{1.452403in}}{\pgfqpoint{2.027980in}{1.448013in}}{\pgfqpoint{2.039030in}{1.448013in}}%
\pgfpathclose%
\pgfusepath{stroke,fill}%
\end{pgfscope}%
\begin{pgfscope}%
\pgfpathrectangle{\pgfqpoint{1.115200in}{0.528000in}}{\pgfqpoint{4.329600in}{3.696000in}} %
\pgfusepath{clip}%
\pgfsetbuttcap%
\pgfsetroundjoin%
\definecolor{currentfill}{rgb}{0.000000,0.000000,1.000000}%
\pgfsetfillcolor{currentfill}%
\pgfsetlinewidth{1.003750pt}%
\definecolor{currentstroke}{rgb}{0.000000,0.000000,1.000000}%
\pgfsetstrokecolor{currentstroke}%
\pgfsetdash{}{0pt}%
\pgfpathmoveto{\pgfqpoint{2.097463in}{1.626023in}}%
\pgfpathcurveto{\pgfqpoint{2.108514in}{1.626023in}}{\pgfqpoint{2.119113in}{1.630413in}}{\pgfqpoint{2.126926in}{1.638226in}}%
\pgfpathcurveto{\pgfqpoint{2.134740in}{1.646040in}}{\pgfqpoint{2.139130in}{1.656639in}}{\pgfqpoint{2.139130in}{1.667689in}}%
\pgfpathcurveto{\pgfqpoint{2.139130in}{1.678739in}}{\pgfqpoint{2.134740in}{1.689338in}}{\pgfqpoint{2.126926in}{1.697152in}}%
\pgfpathcurveto{\pgfqpoint{2.119113in}{1.704966in}}{\pgfqpoint{2.108514in}{1.709356in}}{\pgfqpoint{2.097463in}{1.709356in}}%
\pgfpathcurveto{\pgfqpoint{2.086413in}{1.709356in}}{\pgfqpoint{2.075814in}{1.704966in}}{\pgfqpoint{2.068001in}{1.697152in}}%
\pgfpathcurveto{\pgfqpoint{2.060187in}{1.689338in}}{\pgfqpoint{2.055797in}{1.678739in}}{\pgfqpoint{2.055797in}{1.667689in}}%
\pgfpathcurveto{\pgfqpoint{2.055797in}{1.656639in}}{\pgfqpoint{2.060187in}{1.646040in}}{\pgfqpoint{2.068001in}{1.638226in}}%
\pgfpathcurveto{\pgfqpoint{2.075814in}{1.630413in}}{\pgfqpoint{2.086413in}{1.626023in}}{\pgfqpoint{2.097463in}{1.626023in}}%
\pgfpathclose%
\pgfusepath{stroke,fill}%
\end{pgfscope}%
\begin{pgfscope}%
\pgfpathrectangle{\pgfqpoint{1.115200in}{0.528000in}}{\pgfqpoint{4.329600in}{3.696000in}} %
\pgfusepath{clip}%
\pgfsetbuttcap%
\pgfsetroundjoin%
\definecolor{currentfill}{rgb}{0.000000,0.000000,1.000000}%
\pgfsetfillcolor{currentfill}%
\pgfsetlinewidth{1.003750pt}%
\definecolor{currentstroke}{rgb}{0.000000,0.000000,1.000000}%
\pgfsetstrokecolor{currentstroke}%
\pgfsetdash{}{0pt}%
\pgfpathmoveto{\pgfqpoint{2.208220in}{1.691202in}}%
\pgfpathcurveto{\pgfqpoint{2.219270in}{1.691202in}}{\pgfqpoint{2.229869in}{1.695592in}}{\pgfqpoint{2.237683in}{1.703405in}}%
\pgfpathcurveto{\pgfqpoint{2.245496in}{1.711219in}}{\pgfqpoint{2.249886in}{1.721818in}}{\pgfqpoint{2.249886in}{1.732868in}}%
\pgfpathcurveto{\pgfqpoint{2.249886in}{1.743918in}}{\pgfqpoint{2.245496in}{1.754517in}}{\pgfqpoint{2.237683in}{1.762331in}}%
\pgfpathcurveto{\pgfqpoint{2.229869in}{1.770145in}}{\pgfqpoint{2.219270in}{1.774535in}}{\pgfqpoint{2.208220in}{1.774535in}}%
\pgfpathcurveto{\pgfqpoint{2.197170in}{1.774535in}}{\pgfqpoint{2.186571in}{1.770145in}}{\pgfqpoint{2.178757in}{1.762331in}}%
\pgfpathcurveto{\pgfqpoint{2.170943in}{1.754517in}}{\pgfqpoint{2.166553in}{1.743918in}}{\pgfqpoint{2.166553in}{1.732868in}}%
\pgfpathcurveto{\pgfqpoint{2.166553in}{1.721818in}}{\pgfqpoint{2.170943in}{1.711219in}}{\pgfqpoint{2.178757in}{1.703405in}}%
\pgfpathcurveto{\pgfqpoint{2.186571in}{1.695592in}}{\pgfqpoint{2.197170in}{1.691202in}}{\pgfqpoint{2.208220in}{1.691202in}}%
\pgfpathclose%
\pgfusepath{stroke,fill}%
\end{pgfscope}%
\begin{pgfscope}%
\pgfpathrectangle{\pgfqpoint{1.115200in}{0.528000in}}{\pgfqpoint{4.329600in}{3.696000in}} %
\pgfusepath{clip}%
\pgfsetbuttcap%
\pgfsetroundjoin%
\definecolor{currentfill}{rgb}{0.000000,0.000000,1.000000}%
\pgfsetfillcolor{currentfill}%
\pgfsetlinewidth{1.003750pt}%
\definecolor{currentstroke}{rgb}{0.000000,0.000000,1.000000}%
\pgfsetstrokecolor{currentstroke}%
\pgfsetdash{}{0pt}%
\pgfpathmoveto{\pgfqpoint{2.390590in}{1.609857in}}%
\pgfpathcurveto{\pgfqpoint{2.401640in}{1.609857in}}{\pgfqpoint{2.412239in}{1.614247in}}{\pgfqpoint{2.420052in}{1.622061in}}%
\pgfpathcurveto{\pgfqpoint{2.427866in}{1.629874in}}{\pgfqpoint{2.432256in}{1.640473in}}{\pgfqpoint{2.432256in}{1.651524in}}%
\pgfpathcurveto{\pgfqpoint{2.432256in}{1.662574in}}{\pgfqpoint{2.427866in}{1.673173in}}{\pgfqpoint{2.420052in}{1.680986in}}%
\pgfpathcurveto{\pgfqpoint{2.412239in}{1.688800in}}{\pgfqpoint{2.401640in}{1.693190in}}{\pgfqpoint{2.390590in}{1.693190in}}%
\pgfpathcurveto{\pgfqpoint{2.379539in}{1.693190in}}{\pgfqpoint{2.368940in}{1.688800in}}{\pgfqpoint{2.361127in}{1.680986in}}%
\pgfpathcurveto{\pgfqpoint{2.353313in}{1.673173in}}{\pgfqpoint{2.348923in}{1.662574in}}{\pgfqpoint{2.348923in}{1.651524in}}%
\pgfpathcurveto{\pgfqpoint{2.348923in}{1.640473in}}{\pgfqpoint{2.353313in}{1.629874in}}{\pgfqpoint{2.361127in}{1.622061in}}%
\pgfpathcurveto{\pgfqpoint{2.368940in}{1.614247in}}{\pgfqpoint{2.379539in}{1.609857in}}{\pgfqpoint{2.390590in}{1.609857in}}%
\pgfpathclose%
\pgfusepath{stroke,fill}%
\end{pgfscope}%
\begin{pgfscope}%
\pgfpathrectangle{\pgfqpoint{1.115200in}{0.528000in}}{\pgfqpoint{4.329600in}{3.696000in}} %
\pgfusepath{clip}%
\pgfsetbuttcap%
\pgfsetroundjoin%
\definecolor{currentfill}{rgb}{0.000000,0.000000,1.000000}%
\pgfsetfillcolor{currentfill}%
\pgfsetlinewidth{1.003750pt}%
\definecolor{currentstroke}{rgb}{0.000000,0.000000,1.000000}%
\pgfsetstrokecolor{currentstroke}%
\pgfsetdash{}{0pt}%
\pgfpathmoveto{\pgfqpoint{2.088891in}{1.422241in}}%
\pgfpathcurveto{\pgfqpoint{2.099941in}{1.422241in}}{\pgfqpoint{2.110540in}{1.426632in}}{\pgfqpoint{2.118354in}{1.434445in}}%
\pgfpathcurveto{\pgfqpoint{2.126167in}{1.442259in}}{\pgfqpoint{2.130558in}{1.452858in}}{\pgfqpoint{2.130558in}{1.463908in}}%
\pgfpathcurveto{\pgfqpoint{2.130558in}{1.474958in}}{\pgfqpoint{2.126167in}{1.485557in}}{\pgfqpoint{2.118354in}{1.493371in}}%
\pgfpathcurveto{\pgfqpoint{2.110540in}{1.501184in}}{\pgfqpoint{2.099941in}{1.505575in}}{\pgfqpoint{2.088891in}{1.505575in}}%
\pgfpathcurveto{\pgfqpoint{2.077841in}{1.505575in}}{\pgfqpoint{2.067242in}{1.501184in}}{\pgfqpoint{2.059428in}{1.493371in}}%
\pgfpathcurveto{\pgfqpoint{2.051615in}{1.485557in}}{\pgfqpoint{2.047224in}{1.474958in}}{\pgfqpoint{2.047224in}{1.463908in}}%
\pgfpathcurveto{\pgfqpoint{2.047224in}{1.452858in}}{\pgfqpoint{2.051615in}{1.442259in}}{\pgfqpoint{2.059428in}{1.434445in}}%
\pgfpathcurveto{\pgfqpoint{2.067242in}{1.426632in}}{\pgfqpoint{2.077841in}{1.422241in}}{\pgfqpoint{2.088891in}{1.422241in}}%
\pgfpathclose%
\pgfusepath{stroke,fill}%
\end{pgfscope}%
\begin{pgfscope}%
\pgfpathrectangle{\pgfqpoint{1.115200in}{0.528000in}}{\pgfqpoint{4.329600in}{3.696000in}} %
\pgfusepath{clip}%
\pgfsetbuttcap%
\pgfsetroundjoin%
\definecolor{currentfill}{rgb}{0.000000,0.000000,1.000000}%
\pgfsetfillcolor{currentfill}%
\pgfsetlinewidth{1.003750pt}%
\definecolor{currentstroke}{rgb}{0.000000,0.000000,1.000000}%
\pgfsetstrokecolor{currentstroke}%
\pgfsetdash{}{0pt}%
\pgfpathmoveto{\pgfqpoint{2.047624in}{1.516614in}}%
\pgfpathcurveto{\pgfqpoint{2.058674in}{1.516614in}}{\pgfqpoint{2.069273in}{1.521005in}}{\pgfqpoint{2.077087in}{1.528818in}}%
\pgfpathcurveto{\pgfqpoint{2.084900in}{1.536632in}}{\pgfqpoint{2.089291in}{1.547231in}}{\pgfqpoint{2.089291in}{1.558281in}}%
\pgfpathcurveto{\pgfqpoint{2.089291in}{1.569331in}}{\pgfqpoint{2.084900in}{1.579930in}}{\pgfqpoint{2.077087in}{1.587744in}}%
\pgfpathcurveto{\pgfqpoint{2.069273in}{1.595557in}}{\pgfqpoint{2.058674in}{1.599948in}}{\pgfqpoint{2.047624in}{1.599948in}}%
\pgfpathcurveto{\pgfqpoint{2.036574in}{1.599948in}}{\pgfqpoint{2.025975in}{1.595557in}}{\pgfqpoint{2.018161in}{1.587744in}}%
\pgfpathcurveto{\pgfqpoint{2.010347in}{1.579930in}}{\pgfqpoint{2.005957in}{1.569331in}}{\pgfqpoint{2.005957in}{1.558281in}}%
\pgfpathcurveto{\pgfqpoint{2.005957in}{1.547231in}}{\pgfqpoint{2.010347in}{1.536632in}}{\pgfqpoint{2.018161in}{1.528818in}}%
\pgfpathcurveto{\pgfqpoint{2.025975in}{1.521005in}}{\pgfqpoint{2.036574in}{1.516614in}}{\pgfqpoint{2.047624in}{1.516614in}}%
\pgfpathclose%
\pgfusepath{stroke,fill}%
\end{pgfscope}%
\begin{pgfscope}%
\pgfpathrectangle{\pgfqpoint{1.115200in}{0.528000in}}{\pgfqpoint{4.329600in}{3.696000in}} %
\pgfusepath{clip}%
\pgfsetbuttcap%
\pgfsetroundjoin%
\definecolor{currentfill}{rgb}{0.000000,0.000000,1.000000}%
\pgfsetfillcolor{currentfill}%
\pgfsetlinewidth{1.003750pt}%
\definecolor{currentstroke}{rgb}{0.000000,0.000000,1.000000}%
\pgfsetstrokecolor{currentstroke}%
\pgfsetdash{}{0pt}%
\pgfpathmoveto{\pgfqpoint{2.322569in}{1.657460in}}%
\pgfpathcurveto{\pgfqpoint{2.333619in}{1.657460in}}{\pgfqpoint{2.344218in}{1.661850in}}{\pgfqpoint{2.352031in}{1.669664in}}%
\pgfpathcurveto{\pgfqpoint{2.359845in}{1.677478in}}{\pgfqpoint{2.364235in}{1.688077in}}{\pgfqpoint{2.364235in}{1.699127in}}%
\pgfpathcurveto{\pgfqpoint{2.364235in}{1.710177in}}{\pgfqpoint{2.359845in}{1.720776in}}{\pgfqpoint{2.352031in}{1.728589in}}%
\pgfpathcurveto{\pgfqpoint{2.344218in}{1.736403in}}{\pgfqpoint{2.333619in}{1.740793in}}{\pgfqpoint{2.322569in}{1.740793in}}%
\pgfpathcurveto{\pgfqpoint{2.311518in}{1.740793in}}{\pgfqpoint{2.300919in}{1.736403in}}{\pgfqpoint{2.293106in}{1.728589in}}%
\pgfpathcurveto{\pgfqpoint{2.285292in}{1.720776in}}{\pgfqpoint{2.280902in}{1.710177in}}{\pgfqpoint{2.280902in}{1.699127in}}%
\pgfpathcurveto{\pgfqpoint{2.280902in}{1.688077in}}{\pgfqpoint{2.285292in}{1.677478in}}{\pgfqpoint{2.293106in}{1.669664in}}%
\pgfpathcurveto{\pgfqpoint{2.300919in}{1.661850in}}{\pgfqpoint{2.311518in}{1.657460in}}{\pgfqpoint{2.322569in}{1.657460in}}%
\pgfpathclose%
\pgfusepath{stroke,fill}%
\end{pgfscope}%
\begin{pgfscope}%
\pgfpathrectangle{\pgfqpoint{1.115200in}{0.528000in}}{\pgfqpoint{4.329600in}{3.696000in}} %
\pgfusepath{clip}%
\pgfsetbuttcap%
\pgfsetroundjoin%
\definecolor{currentfill}{rgb}{0.000000,0.000000,1.000000}%
\pgfsetfillcolor{currentfill}%
\pgfsetlinewidth{1.003750pt}%
\definecolor{currentstroke}{rgb}{0.000000,0.000000,1.000000}%
\pgfsetstrokecolor{currentstroke}%
\pgfsetdash{}{0pt}%
\pgfpathmoveto{\pgfqpoint{2.151125in}{1.612676in}}%
\pgfpathcurveto{\pgfqpoint{2.162175in}{1.612676in}}{\pgfqpoint{2.172774in}{1.617066in}}{\pgfqpoint{2.180588in}{1.624880in}}%
\pgfpathcurveto{\pgfqpoint{2.188401in}{1.632694in}}{\pgfqpoint{2.192792in}{1.643293in}}{\pgfqpoint{2.192792in}{1.654343in}}%
\pgfpathcurveto{\pgfqpoint{2.192792in}{1.665393in}}{\pgfqpoint{2.188401in}{1.675992in}}{\pgfqpoint{2.180588in}{1.683806in}}%
\pgfpathcurveto{\pgfqpoint{2.172774in}{1.691619in}}{\pgfqpoint{2.162175in}{1.696009in}}{\pgfqpoint{2.151125in}{1.696009in}}%
\pgfpathcurveto{\pgfqpoint{2.140075in}{1.696009in}}{\pgfqpoint{2.129476in}{1.691619in}}{\pgfqpoint{2.121662in}{1.683806in}}%
\pgfpathcurveto{\pgfqpoint{2.113849in}{1.675992in}}{\pgfqpoint{2.109458in}{1.665393in}}{\pgfqpoint{2.109458in}{1.654343in}}%
\pgfpathcurveto{\pgfqpoint{2.109458in}{1.643293in}}{\pgfqpoint{2.113849in}{1.632694in}}{\pgfqpoint{2.121662in}{1.624880in}}%
\pgfpathcurveto{\pgfqpoint{2.129476in}{1.617066in}}{\pgfqpoint{2.140075in}{1.612676in}}{\pgfqpoint{2.151125in}{1.612676in}}%
\pgfpathclose%
\pgfusepath{stroke,fill}%
\end{pgfscope}%
\begin{pgfscope}%
\pgfpathrectangle{\pgfqpoint{1.115200in}{0.528000in}}{\pgfqpoint{4.329600in}{3.696000in}} %
\pgfusepath{clip}%
\pgfsetbuttcap%
\pgfsetroundjoin%
\definecolor{currentfill}{rgb}{0.000000,0.000000,1.000000}%
\pgfsetfillcolor{currentfill}%
\pgfsetlinewidth{1.003750pt}%
\definecolor{currentstroke}{rgb}{0.000000,0.000000,1.000000}%
\pgfsetstrokecolor{currentstroke}%
\pgfsetdash{}{0pt}%
\pgfpathmoveto{\pgfqpoint{1.954412in}{1.406613in}}%
\pgfpathcurveto{\pgfqpoint{1.965462in}{1.406613in}}{\pgfqpoint{1.976061in}{1.411003in}}{\pgfqpoint{1.983875in}{1.418817in}}%
\pgfpathcurveto{\pgfqpoint{1.991688in}{1.426630in}}{\pgfqpoint{1.996078in}{1.437229in}}{\pgfqpoint{1.996078in}{1.448279in}}%
\pgfpathcurveto{\pgfqpoint{1.996078in}{1.459329in}}{\pgfqpoint{1.991688in}{1.469928in}}{\pgfqpoint{1.983875in}{1.477742in}}%
\pgfpathcurveto{\pgfqpoint{1.976061in}{1.485556in}}{\pgfqpoint{1.965462in}{1.489946in}}{\pgfqpoint{1.954412in}{1.489946in}}%
\pgfpathcurveto{\pgfqpoint{1.943362in}{1.489946in}}{\pgfqpoint{1.932763in}{1.485556in}}{\pgfqpoint{1.924949in}{1.477742in}}%
\pgfpathcurveto{\pgfqpoint{1.917135in}{1.469928in}}{\pgfqpoint{1.912745in}{1.459329in}}{\pgfqpoint{1.912745in}{1.448279in}}%
\pgfpathcurveto{\pgfqpoint{1.912745in}{1.437229in}}{\pgfqpoint{1.917135in}{1.426630in}}{\pgfqpoint{1.924949in}{1.418817in}}%
\pgfpathcurveto{\pgfqpoint{1.932763in}{1.411003in}}{\pgfqpoint{1.943362in}{1.406613in}}{\pgfqpoint{1.954412in}{1.406613in}}%
\pgfpathclose%
\pgfusepath{stroke,fill}%
\end{pgfscope}%
\begin{pgfscope}%
\pgfpathrectangle{\pgfqpoint{1.115200in}{0.528000in}}{\pgfqpoint{4.329600in}{3.696000in}} %
\pgfusepath{clip}%
\pgfsetbuttcap%
\pgfsetroundjoin%
\definecolor{currentfill}{rgb}{0.000000,0.000000,1.000000}%
\pgfsetfillcolor{currentfill}%
\pgfsetlinewidth{1.003750pt}%
\definecolor{currentstroke}{rgb}{0.000000,0.000000,1.000000}%
\pgfsetstrokecolor{currentstroke}%
\pgfsetdash{}{0pt}%
\pgfpathmoveto{\pgfqpoint{2.364207in}{1.613127in}}%
\pgfpathcurveto{\pgfqpoint{2.375257in}{1.613127in}}{\pgfqpoint{2.385856in}{1.617517in}}{\pgfqpoint{2.393670in}{1.625331in}}%
\pgfpathcurveto{\pgfqpoint{2.401484in}{1.633144in}}{\pgfqpoint{2.405874in}{1.643743in}}{\pgfqpoint{2.405874in}{1.654793in}}%
\pgfpathcurveto{\pgfqpoint{2.405874in}{1.665844in}}{\pgfqpoint{2.401484in}{1.676443in}}{\pgfqpoint{2.393670in}{1.684256in}}%
\pgfpathcurveto{\pgfqpoint{2.385856in}{1.692070in}}{\pgfqpoint{2.375257in}{1.696460in}}{\pgfqpoint{2.364207in}{1.696460in}}%
\pgfpathcurveto{\pgfqpoint{2.353157in}{1.696460in}}{\pgfqpoint{2.342558in}{1.692070in}}{\pgfqpoint{2.334745in}{1.684256in}}%
\pgfpathcurveto{\pgfqpoint{2.326931in}{1.676443in}}{\pgfqpoint{2.322541in}{1.665844in}}{\pgfqpoint{2.322541in}{1.654793in}}%
\pgfpathcurveto{\pgfqpoint{2.322541in}{1.643743in}}{\pgfqpoint{2.326931in}{1.633144in}}{\pgfqpoint{2.334745in}{1.625331in}}%
\pgfpathcurveto{\pgfqpoint{2.342558in}{1.617517in}}{\pgfqpoint{2.353157in}{1.613127in}}{\pgfqpoint{2.364207in}{1.613127in}}%
\pgfpathclose%
\pgfusepath{stroke,fill}%
\end{pgfscope}%
\begin{pgfscope}%
\pgfpathrectangle{\pgfqpoint{1.115200in}{0.528000in}}{\pgfqpoint{4.329600in}{3.696000in}} %
\pgfusepath{clip}%
\pgfsetbuttcap%
\pgfsetroundjoin%
\definecolor{currentfill}{rgb}{0.000000,0.000000,1.000000}%
\pgfsetfillcolor{currentfill}%
\pgfsetlinewidth{1.003750pt}%
\definecolor{currentstroke}{rgb}{0.000000,0.000000,1.000000}%
\pgfsetstrokecolor{currentstroke}%
\pgfsetdash{}{0pt}%
\pgfpathmoveto{\pgfqpoint{2.035372in}{1.381613in}}%
\pgfpathcurveto{\pgfqpoint{2.046422in}{1.381613in}}{\pgfqpoint{2.057021in}{1.386003in}}{\pgfqpoint{2.064835in}{1.393817in}}%
\pgfpathcurveto{\pgfqpoint{2.072648in}{1.401631in}}{\pgfqpoint{2.077038in}{1.412230in}}{\pgfqpoint{2.077038in}{1.423280in}}%
\pgfpathcurveto{\pgfqpoint{2.077038in}{1.434330in}}{\pgfqpoint{2.072648in}{1.444929in}}{\pgfqpoint{2.064835in}{1.452743in}}%
\pgfpathcurveto{\pgfqpoint{2.057021in}{1.460556in}}{\pgfqpoint{2.046422in}{1.464946in}}{\pgfqpoint{2.035372in}{1.464946in}}%
\pgfpathcurveto{\pgfqpoint{2.024322in}{1.464946in}}{\pgfqpoint{2.013723in}{1.460556in}}{\pgfqpoint{2.005909in}{1.452743in}}%
\pgfpathcurveto{\pgfqpoint{1.998095in}{1.444929in}}{\pgfqpoint{1.993705in}{1.434330in}}{\pgfqpoint{1.993705in}{1.423280in}}%
\pgfpathcurveto{\pgfqpoint{1.993705in}{1.412230in}}{\pgfqpoint{1.998095in}{1.401631in}}{\pgfqpoint{2.005909in}{1.393817in}}%
\pgfpathcurveto{\pgfqpoint{2.013723in}{1.386003in}}{\pgfqpoint{2.024322in}{1.381613in}}{\pgfqpoint{2.035372in}{1.381613in}}%
\pgfpathclose%
\pgfusepath{stroke,fill}%
\end{pgfscope}%
\begin{pgfscope}%
\pgfpathrectangle{\pgfqpoint{1.115200in}{0.528000in}}{\pgfqpoint{4.329600in}{3.696000in}} %
\pgfusepath{clip}%
\pgfsetbuttcap%
\pgfsetroundjoin%
\definecolor{currentfill}{rgb}{0.000000,0.000000,1.000000}%
\pgfsetfillcolor{currentfill}%
\pgfsetlinewidth{1.003750pt}%
\definecolor{currentstroke}{rgb}{0.000000,0.000000,1.000000}%
\pgfsetstrokecolor{currentstroke}%
\pgfsetdash{}{0pt}%
\pgfpathmoveto{\pgfqpoint{2.082263in}{1.204573in}}%
\pgfpathcurveto{\pgfqpoint{2.093313in}{1.204573in}}{\pgfqpoint{2.103912in}{1.208964in}}{\pgfqpoint{2.111726in}{1.216777in}}%
\pgfpathcurveto{\pgfqpoint{2.119539in}{1.224591in}}{\pgfqpoint{2.123930in}{1.235190in}}{\pgfqpoint{2.123930in}{1.246240in}}%
\pgfpathcurveto{\pgfqpoint{2.123930in}{1.257290in}}{\pgfqpoint{2.119539in}{1.267889in}}{\pgfqpoint{2.111726in}{1.275703in}}%
\pgfpathcurveto{\pgfqpoint{2.103912in}{1.283516in}}{\pgfqpoint{2.093313in}{1.287907in}}{\pgfqpoint{2.082263in}{1.287907in}}%
\pgfpathcurveto{\pgfqpoint{2.071213in}{1.287907in}}{\pgfqpoint{2.060614in}{1.283516in}}{\pgfqpoint{2.052800in}{1.275703in}}%
\pgfpathcurveto{\pgfqpoint{2.044987in}{1.267889in}}{\pgfqpoint{2.040596in}{1.257290in}}{\pgfqpoint{2.040596in}{1.246240in}}%
\pgfpathcurveto{\pgfqpoint{2.040596in}{1.235190in}}{\pgfqpoint{2.044987in}{1.224591in}}{\pgfqpoint{2.052800in}{1.216777in}}%
\pgfpathcurveto{\pgfqpoint{2.060614in}{1.208964in}}{\pgfqpoint{2.071213in}{1.204573in}}{\pgfqpoint{2.082263in}{1.204573in}}%
\pgfpathclose%
\pgfusepath{stroke,fill}%
\end{pgfscope}%
\begin{pgfscope}%
\pgfpathrectangle{\pgfqpoint{1.115200in}{0.528000in}}{\pgfqpoint{4.329600in}{3.696000in}} %
\pgfusepath{clip}%
\pgfsetbuttcap%
\pgfsetroundjoin%
\definecolor{currentfill}{rgb}{0.000000,0.000000,1.000000}%
\pgfsetfillcolor{currentfill}%
\pgfsetlinewidth{1.003750pt}%
\definecolor{currentstroke}{rgb}{0.000000,0.000000,1.000000}%
\pgfsetstrokecolor{currentstroke}%
\pgfsetdash{}{0pt}%
\pgfpathmoveto{\pgfqpoint{2.007490in}{1.569873in}}%
\pgfpathcurveto{\pgfqpoint{2.018541in}{1.569873in}}{\pgfqpoint{2.029140in}{1.574264in}}{\pgfqpoint{2.036953in}{1.582077in}}%
\pgfpathcurveto{\pgfqpoint{2.044767in}{1.589891in}}{\pgfqpoint{2.049157in}{1.600490in}}{\pgfqpoint{2.049157in}{1.611540in}}%
\pgfpathcurveto{\pgfqpoint{2.049157in}{1.622590in}}{\pgfqpoint{2.044767in}{1.633189in}}{\pgfqpoint{2.036953in}{1.641003in}}%
\pgfpathcurveto{\pgfqpoint{2.029140in}{1.648816in}}{\pgfqpoint{2.018541in}{1.653207in}}{\pgfqpoint{2.007490in}{1.653207in}}%
\pgfpathcurveto{\pgfqpoint{1.996440in}{1.653207in}}{\pgfqpoint{1.985841in}{1.648816in}}{\pgfqpoint{1.978028in}{1.641003in}}%
\pgfpathcurveto{\pgfqpoint{1.970214in}{1.633189in}}{\pgfqpoint{1.965824in}{1.622590in}}{\pgfqpoint{1.965824in}{1.611540in}}%
\pgfpathcurveto{\pgfqpoint{1.965824in}{1.600490in}}{\pgfqpoint{1.970214in}{1.589891in}}{\pgfqpoint{1.978028in}{1.582077in}}%
\pgfpathcurveto{\pgfqpoint{1.985841in}{1.574264in}}{\pgfqpoint{1.996440in}{1.569873in}}{\pgfqpoint{2.007490in}{1.569873in}}%
\pgfpathclose%
\pgfusepath{stroke,fill}%
\end{pgfscope}%
\begin{pgfscope}%
\pgfpathrectangle{\pgfqpoint{1.115200in}{0.528000in}}{\pgfqpoint{4.329600in}{3.696000in}} %
\pgfusepath{clip}%
\pgfsetbuttcap%
\pgfsetroundjoin%
\definecolor{currentfill}{rgb}{0.000000,0.000000,1.000000}%
\pgfsetfillcolor{currentfill}%
\pgfsetlinewidth{1.003750pt}%
\definecolor{currentstroke}{rgb}{0.000000,0.000000,1.000000}%
\pgfsetstrokecolor{currentstroke}%
\pgfsetdash{}{0pt}%
\pgfpathmoveto{\pgfqpoint{2.228029in}{1.693004in}}%
\pgfpathcurveto{\pgfqpoint{2.239079in}{1.693004in}}{\pgfqpoint{2.249678in}{1.697394in}}{\pgfqpoint{2.257492in}{1.705207in}}%
\pgfpathcurveto{\pgfqpoint{2.265306in}{1.713021in}}{\pgfqpoint{2.269696in}{1.723620in}}{\pgfqpoint{2.269696in}{1.734670in}}%
\pgfpathcurveto{\pgfqpoint{2.269696in}{1.745720in}}{\pgfqpoint{2.265306in}{1.756319in}}{\pgfqpoint{2.257492in}{1.764133in}}%
\pgfpathcurveto{\pgfqpoint{2.249678in}{1.771947in}}{\pgfqpoint{2.239079in}{1.776337in}}{\pgfqpoint{2.228029in}{1.776337in}}%
\pgfpathcurveto{\pgfqpoint{2.216979in}{1.776337in}}{\pgfqpoint{2.206380in}{1.771947in}}{\pgfqpoint{2.198567in}{1.764133in}}%
\pgfpathcurveto{\pgfqpoint{2.190753in}{1.756319in}}{\pgfqpoint{2.186363in}{1.745720in}}{\pgfqpoint{2.186363in}{1.734670in}}%
\pgfpathcurveto{\pgfqpoint{2.186363in}{1.723620in}}{\pgfqpoint{2.190753in}{1.713021in}}{\pgfqpoint{2.198567in}{1.705207in}}%
\pgfpathcurveto{\pgfqpoint{2.206380in}{1.697394in}}{\pgfqpoint{2.216979in}{1.693004in}}{\pgfqpoint{2.228029in}{1.693004in}}%
\pgfpathclose%
\pgfusepath{stroke,fill}%
\end{pgfscope}%
\begin{pgfscope}%
\pgfpathrectangle{\pgfqpoint{1.115200in}{0.528000in}}{\pgfqpoint{4.329600in}{3.696000in}} %
\pgfusepath{clip}%
\pgfsetbuttcap%
\pgfsetroundjoin%
\definecolor{currentfill}{rgb}{0.000000,0.000000,1.000000}%
\pgfsetfillcolor{currentfill}%
\pgfsetlinewidth{1.003750pt}%
\definecolor{currentstroke}{rgb}{0.000000,0.000000,1.000000}%
\pgfsetstrokecolor{currentstroke}%
\pgfsetdash{}{0pt}%
\pgfpathmoveto{\pgfqpoint{2.341271in}{1.612330in}}%
\pgfpathcurveto{\pgfqpoint{2.352321in}{1.612330in}}{\pgfqpoint{2.362920in}{1.616720in}}{\pgfqpoint{2.370734in}{1.624534in}}%
\pgfpathcurveto{\pgfqpoint{2.378548in}{1.632347in}}{\pgfqpoint{2.382938in}{1.642946in}}{\pgfqpoint{2.382938in}{1.653997in}}%
\pgfpathcurveto{\pgfqpoint{2.382938in}{1.665047in}}{\pgfqpoint{2.378548in}{1.675646in}}{\pgfqpoint{2.370734in}{1.683459in}}%
\pgfpathcurveto{\pgfqpoint{2.362920in}{1.691273in}}{\pgfqpoint{2.352321in}{1.695663in}}{\pgfqpoint{2.341271in}{1.695663in}}%
\pgfpathcurveto{\pgfqpoint{2.330221in}{1.695663in}}{\pgfqpoint{2.319622in}{1.691273in}}{\pgfqpoint{2.311808in}{1.683459in}}%
\pgfpathcurveto{\pgfqpoint{2.303995in}{1.675646in}}{\pgfqpoint{2.299604in}{1.665047in}}{\pgfqpoint{2.299604in}{1.653997in}}%
\pgfpathcurveto{\pgfqpoint{2.299604in}{1.642946in}}{\pgfqpoint{2.303995in}{1.632347in}}{\pgfqpoint{2.311808in}{1.624534in}}%
\pgfpathcurveto{\pgfqpoint{2.319622in}{1.616720in}}{\pgfqpoint{2.330221in}{1.612330in}}{\pgfqpoint{2.341271in}{1.612330in}}%
\pgfpathclose%
\pgfusepath{stroke,fill}%
\end{pgfscope}%
\begin{pgfscope}%
\pgfpathrectangle{\pgfqpoint{1.115200in}{0.528000in}}{\pgfqpoint{4.329600in}{3.696000in}} %
\pgfusepath{clip}%
\pgfsetbuttcap%
\pgfsetroundjoin%
\definecolor{currentfill}{rgb}{0.000000,0.000000,1.000000}%
\pgfsetfillcolor{currentfill}%
\pgfsetlinewidth{1.003750pt}%
\definecolor{currentstroke}{rgb}{0.000000,0.000000,1.000000}%
\pgfsetstrokecolor{currentstroke}%
\pgfsetdash{}{0pt}%
\pgfpathmoveto{\pgfqpoint{2.111720in}{1.775026in}}%
\pgfpathcurveto{\pgfqpoint{2.122770in}{1.775026in}}{\pgfqpoint{2.133369in}{1.779416in}}{\pgfqpoint{2.141183in}{1.787230in}}%
\pgfpathcurveto{\pgfqpoint{2.148996in}{1.795043in}}{\pgfqpoint{2.153387in}{1.805642in}}{\pgfqpoint{2.153387in}{1.816692in}}%
\pgfpathcurveto{\pgfqpoint{2.153387in}{1.827742in}}{\pgfqpoint{2.148996in}{1.838342in}}{\pgfqpoint{2.141183in}{1.846155in}}%
\pgfpathcurveto{\pgfqpoint{2.133369in}{1.853969in}}{\pgfqpoint{2.122770in}{1.858359in}}{\pgfqpoint{2.111720in}{1.858359in}}%
\pgfpathcurveto{\pgfqpoint{2.100670in}{1.858359in}}{\pgfqpoint{2.090071in}{1.853969in}}{\pgfqpoint{2.082257in}{1.846155in}}%
\pgfpathcurveto{\pgfqpoint{2.074444in}{1.838342in}}{\pgfqpoint{2.070053in}{1.827742in}}{\pgfqpoint{2.070053in}{1.816692in}}%
\pgfpathcurveto{\pgfqpoint{2.070053in}{1.805642in}}{\pgfqpoint{2.074444in}{1.795043in}}{\pgfqpoint{2.082257in}{1.787230in}}%
\pgfpathcurveto{\pgfqpoint{2.090071in}{1.779416in}}{\pgfqpoint{2.100670in}{1.775026in}}{\pgfqpoint{2.111720in}{1.775026in}}%
\pgfpathclose%
\pgfusepath{stroke,fill}%
\end{pgfscope}%
\begin{pgfscope}%
\pgfpathrectangle{\pgfqpoint{1.115200in}{0.528000in}}{\pgfqpoint{4.329600in}{3.696000in}} %
\pgfusepath{clip}%
\pgfsetbuttcap%
\pgfsetroundjoin%
\definecolor{currentfill}{rgb}{0.000000,0.000000,1.000000}%
\pgfsetfillcolor{currentfill}%
\pgfsetlinewidth{1.003750pt}%
\definecolor{currentstroke}{rgb}{0.000000,0.000000,1.000000}%
\pgfsetstrokecolor{currentstroke}%
\pgfsetdash{}{0pt}%
\pgfpathmoveto{\pgfqpoint{2.233985in}{1.819457in}}%
\pgfpathcurveto{\pgfqpoint{2.245035in}{1.819457in}}{\pgfqpoint{2.255634in}{1.823847in}}{\pgfqpoint{2.263447in}{1.831661in}}%
\pgfpathcurveto{\pgfqpoint{2.271261in}{1.839474in}}{\pgfqpoint{2.275651in}{1.850073in}}{\pgfqpoint{2.275651in}{1.861123in}}%
\pgfpathcurveto{\pgfqpoint{2.275651in}{1.872174in}}{\pgfqpoint{2.271261in}{1.882773in}}{\pgfqpoint{2.263447in}{1.890586in}}%
\pgfpathcurveto{\pgfqpoint{2.255634in}{1.898400in}}{\pgfqpoint{2.245035in}{1.902790in}}{\pgfqpoint{2.233985in}{1.902790in}}%
\pgfpathcurveto{\pgfqpoint{2.222934in}{1.902790in}}{\pgfqpoint{2.212335in}{1.898400in}}{\pgfqpoint{2.204522in}{1.890586in}}%
\pgfpathcurveto{\pgfqpoint{2.196708in}{1.882773in}}{\pgfqpoint{2.192318in}{1.872174in}}{\pgfqpoint{2.192318in}{1.861123in}}%
\pgfpathcurveto{\pgfqpoint{2.192318in}{1.850073in}}{\pgfqpoint{2.196708in}{1.839474in}}{\pgfqpoint{2.204522in}{1.831661in}}%
\pgfpathcurveto{\pgfqpoint{2.212335in}{1.823847in}}{\pgfqpoint{2.222934in}{1.819457in}}{\pgfqpoint{2.233985in}{1.819457in}}%
\pgfpathclose%
\pgfusepath{stroke,fill}%
\end{pgfscope}%
\begin{pgfscope}%
\pgfpathrectangle{\pgfqpoint{1.115200in}{0.528000in}}{\pgfqpoint{4.329600in}{3.696000in}} %
\pgfusepath{clip}%
\pgfsetbuttcap%
\pgfsetroundjoin%
\definecolor{currentfill}{rgb}{0.000000,0.000000,1.000000}%
\pgfsetfillcolor{currentfill}%
\pgfsetlinewidth{1.003750pt}%
\definecolor{currentstroke}{rgb}{0.000000,0.000000,1.000000}%
\pgfsetstrokecolor{currentstroke}%
\pgfsetdash{}{0pt}%
\pgfpathmoveto{\pgfqpoint{2.096523in}{1.598333in}}%
\pgfpathcurveto{\pgfqpoint{2.107573in}{1.598333in}}{\pgfqpoint{2.118172in}{1.602723in}}{\pgfqpoint{2.125986in}{1.610536in}}%
\pgfpathcurveto{\pgfqpoint{2.133799in}{1.618350in}}{\pgfqpoint{2.138190in}{1.628949in}}{\pgfqpoint{2.138190in}{1.639999in}}%
\pgfpathcurveto{\pgfqpoint{2.138190in}{1.651049in}}{\pgfqpoint{2.133799in}{1.661648in}}{\pgfqpoint{2.125986in}{1.669462in}}%
\pgfpathcurveto{\pgfqpoint{2.118172in}{1.677276in}}{\pgfqpoint{2.107573in}{1.681666in}}{\pgfqpoint{2.096523in}{1.681666in}}%
\pgfpathcurveto{\pgfqpoint{2.085473in}{1.681666in}}{\pgfqpoint{2.074874in}{1.677276in}}{\pgfqpoint{2.067060in}{1.669462in}}%
\pgfpathcurveto{\pgfqpoint{2.059247in}{1.661648in}}{\pgfqpoint{2.054856in}{1.651049in}}{\pgfqpoint{2.054856in}{1.639999in}}%
\pgfpathcurveto{\pgfqpoint{2.054856in}{1.628949in}}{\pgfqpoint{2.059247in}{1.618350in}}{\pgfqpoint{2.067060in}{1.610536in}}%
\pgfpathcurveto{\pgfqpoint{2.074874in}{1.602723in}}{\pgfqpoint{2.085473in}{1.598333in}}{\pgfqpoint{2.096523in}{1.598333in}}%
\pgfpathclose%
\pgfusepath{stroke,fill}%
\end{pgfscope}%
\begin{pgfscope}%
\pgfpathrectangle{\pgfqpoint{1.115200in}{0.528000in}}{\pgfqpoint{4.329600in}{3.696000in}} %
\pgfusepath{clip}%
\pgfsetbuttcap%
\pgfsetroundjoin%
\definecolor{currentfill}{rgb}{0.000000,0.000000,1.000000}%
\pgfsetfillcolor{currentfill}%
\pgfsetlinewidth{1.003750pt}%
\definecolor{currentstroke}{rgb}{0.000000,0.000000,1.000000}%
\pgfsetstrokecolor{currentstroke}%
\pgfsetdash{}{0pt}%
\pgfpathmoveto{\pgfqpoint{2.541319in}{1.440429in}}%
\pgfpathcurveto{\pgfqpoint{2.552369in}{1.440429in}}{\pgfqpoint{2.562968in}{1.444819in}}{\pgfqpoint{2.570781in}{1.452633in}}%
\pgfpathcurveto{\pgfqpoint{2.578595in}{1.460446in}}{\pgfqpoint{2.582985in}{1.471045in}}{\pgfqpoint{2.582985in}{1.482095in}}%
\pgfpathcurveto{\pgfqpoint{2.582985in}{1.493146in}}{\pgfqpoint{2.578595in}{1.503745in}}{\pgfqpoint{2.570781in}{1.511558in}}%
\pgfpathcurveto{\pgfqpoint{2.562968in}{1.519372in}}{\pgfqpoint{2.552369in}{1.523762in}}{\pgfqpoint{2.541319in}{1.523762in}}%
\pgfpathcurveto{\pgfqpoint{2.530269in}{1.523762in}}{\pgfqpoint{2.519670in}{1.519372in}}{\pgfqpoint{2.511856in}{1.511558in}}%
\pgfpathcurveto{\pgfqpoint{2.504042in}{1.503745in}}{\pgfqpoint{2.499652in}{1.493146in}}{\pgfqpoint{2.499652in}{1.482095in}}%
\pgfpathcurveto{\pgfqpoint{2.499652in}{1.471045in}}{\pgfqpoint{2.504042in}{1.460446in}}{\pgfqpoint{2.511856in}{1.452633in}}%
\pgfpathcurveto{\pgfqpoint{2.519670in}{1.444819in}}{\pgfqpoint{2.530269in}{1.440429in}}{\pgfqpoint{2.541319in}{1.440429in}}%
\pgfpathclose%
\pgfusepath{stroke,fill}%
\end{pgfscope}%
\begin{pgfscope}%
\pgfpathrectangle{\pgfqpoint{1.115200in}{0.528000in}}{\pgfqpoint{4.329600in}{3.696000in}} %
\pgfusepath{clip}%
\pgfsetbuttcap%
\pgfsetmiterjoin%
\pgfsetlinewidth{1.003750pt}%
\definecolor{currentstroke}{rgb}{0.000000,0.500000,0.000000}%
\pgfsetstrokecolor{currentstroke}%
\pgfsetdash{{3.700000pt}{1.600000pt}}{0.000000pt}%
\pgfpathmoveto{\pgfqpoint{3.839363in}{2.811342in}}%
\pgfpathcurveto{\pgfqpoint{3.992603in}{2.811342in}}{\pgfqpoint{4.139588in}{2.872225in}}{\pgfqpoint{4.247945in}{2.980582in}}%
\pgfpathcurveto{\pgfqpoint{4.356303in}{3.088939in}}{\pgfqpoint{4.417186in}{3.235924in}}{\pgfqpoint{4.417186in}{3.389165in}}%
\pgfpathcurveto{\pgfqpoint{4.417186in}{3.542405in}}{\pgfqpoint{4.356303in}{3.689390in}}{\pgfqpoint{4.247945in}{3.797747in}}%
\pgfpathcurveto{\pgfqpoint{4.139588in}{3.906104in}}{\pgfqpoint{3.992603in}{3.966987in}}{\pgfqpoint{3.839363in}{3.966987in}}%
\pgfpathcurveto{\pgfqpoint{3.686122in}{3.966987in}}{\pgfqpoint{3.539138in}{3.906104in}}{\pgfqpoint{3.430780in}{3.797747in}}%
\pgfpathcurveto{\pgfqpoint{3.322423in}{3.689390in}}{\pgfqpoint{3.261540in}{3.542405in}}{\pgfqpoint{3.261540in}{3.389165in}}%
\pgfpathcurveto{\pgfqpoint{3.261540in}{3.235924in}}{\pgfqpoint{3.322423in}{3.088939in}}{\pgfqpoint{3.430780in}{2.980582in}}%
\pgfpathcurveto{\pgfqpoint{3.539138in}{2.872225in}}{\pgfqpoint{3.686122in}{2.811342in}}{\pgfqpoint{3.839363in}{2.811342in}}%
\pgfpathclose%
\pgfusepath{stroke}%
\end{pgfscope}%
\begin{pgfscope}%
\pgfpathrectangle{\pgfqpoint{1.115200in}{0.528000in}}{\pgfqpoint{4.329600in}{3.696000in}} %
\pgfusepath{clip}%
\pgfsetbuttcap%
\pgfsetmiterjoin%
\pgfsetlinewidth{1.003750pt}%
\definecolor{currentstroke}{rgb}{1.000000,0.000000,0.000000}%
\pgfsetstrokecolor{currentstroke}%
\pgfsetdash{{3.700000pt}{1.600000pt}}{0.000000pt}%
\pgfpathmoveto{\pgfqpoint{4.680009in}{1.290583in}}%
\pgfpathcurveto{\pgfqpoint{4.851604in}{1.290583in}}{\pgfqpoint{5.016194in}{1.358759in}}{\pgfqpoint{5.137530in}{1.480095in}}%
\pgfpathcurveto{\pgfqpoint{5.258866in}{1.601431in}}{\pgfqpoint{5.327042in}{1.766021in}}{\pgfqpoint{5.327042in}{1.937616in}}%
\pgfpathcurveto{\pgfqpoint{5.327042in}{2.109211in}}{\pgfqpoint{5.258866in}{2.273801in}}{\pgfqpoint{5.137530in}{2.395137in}}%
\pgfpathcurveto{\pgfqpoint{5.016194in}{2.516473in}}{\pgfqpoint{4.851604in}{2.584648in}}{\pgfqpoint{4.680009in}{2.584648in}}%
\pgfpathcurveto{\pgfqpoint{4.508414in}{2.584648in}}{\pgfqpoint{4.343824in}{2.516473in}}{\pgfqpoint{4.222488in}{2.395137in}}%
\pgfpathcurveto{\pgfqpoint{4.101152in}{2.273801in}}{\pgfqpoint{4.032977in}{2.109211in}}{\pgfqpoint{4.032977in}{1.937616in}}%
\pgfpathcurveto{\pgfqpoint{4.032977in}{1.766021in}}{\pgfqpoint{4.101152in}{1.601431in}}{\pgfqpoint{4.222488in}{1.480095in}}%
\pgfpathcurveto{\pgfqpoint{4.343824in}{1.358759in}}{\pgfqpoint{4.508414in}{1.290583in}}{\pgfqpoint{4.680009in}{1.290583in}}%
\pgfpathclose%
\pgfusepath{stroke}%
\end{pgfscope}%
\begin{pgfscope}%
\pgfpathrectangle{\pgfqpoint{1.115200in}{0.528000in}}{\pgfqpoint{4.329600in}{3.696000in}} %
\pgfusepath{clip}%
\pgfsetbuttcap%
\pgfsetmiterjoin%
\pgfsetlinewidth{1.003750pt}%
\definecolor{currentstroke}{rgb}{0.000000,0.000000,1.000000}%
\pgfsetstrokecolor{currentstroke}%
\pgfsetdash{{3.700000pt}{1.600000pt}}{0.000000pt}%
\pgfpathmoveto{\pgfqpoint{2.158818in}{1.136556in}}%
\pgfpathcurveto{\pgfqpoint{2.285864in}{1.136556in}}{\pgfqpoint{2.407724in}{1.187032in}}{\pgfqpoint{2.497559in}{1.276867in}}%
\pgfpathcurveto{\pgfqpoint{2.587395in}{1.366702in}}{\pgfqpoint{2.637871in}{1.488562in}}{\pgfqpoint{2.637871in}{1.615608in}}%
\pgfpathcurveto{\pgfqpoint{2.637871in}{1.742655in}}{\pgfqpoint{2.587395in}{1.864515in}}{\pgfqpoint{2.497559in}{1.954350in}}%
\pgfpathcurveto{\pgfqpoint{2.407724in}{2.044185in}}{\pgfqpoint{2.285864in}{2.094661in}}{\pgfqpoint{2.158818in}{2.094661in}}%
\pgfpathcurveto{\pgfqpoint{2.031772in}{2.094661in}}{\pgfqpoint{1.909912in}{2.044185in}}{\pgfqpoint{1.820077in}{1.954350in}}%
\pgfpathcurveto{\pgfqpoint{1.730242in}{1.864515in}}{\pgfqpoint{1.679766in}{1.742655in}}{\pgfqpoint{1.679766in}{1.615608in}}%
\pgfpathcurveto{\pgfqpoint{1.679766in}{1.488562in}}{\pgfqpoint{1.730242in}{1.366702in}}{\pgfqpoint{1.820077in}{1.276867in}}%
\pgfpathcurveto{\pgfqpoint{1.909912in}{1.187032in}}{\pgfqpoint{2.031772in}{1.136556in}}{\pgfqpoint{2.158818in}{1.136556in}}%
\pgfpathclose%
\pgfusepath{stroke}%
\end{pgfscope}%
\begin{pgfscope}%
\pgfsetbuttcap%
\pgfsetroundjoin%
\definecolor{currentfill}{rgb}{0.000000,0.000000,0.000000}%
\pgfsetfillcolor{currentfill}%
\pgfsetlinewidth{0.803000pt}%
\definecolor{currentstroke}{rgb}{0.000000,0.000000,0.000000}%
\pgfsetstrokecolor{currentstroke}%
\pgfsetdash{}{0pt}%
\pgfsys@defobject{currentmarker}{\pgfqpoint{0.000000in}{-0.048611in}}{\pgfqpoint{0.000000in}{0.000000in}}{%
\pgfpathmoveto{\pgfqpoint{0.000000in}{0.000000in}}%
\pgfpathlineto{\pgfqpoint{0.000000in}{-0.048611in}}%
\pgfusepath{stroke,fill}%
}%
\begin{pgfscope}%
\pgfsys@transformshift{1.115200in}{0.528000in}%
\pgfsys@useobject{currentmarker}{}%
\end{pgfscope}%
\end{pgfscope}%
\begin{pgfscope}%
\pgftext[x=1.115200in,y=0.430778in,,top]{\sffamily\fontsize{10.000000}{12.000000}\selectfont −2.0}%
\end{pgfscope}%
\begin{pgfscope}%
\pgfsetbuttcap%
\pgfsetroundjoin%
\definecolor{currentfill}{rgb}{0.000000,0.000000,0.000000}%
\pgfsetfillcolor{currentfill}%
\pgfsetlinewidth{0.803000pt}%
\definecolor{currentstroke}{rgb}{0.000000,0.000000,0.000000}%
\pgfsetstrokecolor{currentstroke}%
\pgfsetdash{}{0pt}%
\pgfsys@defobject{currentmarker}{\pgfqpoint{0.000000in}{-0.048611in}}{\pgfqpoint{0.000000in}{0.000000in}}{%
\pgfpathmoveto{\pgfqpoint{0.000000in}{0.000000in}}%
\pgfpathlineto{\pgfqpoint{0.000000in}{-0.048611in}}%
\pgfusepath{stroke,fill}%
}%
\begin{pgfscope}%
\pgfsys@transformshift{1.643200in}{0.528000in}%
\pgfsys@useobject{currentmarker}{}%
\end{pgfscope}%
\end{pgfscope}%
\begin{pgfscope}%
\pgftext[x=1.643200in,y=0.430778in,,top]{\sffamily\fontsize{10.000000}{12.000000}\selectfont −1.5}%
\end{pgfscope}%
\begin{pgfscope}%
\pgfsetbuttcap%
\pgfsetroundjoin%
\definecolor{currentfill}{rgb}{0.000000,0.000000,0.000000}%
\pgfsetfillcolor{currentfill}%
\pgfsetlinewidth{0.803000pt}%
\definecolor{currentstroke}{rgb}{0.000000,0.000000,0.000000}%
\pgfsetstrokecolor{currentstroke}%
\pgfsetdash{}{0pt}%
\pgfsys@defobject{currentmarker}{\pgfqpoint{0.000000in}{-0.048611in}}{\pgfqpoint{0.000000in}{0.000000in}}{%
\pgfpathmoveto{\pgfqpoint{0.000000in}{0.000000in}}%
\pgfpathlineto{\pgfqpoint{0.000000in}{-0.048611in}}%
\pgfusepath{stroke,fill}%
}%
\begin{pgfscope}%
\pgfsys@transformshift{2.171200in}{0.528000in}%
\pgfsys@useobject{currentmarker}{}%
\end{pgfscope}%
\end{pgfscope}%
\begin{pgfscope}%
\pgftext[x=2.171200in,y=0.430778in,,top]{\sffamily\fontsize{10.000000}{12.000000}\selectfont −1.0}%
\end{pgfscope}%
\begin{pgfscope}%
\pgfsetbuttcap%
\pgfsetroundjoin%
\definecolor{currentfill}{rgb}{0.000000,0.000000,0.000000}%
\pgfsetfillcolor{currentfill}%
\pgfsetlinewidth{0.803000pt}%
\definecolor{currentstroke}{rgb}{0.000000,0.000000,0.000000}%
\pgfsetstrokecolor{currentstroke}%
\pgfsetdash{}{0pt}%
\pgfsys@defobject{currentmarker}{\pgfqpoint{0.000000in}{-0.048611in}}{\pgfqpoint{0.000000in}{0.000000in}}{%
\pgfpathmoveto{\pgfqpoint{0.000000in}{0.000000in}}%
\pgfpathlineto{\pgfqpoint{0.000000in}{-0.048611in}}%
\pgfusepath{stroke,fill}%
}%
\begin{pgfscope}%
\pgfsys@transformshift{2.699200in}{0.528000in}%
\pgfsys@useobject{currentmarker}{}%
\end{pgfscope}%
\end{pgfscope}%
\begin{pgfscope}%
\pgftext[x=2.699200in,y=0.430778in,,top]{\sffamily\fontsize{10.000000}{12.000000}\selectfont −0.5}%
\end{pgfscope}%
\begin{pgfscope}%
\pgfsetbuttcap%
\pgfsetroundjoin%
\definecolor{currentfill}{rgb}{0.000000,0.000000,0.000000}%
\pgfsetfillcolor{currentfill}%
\pgfsetlinewidth{0.803000pt}%
\definecolor{currentstroke}{rgb}{0.000000,0.000000,0.000000}%
\pgfsetstrokecolor{currentstroke}%
\pgfsetdash{}{0pt}%
\pgfsys@defobject{currentmarker}{\pgfqpoint{0.000000in}{-0.048611in}}{\pgfqpoint{0.000000in}{0.000000in}}{%
\pgfpathmoveto{\pgfqpoint{0.000000in}{0.000000in}}%
\pgfpathlineto{\pgfqpoint{0.000000in}{-0.048611in}}%
\pgfusepath{stroke,fill}%
}%
\begin{pgfscope}%
\pgfsys@transformshift{3.227200in}{0.528000in}%
\pgfsys@useobject{currentmarker}{}%
\end{pgfscope}%
\end{pgfscope}%
\begin{pgfscope}%
\pgftext[x=3.227200in,y=0.430778in,,top]{\sffamily\fontsize{10.000000}{12.000000}\selectfont 0.0}%
\end{pgfscope}%
\begin{pgfscope}%
\pgfsetbuttcap%
\pgfsetroundjoin%
\definecolor{currentfill}{rgb}{0.000000,0.000000,0.000000}%
\pgfsetfillcolor{currentfill}%
\pgfsetlinewidth{0.803000pt}%
\definecolor{currentstroke}{rgb}{0.000000,0.000000,0.000000}%
\pgfsetstrokecolor{currentstroke}%
\pgfsetdash{}{0pt}%
\pgfsys@defobject{currentmarker}{\pgfqpoint{0.000000in}{-0.048611in}}{\pgfqpoint{0.000000in}{0.000000in}}{%
\pgfpathmoveto{\pgfqpoint{0.000000in}{0.000000in}}%
\pgfpathlineto{\pgfqpoint{0.000000in}{-0.048611in}}%
\pgfusepath{stroke,fill}%
}%
\begin{pgfscope}%
\pgfsys@transformshift{3.755200in}{0.528000in}%
\pgfsys@useobject{currentmarker}{}%
\end{pgfscope}%
\end{pgfscope}%
\begin{pgfscope}%
\pgftext[x=3.755200in,y=0.430778in,,top]{\sffamily\fontsize{10.000000}{12.000000}\selectfont 0.5}%
\end{pgfscope}%
\begin{pgfscope}%
\pgfsetbuttcap%
\pgfsetroundjoin%
\definecolor{currentfill}{rgb}{0.000000,0.000000,0.000000}%
\pgfsetfillcolor{currentfill}%
\pgfsetlinewidth{0.803000pt}%
\definecolor{currentstroke}{rgb}{0.000000,0.000000,0.000000}%
\pgfsetstrokecolor{currentstroke}%
\pgfsetdash{}{0pt}%
\pgfsys@defobject{currentmarker}{\pgfqpoint{0.000000in}{-0.048611in}}{\pgfqpoint{0.000000in}{0.000000in}}{%
\pgfpathmoveto{\pgfqpoint{0.000000in}{0.000000in}}%
\pgfpathlineto{\pgfqpoint{0.000000in}{-0.048611in}}%
\pgfusepath{stroke,fill}%
}%
\begin{pgfscope}%
\pgfsys@transformshift{4.283200in}{0.528000in}%
\pgfsys@useobject{currentmarker}{}%
\end{pgfscope}%
\end{pgfscope}%
\begin{pgfscope}%
\pgftext[x=4.283200in,y=0.430778in,,top]{\sffamily\fontsize{10.000000}{12.000000}\selectfont 1.0}%
\end{pgfscope}%
\begin{pgfscope}%
\pgfsetbuttcap%
\pgfsetroundjoin%
\definecolor{currentfill}{rgb}{0.000000,0.000000,0.000000}%
\pgfsetfillcolor{currentfill}%
\pgfsetlinewidth{0.803000pt}%
\definecolor{currentstroke}{rgb}{0.000000,0.000000,0.000000}%
\pgfsetstrokecolor{currentstroke}%
\pgfsetdash{}{0pt}%
\pgfsys@defobject{currentmarker}{\pgfqpoint{0.000000in}{-0.048611in}}{\pgfqpoint{0.000000in}{0.000000in}}{%
\pgfpathmoveto{\pgfqpoint{0.000000in}{0.000000in}}%
\pgfpathlineto{\pgfqpoint{0.000000in}{-0.048611in}}%
\pgfusepath{stroke,fill}%
}%
\begin{pgfscope}%
\pgfsys@transformshift{4.811200in}{0.528000in}%
\pgfsys@useobject{currentmarker}{}%
\end{pgfscope}%
\end{pgfscope}%
\begin{pgfscope}%
\pgftext[x=4.811200in,y=0.430778in,,top]{\sffamily\fontsize{10.000000}{12.000000}\selectfont 1.5}%
\end{pgfscope}%
\begin{pgfscope}%
\pgfsetbuttcap%
\pgfsetroundjoin%
\definecolor{currentfill}{rgb}{0.000000,0.000000,0.000000}%
\pgfsetfillcolor{currentfill}%
\pgfsetlinewidth{0.803000pt}%
\definecolor{currentstroke}{rgb}{0.000000,0.000000,0.000000}%
\pgfsetstrokecolor{currentstroke}%
\pgfsetdash{}{0pt}%
\pgfsys@defobject{currentmarker}{\pgfqpoint{0.000000in}{-0.048611in}}{\pgfqpoint{0.000000in}{0.000000in}}{%
\pgfpathmoveto{\pgfqpoint{0.000000in}{0.000000in}}%
\pgfpathlineto{\pgfqpoint{0.000000in}{-0.048611in}}%
\pgfusepath{stroke,fill}%
}%
\begin{pgfscope}%
\pgfsys@transformshift{5.339200in}{0.528000in}%
\pgfsys@useobject{currentmarker}{}%
\end{pgfscope}%
\end{pgfscope}%
\begin{pgfscope}%
\pgftext[x=5.339200in,y=0.430778in,,top]{\sffamily\fontsize{10.000000}{12.000000}\selectfont 2.0}%
\end{pgfscope}%
\begin{pgfscope}%
\pgftext[x=3.280000in,y=0.240809in,,top]{\sffamily\fontsize{10.000000}{12.000000}\selectfont x}%
\end{pgfscope}%
\begin{pgfscope}%
\pgfsetbuttcap%
\pgfsetroundjoin%
\definecolor{currentfill}{rgb}{0.000000,0.000000,0.000000}%
\pgfsetfillcolor{currentfill}%
\pgfsetlinewidth{0.803000pt}%
\definecolor{currentstroke}{rgb}{0.000000,0.000000,0.000000}%
\pgfsetstrokecolor{currentstroke}%
\pgfsetdash{}{0pt}%
\pgfsys@defobject{currentmarker}{\pgfqpoint{-0.048611in}{0.000000in}}{\pgfqpoint{0.000000in}{0.000000in}}{%
\pgfpathmoveto{\pgfqpoint{0.000000in}{0.000000in}}%
\pgfpathlineto{\pgfqpoint{-0.048611in}{0.000000in}}%
\pgfusepath{stroke,fill}%
}%
\begin{pgfscope}%
\pgfsys@transformshift{1.115200in}{0.528000in}%
\pgfsys@useobject{currentmarker}{}%
\end{pgfscope}%
\end{pgfscope}%
\begin{pgfscope}%
\pgftext[x=0.680725in,y=0.475238in,left,base]{\sffamily\fontsize{10.000000}{12.000000}\selectfont −2.0}%
\end{pgfscope}%
\begin{pgfscope}%
\pgfsetbuttcap%
\pgfsetroundjoin%
\definecolor{currentfill}{rgb}{0.000000,0.000000,0.000000}%
\pgfsetfillcolor{currentfill}%
\pgfsetlinewidth{0.803000pt}%
\definecolor{currentstroke}{rgb}{0.000000,0.000000,0.000000}%
\pgfsetstrokecolor{currentstroke}%
\pgfsetdash{}{0pt}%
\pgfsys@defobject{currentmarker}{\pgfqpoint{-0.048611in}{0.000000in}}{\pgfqpoint{0.000000in}{0.000000in}}{%
\pgfpathmoveto{\pgfqpoint{0.000000in}{0.000000in}}%
\pgfpathlineto{\pgfqpoint{-0.048611in}{0.000000in}}%
\pgfusepath{stroke,fill}%
}%
\begin{pgfscope}%
\pgfsys@transformshift{1.115200in}{1.056000in}%
\pgfsys@useobject{currentmarker}{}%
\end{pgfscope}%
\end{pgfscope}%
\begin{pgfscope}%
\pgftext[x=0.680725in,y=1.003238in,left,base]{\sffamily\fontsize{10.000000}{12.000000}\selectfont −1.5}%
\end{pgfscope}%
\begin{pgfscope}%
\pgfsetbuttcap%
\pgfsetroundjoin%
\definecolor{currentfill}{rgb}{0.000000,0.000000,0.000000}%
\pgfsetfillcolor{currentfill}%
\pgfsetlinewidth{0.803000pt}%
\definecolor{currentstroke}{rgb}{0.000000,0.000000,0.000000}%
\pgfsetstrokecolor{currentstroke}%
\pgfsetdash{}{0pt}%
\pgfsys@defobject{currentmarker}{\pgfqpoint{-0.048611in}{0.000000in}}{\pgfqpoint{0.000000in}{0.000000in}}{%
\pgfpathmoveto{\pgfqpoint{0.000000in}{0.000000in}}%
\pgfpathlineto{\pgfqpoint{-0.048611in}{0.000000in}}%
\pgfusepath{stroke,fill}%
}%
\begin{pgfscope}%
\pgfsys@transformshift{1.115200in}{1.584000in}%
\pgfsys@useobject{currentmarker}{}%
\end{pgfscope}%
\end{pgfscope}%
\begin{pgfscope}%
\pgftext[x=0.680725in,y=1.531238in,left,base]{\sffamily\fontsize{10.000000}{12.000000}\selectfont −1.0}%
\end{pgfscope}%
\begin{pgfscope}%
\pgfsetbuttcap%
\pgfsetroundjoin%
\definecolor{currentfill}{rgb}{0.000000,0.000000,0.000000}%
\pgfsetfillcolor{currentfill}%
\pgfsetlinewidth{0.803000pt}%
\definecolor{currentstroke}{rgb}{0.000000,0.000000,0.000000}%
\pgfsetstrokecolor{currentstroke}%
\pgfsetdash{}{0pt}%
\pgfsys@defobject{currentmarker}{\pgfqpoint{-0.048611in}{0.000000in}}{\pgfqpoint{0.000000in}{0.000000in}}{%
\pgfpathmoveto{\pgfqpoint{0.000000in}{0.000000in}}%
\pgfpathlineto{\pgfqpoint{-0.048611in}{0.000000in}}%
\pgfusepath{stroke,fill}%
}%
\begin{pgfscope}%
\pgfsys@transformshift{1.115200in}{2.112000in}%
\pgfsys@useobject{currentmarker}{}%
\end{pgfscope}%
\end{pgfscope}%
\begin{pgfscope}%
\pgftext[x=0.680725in,y=2.059238in,left,base]{\sffamily\fontsize{10.000000}{12.000000}\selectfont −0.5}%
\end{pgfscope}%
\begin{pgfscope}%
\pgfsetbuttcap%
\pgfsetroundjoin%
\definecolor{currentfill}{rgb}{0.000000,0.000000,0.000000}%
\pgfsetfillcolor{currentfill}%
\pgfsetlinewidth{0.803000pt}%
\definecolor{currentstroke}{rgb}{0.000000,0.000000,0.000000}%
\pgfsetstrokecolor{currentstroke}%
\pgfsetdash{}{0pt}%
\pgfsys@defobject{currentmarker}{\pgfqpoint{-0.048611in}{0.000000in}}{\pgfqpoint{0.000000in}{0.000000in}}{%
\pgfpathmoveto{\pgfqpoint{0.000000in}{0.000000in}}%
\pgfpathlineto{\pgfqpoint{-0.048611in}{0.000000in}}%
\pgfusepath{stroke,fill}%
}%
\begin{pgfscope}%
\pgfsys@transformshift{1.115200in}{2.640000in}%
\pgfsys@useobject{currentmarker}{}%
\end{pgfscope}%
\end{pgfscope}%
\begin{pgfscope}%
\pgftext[x=0.797098in,y=2.587238in,left,base]{\sffamily\fontsize{10.000000}{12.000000}\selectfont 0.0}%
\end{pgfscope}%
\begin{pgfscope}%
\pgfsetbuttcap%
\pgfsetroundjoin%
\definecolor{currentfill}{rgb}{0.000000,0.000000,0.000000}%
\pgfsetfillcolor{currentfill}%
\pgfsetlinewidth{0.803000pt}%
\definecolor{currentstroke}{rgb}{0.000000,0.000000,0.000000}%
\pgfsetstrokecolor{currentstroke}%
\pgfsetdash{}{0pt}%
\pgfsys@defobject{currentmarker}{\pgfqpoint{-0.048611in}{0.000000in}}{\pgfqpoint{0.000000in}{0.000000in}}{%
\pgfpathmoveto{\pgfqpoint{0.000000in}{0.000000in}}%
\pgfpathlineto{\pgfqpoint{-0.048611in}{0.000000in}}%
\pgfusepath{stroke,fill}%
}%
\begin{pgfscope}%
\pgfsys@transformshift{1.115200in}{3.168000in}%
\pgfsys@useobject{currentmarker}{}%
\end{pgfscope}%
\end{pgfscope}%
\begin{pgfscope}%
\pgftext[x=0.797098in,y=3.115238in,left,base]{\sffamily\fontsize{10.000000}{12.000000}\selectfont 0.5}%
\end{pgfscope}%
\begin{pgfscope}%
\pgfsetbuttcap%
\pgfsetroundjoin%
\definecolor{currentfill}{rgb}{0.000000,0.000000,0.000000}%
\pgfsetfillcolor{currentfill}%
\pgfsetlinewidth{0.803000pt}%
\definecolor{currentstroke}{rgb}{0.000000,0.000000,0.000000}%
\pgfsetstrokecolor{currentstroke}%
\pgfsetdash{}{0pt}%
\pgfsys@defobject{currentmarker}{\pgfqpoint{-0.048611in}{0.000000in}}{\pgfqpoint{0.000000in}{0.000000in}}{%
\pgfpathmoveto{\pgfqpoint{0.000000in}{0.000000in}}%
\pgfpathlineto{\pgfqpoint{-0.048611in}{0.000000in}}%
\pgfusepath{stroke,fill}%
}%
\begin{pgfscope}%
\pgfsys@transformshift{1.115200in}{3.696000in}%
\pgfsys@useobject{currentmarker}{}%
\end{pgfscope}%
\end{pgfscope}%
\begin{pgfscope}%
\pgftext[x=0.797098in,y=3.643238in,left,base]{\sffamily\fontsize{10.000000}{12.000000}\selectfont 1.0}%
\end{pgfscope}%
\begin{pgfscope}%
\pgfsetbuttcap%
\pgfsetroundjoin%
\definecolor{currentfill}{rgb}{0.000000,0.000000,0.000000}%
\pgfsetfillcolor{currentfill}%
\pgfsetlinewidth{0.803000pt}%
\definecolor{currentstroke}{rgb}{0.000000,0.000000,0.000000}%
\pgfsetstrokecolor{currentstroke}%
\pgfsetdash{}{0pt}%
\pgfsys@defobject{currentmarker}{\pgfqpoint{-0.048611in}{0.000000in}}{\pgfqpoint{0.000000in}{0.000000in}}{%
\pgfpathmoveto{\pgfqpoint{0.000000in}{0.000000in}}%
\pgfpathlineto{\pgfqpoint{-0.048611in}{0.000000in}}%
\pgfusepath{stroke,fill}%
}%
\begin{pgfscope}%
\pgfsys@transformshift{1.115200in}{4.224000in}%
\pgfsys@useobject{currentmarker}{}%
\end{pgfscope}%
\end{pgfscope}%
\begin{pgfscope}%
\pgftext[x=0.797098in,y=4.171238in,left,base]{\sffamily\fontsize{10.000000}{12.000000}\selectfont 1.5}%
\end{pgfscope}%
\begin{pgfscope}%
\pgftext[x=0.625169in,y=2.376000in,,bottom]{\sffamily\fontsize{10.000000}{12.000000}\selectfont y}%
\end{pgfscope}%
\begin{pgfscope}%
\pgfsetrectcap%
\pgfsetmiterjoin%
\pgfsetlinewidth{0.803000pt}%
\definecolor{currentstroke}{rgb}{0.000000,0.000000,0.000000}%
\pgfsetstrokecolor{currentstroke}%
\pgfsetdash{}{0pt}%
\pgfpathmoveto{\pgfqpoint{1.115200in}{0.528000in}}%
\pgfpathlineto{\pgfqpoint{1.115200in}{4.224000in}}%
\pgfusepath{stroke}%
\end{pgfscope}%
\begin{pgfscope}%
\pgfsetrectcap%
\pgfsetmiterjoin%
\pgfsetlinewidth{0.803000pt}%
\definecolor{currentstroke}{rgb}{0.000000,0.000000,0.000000}%
\pgfsetstrokecolor{currentstroke}%
\pgfsetdash{}{0pt}%
\pgfpathmoveto{\pgfqpoint{5.444800in}{0.528000in}}%
\pgfpathlineto{\pgfqpoint{5.444800in}{4.224000in}}%
\pgfusepath{stroke}%
\end{pgfscope}%
\begin{pgfscope}%
\pgfsetrectcap%
\pgfsetmiterjoin%
\pgfsetlinewidth{0.803000pt}%
\definecolor{currentstroke}{rgb}{0.000000,0.000000,0.000000}%
\pgfsetstrokecolor{currentstroke}%
\pgfsetdash{}{0pt}%
\pgfpathmoveto{\pgfqpoint{1.115200in}{0.528000in}}%
\pgfpathlineto{\pgfqpoint{5.444800in}{0.528000in}}%
\pgfusepath{stroke}%
\end{pgfscope}%
\begin{pgfscope}%
\pgfsetrectcap%
\pgfsetmiterjoin%
\pgfsetlinewidth{0.803000pt}%
\definecolor{currentstroke}{rgb}{0.000000,0.000000,0.000000}%
\pgfsetstrokecolor{currentstroke}%
\pgfsetdash{}{0pt}%
\pgfpathmoveto{\pgfqpoint{1.115200in}{4.224000in}}%
\pgfpathlineto{\pgfqpoint{5.444800in}{4.224000in}}%
\pgfusepath{stroke}%
\end{pgfscope}%
\end{pgfpicture}%
\makeatother%
\endgroup%
}
	\caption{Klasterovanje}
	\label{fig:klaster}
\end{figure}

\section{Dizajn sistema za mašinsko učenje}

Okvirno, koraci u rešavanju problema su sledeći \cite{janicic_vi}:
\begin{itemize}
	\item Prepoznavanje problema mašinskog učenja (nadgledano učenje, \\nenadgledano učenje, učenje potkrepljivanjem);
	\item Prikupljanje i obrada podataka, zajedno sa odabirom atributa;
	\item Odabir skupa dopustivih modela;
	\item Odabir algoritma učenja; moguće je odabrati postojeći algoritam ili razviti neki novi koji bolje odgovara problemu
	\item Izbor mere kvaliteta učenja;
	\item Obuka, evaluacija i, ukoliko je neophodno, ponavljanje nekog od prethodnih koraka radi unapređenja naučenog modela 
\end{itemize}

Prilikom odabira modela treba imati na umu vrstu problema koja se rešava, količinu podataka, zakonitosti koje važe u podacima, itd.


\subsection{Podaci}

Mašinsko učenje bavi se generalizacijom nad nepoznatim objektima na osnovu već viđenih objekata. Pod pojmom objekta misli se na pojedinačni podatak koji sistem vidi. Koriste se još i izrazi primerak i instanca. Vrednosti podataka pripadaju nekom unapred zadatom skupu. Podaci mogu biti različitog tipa: numerički ili kategorički. Skupovi koji određuju vrednosti kojima se instance određuju nisu unapred zadati i neophodno ih je odrediti na način pogodan za rešavanje konkretnog problema. Na primer, ukoliko je neophodno razvrstati slike životinja i biljaka na te dve kategorije, informacija o količini zelene boje na slici može biti prilično korisna, dok pri razvrstavanju vrste biljaka u zavisnosti od lista ovaj podatak skoro nije upotrebljiv (ali podatak o nijansi zelene boje može biti). Dakle, dobar izbor atributa imaće veliki uticaj na kasnije korake učenja. Podaci se sistemu daju kao vektori atributa. Nije uvek neophodno vršiti ekstrakciju atributa. Na primer, neuronske mreže u stanju su da uče nad sirovim podacima.
\par

Podaci se neretko pre slanja sistemu obrađuju na neki način; ovaj postupak zove se pretprocesiranje. Postoje mnogi razlozi za pretprocesiranje a glavni cilj jeste da se dobiju objekti nad kojima učenje može da se vrši, imajući u vidu zahteve algoritama učenja. Međutim, i to zavisi od problema. Nekada će nepotpuni objekti, podaci koji ne sadrže sve informacije neophodne za učenje, biti izbačeni iz skupa podataka koji se razmatra, a u nekom drugom slučaju, i oni će biti korišćeni. Primeri pretprocesiranja su pretvaranje slike koja je u boji u crno beli zapis, normalizacija, umetanje nedostajućih vrednosti, itd.

\subsection{Evaluacija modela}

Nakon obučavanja (treniranja), neophodno je izvršiti evaluaciju dobijenog modela. Na koji god način se ovo izvršava, podaci korišćeni za obučavanje ne smeju se koristiti za evaluaciju modela. Često se pribegava podeli podataka na skupove za obučavanje i za testiranje. Skup za obučavanje obično iznosi dve trećine skupa ukupnih podataka. No, kako različite podele skupa mogu izazvati dobijanje različitih modela, slučajno deljenje nije najbolji izbor.  Često korišćena tehnika jeste unakrsna validacija. Ovaj pristup podrazumeva podelu skupa podataka $D$ na $K$ podskupova približno jednake veličine, $S_i$ za $i=1,...,K$. Tada se za svako $i$ model trenira na skupu $D \setminus S_i$ a evaluacija se vrši pomoću podataka iz $S_i$. Posle izvedenog postupka za sve $i$, kao konačna ocena uzima se prosečna ocena svakog od $K$ treniranja i evaluacija modela. Za vrednosti $K$ uobičajeno se uzimaju vrednosti 5 ili 10. Ovaj metod vodi pouzdanijoj oceni kvaliteta modela.

\section{Problemi pri mašinskom učenju}


Kao što je podrazumevano pri pomenu pojma generalizacije, nije dovoljno odrediti funkciju koja dobro određuje izlazne vrednosti na osnovu promenljivih nad kojima se uči već je poželjno i novim ulaznim podacima dodeliti tačnu izlaznu vrednost. Odavde se može videti da je primer lošeg sistema za mašinsko učenje onaj sistem koji će izuzetno dobro naučiti da preslikava ulazne vrednosti iz skupa za učenje u odgovarajuće izlazne vrednosti ali u situaciji kada se iz tog skupa izađe neće davati zadovoljavajuće rezultate. Ovaj problem ima svoje ime: preprilagođavanje. Postoji i problem potprilagođavanja, koji podrazumeva da se sistem nije dovoljno prilagodio podacima. I preprilagođavanje i potprilagođavanje predstavljaju veliki problem ukoliko do njih dođe. Primer preprilagođavanja može se videti na slici \ref{fig:odabir}, koja prikazuje razliku između dva modela iz skupa dopustivih modela za linearnu regresiju polinomom nad 10 različitih tačaka. Na levom delu slike prikazan je polinom reda 1 (prava) a na desnom delu prikazan je polinom reda 10. Iako će polinom reda 10 savršeno opisivati 10 tačaka sa slike, prava će verovatno bolje generalizovati nad novim podacima.
\begin{figure}
	\centering
	\resizebox{.45\linewidth}{!}{%% Creator: Matplotlib, PGF backend
%%
%% To include the figure in your LaTeX document, write
%%   \input{<filename>.pgf}
%%
%% Make sure the required packages are loaded in your preamble
%%   \usepackage{pgf}
%%
%% Figures using additional raster images can only be included by \input if
%% they are in the same directory as the main LaTeX file. For loading figures
%% from other directories you can use the `import` package
%%   \usepackage{import}
%% and then include the figures with
%%   \import{<path to file>}{<filename>.pgf}
%%
%% Matplotlib used the following preamble
%%   \usepackage{fontspec}
%%   \setmainfont{Bitstream Vera Serif}
%%   \setsansfont{Bitstream Vera Sans}
%%   \setmonofont{Bitstream Vera Sans Mono}
%%
\begingroup%
\makeatletter%
\begin{pgfpicture}%
\pgfpathrectangle{\pgfpointorigin}{\pgfqpoint{8.000000in}{6.000000in}}%
\pgfusepath{use as bounding box, clip}%
\begin{pgfscope}%
\pgfsetbuttcap%
\pgfsetmiterjoin%
\definecolor{currentfill}{rgb}{1.000000,1.000000,1.000000}%
\pgfsetfillcolor{currentfill}%
\pgfsetlinewidth{0.000000pt}%
\definecolor{currentstroke}{rgb}{1.000000,1.000000,1.000000}%
\pgfsetstrokecolor{currentstroke}%
\pgfsetdash{}{0pt}%
\pgfpathmoveto{\pgfqpoint{0.000000in}{0.000000in}}%
\pgfpathlineto{\pgfqpoint{8.000000in}{0.000000in}}%
\pgfpathlineto{\pgfqpoint{8.000000in}{6.000000in}}%
\pgfpathlineto{\pgfqpoint{0.000000in}{6.000000in}}%
\pgfpathclose%
\pgfusepath{fill}%
\end{pgfscope}%
\begin{pgfscope}%
\pgfsetbuttcap%
\pgfsetmiterjoin%
\definecolor{currentfill}{rgb}{1.000000,1.000000,1.000000}%
\pgfsetfillcolor{currentfill}%
\pgfsetlinewidth{0.000000pt}%
\definecolor{currentstroke}{rgb}{0.000000,0.000000,0.000000}%
\pgfsetstrokecolor{currentstroke}%
\pgfsetstrokeopacity{0.000000}%
\pgfsetdash{}{0pt}%
\pgfpathmoveto{\pgfqpoint{1.000000in}{0.600000in}}%
\pgfpathlineto{\pgfqpoint{7.200000in}{0.600000in}}%
\pgfpathlineto{\pgfqpoint{7.200000in}{5.400000in}}%
\pgfpathlineto{\pgfqpoint{1.000000in}{5.400000in}}%
\pgfpathclose%
\pgfusepath{fill}%
\end{pgfscope}%
\begin{pgfscope}%
\pgfpathrectangle{\pgfqpoint{1.000000in}{0.600000in}}{\pgfqpoint{6.200000in}{4.800000in}} %
\pgfusepath{clip}%
\pgfsetbuttcap%
\pgfsetroundjoin%
\definecolor{currentfill}{rgb}{0.000000,0.500000,0.000000}%
\pgfsetfillcolor{currentfill}%
\pgfsetlinewidth{1.003750pt}%
\definecolor{currentstroke}{rgb}{0.000000,0.500000,0.000000}%
\pgfsetstrokecolor{currentstroke}%
\pgfsetdash{}{0pt}%
\pgfpathmoveto{\pgfqpoint{1.442857in}{0.663119in}}%
\pgfpathcurveto{\pgfqpoint{1.459330in}{0.663119in}}{\pgfqpoint{1.475130in}{0.669664in}}{\pgfqpoint{1.486778in}{0.681312in}}%
\pgfpathcurveto{\pgfqpoint{1.498426in}{0.692960in}}{\pgfqpoint{1.504970in}{0.708760in}}{\pgfqpoint{1.504970in}{0.725232in}}%
\pgfpathcurveto{\pgfqpoint{1.504970in}{0.741705in}}{\pgfqpoint{1.498426in}{0.757505in}}{\pgfqpoint{1.486778in}{0.769153in}}%
\pgfpathcurveto{\pgfqpoint{1.475130in}{0.780801in}}{\pgfqpoint{1.459330in}{0.787345in}}{\pgfqpoint{1.442857in}{0.787345in}}%
\pgfpathcurveto{\pgfqpoint{1.426385in}{0.787345in}}{\pgfqpoint{1.410584in}{0.780801in}}{\pgfqpoint{1.398937in}{0.769153in}}%
\pgfpathcurveto{\pgfqpoint{1.387289in}{0.757505in}}{\pgfqpoint{1.380744in}{0.741705in}}{\pgfqpoint{1.380744in}{0.725232in}}%
\pgfpathcurveto{\pgfqpoint{1.380744in}{0.708760in}}{\pgfqpoint{1.387289in}{0.692960in}}{\pgfqpoint{1.398937in}{0.681312in}}%
\pgfpathcurveto{\pgfqpoint{1.410584in}{0.669664in}}{\pgfqpoint{1.426385in}{0.663119in}}{\pgfqpoint{1.442857in}{0.663119in}}%
\pgfpathclose%
\pgfusepath{stroke,fill}%
\end{pgfscope}%
\begin{pgfscope}%
\pgfpathrectangle{\pgfqpoint{1.000000in}{0.600000in}}{\pgfqpoint{6.200000in}{4.800000in}} %
\pgfusepath{clip}%
\pgfsetbuttcap%
\pgfsetroundjoin%
\definecolor{currentfill}{rgb}{0.000000,0.500000,0.000000}%
\pgfsetfillcolor{currentfill}%
\pgfsetlinewidth{1.003750pt}%
\definecolor{currentstroke}{rgb}{0.000000,0.500000,0.000000}%
\pgfsetstrokecolor{currentstroke}%
\pgfsetdash{}{0pt}%
\pgfpathmoveto{\pgfqpoint{2.057937in}{1.339529in}}%
\pgfpathcurveto{\pgfqpoint{2.074409in}{1.339529in}}{\pgfqpoint{2.090209in}{1.346074in}}{\pgfqpoint{2.101857in}{1.357722in}}%
\pgfpathcurveto{\pgfqpoint{2.113505in}{1.369370in}}{\pgfqpoint{2.120050in}{1.385170in}}{\pgfqpoint{2.120050in}{1.401642in}}%
\pgfpathcurveto{\pgfqpoint{2.120050in}{1.418115in}}{\pgfqpoint{2.113505in}{1.433915in}}{\pgfqpoint{2.101857in}{1.445563in}}%
\pgfpathcurveto{\pgfqpoint{2.090209in}{1.457211in}}{\pgfqpoint{2.074409in}{1.463755in}}{\pgfqpoint{2.057937in}{1.463755in}}%
\pgfpathcurveto{\pgfqpoint{2.041464in}{1.463755in}}{\pgfqpoint{2.025664in}{1.457211in}}{\pgfqpoint{2.014016in}{1.445563in}}%
\pgfpathcurveto{\pgfqpoint{2.002368in}{1.433915in}}{\pgfqpoint{1.995824in}{1.418115in}}{\pgfqpoint{1.995824in}{1.401642in}}%
\pgfpathcurveto{\pgfqpoint{1.995824in}{1.385170in}}{\pgfqpoint{2.002368in}{1.369370in}}{\pgfqpoint{2.014016in}{1.357722in}}%
\pgfpathcurveto{\pgfqpoint{2.025664in}{1.346074in}}{\pgfqpoint{2.041464in}{1.339529in}}{\pgfqpoint{2.057937in}{1.339529in}}%
\pgfpathclose%
\pgfusepath{stroke,fill}%
\end{pgfscope}%
\begin{pgfscope}%
\pgfpathrectangle{\pgfqpoint{1.000000in}{0.600000in}}{\pgfqpoint{6.200000in}{4.800000in}} %
\pgfusepath{clip}%
\pgfsetbuttcap%
\pgfsetroundjoin%
\definecolor{currentfill}{rgb}{0.000000,0.500000,0.000000}%
\pgfsetfillcolor{currentfill}%
\pgfsetlinewidth{1.003750pt}%
\definecolor{currentstroke}{rgb}{0.000000,0.500000,0.000000}%
\pgfsetstrokecolor{currentstroke}%
\pgfsetdash{}{0pt}%
\pgfpathmoveto{\pgfqpoint{2.673016in}{1.632439in}}%
\pgfpathcurveto{\pgfqpoint{2.689488in}{1.632439in}}{\pgfqpoint{2.705289in}{1.638984in}}{\pgfqpoint{2.716936in}{1.650632in}}%
\pgfpathcurveto{\pgfqpoint{2.728584in}{1.662279in}}{\pgfqpoint{2.735129in}{1.678080in}}{\pgfqpoint{2.735129in}{1.694552in}}%
\pgfpathcurveto{\pgfqpoint{2.735129in}{1.711025in}}{\pgfqpoint{2.728584in}{1.726825in}}{\pgfqpoint{2.716936in}{1.738473in}}%
\pgfpathcurveto{\pgfqpoint{2.705289in}{1.750120in}}{\pgfqpoint{2.689488in}{1.756665in}}{\pgfqpoint{2.673016in}{1.756665in}}%
\pgfpathcurveto{\pgfqpoint{2.656543in}{1.756665in}}{\pgfqpoint{2.640743in}{1.750120in}}{\pgfqpoint{2.629095in}{1.738473in}}%
\pgfpathcurveto{\pgfqpoint{2.617447in}{1.726825in}}{\pgfqpoint{2.610903in}{1.711025in}}{\pgfqpoint{2.610903in}{1.694552in}}%
\pgfpathcurveto{\pgfqpoint{2.610903in}{1.678080in}}{\pgfqpoint{2.617447in}{1.662279in}}{\pgfqpoint{2.629095in}{1.650632in}}%
\pgfpathcurveto{\pgfqpoint{2.640743in}{1.638984in}}{\pgfqpoint{2.656543in}{1.632439in}}{\pgfqpoint{2.673016in}{1.632439in}}%
\pgfpathclose%
\pgfusepath{stroke,fill}%
\end{pgfscope}%
\begin{pgfscope}%
\pgfpathrectangle{\pgfqpoint{1.000000in}{0.600000in}}{\pgfqpoint{6.200000in}{4.800000in}} %
\pgfusepath{clip}%
\pgfsetbuttcap%
\pgfsetroundjoin%
\definecolor{currentfill}{rgb}{0.000000,0.500000,0.000000}%
\pgfsetfillcolor{currentfill}%
\pgfsetlinewidth{1.003750pt}%
\definecolor{currentstroke}{rgb}{0.000000,0.500000,0.000000}%
\pgfsetstrokecolor{currentstroke}%
\pgfsetdash{}{0pt}%
\pgfpathmoveto{\pgfqpoint{3.288095in}{2.327298in}}%
\pgfpathcurveto{\pgfqpoint{3.304568in}{2.327298in}}{\pgfqpoint{3.320368in}{2.333843in}}{\pgfqpoint{3.332016in}{2.345491in}}%
\pgfpathcurveto{\pgfqpoint{3.343664in}{2.357138in}}{\pgfqpoint{3.350208in}{2.372939in}}{\pgfqpoint{3.350208in}{2.389411in}}%
\pgfpathcurveto{\pgfqpoint{3.350208in}{2.405884in}}{\pgfqpoint{3.343664in}{2.421684in}}{\pgfqpoint{3.332016in}{2.433332in}}%
\pgfpathcurveto{\pgfqpoint{3.320368in}{2.444979in}}{\pgfqpoint{3.304568in}{2.451524in}}{\pgfqpoint{3.288095in}{2.451524in}}%
\pgfpathcurveto{\pgfqpoint{3.271623in}{2.451524in}}{\pgfqpoint{3.255823in}{2.444979in}}{\pgfqpoint{3.244175in}{2.433332in}}%
\pgfpathcurveto{\pgfqpoint{3.232527in}{2.421684in}}{\pgfqpoint{3.225982in}{2.405884in}}{\pgfqpoint{3.225982in}{2.389411in}}%
\pgfpathcurveto{\pgfqpoint{3.225982in}{2.372939in}}{\pgfqpoint{3.232527in}{2.357138in}}{\pgfqpoint{3.244175in}{2.345491in}}%
\pgfpathcurveto{\pgfqpoint{3.255823in}{2.333843in}}{\pgfqpoint{3.271623in}{2.327298in}}{\pgfqpoint{3.288095in}{2.327298in}}%
\pgfpathclose%
\pgfusepath{stroke,fill}%
\end{pgfscope}%
\begin{pgfscope}%
\pgfpathrectangle{\pgfqpoint{1.000000in}{0.600000in}}{\pgfqpoint{6.200000in}{4.800000in}} %
\pgfusepath{clip}%
\pgfsetbuttcap%
\pgfsetroundjoin%
\definecolor{currentfill}{rgb}{0.000000,0.500000,0.000000}%
\pgfsetfillcolor{currentfill}%
\pgfsetlinewidth{1.003750pt}%
\definecolor{currentstroke}{rgb}{0.000000,0.500000,0.000000}%
\pgfsetstrokecolor{currentstroke}%
\pgfsetdash{}{0pt}%
\pgfpathmoveto{\pgfqpoint{3.903175in}{2.688513in}}%
\pgfpathcurveto{\pgfqpoint{3.919647in}{2.688513in}}{\pgfqpoint{3.935447in}{2.695057in}}{\pgfqpoint{3.947095in}{2.706705in}}%
\pgfpathcurveto{\pgfqpoint{3.958743in}{2.718353in}}{\pgfqpoint{3.965288in}{2.734153in}}{\pgfqpoint{3.965288in}{2.750626in}}%
\pgfpathcurveto{\pgfqpoint{3.965288in}{2.767098in}}{\pgfqpoint{3.958743in}{2.782898in}}{\pgfqpoint{3.947095in}{2.794546in}}%
\pgfpathcurveto{\pgfqpoint{3.935447in}{2.806194in}}{\pgfqpoint{3.919647in}{2.812739in}}{\pgfqpoint{3.903175in}{2.812739in}}%
\pgfpathcurveto{\pgfqpoint{3.886702in}{2.812739in}}{\pgfqpoint{3.870902in}{2.806194in}}{\pgfqpoint{3.859254in}{2.794546in}}%
\pgfpathcurveto{\pgfqpoint{3.847606in}{2.782898in}}{\pgfqpoint{3.841062in}{2.767098in}}{\pgfqpoint{3.841062in}{2.750626in}}%
\pgfpathcurveto{\pgfqpoint{3.841062in}{2.734153in}}{\pgfqpoint{3.847606in}{2.718353in}}{\pgfqpoint{3.859254in}{2.706705in}}%
\pgfpathcurveto{\pgfqpoint{3.870902in}{2.695057in}}{\pgfqpoint{3.886702in}{2.688513in}}{\pgfqpoint{3.903175in}{2.688513in}}%
\pgfpathclose%
\pgfusepath{stroke,fill}%
\end{pgfscope}%
\begin{pgfscope}%
\pgfpathrectangle{\pgfqpoint{1.000000in}{0.600000in}}{\pgfqpoint{6.200000in}{4.800000in}} %
\pgfusepath{clip}%
\pgfsetbuttcap%
\pgfsetroundjoin%
\definecolor{currentfill}{rgb}{0.000000,0.500000,0.000000}%
\pgfsetfillcolor{currentfill}%
\pgfsetlinewidth{1.003750pt}%
\definecolor{currentstroke}{rgb}{0.000000,0.500000,0.000000}%
\pgfsetstrokecolor{currentstroke}%
\pgfsetdash{}{0pt}%
\pgfpathmoveto{\pgfqpoint{4.518254in}{3.211433in}}%
\pgfpathcurveto{\pgfqpoint{4.534727in}{3.211433in}}{\pgfqpoint{4.550527in}{3.217978in}}{\pgfqpoint{4.562174in}{3.229626in}}%
\pgfpathcurveto{\pgfqpoint{4.573822in}{3.241274in}}{\pgfqpoint{4.580367in}{3.257074in}}{\pgfqpoint{4.580367in}{3.273546in}}%
\pgfpathcurveto{\pgfqpoint{4.580367in}{3.290019in}}{\pgfqpoint{4.573822in}{3.305819in}}{\pgfqpoint{4.562174in}{3.317467in}}%
\pgfpathcurveto{\pgfqpoint{4.550527in}{3.329115in}}{\pgfqpoint{4.534727in}{3.335659in}}{\pgfqpoint{4.518254in}{3.335659in}}%
\pgfpathcurveto{\pgfqpoint{4.501781in}{3.335659in}}{\pgfqpoint{4.485981in}{3.329115in}}{\pgfqpoint{4.474333in}{3.317467in}}%
\pgfpathcurveto{\pgfqpoint{4.462686in}{3.305819in}}{\pgfqpoint{4.456141in}{3.290019in}}{\pgfqpoint{4.456141in}{3.273546in}}%
\pgfpathcurveto{\pgfqpoint{4.456141in}{3.257074in}}{\pgfqpoint{4.462686in}{3.241274in}}{\pgfqpoint{4.474333in}{3.229626in}}%
\pgfpathcurveto{\pgfqpoint{4.485981in}{3.217978in}}{\pgfqpoint{4.501781in}{3.211433in}}{\pgfqpoint{4.518254in}{3.211433in}}%
\pgfpathclose%
\pgfusepath{stroke,fill}%
\end{pgfscope}%
\begin{pgfscope}%
\pgfpathrectangle{\pgfqpoint{1.000000in}{0.600000in}}{\pgfqpoint{6.200000in}{4.800000in}} %
\pgfusepath{clip}%
\pgfsetbuttcap%
\pgfsetroundjoin%
\definecolor{currentfill}{rgb}{0.000000,0.500000,0.000000}%
\pgfsetfillcolor{currentfill}%
\pgfsetlinewidth{1.003750pt}%
\definecolor{currentstroke}{rgb}{0.000000,0.500000,0.000000}%
\pgfsetstrokecolor{currentstroke}%
\pgfsetdash{}{0pt}%
\pgfpathmoveto{\pgfqpoint{5.133333in}{3.743361in}}%
\pgfpathcurveto{\pgfqpoint{5.149806in}{3.743361in}}{\pgfqpoint{5.165606in}{3.749906in}}{\pgfqpoint{5.177254in}{3.761553in}}%
\pgfpathcurveto{\pgfqpoint{5.188902in}{3.773201in}}{\pgfqpoint{5.195446in}{3.789001in}}{\pgfqpoint{5.195446in}{3.805474in}}%
\pgfpathcurveto{\pgfqpoint{5.195446in}{3.821946in}}{\pgfqpoint{5.188902in}{3.837747in}}{\pgfqpoint{5.177254in}{3.849394in}}%
\pgfpathcurveto{\pgfqpoint{5.165606in}{3.861042in}}{\pgfqpoint{5.149806in}{3.867587in}}{\pgfqpoint{5.133333in}{3.867587in}}%
\pgfpathcurveto{\pgfqpoint{5.116861in}{3.867587in}}{\pgfqpoint{5.101061in}{3.861042in}}{\pgfqpoint{5.089413in}{3.849394in}}%
\pgfpathcurveto{\pgfqpoint{5.077765in}{3.837747in}}{\pgfqpoint{5.071220in}{3.821946in}}{\pgfqpoint{5.071220in}{3.805474in}}%
\pgfpathcurveto{\pgfqpoint{5.071220in}{3.789001in}}{\pgfqpoint{5.077765in}{3.773201in}}{\pgfqpoint{5.089413in}{3.761553in}}%
\pgfpathcurveto{\pgfqpoint{5.101061in}{3.749906in}}{\pgfqpoint{5.116861in}{3.743361in}}{\pgfqpoint{5.133333in}{3.743361in}}%
\pgfpathclose%
\pgfusepath{stroke,fill}%
\end{pgfscope}%
\begin{pgfscope}%
\pgfpathrectangle{\pgfqpoint{1.000000in}{0.600000in}}{\pgfqpoint{6.200000in}{4.800000in}} %
\pgfusepath{clip}%
\pgfsetbuttcap%
\pgfsetroundjoin%
\definecolor{currentfill}{rgb}{0.000000,0.500000,0.000000}%
\pgfsetfillcolor{currentfill}%
\pgfsetlinewidth{1.003750pt}%
\definecolor{currentstroke}{rgb}{0.000000,0.500000,0.000000}%
\pgfsetstrokecolor{currentstroke}%
\pgfsetdash{}{0pt}%
\pgfpathmoveto{\pgfqpoint{5.748413in}{4.157951in}}%
\pgfpathcurveto{\pgfqpoint{5.764885in}{4.157951in}}{\pgfqpoint{5.780685in}{4.164496in}}{\pgfqpoint{5.792333in}{4.176144in}}%
\pgfpathcurveto{\pgfqpoint{5.803981in}{4.187792in}}{\pgfqpoint{5.810526in}{4.203592in}}{\pgfqpoint{5.810526in}{4.220064in}}%
\pgfpathcurveto{\pgfqpoint{5.810526in}{4.236537in}}{\pgfqpoint{5.803981in}{4.252337in}}{\pgfqpoint{5.792333in}{4.263985in}}%
\pgfpathcurveto{\pgfqpoint{5.780685in}{4.275633in}}{\pgfqpoint{5.764885in}{4.282177in}}{\pgfqpoint{5.748413in}{4.282177in}}%
\pgfpathcurveto{\pgfqpoint{5.731940in}{4.282177in}}{\pgfqpoint{5.716140in}{4.275633in}}{\pgfqpoint{5.704492in}{4.263985in}}%
\pgfpathcurveto{\pgfqpoint{5.692844in}{4.252337in}}{\pgfqpoint{5.686300in}{4.236537in}}{\pgfqpoint{5.686300in}{4.220064in}}%
\pgfpathcurveto{\pgfqpoint{5.686300in}{4.203592in}}{\pgfqpoint{5.692844in}{4.187792in}}{\pgfqpoint{5.704492in}{4.176144in}}%
\pgfpathcurveto{\pgfqpoint{5.716140in}{4.164496in}}{\pgfqpoint{5.731940in}{4.157951in}}{\pgfqpoint{5.748413in}{4.157951in}}%
\pgfpathclose%
\pgfusepath{stroke,fill}%
\end{pgfscope}%
\begin{pgfscope}%
\pgfpathrectangle{\pgfqpoint{1.000000in}{0.600000in}}{\pgfqpoint{6.200000in}{4.800000in}} %
\pgfusepath{clip}%
\pgfsetbuttcap%
\pgfsetroundjoin%
\definecolor{currentfill}{rgb}{0.000000,0.500000,0.000000}%
\pgfsetfillcolor{currentfill}%
\pgfsetlinewidth{1.003750pt}%
\definecolor{currentstroke}{rgb}{0.000000,0.500000,0.000000}%
\pgfsetstrokecolor{currentstroke}%
\pgfsetdash{}{0pt}%
\pgfpathmoveto{\pgfqpoint{6.363492in}{4.580100in}}%
\pgfpathcurveto{\pgfqpoint{6.379965in}{4.580100in}}{\pgfqpoint{6.395765in}{4.586645in}}{\pgfqpoint{6.407413in}{4.598293in}}%
\pgfpathcurveto{\pgfqpoint{6.419060in}{4.609941in}}{\pgfqpoint{6.425605in}{4.625741in}}{\pgfqpoint{6.425605in}{4.642213in}}%
\pgfpathcurveto{\pgfqpoint{6.425605in}{4.658686in}}{\pgfqpoint{6.419060in}{4.674486in}}{\pgfqpoint{6.407413in}{4.686134in}}%
\pgfpathcurveto{\pgfqpoint{6.395765in}{4.697782in}}{\pgfqpoint{6.379965in}{4.704326in}}{\pgfqpoint{6.363492in}{4.704326in}}%
\pgfpathcurveto{\pgfqpoint{6.347020in}{4.704326in}}{\pgfqpoint{6.331219in}{4.697782in}}{\pgfqpoint{6.319572in}{4.686134in}}%
\pgfpathcurveto{\pgfqpoint{6.307924in}{4.674486in}}{\pgfqpoint{6.301379in}{4.658686in}}{\pgfqpoint{6.301379in}{4.642213in}}%
\pgfpathcurveto{\pgfqpoint{6.301379in}{4.625741in}}{\pgfqpoint{6.307924in}{4.609941in}}{\pgfqpoint{6.319572in}{4.598293in}}%
\pgfpathcurveto{\pgfqpoint{6.331219in}{4.586645in}}{\pgfqpoint{6.347020in}{4.580100in}}{\pgfqpoint{6.363492in}{4.580100in}}%
\pgfpathclose%
\pgfusepath{stroke,fill}%
\end{pgfscope}%
\begin{pgfscope}%
\pgfpathrectangle{\pgfqpoint{1.000000in}{0.600000in}}{\pgfqpoint{6.200000in}{4.800000in}} %
\pgfusepath{clip}%
\pgfsetbuttcap%
\pgfsetroundjoin%
\definecolor{currentfill}{rgb}{0.000000,0.500000,0.000000}%
\pgfsetfillcolor{currentfill}%
\pgfsetlinewidth{1.003750pt}%
\definecolor{currentstroke}{rgb}{0.000000,0.500000,0.000000}%
\pgfsetstrokecolor{currentstroke}%
\pgfsetdash{}{0pt}%
\pgfpathmoveto{\pgfqpoint{6.978571in}{5.035125in}}%
\pgfpathcurveto{\pgfqpoint{6.995044in}{5.035125in}}{\pgfqpoint{7.010844in}{5.041670in}}{\pgfqpoint{7.022492in}{5.053318in}}%
\pgfpathcurveto{\pgfqpoint{7.034140in}{5.064966in}}{\pgfqpoint{7.040684in}{5.080766in}}{\pgfqpoint{7.040684in}{5.097238in}}%
\pgfpathcurveto{\pgfqpoint{7.040684in}{5.113711in}}{\pgfqpoint{7.034140in}{5.129511in}}{\pgfqpoint{7.022492in}{5.141159in}}%
\pgfpathcurveto{\pgfqpoint{7.010844in}{5.152807in}}{\pgfqpoint{6.995044in}{5.159351in}}{\pgfqpoint{6.978571in}{5.159351in}}%
\pgfpathcurveto{\pgfqpoint{6.962099in}{5.159351in}}{\pgfqpoint{6.946299in}{5.152807in}}{\pgfqpoint{6.934651in}{5.141159in}}%
\pgfpathcurveto{\pgfqpoint{6.923003in}{5.129511in}}{\pgfqpoint{6.916458in}{5.113711in}}{\pgfqpoint{6.916458in}{5.097238in}}%
\pgfpathcurveto{\pgfqpoint{6.916458in}{5.080766in}}{\pgfqpoint{6.923003in}{5.064966in}}{\pgfqpoint{6.934651in}{5.053318in}}%
\pgfpathcurveto{\pgfqpoint{6.946299in}{5.041670in}}{\pgfqpoint{6.962099in}{5.035125in}}{\pgfqpoint{6.978571in}{5.035125in}}%
\pgfpathclose%
\pgfusepath{stroke,fill}%
\end{pgfscope}%
\begin{pgfscope}%
\pgfpathrectangle{\pgfqpoint{1.000000in}{0.600000in}}{\pgfqpoint{6.200000in}{4.800000in}} %
\pgfusepath{clip}%
\pgfsetrectcap%
\pgfsetroundjoin%
\pgfsetlinewidth{2.509375pt}%
\definecolor{currentstroke}{rgb}{1.000000,0.000000,0.000000}%
\pgfsetstrokecolor{currentstroke}%
\pgfsetdash{}{0pt}%
\pgfpathmoveto{\pgfqpoint{1.442857in}{0.833708in}}%
\pgfpathlineto{\pgfqpoint{1.555831in}{0.922128in}}%
\pgfpathlineto{\pgfqpoint{1.668805in}{1.010548in}}%
\pgfpathlineto{\pgfqpoint{1.781778in}{1.098968in}}%
\pgfpathlineto{\pgfqpoint{1.894752in}{1.187388in}}%
\pgfpathlineto{\pgfqpoint{2.007726in}{1.275808in}}%
\pgfpathlineto{\pgfqpoint{2.120700in}{1.364228in}}%
\pgfpathlineto{\pgfqpoint{2.233673in}{1.452648in}}%
\pgfpathlineto{\pgfqpoint{2.346647in}{1.541068in}}%
\pgfpathlineto{\pgfqpoint{2.459621in}{1.629488in}}%
\pgfpathlineto{\pgfqpoint{2.572595in}{1.717909in}}%
\pgfpathlineto{\pgfqpoint{2.685569in}{1.806329in}}%
\pgfpathlineto{\pgfqpoint{2.798542in}{1.894749in}}%
\pgfpathlineto{\pgfqpoint{2.911516in}{1.983169in}}%
\pgfpathlineto{\pgfqpoint{3.024490in}{2.071589in}}%
\pgfpathlineto{\pgfqpoint{3.137464in}{2.160009in}}%
\pgfpathlineto{\pgfqpoint{3.250437in}{2.248429in}}%
\pgfpathlineto{\pgfqpoint{3.363411in}{2.336849in}}%
\pgfpathlineto{\pgfqpoint{3.476385in}{2.425269in}}%
\pgfpathlineto{\pgfqpoint{3.589359in}{2.513689in}}%
\pgfpathlineto{\pgfqpoint{3.702332in}{2.602110in}}%
\pgfpathlineto{\pgfqpoint{3.815306in}{2.690530in}}%
\pgfpathlineto{\pgfqpoint{3.928280in}{2.778950in}}%
\pgfpathlineto{\pgfqpoint{4.041254in}{2.867370in}}%
\pgfpathlineto{\pgfqpoint{4.154227in}{2.955790in}}%
\pgfpathlineto{\pgfqpoint{4.267201in}{3.044210in}}%
\pgfpathlineto{\pgfqpoint{4.380175in}{3.132630in}}%
\pgfpathlineto{\pgfqpoint{4.493149in}{3.221050in}}%
\pgfpathlineto{\pgfqpoint{4.606122in}{3.309470in}}%
\pgfpathlineto{\pgfqpoint{4.719096in}{3.397890in}}%
\pgfpathlineto{\pgfqpoint{4.832070in}{3.486311in}}%
\pgfpathlineto{\pgfqpoint{4.945044in}{3.574731in}}%
\pgfpathlineto{\pgfqpoint{5.058017in}{3.663151in}}%
\pgfpathlineto{\pgfqpoint{5.170991in}{3.751571in}}%
\pgfpathlineto{\pgfqpoint{5.283965in}{3.839991in}}%
\pgfpathlineto{\pgfqpoint{5.396939in}{3.928411in}}%
\pgfpathlineto{\pgfqpoint{5.509913in}{4.016831in}}%
\pgfpathlineto{\pgfqpoint{5.622886in}{4.105251in}}%
\pgfpathlineto{\pgfqpoint{5.735860in}{4.193671in}}%
\pgfpathlineto{\pgfqpoint{5.848834in}{4.282091in}}%
\pgfpathlineto{\pgfqpoint{5.961808in}{4.370512in}}%
\pgfpathlineto{\pgfqpoint{6.074781in}{4.458932in}}%
\pgfpathlineto{\pgfqpoint{6.187755in}{4.547352in}}%
\pgfpathlineto{\pgfqpoint{6.300729in}{4.635772in}}%
\pgfpathlineto{\pgfqpoint{6.413703in}{4.724192in}}%
\pgfpathlineto{\pgfqpoint{6.526676in}{4.812612in}}%
\pgfpathlineto{\pgfqpoint{6.639650in}{4.901032in}}%
\pgfpathlineto{\pgfqpoint{6.752624in}{4.989452in}}%
\pgfpathlineto{\pgfqpoint{6.865598in}{5.077872in}}%
\pgfpathlineto{\pgfqpoint{6.978571in}{5.166292in}}%
\pgfusepath{stroke}%
\end{pgfscope}%
\begin{pgfscope}%
\pgfsetrectcap%
\pgfsetmiterjoin%
\pgfsetlinewidth{1.003750pt}%
\definecolor{currentstroke}{rgb}{0.000000,0.000000,0.000000}%
\pgfsetstrokecolor{currentstroke}%
\pgfsetdash{}{0pt}%
\pgfpathmoveto{\pgfqpoint{7.200000in}{0.600000in}}%
\pgfpathlineto{\pgfqpoint{7.200000in}{5.400000in}}%
\pgfusepath{stroke}%
\end{pgfscope}%
\begin{pgfscope}%
\pgfsetrectcap%
\pgfsetmiterjoin%
\pgfsetlinewidth{1.003750pt}%
\definecolor{currentstroke}{rgb}{0.000000,0.000000,0.000000}%
\pgfsetstrokecolor{currentstroke}%
\pgfsetdash{}{0pt}%
\pgfpathmoveto{\pgfqpoint{1.000000in}{0.600000in}}%
\pgfpathlineto{\pgfqpoint{7.200000in}{0.600000in}}%
\pgfusepath{stroke}%
\end{pgfscope}%
\begin{pgfscope}%
\pgfsetrectcap%
\pgfsetmiterjoin%
\pgfsetlinewidth{1.003750pt}%
\definecolor{currentstroke}{rgb}{0.000000,0.000000,0.000000}%
\pgfsetstrokecolor{currentstroke}%
\pgfsetdash{}{0pt}%
\pgfpathmoveto{\pgfqpoint{1.000000in}{0.600000in}}%
\pgfpathlineto{\pgfqpoint{1.000000in}{5.400000in}}%
\pgfusepath{stroke}%
\end{pgfscope}%
\begin{pgfscope}%
\pgfsetrectcap%
\pgfsetmiterjoin%
\pgfsetlinewidth{1.003750pt}%
\definecolor{currentstroke}{rgb}{0.000000,0.000000,0.000000}%
\pgfsetstrokecolor{currentstroke}%
\pgfsetdash{}{0pt}%
\pgfpathmoveto{\pgfqpoint{1.000000in}{5.400000in}}%
\pgfpathlineto{\pgfqpoint{7.200000in}{5.400000in}}%
\pgfusepath{stroke}%
\end{pgfscope}%
\begin{pgfscope}%
\pgfsetbuttcap%
\pgfsetroundjoin%
\definecolor{currentfill}{rgb}{0.000000,0.000000,0.000000}%
\pgfsetfillcolor{currentfill}%
\pgfsetlinewidth{0.501875pt}%
\definecolor{currentstroke}{rgb}{0.000000,0.000000,0.000000}%
\pgfsetstrokecolor{currentstroke}%
\pgfsetdash{}{0pt}%
\pgfsys@defobject{currentmarker}{\pgfqpoint{0.000000in}{0.000000in}}{\pgfqpoint{0.000000in}{0.055556in}}{%
\pgfpathmoveto{\pgfqpoint{0.000000in}{0.000000in}}%
\pgfpathlineto{\pgfqpoint{0.000000in}{0.055556in}}%
\pgfusepath{stroke,fill}%
}%
\begin{pgfscope}%
\pgfsys@transformshift{1.442857in}{0.600000in}%
\pgfsys@useobject{currentmarker}{}%
\end{pgfscope}%
\end{pgfscope}%
\begin{pgfscope}%
\pgfsetbuttcap%
\pgfsetroundjoin%
\definecolor{currentfill}{rgb}{0.000000,0.000000,0.000000}%
\pgfsetfillcolor{currentfill}%
\pgfsetlinewidth{0.501875pt}%
\definecolor{currentstroke}{rgb}{0.000000,0.000000,0.000000}%
\pgfsetstrokecolor{currentstroke}%
\pgfsetdash{}{0pt}%
\pgfsys@defobject{currentmarker}{\pgfqpoint{0.000000in}{-0.055556in}}{\pgfqpoint{0.000000in}{0.000000in}}{%
\pgfpathmoveto{\pgfqpoint{0.000000in}{0.000000in}}%
\pgfpathlineto{\pgfqpoint{0.000000in}{-0.055556in}}%
\pgfusepath{stroke,fill}%
}%
\begin{pgfscope}%
\pgfsys@transformshift{1.442857in}{5.400000in}%
\pgfsys@useobject{currentmarker}{}%
\end{pgfscope}%
\end{pgfscope}%
\begin{pgfscope}%
\pgftext[x=1.442857in,y=0.544444in,,top]{\sffamily\fontsize{12.000000}{14.400000}\selectfont 0}%
\end{pgfscope}%
\begin{pgfscope}%
\pgfsetbuttcap%
\pgfsetroundjoin%
\definecolor{currentfill}{rgb}{0.000000,0.000000,0.000000}%
\pgfsetfillcolor{currentfill}%
\pgfsetlinewidth{0.501875pt}%
\definecolor{currentstroke}{rgb}{0.000000,0.000000,0.000000}%
\pgfsetstrokecolor{currentstroke}%
\pgfsetdash{}{0pt}%
\pgfsys@defobject{currentmarker}{\pgfqpoint{0.000000in}{0.000000in}}{\pgfqpoint{0.000000in}{0.055556in}}{%
\pgfpathmoveto{\pgfqpoint{0.000000in}{0.000000in}}%
\pgfpathlineto{\pgfqpoint{0.000000in}{0.055556in}}%
\pgfusepath{stroke,fill}%
}%
\begin{pgfscope}%
\pgfsys@transformshift{2.550000in}{0.600000in}%
\pgfsys@useobject{currentmarker}{}%
\end{pgfscope}%
\end{pgfscope}%
\begin{pgfscope}%
\pgfsetbuttcap%
\pgfsetroundjoin%
\definecolor{currentfill}{rgb}{0.000000,0.000000,0.000000}%
\pgfsetfillcolor{currentfill}%
\pgfsetlinewidth{0.501875pt}%
\definecolor{currentstroke}{rgb}{0.000000,0.000000,0.000000}%
\pgfsetstrokecolor{currentstroke}%
\pgfsetdash{}{0pt}%
\pgfsys@defobject{currentmarker}{\pgfqpoint{0.000000in}{-0.055556in}}{\pgfqpoint{0.000000in}{0.000000in}}{%
\pgfpathmoveto{\pgfqpoint{0.000000in}{0.000000in}}%
\pgfpathlineto{\pgfqpoint{0.000000in}{-0.055556in}}%
\pgfusepath{stroke,fill}%
}%
\begin{pgfscope}%
\pgfsys@transformshift{2.550000in}{5.400000in}%
\pgfsys@useobject{currentmarker}{}%
\end{pgfscope}%
\end{pgfscope}%
\begin{pgfscope}%
\pgftext[x=2.550000in,y=0.544444in,,top]{\sffamily\fontsize{12.000000}{14.400000}\selectfont 1}%
\end{pgfscope}%
\begin{pgfscope}%
\pgfsetbuttcap%
\pgfsetroundjoin%
\definecolor{currentfill}{rgb}{0.000000,0.000000,0.000000}%
\pgfsetfillcolor{currentfill}%
\pgfsetlinewidth{0.501875pt}%
\definecolor{currentstroke}{rgb}{0.000000,0.000000,0.000000}%
\pgfsetstrokecolor{currentstroke}%
\pgfsetdash{}{0pt}%
\pgfsys@defobject{currentmarker}{\pgfqpoint{0.000000in}{0.000000in}}{\pgfqpoint{0.000000in}{0.055556in}}{%
\pgfpathmoveto{\pgfqpoint{0.000000in}{0.000000in}}%
\pgfpathlineto{\pgfqpoint{0.000000in}{0.055556in}}%
\pgfusepath{stroke,fill}%
}%
\begin{pgfscope}%
\pgfsys@transformshift{3.657143in}{0.600000in}%
\pgfsys@useobject{currentmarker}{}%
\end{pgfscope}%
\end{pgfscope}%
\begin{pgfscope}%
\pgfsetbuttcap%
\pgfsetroundjoin%
\definecolor{currentfill}{rgb}{0.000000,0.000000,0.000000}%
\pgfsetfillcolor{currentfill}%
\pgfsetlinewidth{0.501875pt}%
\definecolor{currentstroke}{rgb}{0.000000,0.000000,0.000000}%
\pgfsetstrokecolor{currentstroke}%
\pgfsetdash{}{0pt}%
\pgfsys@defobject{currentmarker}{\pgfqpoint{0.000000in}{-0.055556in}}{\pgfqpoint{0.000000in}{0.000000in}}{%
\pgfpathmoveto{\pgfqpoint{0.000000in}{0.000000in}}%
\pgfpathlineto{\pgfqpoint{0.000000in}{-0.055556in}}%
\pgfusepath{stroke,fill}%
}%
\begin{pgfscope}%
\pgfsys@transformshift{3.657143in}{5.400000in}%
\pgfsys@useobject{currentmarker}{}%
\end{pgfscope}%
\end{pgfscope}%
\begin{pgfscope}%
\pgftext[x=3.657143in,y=0.544444in,,top]{\sffamily\fontsize{12.000000}{14.400000}\selectfont 2}%
\end{pgfscope}%
\begin{pgfscope}%
\pgfsetbuttcap%
\pgfsetroundjoin%
\definecolor{currentfill}{rgb}{0.000000,0.000000,0.000000}%
\pgfsetfillcolor{currentfill}%
\pgfsetlinewidth{0.501875pt}%
\definecolor{currentstroke}{rgb}{0.000000,0.000000,0.000000}%
\pgfsetstrokecolor{currentstroke}%
\pgfsetdash{}{0pt}%
\pgfsys@defobject{currentmarker}{\pgfqpoint{0.000000in}{0.000000in}}{\pgfqpoint{0.000000in}{0.055556in}}{%
\pgfpathmoveto{\pgfqpoint{0.000000in}{0.000000in}}%
\pgfpathlineto{\pgfqpoint{0.000000in}{0.055556in}}%
\pgfusepath{stroke,fill}%
}%
\begin{pgfscope}%
\pgfsys@transformshift{4.764286in}{0.600000in}%
\pgfsys@useobject{currentmarker}{}%
\end{pgfscope}%
\end{pgfscope}%
\begin{pgfscope}%
\pgfsetbuttcap%
\pgfsetroundjoin%
\definecolor{currentfill}{rgb}{0.000000,0.000000,0.000000}%
\pgfsetfillcolor{currentfill}%
\pgfsetlinewidth{0.501875pt}%
\definecolor{currentstroke}{rgb}{0.000000,0.000000,0.000000}%
\pgfsetstrokecolor{currentstroke}%
\pgfsetdash{}{0pt}%
\pgfsys@defobject{currentmarker}{\pgfqpoint{0.000000in}{-0.055556in}}{\pgfqpoint{0.000000in}{0.000000in}}{%
\pgfpathmoveto{\pgfqpoint{0.000000in}{0.000000in}}%
\pgfpathlineto{\pgfqpoint{0.000000in}{-0.055556in}}%
\pgfusepath{stroke,fill}%
}%
\begin{pgfscope}%
\pgfsys@transformshift{4.764286in}{5.400000in}%
\pgfsys@useobject{currentmarker}{}%
\end{pgfscope}%
\end{pgfscope}%
\begin{pgfscope}%
\pgftext[x=4.764286in,y=0.544444in,,top]{\sffamily\fontsize{12.000000}{14.400000}\selectfont 3}%
\end{pgfscope}%
\begin{pgfscope}%
\pgfsetbuttcap%
\pgfsetroundjoin%
\definecolor{currentfill}{rgb}{0.000000,0.000000,0.000000}%
\pgfsetfillcolor{currentfill}%
\pgfsetlinewidth{0.501875pt}%
\definecolor{currentstroke}{rgb}{0.000000,0.000000,0.000000}%
\pgfsetstrokecolor{currentstroke}%
\pgfsetdash{}{0pt}%
\pgfsys@defobject{currentmarker}{\pgfqpoint{0.000000in}{0.000000in}}{\pgfqpoint{0.000000in}{0.055556in}}{%
\pgfpathmoveto{\pgfqpoint{0.000000in}{0.000000in}}%
\pgfpathlineto{\pgfqpoint{0.000000in}{0.055556in}}%
\pgfusepath{stroke,fill}%
}%
\begin{pgfscope}%
\pgfsys@transformshift{5.871429in}{0.600000in}%
\pgfsys@useobject{currentmarker}{}%
\end{pgfscope}%
\end{pgfscope}%
\begin{pgfscope}%
\pgfsetbuttcap%
\pgfsetroundjoin%
\definecolor{currentfill}{rgb}{0.000000,0.000000,0.000000}%
\pgfsetfillcolor{currentfill}%
\pgfsetlinewidth{0.501875pt}%
\definecolor{currentstroke}{rgb}{0.000000,0.000000,0.000000}%
\pgfsetstrokecolor{currentstroke}%
\pgfsetdash{}{0pt}%
\pgfsys@defobject{currentmarker}{\pgfqpoint{0.000000in}{-0.055556in}}{\pgfqpoint{0.000000in}{0.000000in}}{%
\pgfpathmoveto{\pgfqpoint{0.000000in}{0.000000in}}%
\pgfpathlineto{\pgfqpoint{0.000000in}{-0.055556in}}%
\pgfusepath{stroke,fill}%
}%
\begin{pgfscope}%
\pgfsys@transformshift{5.871429in}{5.400000in}%
\pgfsys@useobject{currentmarker}{}%
\end{pgfscope}%
\end{pgfscope}%
\begin{pgfscope}%
\pgftext[x=5.871429in,y=0.544444in,,top]{\sffamily\fontsize{12.000000}{14.400000}\selectfont 4}%
\end{pgfscope}%
\begin{pgfscope}%
\pgfsetbuttcap%
\pgfsetroundjoin%
\definecolor{currentfill}{rgb}{0.000000,0.000000,0.000000}%
\pgfsetfillcolor{currentfill}%
\pgfsetlinewidth{0.501875pt}%
\definecolor{currentstroke}{rgb}{0.000000,0.000000,0.000000}%
\pgfsetstrokecolor{currentstroke}%
\pgfsetdash{}{0pt}%
\pgfsys@defobject{currentmarker}{\pgfqpoint{0.000000in}{0.000000in}}{\pgfqpoint{0.000000in}{0.055556in}}{%
\pgfpathmoveto{\pgfqpoint{0.000000in}{0.000000in}}%
\pgfpathlineto{\pgfqpoint{0.000000in}{0.055556in}}%
\pgfusepath{stroke,fill}%
}%
\begin{pgfscope}%
\pgfsys@transformshift{6.978571in}{0.600000in}%
\pgfsys@useobject{currentmarker}{}%
\end{pgfscope}%
\end{pgfscope}%
\begin{pgfscope}%
\pgfsetbuttcap%
\pgfsetroundjoin%
\definecolor{currentfill}{rgb}{0.000000,0.000000,0.000000}%
\pgfsetfillcolor{currentfill}%
\pgfsetlinewidth{0.501875pt}%
\definecolor{currentstroke}{rgb}{0.000000,0.000000,0.000000}%
\pgfsetstrokecolor{currentstroke}%
\pgfsetdash{}{0pt}%
\pgfsys@defobject{currentmarker}{\pgfqpoint{0.000000in}{-0.055556in}}{\pgfqpoint{0.000000in}{0.000000in}}{%
\pgfpathmoveto{\pgfqpoint{0.000000in}{0.000000in}}%
\pgfpathlineto{\pgfqpoint{0.000000in}{-0.055556in}}%
\pgfusepath{stroke,fill}%
}%
\begin{pgfscope}%
\pgfsys@transformshift{6.978571in}{5.400000in}%
\pgfsys@useobject{currentmarker}{}%
\end{pgfscope}%
\end{pgfscope}%
\begin{pgfscope}%
\pgftext[x=6.978571in,y=0.544444in,,top]{\sffamily\fontsize{12.000000}{14.400000}\selectfont 5}%
\end{pgfscope}%
\begin{pgfscope}%
\pgftext[x=4.100000in,y=0.313705in,,top]{\sffamily\fontsize{12.000000}{14.400000}\selectfont x}%
\end{pgfscope}%
\begin{pgfscope}%
\pgfsetbuttcap%
\pgfsetroundjoin%
\definecolor{currentfill}{rgb}{0.000000,0.000000,0.000000}%
\pgfsetfillcolor{currentfill}%
\pgfsetlinewidth{0.501875pt}%
\definecolor{currentstroke}{rgb}{0.000000,0.000000,0.000000}%
\pgfsetstrokecolor{currentstroke}%
\pgfsetdash{}{0pt}%
\pgfsys@defobject{currentmarker}{\pgfqpoint{0.000000in}{0.000000in}}{\pgfqpoint{0.055556in}{0.000000in}}{%
\pgfpathmoveto{\pgfqpoint{0.000000in}{0.000000in}}%
\pgfpathlineto{\pgfqpoint{0.055556in}{0.000000in}}%
\pgfusepath{stroke,fill}%
}%
\begin{pgfscope}%
\pgfsys@transformshift{1.000000in}{0.660380in}%
\pgfsys@useobject{currentmarker}{}%
\end{pgfscope}%
\end{pgfscope}%
\begin{pgfscope}%
\pgfsetbuttcap%
\pgfsetroundjoin%
\definecolor{currentfill}{rgb}{0.000000,0.000000,0.000000}%
\pgfsetfillcolor{currentfill}%
\pgfsetlinewidth{0.501875pt}%
\definecolor{currentstroke}{rgb}{0.000000,0.000000,0.000000}%
\pgfsetstrokecolor{currentstroke}%
\pgfsetdash{}{0pt}%
\pgfsys@defobject{currentmarker}{\pgfqpoint{-0.055556in}{0.000000in}}{\pgfqpoint{0.000000in}{0.000000in}}{%
\pgfpathmoveto{\pgfqpoint{0.000000in}{0.000000in}}%
\pgfpathlineto{\pgfqpoint{-0.055556in}{0.000000in}}%
\pgfusepath{stroke,fill}%
}%
\begin{pgfscope}%
\pgfsys@transformshift{7.200000in}{0.660380in}%
\pgfsys@useobject{currentmarker}{}%
\end{pgfscope}%
\end{pgfscope}%
\begin{pgfscope}%
\pgftext[x=0.944444in,y=0.660380in,right,]{\sffamily\fontsize{12.000000}{14.400000}\selectfont 1}%
\end{pgfscope}%
\begin{pgfscope}%
\pgfsetbuttcap%
\pgfsetroundjoin%
\definecolor{currentfill}{rgb}{0.000000,0.000000,0.000000}%
\pgfsetfillcolor{currentfill}%
\pgfsetlinewidth{0.501875pt}%
\definecolor{currentstroke}{rgb}{0.000000,0.000000,0.000000}%
\pgfsetstrokecolor{currentstroke}%
\pgfsetdash{}{0pt}%
\pgfsys@defobject{currentmarker}{\pgfqpoint{0.000000in}{0.000000in}}{\pgfqpoint{0.055556in}{0.000000in}}{%
\pgfpathmoveto{\pgfqpoint{0.000000in}{0.000000in}}%
\pgfpathlineto{\pgfqpoint{0.055556in}{0.000000in}}%
\pgfusepath{stroke,fill}%
}%
\begin{pgfscope}%
\pgfsys@transformshift{1.000000in}{1.244650in}%
\pgfsys@useobject{currentmarker}{}%
\end{pgfscope}%
\end{pgfscope}%
\begin{pgfscope}%
\pgfsetbuttcap%
\pgfsetroundjoin%
\definecolor{currentfill}{rgb}{0.000000,0.000000,0.000000}%
\pgfsetfillcolor{currentfill}%
\pgfsetlinewidth{0.501875pt}%
\definecolor{currentstroke}{rgb}{0.000000,0.000000,0.000000}%
\pgfsetstrokecolor{currentstroke}%
\pgfsetdash{}{0pt}%
\pgfsys@defobject{currentmarker}{\pgfqpoint{-0.055556in}{0.000000in}}{\pgfqpoint{0.000000in}{0.000000in}}{%
\pgfpathmoveto{\pgfqpoint{0.000000in}{0.000000in}}%
\pgfpathlineto{\pgfqpoint{-0.055556in}{0.000000in}}%
\pgfusepath{stroke,fill}%
}%
\begin{pgfscope}%
\pgfsys@transformshift{7.200000in}{1.244650in}%
\pgfsys@useobject{currentmarker}{}%
\end{pgfscope}%
\end{pgfscope}%
\begin{pgfscope}%
\pgftext[x=0.944444in,y=1.244650in,right,]{\sffamily\fontsize{12.000000}{14.400000}\selectfont 2}%
\end{pgfscope}%
\begin{pgfscope}%
\pgfsetbuttcap%
\pgfsetroundjoin%
\definecolor{currentfill}{rgb}{0.000000,0.000000,0.000000}%
\pgfsetfillcolor{currentfill}%
\pgfsetlinewidth{0.501875pt}%
\definecolor{currentstroke}{rgb}{0.000000,0.000000,0.000000}%
\pgfsetstrokecolor{currentstroke}%
\pgfsetdash{}{0pt}%
\pgfsys@defobject{currentmarker}{\pgfqpoint{0.000000in}{0.000000in}}{\pgfqpoint{0.055556in}{0.000000in}}{%
\pgfpathmoveto{\pgfqpoint{0.000000in}{0.000000in}}%
\pgfpathlineto{\pgfqpoint{0.055556in}{0.000000in}}%
\pgfusepath{stroke,fill}%
}%
\begin{pgfscope}%
\pgfsys@transformshift{1.000000in}{1.828919in}%
\pgfsys@useobject{currentmarker}{}%
\end{pgfscope}%
\end{pgfscope}%
\begin{pgfscope}%
\pgfsetbuttcap%
\pgfsetroundjoin%
\definecolor{currentfill}{rgb}{0.000000,0.000000,0.000000}%
\pgfsetfillcolor{currentfill}%
\pgfsetlinewidth{0.501875pt}%
\definecolor{currentstroke}{rgb}{0.000000,0.000000,0.000000}%
\pgfsetstrokecolor{currentstroke}%
\pgfsetdash{}{0pt}%
\pgfsys@defobject{currentmarker}{\pgfqpoint{-0.055556in}{0.000000in}}{\pgfqpoint{0.000000in}{0.000000in}}{%
\pgfpathmoveto{\pgfqpoint{0.000000in}{0.000000in}}%
\pgfpathlineto{\pgfqpoint{-0.055556in}{0.000000in}}%
\pgfusepath{stroke,fill}%
}%
\begin{pgfscope}%
\pgfsys@transformshift{7.200000in}{1.828919in}%
\pgfsys@useobject{currentmarker}{}%
\end{pgfscope}%
\end{pgfscope}%
\begin{pgfscope}%
\pgftext[x=0.944444in,y=1.828919in,right,]{\sffamily\fontsize{12.000000}{14.400000}\selectfont 3}%
\end{pgfscope}%
\begin{pgfscope}%
\pgfsetbuttcap%
\pgfsetroundjoin%
\definecolor{currentfill}{rgb}{0.000000,0.000000,0.000000}%
\pgfsetfillcolor{currentfill}%
\pgfsetlinewidth{0.501875pt}%
\definecolor{currentstroke}{rgb}{0.000000,0.000000,0.000000}%
\pgfsetstrokecolor{currentstroke}%
\pgfsetdash{}{0pt}%
\pgfsys@defobject{currentmarker}{\pgfqpoint{0.000000in}{0.000000in}}{\pgfqpoint{0.055556in}{0.000000in}}{%
\pgfpathmoveto{\pgfqpoint{0.000000in}{0.000000in}}%
\pgfpathlineto{\pgfqpoint{0.055556in}{0.000000in}}%
\pgfusepath{stroke,fill}%
}%
\begin{pgfscope}%
\pgfsys@transformshift{1.000000in}{2.413188in}%
\pgfsys@useobject{currentmarker}{}%
\end{pgfscope}%
\end{pgfscope}%
\begin{pgfscope}%
\pgfsetbuttcap%
\pgfsetroundjoin%
\definecolor{currentfill}{rgb}{0.000000,0.000000,0.000000}%
\pgfsetfillcolor{currentfill}%
\pgfsetlinewidth{0.501875pt}%
\definecolor{currentstroke}{rgb}{0.000000,0.000000,0.000000}%
\pgfsetstrokecolor{currentstroke}%
\pgfsetdash{}{0pt}%
\pgfsys@defobject{currentmarker}{\pgfqpoint{-0.055556in}{0.000000in}}{\pgfqpoint{0.000000in}{0.000000in}}{%
\pgfpathmoveto{\pgfqpoint{0.000000in}{0.000000in}}%
\pgfpathlineto{\pgfqpoint{-0.055556in}{0.000000in}}%
\pgfusepath{stroke,fill}%
}%
\begin{pgfscope}%
\pgfsys@transformshift{7.200000in}{2.413188in}%
\pgfsys@useobject{currentmarker}{}%
\end{pgfscope}%
\end{pgfscope}%
\begin{pgfscope}%
\pgftext[x=0.944444in,y=2.413188in,right,]{\sffamily\fontsize{12.000000}{14.400000}\selectfont 4}%
\end{pgfscope}%
\begin{pgfscope}%
\pgfsetbuttcap%
\pgfsetroundjoin%
\definecolor{currentfill}{rgb}{0.000000,0.000000,0.000000}%
\pgfsetfillcolor{currentfill}%
\pgfsetlinewidth{0.501875pt}%
\definecolor{currentstroke}{rgb}{0.000000,0.000000,0.000000}%
\pgfsetstrokecolor{currentstroke}%
\pgfsetdash{}{0pt}%
\pgfsys@defobject{currentmarker}{\pgfqpoint{0.000000in}{0.000000in}}{\pgfqpoint{0.055556in}{0.000000in}}{%
\pgfpathmoveto{\pgfqpoint{0.000000in}{0.000000in}}%
\pgfpathlineto{\pgfqpoint{0.055556in}{0.000000in}}%
\pgfusepath{stroke,fill}%
}%
\begin{pgfscope}%
\pgfsys@transformshift{1.000000in}{2.997457in}%
\pgfsys@useobject{currentmarker}{}%
\end{pgfscope}%
\end{pgfscope}%
\begin{pgfscope}%
\pgfsetbuttcap%
\pgfsetroundjoin%
\definecolor{currentfill}{rgb}{0.000000,0.000000,0.000000}%
\pgfsetfillcolor{currentfill}%
\pgfsetlinewidth{0.501875pt}%
\definecolor{currentstroke}{rgb}{0.000000,0.000000,0.000000}%
\pgfsetstrokecolor{currentstroke}%
\pgfsetdash{}{0pt}%
\pgfsys@defobject{currentmarker}{\pgfqpoint{-0.055556in}{0.000000in}}{\pgfqpoint{0.000000in}{0.000000in}}{%
\pgfpathmoveto{\pgfqpoint{0.000000in}{0.000000in}}%
\pgfpathlineto{\pgfqpoint{-0.055556in}{0.000000in}}%
\pgfusepath{stroke,fill}%
}%
\begin{pgfscope}%
\pgfsys@transformshift{7.200000in}{2.997457in}%
\pgfsys@useobject{currentmarker}{}%
\end{pgfscope}%
\end{pgfscope}%
\begin{pgfscope}%
\pgftext[x=0.944444in,y=2.997457in,right,]{\sffamily\fontsize{12.000000}{14.400000}\selectfont 5}%
\end{pgfscope}%
\begin{pgfscope}%
\pgfsetbuttcap%
\pgfsetroundjoin%
\definecolor{currentfill}{rgb}{0.000000,0.000000,0.000000}%
\pgfsetfillcolor{currentfill}%
\pgfsetlinewidth{0.501875pt}%
\definecolor{currentstroke}{rgb}{0.000000,0.000000,0.000000}%
\pgfsetstrokecolor{currentstroke}%
\pgfsetdash{}{0pt}%
\pgfsys@defobject{currentmarker}{\pgfqpoint{0.000000in}{0.000000in}}{\pgfqpoint{0.055556in}{0.000000in}}{%
\pgfpathmoveto{\pgfqpoint{0.000000in}{0.000000in}}%
\pgfpathlineto{\pgfqpoint{0.055556in}{0.000000in}}%
\pgfusepath{stroke,fill}%
}%
\begin{pgfscope}%
\pgfsys@transformshift{1.000000in}{3.581726in}%
\pgfsys@useobject{currentmarker}{}%
\end{pgfscope}%
\end{pgfscope}%
\begin{pgfscope}%
\pgfsetbuttcap%
\pgfsetroundjoin%
\definecolor{currentfill}{rgb}{0.000000,0.000000,0.000000}%
\pgfsetfillcolor{currentfill}%
\pgfsetlinewidth{0.501875pt}%
\definecolor{currentstroke}{rgb}{0.000000,0.000000,0.000000}%
\pgfsetstrokecolor{currentstroke}%
\pgfsetdash{}{0pt}%
\pgfsys@defobject{currentmarker}{\pgfqpoint{-0.055556in}{0.000000in}}{\pgfqpoint{0.000000in}{0.000000in}}{%
\pgfpathmoveto{\pgfqpoint{0.000000in}{0.000000in}}%
\pgfpathlineto{\pgfqpoint{-0.055556in}{0.000000in}}%
\pgfusepath{stroke,fill}%
}%
\begin{pgfscope}%
\pgfsys@transformshift{7.200000in}{3.581726in}%
\pgfsys@useobject{currentmarker}{}%
\end{pgfscope}%
\end{pgfscope}%
\begin{pgfscope}%
\pgftext[x=0.944444in,y=3.581726in,right,]{\sffamily\fontsize{12.000000}{14.400000}\selectfont 6}%
\end{pgfscope}%
\begin{pgfscope}%
\pgfsetbuttcap%
\pgfsetroundjoin%
\definecolor{currentfill}{rgb}{0.000000,0.000000,0.000000}%
\pgfsetfillcolor{currentfill}%
\pgfsetlinewidth{0.501875pt}%
\definecolor{currentstroke}{rgb}{0.000000,0.000000,0.000000}%
\pgfsetstrokecolor{currentstroke}%
\pgfsetdash{}{0pt}%
\pgfsys@defobject{currentmarker}{\pgfqpoint{0.000000in}{0.000000in}}{\pgfqpoint{0.055556in}{0.000000in}}{%
\pgfpathmoveto{\pgfqpoint{0.000000in}{0.000000in}}%
\pgfpathlineto{\pgfqpoint{0.055556in}{0.000000in}}%
\pgfusepath{stroke,fill}%
}%
\begin{pgfscope}%
\pgfsys@transformshift{1.000000in}{4.165995in}%
\pgfsys@useobject{currentmarker}{}%
\end{pgfscope}%
\end{pgfscope}%
\begin{pgfscope}%
\pgfsetbuttcap%
\pgfsetroundjoin%
\definecolor{currentfill}{rgb}{0.000000,0.000000,0.000000}%
\pgfsetfillcolor{currentfill}%
\pgfsetlinewidth{0.501875pt}%
\definecolor{currentstroke}{rgb}{0.000000,0.000000,0.000000}%
\pgfsetstrokecolor{currentstroke}%
\pgfsetdash{}{0pt}%
\pgfsys@defobject{currentmarker}{\pgfqpoint{-0.055556in}{0.000000in}}{\pgfqpoint{0.000000in}{0.000000in}}{%
\pgfpathmoveto{\pgfqpoint{0.000000in}{0.000000in}}%
\pgfpathlineto{\pgfqpoint{-0.055556in}{0.000000in}}%
\pgfusepath{stroke,fill}%
}%
\begin{pgfscope}%
\pgfsys@transformshift{7.200000in}{4.165995in}%
\pgfsys@useobject{currentmarker}{}%
\end{pgfscope}%
\end{pgfscope}%
\begin{pgfscope}%
\pgftext[x=0.944444in,y=4.165995in,right,]{\sffamily\fontsize{12.000000}{14.400000}\selectfont 7}%
\end{pgfscope}%
\begin{pgfscope}%
\pgfsetbuttcap%
\pgfsetroundjoin%
\definecolor{currentfill}{rgb}{0.000000,0.000000,0.000000}%
\pgfsetfillcolor{currentfill}%
\pgfsetlinewidth{0.501875pt}%
\definecolor{currentstroke}{rgb}{0.000000,0.000000,0.000000}%
\pgfsetstrokecolor{currentstroke}%
\pgfsetdash{}{0pt}%
\pgfsys@defobject{currentmarker}{\pgfqpoint{0.000000in}{0.000000in}}{\pgfqpoint{0.055556in}{0.000000in}}{%
\pgfpathmoveto{\pgfqpoint{0.000000in}{0.000000in}}%
\pgfpathlineto{\pgfqpoint{0.055556in}{0.000000in}}%
\pgfusepath{stroke,fill}%
}%
\begin{pgfscope}%
\pgfsys@transformshift{1.000000in}{4.750264in}%
\pgfsys@useobject{currentmarker}{}%
\end{pgfscope}%
\end{pgfscope}%
\begin{pgfscope}%
\pgfsetbuttcap%
\pgfsetroundjoin%
\definecolor{currentfill}{rgb}{0.000000,0.000000,0.000000}%
\pgfsetfillcolor{currentfill}%
\pgfsetlinewidth{0.501875pt}%
\definecolor{currentstroke}{rgb}{0.000000,0.000000,0.000000}%
\pgfsetstrokecolor{currentstroke}%
\pgfsetdash{}{0pt}%
\pgfsys@defobject{currentmarker}{\pgfqpoint{-0.055556in}{0.000000in}}{\pgfqpoint{0.000000in}{0.000000in}}{%
\pgfpathmoveto{\pgfqpoint{0.000000in}{0.000000in}}%
\pgfpathlineto{\pgfqpoint{-0.055556in}{0.000000in}}%
\pgfusepath{stroke,fill}%
}%
\begin{pgfscope}%
\pgfsys@transformshift{7.200000in}{4.750264in}%
\pgfsys@useobject{currentmarker}{}%
\end{pgfscope}%
\end{pgfscope}%
\begin{pgfscope}%
\pgftext[x=0.944444in,y=4.750264in,right,]{\sffamily\fontsize{12.000000}{14.400000}\selectfont 8}%
\end{pgfscope}%
\begin{pgfscope}%
\pgfsetbuttcap%
\pgfsetroundjoin%
\definecolor{currentfill}{rgb}{0.000000,0.000000,0.000000}%
\pgfsetfillcolor{currentfill}%
\pgfsetlinewidth{0.501875pt}%
\definecolor{currentstroke}{rgb}{0.000000,0.000000,0.000000}%
\pgfsetstrokecolor{currentstroke}%
\pgfsetdash{}{0pt}%
\pgfsys@defobject{currentmarker}{\pgfqpoint{0.000000in}{0.000000in}}{\pgfqpoint{0.055556in}{0.000000in}}{%
\pgfpathmoveto{\pgfqpoint{0.000000in}{0.000000in}}%
\pgfpathlineto{\pgfqpoint{0.055556in}{0.000000in}}%
\pgfusepath{stroke,fill}%
}%
\begin{pgfscope}%
\pgfsys@transformshift{1.000000in}{5.334533in}%
\pgfsys@useobject{currentmarker}{}%
\end{pgfscope}%
\end{pgfscope}%
\begin{pgfscope}%
\pgfsetbuttcap%
\pgfsetroundjoin%
\definecolor{currentfill}{rgb}{0.000000,0.000000,0.000000}%
\pgfsetfillcolor{currentfill}%
\pgfsetlinewidth{0.501875pt}%
\definecolor{currentstroke}{rgb}{0.000000,0.000000,0.000000}%
\pgfsetstrokecolor{currentstroke}%
\pgfsetdash{}{0pt}%
\pgfsys@defobject{currentmarker}{\pgfqpoint{-0.055556in}{0.000000in}}{\pgfqpoint{0.000000in}{0.000000in}}{%
\pgfpathmoveto{\pgfqpoint{0.000000in}{0.000000in}}%
\pgfpathlineto{\pgfqpoint{-0.055556in}{0.000000in}}%
\pgfusepath{stroke,fill}%
}%
\begin{pgfscope}%
\pgfsys@transformshift{7.200000in}{5.334533in}%
\pgfsys@useobject{currentmarker}{}%
\end{pgfscope}%
\end{pgfscope}%
\begin{pgfscope}%
\pgftext[x=0.944444in,y=5.334533in,right,]{\sffamily\fontsize{12.000000}{14.400000}\selectfont 9}%
\end{pgfscope}%
\begin{pgfscope}%
\pgftext[x=0.768962in,y=3.000000in,,bottom]{\sffamily\fontsize{12.000000}{14.400000}\selectfont y}%
\end{pgfscope}%
\end{pgfpicture}%
\makeatother%
\endgroup%
}
	\resizebox{.45\linewidth}{!}{%% Creator: Matplotlib, PGF backend
%%
%% To include the figure in your LaTeX document, write
%%   \input{<filename>.pgf}
%%
%% Make sure the required packages are loaded in your preamble
%%   \usepackage{pgf}
%%
%% Figures using additional raster images can only be included by \input if
%% they are in the same directory as the main LaTeX file. For loading figures
%% from other directories you can use the `import` package
%%   \usepackage{import}
%% and then include the figures with
%%   \import{<path to file>}{<filename>.pgf}
%%
%% Matplotlib used the following preamble
%%   \usepackage{fontspec}
%%   \setmainfont{Bitstream Vera Serif}
%%   \setsansfont{Bitstream Vera Sans}
%%   \setmonofont{Bitstream Vera Sans Mono}
%%
\begingroup%
\makeatletter%
\begin{pgfpicture}%
\pgfpathrectangle{\pgfpointorigin}{\pgfqpoint{8.000000in}{6.000000in}}%
\pgfusepath{use as bounding box, clip}%
\begin{pgfscope}%
\pgfsetbuttcap%
\pgfsetmiterjoin%
\definecolor{currentfill}{rgb}{1.000000,1.000000,1.000000}%
\pgfsetfillcolor{currentfill}%
\pgfsetlinewidth{0.000000pt}%
\definecolor{currentstroke}{rgb}{1.000000,1.000000,1.000000}%
\pgfsetstrokecolor{currentstroke}%
\pgfsetdash{}{0pt}%
\pgfpathmoveto{\pgfqpoint{0.000000in}{0.000000in}}%
\pgfpathlineto{\pgfqpoint{8.000000in}{0.000000in}}%
\pgfpathlineto{\pgfqpoint{8.000000in}{6.000000in}}%
\pgfpathlineto{\pgfqpoint{0.000000in}{6.000000in}}%
\pgfpathclose%
\pgfusepath{fill}%
\end{pgfscope}%
\begin{pgfscope}%
\pgfsetbuttcap%
\pgfsetmiterjoin%
\definecolor{currentfill}{rgb}{1.000000,1.000000,1.000000}%
\pgfsetfillcolor{currentfill}%
\pgfsetlinewidth{0.000000pt}%
\definecolor{currentstroke}{rgb}{0.000000,0.000000,0.000000}%
\pgfsetstrokecolor{currentstroke}%
\pgfsetstrokeopacity{0.000000}%
\pgfsetdash{}{0pt}%
\pgfpathmoveto{\pgfqpoint{1.000000in}{0.600000in}}%
\pgfpathlineto{\pgfqpoint{7.200000in}{0.600000in}}%
\pgfpathlineto{\pgfqpoint{7.200000in}{5.400000in}}%
\pgfpathlineto{\pgfqpoint{1.000000in}{5.400000in}}%
\pgfpathclose%
\pgfusepath{fill}%
\end{pgfscope}%
\begin{pgfscope}%
\pgfpathrectangle{\pgfqpoint{1.000000in}{0.600000in}}{\pgfqpoint{6.200000in}{4.800000in}} %
\pgfusepath{clip}%
\pgfsetbuttcap%
\pgfsetroundjoin%
\definecolor{currentfill}{rgb}{0.000000,0.500000,0.000000}%
\pgfsetfillcolor{currentfill}%
\pgfsetlinewidth{1.003750pt}%
\definecolor{currentstroke}{rgb}{0.000000,0.500000,0.000000}%
\pgfsetstrokecolor{currentstroke}%
\pgfsetdash{}{0pt}%
\pgfpathmoveto{\pgfqpoint{1.442857in}{0.663119in}}%
\pgfpathcurveto{\pgfqpoint{1.459330in}{0.663119in}}{\pgfqpoint{1.475130in}{0.669664in}}{\pgfqpoint{1.486778in}{0.681312in}}%
\pgfpathcurveto{\pgfqpoint{1.498426in}{0.692960in}}{\pgfqpoint{1.504970in}{0.708760in}}{\pgfqpoint{1.504970in}{0.725232in}}%
\pgfpathcurveto{\pgfqpoint{1.504970in}{0.741705in}}{\pgfqpoint{1.498426in}{0.757505in}}{\pgfqpoint{1.486778in}{0.769153in}}%
\pgfpathcurveto{\pgfqpoint{1.475130in}{0.780801in}}{\pgfqpoint{1.459330in}{0.787345in}}{\pgfqpoint{1.442857in}{0.787345in}}%
\pgfpathcurveto{\pgfqpoint{1.426385in}{0.787345in}}{\pgfqpoint{1.410584in}{0.780801in}}{\pgfqpoint{1.398937in}{0.769153in}}%
\pgfpathcurveto{\pgfqpoint{1.387289in}{0.757505in}}{\pgfqpoint{1.380744in}{0.741705in}}{\pgfqpoint{1.380744in}{0.725232in}}%
\pgfpathcurveto{\pgfqpoint{1.380744in}{0.708760in}}{\pgfqpoint{1.387289in}{0.692960in}}{\pgfqpoint{1.398937in}{0.681312in}}%
\pgfpathcurveto{\pgfqpoint{1.410584in}{0.669664in}}{\pgfqpoint{1.426385in}{0.663119in}}{\pgfqpoint{1.442857in}{0.663119in}}%
\pgfpathclose%
\pgfusepath{stroke,fill}%
\end{pgfscope}%
\begin{pgfscope}%
\pgfpathrectangle{\pgfqpoint{1.000000in}{0.600000in}}{\pgfqpoint{6.200000in}{4.800000in}} %
\pgfusepath{clip}%
\pgfsetbuttcap%
\pgfsetroundjoin%
\definecolor{currentfill}{rgb}{0.000000,0.500000,0.000000}%
\pgfsetfillcolor{currentfill}%
\pgfsetlinewidth{1.003750pt}%
\definecolor{currentstroke}{rgb}{0.000000,0.500000,0.000000}%
\pgfsetstrokecolor{currentstroke}%
\pgfsetdash{}{0pt}%
\pgfpathmoveto{\pgfqpoint{2.057937in}{1.339529in}}%
\pgfpathcurveto{\pgfqpoint{2.074409in}{1.339529in}}{\pgfqpoint{2.090209in}{1.346074in}}{\pgfqpoint{2.101857in}{1.357722in}}%
\pgfpathcurveto{\pgfqpoint{2.113505in}{1.369370in}}{\pgfqpoint{2.120050in}{1.385170in}}{\pgfqpoint{2.120050in}{1.401642in}}%
\pgfpathcurveto{\pgfqpoint{2.120050in}{1.418115in}}{\pgfqpoint{2.113505in}{1.433915in}}{\pgfqpoint{2.101857in}{1.445563in}}%
\pgfpathcurveto{\pgfqpoint{2.090209in}{1.457211in}}{\pgfqpoint{2.074409in}{1.463755in}}{\pgfqpoint{2.057937in}{1.463755in}}%
\pgfpathcurveto{\pgfqpoint{2.041464in}{1.463755in}}{\pgfqpoint{2.025664in}{1.457211in}}{\pgfqpoint{2.014016in}{1.445563in}}%
\pgfpathcurveto{\pgfqpoint{2.002368in}{1.433915in}}{\pgfqpoint{1.995824in}{1.418115in}}{\pgfqpoint{1.995824in}{1.401642in}}%
\pgfpathcurveto{\pgfqpoint{1.995824in}{1.385170in}}{\pgfqpoint{2.002368in}{1.369370in}}{\pgfqpoint{2.014016in}{1.357722in}}%
\pgfpathcurveto{\pgfqpoint{2.025664in}{1.346074in}}{\pgfqpoint{2.041464in}{1.339529in}}{\pgfqpoint{2.057937in}{1.339529in}}%
\pgfpathclose%
\pgfusepath{stroke,fill}%
\end{pgfscope}%
\begin{pgfscope}%
\pgfpathrectangle{\pgfqpoint{1.000000in}{0.600000in}}{\pgfqpoint{6.200000in}{4.800000in}} %
\pgfusepath{clip}%
\pgfsetbuttcap%
\pgfsetroundjoin%
\definecolor{currentfill}{rgb}{0.000000,0.500000,0.000000}%
\pgfsetfillcolor{currentfill}%
\pgfsetlinewidth{1.003750pt}%
\definecolor{currentstroke}{rgb}{0.000000,0.500000,0.000000}%
\pgfsetstrokecolor{currentstroke}%
\pgfsetdash{}{0pt}%
\pgfpathmoveto{\pgfqpoint{2.673016in}{1.632439in}}%
\pgfpathcurveto{\pgfqpoint{2.689488in}{1.632439in}}{\pgfqpoint{2.705289in}{1.638984in}}{\pgfqpoint{2.716936in}{1.650632in}}%
\pgfpathcurveto{\pgfqpoint{2.728584in}{1.662279in}}{\pgfqpoint{2.735129in}{1.678080in}}{\pgfqpoint{2.735129in}{1.694552in}}%
\pgfpathcurveto{\pgfqpoint{2.735129in}{1.711025in}}{\pgfqpoint{2.728584in}{1.726825in}}{\pgfqpoint{2.716936in}{1.738473in}}%
\pgfpathcurveto{\pgfqpoint{2.705289in}{1.750120in}}{\pgfqpoint{2.689488in}{1.756665in}}{\pgfqpoint{2.673016in}{1.756665in}}%
\pgfpathcurveto{\pgfqpoint{2.656543in}{1.756665in}}{\pgfqpoint{2.640743in}{1.750120in}}{\pgfqpoint{2.629095in}{1.738473in}}%
\pgfpathcurveto{\pgfqpoint{2.617447in}{1.726825in}}{\pgfqpoint{2.610903in}{1.711025in}}{\pgfqpoint{2.610903in}{1.694552in}}%
\pgfpathcurveto{\pgfqpoint{2.610903in}{1.678080in}}{\pgfqpoint{2.617447in}{1.662279in}}{\pgfqpoint{2.629095in}{1.650632in}}%
\pgfpathcurveto{\pgfqpoint{2.640743in}{1.638984in}}{\pgfqpoint{2.656543in}{1.632439in}}{\pgfqpoint{2.673016in}{1.632439in}}%
\pgfpathclose%
\pgfusepath{stroke,fill}%
\end{pgfscope}%
\begin{pgfscope}%
\pgfpathrectangle{\pgfqpoint{1.000000in}{0.600000in}}{\pgfqpoint{6.200000in}{4.800000in}} %
\pgfusepath{clip}%
\pgfsetbuttcap%
\pgfsetroundjoin%
\definecolor{currentfill}{rgb}{0.000000,0.500000,0.000000}%
\pgfsetfillcolor{currentfill}%
\pgfsetlinewidth{1.003750pt}%
\definecolor{currentstroke}{rgb}{0.000000,0.500000,0.000000}%
\pgfsetstrokecolor{currentstroke}%
\pgfsetdash{}{0pt}%
\pgfpathmoveto{\pgfqpoint{3.288095in}{2.327298in}}%
\pgfpathcurveto{\pgfqpoint{3.304568in}{2.327298in}}{\pgfqpoint{3.320368in}{2.333843in}}{\pgfqpoint{3.332016in}{2.345491in}}%
\pgfpathcurveto{\pgfqpoint{3.343664in}{2.357138in}}{\pgfqpoint{3.350208in}{2.372939in}}{\pgfqpoint{3.350208in}{2.389411in}}%
\pgfpathcurveto{\pgfqpoint{3.350208in}{2.405884in}}{\pgfqpoint{3.343664in}{2.421684in}}{\pgfqpoint{3.332016in}{2.433332in}}%
\pgfpathcurveto{\pgfqpoint{3.320368in}{2.444979in}}{\pgfqpoint{3.304568in}{2.451524in}}{\pgfqpoint{3.288095in}{2.451524in}}%
\pgfpathcurveto{\pgfqpoint{3.271623in}{2.451524in}}{\pgfqpoint{3.255823in}{2.444979in}}{\pgfqpoint{3.244175in}{2.433332in}}%
\pgfpathcurveto{\pgfqpoint{3.232527in}{2.421684in}}{\pgfqpoint{3.225982in}{2.405884in}}{\pgfqpoint{3.225982in}{2.389411in}}%
\pgfpathcurveto{\pgfqpoint{3.225982in}{2.372939in}}{\pgfqpoint{3.232527in}{2.357138in}}{\pgfqpoint{3.244175in}{2.345491in}}%
\pgfpathcurveto{\pgfqpoint{3.255823in}{2.333843in}}{\pgfqpoint{3.271623in}{2.327298in}}{\pgfqpoint{3.288095in}{2.327298in}}%
\pgfpathclose%
\pgfusepath{stroke,fill}%
\end{pgfscope}%
\begin{pgfscope}%
\pgfpathrectangle{\pgfqpoint{1.000000in}{0.600000in}}{\pgfqpoint{6.200000in}{4.800000in}} %
\pgfusepath{clip}%
\pgfsetbuttcap%
\pgfsetroundjoin%
\definecolor{currentfill}{rgb}{0.000000,0.500000,0.000000}%
\pgfsetfillcolor{currentfill}%
\pgfsetlinewidth{1.003750pt}%
\definecolor{currentstroke}{rgb}{0.000000,0.500000,0.000000}%
\pgfsetstrokecolor{currentstroke}%
\pgfsetdash{}{0pt}%
\pgfpathmoveto{\pgfqpoint{3.903175in}{2.688513in}}%
\pgfpathcurveto{\pgfqpoint{3.919647in}{2.688513in}}{\pgfqpoint{3.935447in}{2.695057in}}{\pgfqpoint{3.947095in}{2.706705in}}%
\pgfpathcurveto{\pgfqpoint{3.958743in}{2.718353in}}{\pgfqpoint{3.965288in}{2.734153in}}{\pgfqpoint{3.965288in}{2.750626in}}%
\pgfpathcurveto{\pgfqpoint{3.965288in}{2.767098in}}{\pgfqpoint{3.958743in}{2.782898in}}{\pgfqpoint{3.947095in}{2.794546in}}%
\pgfpathcurveto{\pgfqpoint{3.935447in}{2.806194in}}{\pgfqpoint{3.919647in}{2.812739in}}{\pgfqpoint{3.903175in}{2.812739in}}%
\pgfpathcurveto{\pgfqpoint{3.886702in}{2.812739in}}{\pgfqpoint{3.870902in}{2.806194in}}{\pgfqpoint{3.859254in}{2.794546in}}%
\pgfpathcurveto{\pgfqpoint{3.847606in}{2.782898in}}{\pgfqpoint{3.841062in}{2.767098in}}{\pgfqpoint{3.841062in}{2.750626in}}%
\pgfpathcurveto{\pgfqpoint{3.841062in}{2.734153in}}{\pgfqpoint{3.847606in}{2.718353in}}{\pgfqpoint{3.859254in}{2.706705in}}%
\pgfpathcurveto{\pgfqpoint{3.870902in}{2.695057in}}{\pgfqpoint{3.886702in}{2.688513in}}{\pgfqpoint{3.903175in}{2.688513in}}%
\pgfpathclose%
\pgfusepath{stroke,fill}%
\end{pgfscope}%
\begin{pgfscope}%
\pgfpathrectangle{\pgfqpoint{1.000000in}{0.600000in}}{\pgfqpoint{6.200000in}{4.800000in}} %
\pgfusepath{clip}%
\pgfsetbuttcap%
\pgfsetroundjoin%
\definecolor{currentfill}{rgb}{0.000000,0.500000,0.000000}%
\pgfsetfillcolor{currentfill}%
\pgfsetlinewidth{1.003750pt}%
\definecolor{currentstroke}{rgb}{0.000000,0.500000,0.000000}%
\pgfsetstrokecolor{currentstroke}%
\pgfsetdash{}{0pt}%
\pgfpathmoveto{\pgfqpoint{4.518254in}{3.211433in}}%
\pgfpathcurveto{\pgfqpoint{4.534727in}{3.211433in}}{\pgfqpoint{4.550527in}{3.217978in}}{\pgfqpoint{4.562174in}{3.229626in}}%
\pgfpathcurveto{\pgfqpoint{4.573822in}{3.241274in}}{\pgfqpoint{4.580367in}{3.257074in}}{\pgfqpoint{4.580367in}{3.273546in}}%
\pgfpathcurveto{\pgfqpoint{4.580367in}{3.290019in}}{\pgfqpoint{4.573822in}{3.305819in}}{\pgfqpoint{4.562174in}{3.317467in}}%
\pgfpathcurveto{\pgfqpoint{4.550527in}{3.329115in}}{\pgfqpoint{4.534727in}{3.335659in}}{\pgfqpoint{4.518254in}{3.335659in}}%
\pgfpathcurveto{\pgfqpoint{4.501781in}{3.335659in}}{\pgfqpoint{4.485981in}{3.329115in}}{\pgfqpoint{4.474333in}{3.317467in}}%
\pgfpathcurveto{\pgfqpoint{4.462686in}{3.305819in}}{\pgfqpoint{4.456141in}{3.290019in}}{\pgfqpoint{4.456141in}{3.273546in}}%
\pgfpathcurveto{\pgfqpoint{4.456141in}{3.257074in}}{\pgfqpoint{4.462686in}{3.241274in}}{\pgfqpoint{4.474333in}{3.229626in}}%
\pgfpathcurveto{\pgfqpoint{4.485981in}{3.217978in}}{\pgfqpoint{4.501781in}{3.211433in}}{\pgfqpoint{4.518254in}{3.211433in}}%
\pgfpathclose%
\pgfusepath{stroke,fill}%
\end{pgfscope}%
\begin{pgfscope}%
\pgfpathrectangle{\pgfqpoint{1.000000in}{0.600000in}}{\pgfqpoint{6.200000in}{4.800000in}} %
\pgfusepath{clip}%
\pgfsetbuttcap%
\pgfsetroundjoin%
\definecolor{currentfill}{rgb}{0.000000,0.500000,0.000000}%
\pgfsetfillcolor{currentfill}%
\pgfsetlinewidth{1.003750pt}%
\definecolor{currentstroke}{rgb}{0.000000,0.500000,0.000000}%
\pgfsetstrokecolor{currentstroke}%
\pgfsetdash{}{0pt}%
\pgfpathmoveto{\pgfqpoint{5.133333in}{3.743361in}}%
\pgfpathcurveto{\pgfqpoint{5.149806in}{3.743361in}}{\pgfqpoint{5.165606in}{3.749906in}}{\pgfqpoint{5.177254in}{3.761553in}}%
\pgfpathcurveto{\pgfqpoint{5.188902in}{3.773201in}}{\pgfqpoint{5.195446in}{3.789001in}}{\pgfqpoint{5.195446in}{3.805474in}}%
\pgfpathcurveto{\pgfqpoint{5.195446in}{3.821946in}}{\pgfqpoint{5.188902in}{3.837747in}}{\pgfqpoint{5.177254in}{3.849394in}}%
\pgfpathcurveto{\pgfqpoint{5.165606in}{3.861042in}}{\pgfqpoint{5.149806in}{3.867587in}}{\pgfqpoint{5.133333in}{3.867587in}}%
\pgfpathcurveto{\pgfqpoint{5.116861in}{3.867587in}}{\pgfqpoint{5.101061in}{3.861042in}}{\pgfqpoint{5.089413in}{3.849394in}}%
\pgfpathcurveto{\pgfqpoint{5.077765in}{3.837747in}}{\pgfqpoint{5.071220in}{3.821946in}}{\pgfqpoint{5.071220in}{3.805474in}}%
\pgfpathcurveto{\pgfqpoint{5.071220in}{3.789001in}}{\pgfqpoint{5.077765in}{3.773201in}}{\pgfqpoint{5.089413in}{3.761553in}}%
\pgfpathcurveto{\pgfqpoint{5.101061in}{3.749906in}}{\pgfqpoint{5.116861in}{3.743361in}}{\pgfqpoint{5.133333in}{3.743361in}}%
\pgfpathclose%
\pgfusepath{stroke,fill}%
\end{pgfscope}%
\begin{pgfscope}%
\pgfpathrectangle{\pgfqpoint{1.000000in}{0.600000in}}{\pgfqpoint{6.200000in}{4.800000in}} %
\pgfusepath{clip}%
\pgfsetbuttcap%
\pgfsetroundjoin%
\definecolor{currentfill}{rgb}{0.000000,0.500000,0.000000}%
\pgfsetfillcolor{currentfill}%
\pgfsetlinewidth{1.003750pt}%
\definecolor{currentstroke}{rgb}{0.000000,0.500000,0.000000}%
\pgfsetstrokecolor{currentstroke}%
\pgfsetdash{}{0pt}%
\pgfpathmoveto{\pgfqpoint{5.748413in}{4.157951in}}%
\pgfpathcurveto{\pgfqpoint{5.764885in}{4.157951in}}{\pgfqpoint{5.780685in}{4.164496in}}{\pgfqpoint{5.792333in}{4.176144in}}%
\pgfpathcurveto{\pgfqpoint{5.803981in}{4.187792in}}{\pgfqpoint{5.810526in}{4.203592in}}{\pgfqpoint{5.810526in}{4.220064in}}%
\pgfpathcurveto{\pgfqpoint{5.810526in}{4.236537in}}{\pgfqpoint{5.803981in}{4.252337in}}{\pgfqpoint{5.792333in}{4.263985in}}%
\pgfpathcurveto{\pgfqpoint{5.780685in}{4.275633in}}{\pgfqpoint{5.764885in}{4.282177in}}{\pgfqpoint{5.748413in}{4.282177in}}%
\pgfpathcurveto{\pgfqpoint{5.731940in}{4.282177in}}{\pgfqpoint{5.716140in}{4.275633in}}{\pgfqpoint{5.704492in}{4.263985in}}%
\pgfpathcurveto{\pgfqpoint{5.692844in}{4.252337in}}{\pgfqpoint{5.686300in}{4.236537in}}{\pgfqpoint{5.686300in}{4.220064in}}%
\pgfpathcurveto{\pgfqpoint{5.686300in}{4.203592in}}{\pgfqpoint{5.692844in}{4.187792in}}{\pgfqpoint{5.704492in}{4.176144in}}%
\pgfpathcurveto{\pgfqpoint{5.716140in}{4.164496in}}{\pgfqpoint{5.731940in}{4.157951in}}{\pgfqpoint{5.748413in}{4.157951in}}%
\pgfpathclose%
\pgfusepath{stroke,fill}%
\end{pgfscope}%
\begin{pgfscope}%
\pgfpathrectangle{\pgfqpoint{1.000000in}{0.600000in}}{\pgfqpoint{6.200000in}{4.800000in}} %
\pgfusepath{clip}%
\pgfsetbuttcap%
\pgfsetroundjoin%
\definecolor{currentfill}{rgb}{0.000000,0.500000,0.000000}%
\pgfsetfillcolor{currentfill}%
\pgfsetlinewidth{1.003750pt}%
\definecolor{currentstroke}{rgb}{0.000000,0.500000,0.000000}%
\pgfsetstrokecolor{currentstroke}%
\pgfsetdash{}{0pt}%
\pgfpathmoveto{\pgfqpoint{6.363492in}{4.580100in}}%
\pgfpathcurveto{\pgfqpoint{6.379965in}{4.580100in}}{\pgfqpoint{6.395765in}{4.586645in}}{\pgfqpoint{6.407413in}{4.598293in}}%
\pgfpathcurveto{\pgfqpoint{6.419060in}{4.609941in}}{\pgfqpoint{6.425605in}{4.625741in}}{\pgfqpoint{6.425605in}{4.642213in}}%
\pgfpathcurveto{\pgfqpoint{6.425605in}{4.658686in}}{\pgfqpoint{6.419060in}{4.674486in}}{\pgfqpoint{6.407413in}{4.686134in}}%
\pgfpathcurveto{\pgfqpoint{6.395765in}{4.697782in}}{\pgfqpoint{6.379965in}{4.704326in}}{\pgfqpoint{6.363492in}{4.704326in}}%
\pgfpathcurveto{\pgfqpoint{6.347020in}{4.704326in}}{\pgfqpoint{6.331219in}{4.697782in}}{\pgfqpoint{6.319572in}{4.686134in}}%
\pgfpathcurveto{\pgfqpoint{6.307924in}{4.674486in}}{\pgfqpoint{6.301379in}{4.658686in}}{\pgfqpoint{6.301379in}{4.642213in}}%
\pgfpathcurveto{\pgfqpoint{6.301379in}{4.625741in}}{\pgfqpoint{6.307924in}{4.609941in}}{\pgfqpoint{6.319572in}{4.598293in}}%
\pgfpathcurveto{\pgfqpoint{6.331219in}{4.586645in}}{\pgfqpoint{6.347020in}{4.580100in}}{\pgfqpoint{6.363492in}{4.580100in}}%
\pgfpathclose%
\pgfusepath{stroke,fill}%
\end{pgfscope}%
\begin{pgfscope}%
\pgfpathrectangle{\pgfqpoint{1.000000in}{0.600000in}}{\pgfqpoint{6.200000in}{4.800000in}} %
\pgfusepath{clip}%
\pgfsetbuttcap%
\pgfsetroundjoin%
\definecolor{currentfill}{rgb}{0.000000,0.500000,0.000000}%
\pgfsetfillcolor{currentfill}%
\pgfsetlinewidth{1.003750pt}%
\definecolor{currentstroke}{rgb}{0.000000,0.500000,0.000000}%
\pgfsetstrokecolor{currentstroke}%
\pgfsetdash{}{0pt}%
\pgfpathmoveto{\pgfqpoint{6.978571in}{5.035125in}}%
\pgfpathcurveto{\pgfqpoint{6.995044in}{5.035125in}}{\pgfqpoint{7.010844in}{5.041670in}}{\pgfqpoint{7.022492in}{5.053318in}}%
\pgfpathcurveto{\pgfqpoint{7.034140in}{5.064966in}}{\pgfqpoint{7.040684in}{5.080766in}}{\pgfqpoint{7.040684in}{5.097238in}}%
\pgfpathcurveto{\pgfqpoint{7.040684in}{5.113711in}}{\pgfqpoint{7.034140in}{5.129511in}}{\pgfqpoint{7.022492in}{5.141159in}}%
\pgfpathcurveto{\pgfqpoint{7.010844in}{5.152807in}}{\pgfqpoint{6.995044in}{5.159351in}}{\pgfqpoint{6.978571in}{5.159351in}}%
\pgfpathcurveto{\pgfqpoint{6.962099in}{5.159351in}}{\pgfqpoint{6.946299in}{5.152807in}}{\pgfqpoint{6.934651in}{5.141159in}}%
\pgfpathcurveto{\pgfqpoint{6.923003in}{5.129511in}}{\pgfqpoint{6.916458in}{5.113711in}}{\pgfqpoint{6.916458in}{5.097238in}}%
\pgfpathcurveto{\pgfqpoint{6.916458in}{5.080766in}}{\pgfqpoint{6.923003in}{5.064966in}}{\pgfqpoint{6.934651in}{5.053318in}}%
\pgfpathcurveto{\pgfqpoint{6.946299in}{5.041670in}}{\pgfqpoint{6.962099in}{5.035125in}}{\pgfqpoint{6.978571in}{5.035125in}}%
\pgfpathclose%
\pgfusepath{stroke,fill}%
\end{pgfscope}%
\begin{pgfscope}%
\pgfpathrectangle{\pgfqpoint{1.000000in}{0.600000in}}{\pgfqpoint{6.200000in}{4.800000in}} %
\pgfusepath{clip}%
\pgfsetrectcap%
\pgfsetroundjoin%
\pgfsetlinewidth{2.509375pt}%
\definecolor{currentstroke}{rgb}{1.000000,0.000000,0.000000}%
\pgfsetstrokecolor{currentstroke}%
\pgfsetdash{}{0pt}%
\pgfpathmoveto{\pgfqpoint{1.442857in}{0.725232in}}%
\pgfpathlineto{\pgfqpoint{1.555831in}{1.416418in}}%
\pgfpathlineto{\pgfqpoint{1.668805in}{1.660708in}}%
\pgfpathlineto{\pgfqpoint{1.781778in}{1.662975in}}%
\pgfpathlineto{\pgfqpoint{1.894752in}{1.561820in}}%
\pgfpathlineto{\pgfqpoint{2.007726in}{1.444758in}}%
\pgfpathlineto{\pgfqpoint{2.120700in}{1.361175in}}%
\pgfpathlineto{\pgfqpoint{2.233673in}{1.333160in}}%
\pgfpathlineto{\pgfqpoint{2.346647in}{1.364382in}}%
\pgfpathlineto{\pgfqpoint{2.459621in}{1.447158in}}%
\pgfpathlineto{\pgfqpoint{2.572595in}{1.567893in}}%
\pgfpathlineto{\pgfqpoint{2.685569in}{1.711102in}}%
\pgfpathlineto{\pgfqpoint{2.798542in}{1.862194in}}%
\pgfpathlineto{\pgfqpoint{2.911516in}{2.009215in}}%
\pgfpathlineto{\pgfqpoint{3.024490in}{2.143744in}}%
\pgfpathlineto{\pgfqpoint{3.137464in}{2.261127in}}%
\pgfpathlineto{\pgfqpoint{3.250437in}{2.360215in}}%
\pgfpathlineto{\pgfqpoint{3.363411in}{2.442753in}}%
\pgfpathlineto{\pgfqpoint{3.476385in}{2.512587in}}%
\pgfpathlineto{\pgfqpoint{3.589359in}{2.574770in}}%
\pgfpathlineto{\pgfqpoint{3.702332in}{2.634707in}}%
\pgfpathlineto{\pgfqpoint{3.815306in}{2.697387in}}%
\pgfpathlineto{\pgfqpoint{3.928280in}{2.766783in}}%
\pgfpathlineto{\pgfqpoint{4.041254in}{2.845441in}}%
\pgfpathlineto{\pgfqpoint{4.154227in}{2.934293in}}%
\pgfpathlineto{\pgfqpoint{4.267201in}{3.032671in}}%
\pgfpathlineto{\pgfqpoint{4.380175in}{3.138526in}}%
\pgfpathlineto{\pgfqpoint{4.493149in}{3.248799in}}%
\pgfpathlineto{\pgfqpoint{4.606122in}{3.359905in}}%
\pgfpathlineto{\pgfqpoint{4.719096in}{3.468277in}}%
\pgfpathlineto{\pgfqpoint{4.832070in}{3.570897in}}%
\pgfpathlineto{\pgfqpoint{4.945044in}{3.665759in}}%
\pgfpathlineto{\pgfqpoint{5.058017in}{3.752194in}}%
\pgfpathlineto{\pgfqpoint{5.170991in}{3.831008in}}%
\pgfpathlineto{\pgfqpoint{5.283965in}{3.904397in}}%
\pgfpathlineto{\pgfqpoint{5.396939in}{3.975593in}}%
\pgfpathlineto{\pgfqpoint{5.509913in}{4.048286in}}%
\pgfpathlineto{\pgfqpoint{5.622886in}{4.125815in}}%
\pgfpathlineto{\pgfqpoint{5.735860in}{4.210239in}}%
\pgfpathlineto{\pgfqpoint{5.848834in}{4.301396in}}%
\pgfpathlineto{\pgfqpoint{5.961808in}{4.396140in}}%
\pgfpathlineto{\pgfqpoint{6.074781in}{4.488000in}}%
\pgfpathlineto{\pgfqpoint{6.187755in}{4.567590in}}%
\pgfpathlineto{\pgfqpoint{6.300729in}{4.624183in}}%
\pgfpathlineto{\pgfqpoint{6.413703in}{4.648937in}}%
\pgfpathlineto{\pgfqpoint{6.526676in}{4.640410in}}%
\pgfpathlineto{\pgfqpoint{6.639650in}{4.613081in}}%
\pgfpathlineto{\pgfqpoint{6.752624in}{4.609733in}}%
\pgfpathlineto{\pgfqpoint{6.865598in}{4.718715in}}%
\pgfpathlineto{\pgfqpoint{6.978571in}{5.097238in}}%
\pgfusepath{stroke}%
\end{pgfscope}%
\begin{pgfscope}%
\pgfsetrectcap%
\pgfsetmiterjoin%
\pgfsetlinewidth{1.003750pt}%
\definecolor{currentstroke}{rgb}{0.000000,0.000000,0.000000}%
\pgfsetstrokecolor{currentstroke}%
\pgfsetdash{}{0pt}%
\pgfpathmoveto{\pgfqpoint{1.000000in}{5.400000in}}%
\pgfpathlineto{\pgfqpoint{7.200000in}{5.400000in}}%
\pgfusepath{stroke}%
\end{pgfscope}%
\begin{pgfscope}%
\pgfsetrectcap%
\pgfsetmiterjoin%
\pgfsetlinewidth{1.003750pt}%
\definecolor{currentstroke}{rgb}{0.000000,0.000000,0.000000}%
\pgfsetstrokecolor{currentstroke}%
\pgfsetdash{}{0pt}%
\pgfpathmoveto{\pgfqpoint{1.000000in}{0.600000in}}%
\pgfpathlineto{\pgfqpoint{1.000000in}{5.400000in}}%
\pgfusepath{stroke}%
\end{pgfscope}%
\begin{pgfscope}%
\pgfsetrectcap%
\pgfsetmiterjoin%
\pgfsetlinewidth{1.003750pt}%
\definecolor{currentstroke}{rgb}{0.000000,0.000000,0.000000}%
\pgfsetstrokecolor{currentstroke}%
\pgfsetdash{}{0pt}%
\pgfpathmoveto{\pgfqpoint{1.000000in}{0.600000in}}%
\pgfpathlineto{\pgfqpoint{7.200000in}{0.600000in}}%
\pgfusepath{stroke}%
\end{pgfscope}%
\begin{pgfscope}%
\pgfsetrectcap%
\pgfsetmiterjoin%
\pgfsetlinewidth{1.003750pt}%
\definecolor{currentstroke}{rgb}{0.000000,0.000000,0.000000}%
\pgfsetstrokecolor{currentstroke}%
\pgfsetdash{}{0pt}%
\pgfpathmoveto{\pgfqpoint{7.200000in}{0.600000in}}%
\pgfpathlineto{\pgfqpoint{7.200000in}{5.400000in}}%
\pgfusepath{stroke}%
\end{pgfscope}%
\begin{pgfscope}%
\pgfsetbuttcap%
\pgfsetroundjoin%
\definecolor{currentfill}{rgb}{0.000000,0.000000,0.000000}%
\pgfsetfillcolor{currentfill}%
\pgfsetlinewidth{0.501875pt}%
\definecolor{currentstroke}{rgb}{0.000000,0.000000,0.000000}%
\pgfsetstrokecolor{currentstroke}%
\pgfsetdash{}{0pt}%
\pgfsys@defobject{currentmarker}{\pgfqpoint{0.000000in}{0.000000in}}{\pgfqpoint{0.000000in}{0.055556in}}{%
\pgfpathmoveto{\pgfqpoint{0.000000in}{0.000000in}}%
\pgfpathlineto{\pgfqpoint{0.000000in}{0.055556in}}%
\pgfusepath{stroke,fill}%
}%
\begin{pgfscope}%
\pgfsys@transformshift{1.442857in}{0.600000in}%
\pgfsys@useobject{currentmarker}{}%
\end{pgfscope}%
\end{pgfscope}%
\begin{pgfscope}%
\pgfsetbuttcap%
\pgfsetroundjoin%
\definecolor{currentfill}{rgb}{0.000000,0.000000,0.000000}%
\pgfsetfillcolor{currentfill}%
\pgfsetlinewidth{0.501875pt}%
\definecolor{currentstroke}{rgb}{0.000000,0.000000,0.000000}%
\pgfsetstrokecolor{currentstroke}%
\pgfsetdash{}{0pt}%
\pgfsys@defobject{currentmarker}{\pgfqpoint{0.000000in}{-0.055556in}}{\pgfqpoint{0.000000in}{0.000000in}}{%
\pgfpathmoveto{\pgfqpoint{0.000000in}{0.000000in}}%
\pgfpathlineto{\pgfqpoint{0.000000in}{-0.055556in}}%
\pgfusepath{stroke,fill}%
}%
\begin{pgfscope}%
\pgfsys@transformshift{1.442857in}{5.400000in}%
\pgfsys@useobject{currentmarker}{}%
\end{pgfscope}%
\end{pgfscope}%
\begin{pgfscope}%
\pgftext[x=1.442857in,y=0.544444in,,top]{\sffamily\fontsize{12.000000}{14.400000}\selectfont 0}%
\end{pgfscope}%
\begin{pgfscope}%
\pgfsetbuttcap%
\pgfsetroundjoin%
\definecolor{currentfill}{rgb}{0.000000,0.000000,0.000000}%
\pgfsetfillcolor{currentfill}%
\pgfsetlinewidth{0.501875pt}%
\definecolor{currentstroke}{rgb}{0.000000,0.000000,0.000000}%
\pgfsetstrokecolor{currentstroke}%
\pgfsetdash{}{0pt}%
\pgfsys@defobject{currentmarker}{\pgfqpoint{0.000000in}{0.000000in}}{\pgfqpoint{0.000000in}{0.055556in}}{%
\pgfpathmoveto{\pgfqpoint{0.000000in}{0.000000in}}%
\pgfpathlineto{\pgfqpoint{0.000000in}{0.055556in}}%
\pgfusepath{stroke,fill}%
}%
\begin{pgfscope}%
\pgfsys@transformshift{2.550000in}{0.600000in}%
\pgfsys@useobject{currentmarker}{}%
\end{pgfscope}%
\end{pgfscope}%
\begin{pgfscope}%
\pgfsetbuttcap%
\pgfsetroundjoin%
\definecolor{currentfill}{rgb}{0.000000,0.000000,0.000000}%
\pgfsetfillcolor{currentfill}%
\pgfsetlinewidth{0.501875pt}%
\definecolor{currentstroke}{rgb}{0.000000,0.000000,0.000000}%
\pgfsetstrokecolor{currentstroke}%
\pgfsetdash{}{0pt}%
\pgfsys@defobject{currentmarker}{\pgfqpoint{0.000000in}{-0.055556in}}{\pgfqpoint{0.000000in}{0.000000in}}{%
\pgfpathmoveto{\pgfqpoint{0.000000in}{0.000000in}}%
\pgfpathlineto{\pgfqpoint{0.000000in}{-0.055556in}}%
\pgfusepath{stroke,fill}%
}%
\begin{pgfscope}%
\pgfsys@transformshift{2.550000in}{5.400000in}%
\pgfsys@useobject{currentmarker}{}%
\end{pgfscope}%
\end{pgfscope}%
\begin{pgfscope}%
\pgftext[x=2.550000in,y=0.544444in,,top]{\sffamily\fontsize{12.000000}{14.400000}\selectfont 1}%
\end{pgfscope}%
\begin{pgfscope}%
\pgfsetbuttcap%
\pgfsetroundjoin%
\definecolor{currentfill}{rgb}{0.000000,0.000000,0.000000}%
\pgfsetfillcolor{currentfill}%
\pgfsetlinewidth{0.501875pt}%
\definecolor{currentstroke}{rgb}{0.000000,0.000000,0.000000}%
\pgfsetstrokecolor{currentstroke}%
\pgfsetdash{}{0pt}%
\pgfsys@defobject{currentmarker}{\pgfqpoint{0.000000in}{0.000000in}}{\pgfqpoint{0.000000in}{0.055556in}}{%
\pgfpathmoveto{\pgfqpoint{0.000000in}{0.000000in}}%
\pgfpathlineto{\pgfqpoint{0.000000in}{0.055556in}}%
\pgfusepath{stroke,fill}%
}%
\begin{pgfscope}%
\pgfsys@transformshift{3.657143in}{0.600000in}%
\pgfsys@useobject{currentmarker}{}%
\end{pgfscope}%
\end{pgfscope}%
\begin{pgfscope}%
\pgfsetbuttcap%
\pgfsetroundjoin%
\definecolor{currentfill}{rgb}{0.000000,0.000000,0.000000}%
\pgfsetfillcolor{currentfill}%
\pgfsetlinewidth{0.501875pt}%
\definecolor{currentstroke}{rgb}{0.000000,0.000000,0.000000}%
\pgfsetstrokecolor{currentstroke}%
\pgfsetdash{}{0pt}%
\pgfsys@defobject{currentmarker}{\pgfqpoint{0.000000in}{-0.055556in}}{\pgfqpoint{0.000000in}{0.000000in}}{%
\pgfpathmoveto{\pgfqpoint{0.000000in}{0.000000in}}%
\pgfpathlineto{\pgfqpoint{0.000000in}{-0.055556in}}%
\pgfusepath{stroke,fill}%
}%
\begin{pgfscope}%
\pgfsys@transformshift{3.657143in}{5.400000in}%
\pgfsys@useobject{currentmarker}{}%
\end{pgfscope}%
\end{pgfscope}%
\begin{pgfscope}%
\pgftext[x=3.657143in,y=0.544444in,,top]{\sffamily\fontsize{12.000000}{14.400000}\selectfont 2}%
\end{pgfscope}%
\begin{pgfscope}%
\pgfsetbuttcap%
\pgfsetroundjoin%
\definecolor{currentfill}{rgb}{0.000000,0.000000,0.000000}%
\pgfsetfillcolor{currentfill}%
\pgfsetlinewidth{0.501875pt}%
\definecolor{currentstroke}{rgb}{0.000000,0.000000,0.000000}%
\pgfsetstrokecolor{currentstroke}%
\pgfsetdash{}{0pt}%
\pgfsys@defobject{currentmarker}{\pgfqpoint{0.000000in}{0.000000in}}{\pgfqpoint{0.000000in}{0.055556in}}{%
\pgfpathmoveto{\pgfqpoint{0.000000in}{0.000000in}}%
\pgfpathlineto{\pgfqpoint{0.000000in}{0.055556in}}%
\pgfusepath{stroke,fill}%
}%
\begin{pgfscope}%
\pgfsys@transformshift{4.764286in}{0.600000in}%
\pgfsys@useobject{currentmarker}{}%
\end{pgfscope}%
\end{pgfscope}%
\begin{pgfscope}%
\pgfsetbuttcap%
\pgfsetroundjoin%
\definecolor{currentfill}{rgb}{0.000000,0.000000,0.000000}%
\pgfsetfillcolor{currentfill}%
\pgfsetlinewidth{0.501875pt}%
\definecolor{currentstroke}{rgb}{0.000000,0.000000,0.000000}%
\pgfsetstrokecolor{currentstroke}%
\pgfsetdash{}{0pt}%
\pgfsys@defobject{currentmarker}{\pgfqpoint{0.000000in}{-0.055556in}}{\pgfqpoint{0.000000in}{0.000000in}}{%
\pgfpathmoveto{\pgfqpoint{0.000000in}{0.000000in}}%
\pgfpathlineto{\pgfqpoint{0.000000in}{-0.055556in}}%
\pgfusepath{stroke,fill}%
}%
\begin{pgfscope}%
\pgfsys@transformshift{4.764286in}{5.400000in}%
\pgfsys@useobject{currentmarker}{}%
\end{pgfscope}%
\end{pgfscope}%
\begin{pgfscope}%
\pgftext[x=4.764286in,y=0.544444in,,top]{\sffamily\fontsize{12.000000}{14.400000}\selectfont 3}%
\end{pgfscope}%
\begin{pgfscope}%
\pgfsetbuttcap%
\pgfsetroundjoin%
\definecolor{currentfill}{rgb}{0.000000,0.000000,0.000000}%
\pgfsetfillcolor{currentfill}%
\pgfsetlinewidth{0.501875pt}%
\definecolor{currentstroke}{rgb}{0.000000,0.000000,0.000000}%
\pgfsetstrokecolor{currentstroke}%
\pgfsetdash{}{0pt}%
\pgfsys@defobject{currentmarker}{\pgfqpoint{0.000000in}{0.000000in}}{\pgfqpoint{0.000000in}{0.055556in}}{%
\pgfpathmoveto{\pgfqpoint{0.000000in}{0.000000in}}%
\pgfpathlineto{\pgfqpoint{0.000000in}{0.055556in}}%
\pgfusepath{stroke,fill}%
}%
\begin{pgfscope}%
\pgfsys@transformshift{5.871429in}{0.600000in}%
\pgfsys@useobject{currentmarker}{}%
\end{pgfscope}%
\end{pgfscope}%
\begin{pgfscope}%
\pgfsetbuttcap%
\pgfsetroundjoin%
\definecolor{currentfill}{rgb}{0.000000,0.000000,0.000000}%
\pgfsetfillcolor{currentfill}%
\pgfsetlinewidth{0.501875pt}%
\definecolor{currentstroke}{rgb}{0.000000,0.000000,0.000000}%
\pgfsetstrokecolor{currentstroke}%
\pgfsetdash{}{0pt}%
\pgfsys@defobject{currentmarker}{\pgfqpoint{0.000000in}{-0.055556in}}{\pgfqpoint{0.000000in}{0.000000in}}{%
\pgfpathmoveto{\pgfqpoint{0.000000in}{0.000000in}}%
\pgfpathlineto{\pgfqpoint{0.000000in}{-0.055556in}}%
\pgfusepath{stroke,fill}%
}%
\begin{pgfscope}%
\pgfsys@transformshift{5.871429in}{5.400000in}%
\pgfsys@useobject{currentmarker}{}%
\end{pgfscope}%
\end{pgfscope}%
\begin{pgfscope}%
\pgftext[x=5.871429in,y=0.544444in,,top]{\sffamily\fontsize{12.000000}{14.400000}\selectfont 4}%
\end{pgfscope}%
\begin{pgfscope}%
\pgfsetbuttcap%
\pgfsetroundjoin%
\definecolor{currentfill}{rgb}{0.000000,0.000000,0.000000}%
\pgfsetfillcolor{currentfill}%
\pgfsetlinewidth{0.501875pt}%
\definecolor{currentstroke}{rgb}{0.000000,0.000000,0.000000}%
\pgfsetstrokecolor{currentstroke}%
\pgfsetdash{}{0pt}%
\pgfsys@defobject{currentmarker}{\pgfqpoint{0.000000in}{0.000000in}}{\pgfqpoint{0.000000in}{0.055556in}}{%
\pgfpathmoveto{\pgfqpoint{0.000000in}{0.000000in}}%
\pgfpathlineto{\pgfqpoint{0.000000in}{0.055556in}}%
\pgfusepath{stroke,fill}%
}%
\begin{pgfscope}%
\pgfsys@transformshift{6.978571in}{0.600000in}%
\pgfsys@useobject{currentmarker}{}%
\end{pgfscope}%
\end{pgfscope}%
\begin{pgfscope}%
\pgfsetbuttcap%
\pgfsetroundjoin%
\definecolor{currentfill}{rgb}{0.000000,0.000000,0.000000}%
\pgfsetfillcolor{currentfill}%
\pgfsetlinewidth{0.501875pt}%
\definecolor{currentstroke}{rgb}{0.000000,0.000000,0.000000}%
\pgfsetstrokecolor{currentstroke}%
\pgfsetdash{}{0pt}%
\pgfsys@defobject{currentmarker}{\pgfqpoint{0.000000in}{-0.055556in}}{\pgfqpoint{0.000000in}{0.000000in}}{%
\pgfpathmoveto{\pgfqpoint{0.000000in}{0.000000in}}%
\pgfpathlineto{\pgfqpoint{0.000000in}{-0.055556in}}%
\pgfusepath{stroke,fill}%
}%
\begin{pgfscope}%
\pgfsys@transformshift{6.978571in}{5.400000in}%
\pgfsys@useobject{currentmarker}{}%
\end{pgfscope}%
\end{pgfscope}%
\begin{pgfscope}%
\pgftext[x=6.978571in,y=0.544444in,,top]{\sffamily\fontsize{12.000000}{14.400000}\selectfont 5}%
\end{pgfscope}%
\begin{pgfscope}%
\pgftext[x=4.100000in,y=0.313705in,,top]{\sffamily\fontsize{12.000000}{14.400000}\selectfont x}%
\end{pgfscope}%
\begin{pgfscope}%
\pgfsetbuttcap%
\pgfsetroundjoin%
\definecolor{currentfill}{rgb}{0.000000,0.000000,0.000000}%
\pgfsetfillcolor{currentfill}%
\pgfsetlinewidth{0.501875pt}%
\definecolor{currentstroke}{rgb}{0.000000,0.000000,0.000000}%
\pgfsetstrokecolor{currentstroke}%
\pgfsetdash{}{0pt}%
\pgfsys@defobject{currentmarker}{\pgfqpoint{0.000000in}{0.000000in}}{\pgfqpoint{0.055556in}{0.000000in}}{%
\pgfpathmoveto{\pgfqpoint{0.000000in}{0.000000in}}%
\pgfpathlineto{\pgfqpoint{0.055556in}{0.000000in}}%
\pgfusepath{stroke,fill}%
}%
\begin{pgfscope}%
\pgfsys@transformshift{1.000000in}{0.660380in}%
\pgfsys@useobject{currentmarker}{}%
\end{pgfscope}%
\end{pgfscope}%
\begin{pgfscope}%
\pgfsetbuttcap%
\pgfsetroundjoin%
\definecolor{currentfill}{rgb}{0.000000,0.000000,0.000000}%
\pgfsetfillcolor{currentfill}%
\pgfsetlinewidth{0.501875pt}%
\definecolor{currentstroke}{rgb}{0.000000,0.000000,0.000000}%
\pgfsetstrokecolor{currentstroke}%
\pgfsetdash{}{0pt}%
\pgfsys@defobject{currentmarker}{\pgfqpoint{-0.055556in}{0.000000in}}{\pgfqpoint{0.000000in}{0.000000in}}{%
\pgfpathmoveto{\pgfqpoint{0.000000in}{0.000000in}}%
\pgfpathlineto{\pgfqpoint{-0.055556in}{0.000000in}}%
\pgfusepath{stroke,fill}%
}%
\begin{pgfscope}%
\pgfsys@transformshift{7.200000in}{0.660380in}%
\pgfsys@useobject{currentmarker}{}%
\end{pgfscope}%
\end{pgfscope}%
\begin{pgfscope}%
\pgftext[x=0.944444in,y=0.660380in,right,]{\sffamily\fontsize{12.000000}{14.400000}\selectfont 1}%
\end{pgfscope}%
\begin{pgfscope}%
\pgfsetbuttcap%
\pgfsetroundjoin%
\definecolor{currentfill}{rgb}{0.000000,0.000000,0.000000}%
\pgfsetfillcolor{currentfill}%
\pgfsetlinewidth{0.501875pt}%
\definecolor{currentstroke}{rgb}{0.000000,0.000000,0.000000}%
\pgfsetstrokecolor{currentstroke}%
\pgfsetdash{}{0pt}%
\pgfsys@defobject{currentmarker}{\pgfqpoint{0.000000in}{0.000000in}}{\pgfqpoint{0.055556in}{0.000000in}}{%
\pgfpathmoveto{\pgfqpoint{0.000000in}{0.000000in}}%
\pgfpathlineto{\pgfqpoint{0.055556in}{0.000000in}}%
\pgfusepath{stroke,fill}%
}%
\begin{pgfscope}%
\pgfsys@transformshift{1.000000in}{1.244650in}%
\pgfsys@useobject{currentmarker}{}%
\end{pgfscope}%
\end{pgfscope}%
\begin{pgfscope}%
\pgfsetbuttcap%
\pgfsetroundjoin%
\definecolor{currentfill}{rgb}{0.000000,0.000000,0.000000}%
\pgfsetfillcolor{currentfill}%
\pgfsetlinewidth{0.501875pt}%
\definecolor{currentstroke}{rgb}{0.000000,0.000000,0.000000}%
\pgfsetstrokecolor{currentstroke}%
\pgfsetdash{}{0pt}%
\pgfsys@defobject{currentmarker}{\pgfqpoint{-0.055556in}{0.000000in}}{\pgfqpoint{0.000000in}{0.000000in}}{%
\pgfpathmoveto{\pgfqpoint{0.000000in}{0.000000in}}%
\pgfpathlineto{\pgfqpoint{-0.055556in}{0.000000in}}%
\pgfusepath{stroke,fill}%
}%
\begin{pgfscope}%
\pgfsys@transformshift{7.200000in}{1.244650in}%
\pgfsys@useobject{currentmarker}{}%
\end{pgfscope}%
\end{pgfscope}%
\begin{pgfscope}%
\pgftext[x=0.944444in,y=1.244650in,right,]{\sffamily\fontsize{12.000000}{14.400000}\selectfont 2}%
\end{pgfscope}%
\begin{pgfscope}%
\pgfsetbuttcap%
\pgfsetroundjoin%
\definecolor{currentfill}{rgb}{0.000000,0.000000,0.000000}%
\pgfsetfillcolor{currentfill}%
\pgfsetlinewidth{0.501875pt}%
\definecolor{currentstroke}{rgb}{0.000000,0.000000,0.000000}%
\pgfsetstrokecolor{currentstroke}%
\pgfsetdash{}{0pt}%
\pgfsys@defobject{currentmarker}{\pgfqpoint{0.000000in}{0.000000in}}{\pgfqpoint{0.055556in}{0.000000in}}{%
\pgfpathmoveto{\pgfqpoint{0.000000in}{0.000000in}}%
\pgfpathlineto{\pgfqpoint{0.055556in}{0.000000in}}%
\pgfusepath{stroke,fill}%
}%
\begin{pgfscope}%
\pgfsys@transformshift{1.000000in}{1.828919in}%
\pgfsys@useobject{currentmarker}{}%
\end{pgfscope}%
\end{pgfscope}%
\begin{pgfscope}%
\pgfsetbuttcap%
\pgfsetroundjoin%
\definecolor{currentfill}{rgb}{0.000000,0.000000,0.000000}%
\pgfsetfillcolor{currentfill}%
\pgfsetlinewidth{0.501875pt}%
\definecolor{currentstroke}{rgb}{0.000000,0.000000,0.000000}%
\pgfsetstrokecolor{currentstroke}%
\pgfsetdash{}{0pt}%
\pgfsys@defobject{currentmarker}{\pgfqpoint{-0.055556in}{0.000000in}}{\pgfqpoint{0.000000in}{0.000000in}}{%
\pgfpathmoveto{\pgfqpoint{0.000000in}{0.000000in}}%
\pgfpathlineto{\pgfqpoint{-0.055556in}{0.000000in}}%
\pgfusepath{stroke,fill}%
}%
\begin{pgfscope}%
\pgfsys@transformshift{7.200000in}{1.828919in}%
\pgfsys@useobject{currentmarker}{}%
\end{pgfscope}%
\end{pgfscope}%
\begin{pgfscope}%
\pgftext[x=0.944444in,y=1.828919in,right,]{\sffamily\fontsize{12.000000}{14.400000}\selectfont 3}%
\end{pgfscope}%
\begin{pgfscope}%
\pgfsetbuttcap%
\pgfsetroundjoin%
\definecolor{currentfill}{rgb}{0.000000,0.000000,0.000000}%
\pgfsetfillcolor{currentfill}%
\pgfsetlinewidth{0.501875pt}%
\definecolor{currentstroke}{rgb}{0.000000,0.000000,0.000000}%
\pgfsetstrokecolor{currentstroke}%
\pgfsetdash{}{0pt}%
\pgfsys@defobject{currentmarker}{\pgfqpoint{0.000000in}{0.000000in}}{\pgfqpoint{0.055556in}{0.000000in}}{%
\pgfpathmoveto{\pgfqpoint{0.000000in}{0.000000in}}%
\pgfpathlineto{\pgfqpoint{0.055556in}{0.000000in}}%
\pgfusepath{stroke,fill}%
}%
\begin{pgfscope}%
\pgfsys@transformshift{1.000000in}{2.413188in}%
\pgfsys@useobject{currentmarker}{}%
\end{pgfscope}%
\end{pgfscope}%
\begin{pgfscope}%
\pgfsetbuttcap%
\pgfsetroundjoin%
\definecolor{currentfill}{rgb}{0.000000,0.000000,0.000000}%
\pgfsetfillcolor{currentfill}%
\pgfsetlinewidth{0.501875pt}%
\definecolor{currentstroke}{rgb}{0.000000,0.000000,0.000000}%
\pgfsetstrokecolor{currentstroke}%
\pgfsetdash{}{0pt}%
\pgfsys@defobject{currentmarker}{\pgfqpoint{-0.055556in}{0.000000in}}{\pgfqpoint{0.000000in}{0.000000in}}{%
\pgfpathmoveto{\pgfqpoint{0.000000in}{0.000000in}}%
\pgfpathlineto{\pgfqpoint{-0.055556in}{0.000000in}}%
\pgfusepath{stroke,fill}%
}%
\begin{pgfscope}%
\pgfsys@transformshift{7.200000in}{2.413188in}%
\pgfsys@useobject{currentmarker}{}%
\end{pgfscope}%
\end{pgfscope}%
\begin{pgfscope}%
\pgftext[x=0.944444in,y=2.413188in,right,]{\sffamily\fontsize{12.000000}{14.400000}\selectfont 4}%
\end{pgfscope}%
\begin{pgfscope}%
\pgfsetbuttcap%
\pgfsetroundjoin%
\definecolor{currentfill}{rgb}{0.000000,0.000000,0.000000}%
\pgfsetfillcolor{currentfill}%
\pgfsetlinewidth{0.501875pt}%
\definecolor{currentstroke}{rgb}{0.000000,0.000000,0.000000}%
\pgfsetstrokecolor{currentstroke}%
\pgfsetdash{}{0pt}%
\pgfsys@defobject{currentmarker}{\pgfqpoint{0.000000in}{0.000000in}}{\pgfqpoint{0.055556in}{0.000000in}}{%
\pgfpathmoveto{\pgfqpoint{0.000000in}{0.000000in}}%
\pgfpathlineto{\pgfqpoint{0.055556in}{0.000000in}}%
\pgfusepath{stroke,fill}%
}%
\begin{pgfscope}%
\pgfsys@transformshift{1.000000in}{2.997457in}%
\pgfsys@useobject{currentmarker}{}%
\end{pgfscope}%
\end{pgfscope}%
\begin{pgfscope}%
\pgfsetbuttcap%
\pgfsetroundjoin%
\definecolor{currentfill}{rgb}{0.000000,0.000000,0.000000}%
\pgfsetfillcolor{currentfill}%
\pgfsetlinewidth{0.501875pt}%
\definecolor{currentstroke}{rgb}{0.000000,0.000000,0.000000}%
\pgfsetstrokecolor{currentstroke}%
\pgfsetdash{}{0pt}%
\pgfsys@defobject{currentmarker}{\pgfqpoint{-0.055556in}{0.000000in}}{\pgfqpoint{0.000000in}{0.000000in}}{%
\pgfpathmoveto{\pgfqpoint{0.000000in}{0.000000in}}%
\pgfpathlineto{\pgfqpoint{-0.055556in}{0.000000in}}%
\pgfusepath{stroke,fill}%
}%
\begin{pgfscope}%
\pgfsys@transformshift{7.200000in}{2.997457in}%
\pgfsys@useobject{currentmarker}{}%
\end{pgfscope}%
\end{pgfscope}%
\begin{pgfscope}%
\pgftext[x=0.944444in,y=2.997457in,right,]{\sffamily\fontsize{12.000000}{14.400000}\selectfont 5}%
\end{pgfscope}%
\begin{pgfscope}%
\pgfsetbuttcap%
\pgfsetroundjoin%
\definecolor{currentfill}{rgb}{0.000000,0.000000,0.000000}%
\pgfsetfillcolor{currentfill}%
\pgfsetlinewidth{0.501875pt}%
\definecolor{currentstroke}{rgb}{0.000000,0.000000,0.000000}%
\pgfsetstrokecolor{currentstroke}%
\pgfsetdash{}{0pt}%
\pgfsys@defobject{currentmarker}{\pgfqpoint{0.000000in}{0.000000in}}{\pgfqpoint{0.055556in}{0.000000in}}{%
\pgfpathmoveto{\pgfqpoint{0.000000in}{0.000000in}}%
\pgfpathlineto{\pgfqpoint{0.055556in}{0.000000in}}%
\pgfusepath{stroke,fill}%
}%
\begin{pgfscope}%
\pgfsys@transformshift{1.000000in}{3.581726in}%
\pgfsys@useobject{currentmarker}{}%
\end{pgfscope}%
\end{pgfscope}%
\begin{pgfscope}%
\pgfsetbuttcap%
\pgfsetroundjoin%
\definecolor{currentfill}{rgb}{0.000000,0.000000,0.000000}%
\pgfsetfillcolor{currentfill}%
\pgfsetlinewidth{0.501875pt}%
\definecolor{currentstroke}{rgb}{0.000000,0.000000,0.000000}%
\pgfsetstrokecolor{currentstroke}%
\pgfsetdash{}{0pt}%
\pgfsys@defobject{currentmarker}{\pgfqpoint{-0.055556in}{0.000000in}}{\pgfqpoint{0.000000in}{0.000000in}}{%
\pgfpathmoveto{\pgfqpoint{0.000000in}{0.000000in}}%
\pgfpathlineto{\pgfqpoint{-0.055556in}{0.000000in}}%
\pgfusepath{stroke,fill}%
}%
\begin{pgfscope}%
\pgfsys@transformshift{7.200000in}{3.581726in}%
\pgfsys@useobject{currentmarker}{}%
\end{pgfscope}%
\end{pgfscope}%
\begin{pgfscope}%
\pgftext[x=0.944444in,y=3.581726in,right,]{\sffamily\fontsize{12.000000}{14.400000}\selectfont 6}%
\end{pgfscope}%
\begin{pgfscope}%
\pgfsetbuttcap%
\pgfsetroundjoin%
\definecolor{currentfill}{rgb}{0.000000,0.000000,0.000000}%
\pgfsetfillcolor{currentfill}%
\pgfsetlinewidth{0.501875pt}%
\definecolor{currentstroke}{rgb}{0.000000,0.000000,0.000000}%
\pgfsetstrokecolor{currentstroke}%
\pgfsetdash{}{0pt}%
\pgfsys@defobject{currentmarker}{\pgfqpoint{0.000000in}{0.000000in}}{\pgfqpoint{0.055556in}{0.000000in}}{%
\pgfpathmoveto{\pgfqpoint{0.000000in}{0.000000in}}%
\pgfpathlineto{\pgfqpoint{0.055556in}{0.000000in}}%
\pgfusepath{stroke,fill}%
}%
\begin{pgfscope}%
\pgfsys@transformshift{1.000000in}{4.165995in}%
\pgfsys@useobject{currentmarker}{}%
\end{pgfscope}%
\end{pgfscope}%
\begin{pgfscope}%
\pgfsetbuttcap%
\pgfsetroundjoin%
\definecolor{currentfill}{rgb}{0.000000,0.000000,0.000000}%
\pgfsetfillcolor{currentfill}%
\pgfsetlinewidth{0.501875pt}%
\definecolor{currentstroke}{rgb}{0.000000,0.000000,0.000000}%
\pgfsetstrokecolor{currentstroke}%
\pgfsetdash{}{0pt}%
\pgfsys@defobject{currentmarker}{\pgfqpoint{-0.055556in}{0.000000in}}{\pgfqpoint{0.000000in}{0.000000in}}{%
\pgfpathmoveto{\pgfqpoint{0.000000in}{0.000000in}}%
\pgfpathlineto{\pgfqpoint{-0.055556in}{0.000000in}}%
\pgfusepath{stroke,fill}%
}%
\begin{pgfscope}%
\pgfsys@transformshift{7.200000in}{4.165995in}%
\pgfsys@useobject{currentmarker}{}%
\end{pgfscope}%
\end{pgfscope}%
\begin{pgfscope}%
\pgftext[x=0.944444in,y=4.165995in,right,]{\sffamily\fontsize{12.000000}{14.400000}\selectfont 7}%
\end{pgfscope}%
\begin{pgfscope}%
\pgfsetbuttcap%
\pgfsetroundjoin%
\definecolor{currentfill}{rgb}{0.000000,0.000000,0.000000}%
\pgfsetfillcolor{currentfill}%
\pgfsetlinewidth{0.501875pt}%
\definecolor{currentstroke}{rgb}{0.000000,0.000000,0.000000}%
\pgfsetstrokecolor{currentstroke}%
\pgfsetdash{}{0pt}%
\pgfsys@defobject{currentmarker}{\pgfqpoint{0.000000in}{0.000000in}}{\pgfqpoint{0.055556in}{0.000000in}}{%
\pgfpathmoveto{\pgfqpoint{0.000000in}{0.000000in}}%
\pgfpathlineto{\pgfqpoint{0.055556in}{0.000000in}}%
\pgfusepath{stroke,fill}%
}%
\begin{pgfscope}%
\pgfsys@transformshift{1.000000in}{4.750264in}%
\pgfsys@useobject{currentmarker}{}%
\end{pgfscope}%
\end{pgfscope}%
\begin{pgfscope}%
\pgfsetbuttcap%
\pgfsetroundjoin%
\definecolor{currentfill}{rgb}{0.000000,0.000000,0.000000}%
\pgfsetfillcolor{currentfill}%
\pgfsetlinewidth{0.501875pt}%
\definecolor{currentstroke}{rgb}{0.000000,0.000000,0.000000}%
\pgfsetstrokecolor{currentstroke}%
\pgfsetdash{}{0pt}%
\pgfsys@defobject{currentmarker}{\pgfqpoint{-0.055556in}{0.000000in}}{\pgfqpoint{0.000000in}{0.000000in}}{%
\pgfpathmoveto{\pgfqpoint{0.000000in}{0.000000in}}%
\pgfpathlineto{\pgfqpoint{-0.055556in}{0.000000in}}%
\pgfusepath{stroke,fill}%
}%
\begin{pgfscope}%
\pgfsys@transformshift{7.200000in}{4.750264in}%
\pgfsys@useobject{currentmarker}{}%
\end{pgfscope}%
\end{pgfscope}%
\begin{pgfscope}%
\pgftext[x=0.944444in,y=4.750264in,right,]{\sffamily\fontsize{12.000000}{14.400000}\selectfont 8}%
\end{pgfscope}%
\begin{pgfscope}%
\pgfsetbuttcap%
\pgfsetroundjoin%
\definecolor{currentfill}{rgb}{0.000000,0.000000,0.000000}%
\pgfsetfillcolor{currentfill}%
\pgfsetlinewidth{0.501875pt}%
\definecolor{currentstroke}{rgb}{0.000000,0.000000,0.000000}%
\pgfsetstrokecolor{currentstroke}%
\pgfsetdash{}{0pt}%
\pgfsys@defobject{currentmarker}{\pgfqpoint{0.000000in}{0.000000in}}{\pgfqpoint{0.055556in}{0.000000in}}{%
\pgfpathmoveto{\pgfqpoint{0.000000in}{0.000000in}}%
\pgfpathlineto{\pgfqpoint{0.055556in}{0.000000in}}%
\pgfusepath{stroke,fill}%
}%
\begin{pgfscope}%
\pgfsys@transformshift{1.000000in}{5.334533in}%
\pgfsys@useobject{currentmarker}{}%
\end{pgfscope}%
\end{pgfscope}%
\begin{pgfscope}%
\pgfsetbuttcap%
\pgfsetroundjoin%
\definecolor{currentfill}{rgb}{0.000000,0.000000,0.000000}%
\pgfsetfillcolor{currentfill}%
\pgfsetlinewidth{0.501875pt}%
\definecolor{currentstroke}{rgb}{0.000000,0.000000,0.000000}%
\pgfsetstrokecolor{currentstroke}%
\pgfsetdash{}{0pt}%
\pgfsys@defobject{currentmarker}{\pgfqpoint{-0.055556in}{0.000000in}}{\pgfqpoint{0.000000in}{0.000000in}}{%
\pgfpathmoveto{\pgfqpoint{0.000000in}{0.000000in}}%
\pgfpathlineto{\pgfqpoint{-0.055556in}{0.000000in}}%
\pgfusepath{stroke,fill}%
}%
\begin{pgfscope}%
\pgfsys@transformshift{7.200000in}{5.334533in}%
\pgfsys@useobject{currentmarker}{}%
\end{pgfscope}%
\end{pgfscope}%
\begin{pgfscope}%
\pgftext[x=0.944444in,y=5.334533in,right,]{\sffamily\fontsize{12.000000}{14.400000}\selectfont 9}%
\end{pgfscope}%
\begin{pgfscope}%
\pgftext[x=0.768962in,y=3.000000in,,bottom]{\sffamily\fontsize{12.000000}{14.400000}\selectfont y}%
\end{pgfscope}%
\end{pgfpicture}%
\makeatother%
\endgroup%
}
	\caption{Primer odabira modela pri linearnoj regresiji polinomom}
	\label{fig:odabir}
\end{figure}

\par
Na još jedan od mogućih problema nailazi se u slučaju neprikladnih podataka. Nekada ulazni atributi ne daju dovoljno informacija o izlaznim. Takođe, moguće je da podataka jednostavno nema dovoljno. U ovom slučaju, sistem ne dobija dovoljno bogat skup informacija kako bi uspešno izvršio generalizaciju. S druge strane, moguće je da postoji prevelika količina podataka. Tada se pribegava pažljivom odabiru podataka koji se koriste za učenje ali ovo u opštem slučaju  treba izbegavati jer su podaci izuzetno vredan element procesa mašinskog učenja. Još jedan problem vezan za podatke može biti njihova nepotpunost. Na primer, moguće je da u nekim instancama postoje nedostajuće vrednosti atributa. \par
Kako je najčešče potrebno pretprocesirati podatke u sklopu procesa mašinskog učenja, moguće je da u ovom postupku dođe do greške. Primera radi, prilikom rada sa konvolutivnim neuronskim mrežama, o kojima će biti reči u jednoj od narednih glava, nekada se slike  u boji pretvaraju u crno-bele. Ako se primeni transformacija koja onemogućuje razlikovanje objekata koji su različiti u početnoj slici a razlikovanje je neophodno za ispravno učenje, tada proces treniranja neće teći kako je planirano. \par


% [ONO ZA RGB U GRAY U DQN BI BILO SOLIDNO OVDE UBACITI] \par
Problem može da nastane i ukoliko nije odabran pravi algoritam za učenje, ukoliko se loše pristupilo procesu optimizacije, prilikom lošeg procesa evaluacije i, naravno, prilikom loše implementacije algoritma. Sve ove prepreke često je moguće prevazići ali je jasno da je neophodno biti izuzetno pažljiv prilikom celog procesa mašinskog učenja.
\chapter{Neuronske mreže}
\label{ch:nn}

Neuronske mreže (eng. neural networks) predstavljaju danas izuzetno popularan vid mašinskog učenja. Ovi modeli izuzetno su fleksibilni i imaju široku primenu.  Koriste se za prepoznavanje govora, prevođenje, prepoznavanje oblika na slikama, upravljanje vozilima, uspostavljanje dijagnoza u medicini, igranje igara itd. [Neuronskim mrežama može se aproksimirati proizvoljna neprekidna funkcija.] Pun naziv je veštačka neuronska mreža (eng. artificial neural network, skr. ANN) jer se ovakvi modeli idejno zasnivaju na načinu na koji mozak funkcioniše. Osnovne gradivne jedinice, neuroni, zasnovani su na neuronima u mozgu, dok veze između njih predstavljaju sinapse\footnote{Sinapsa je struktura koja omogućuje komunikaciju između neurona.}. Te veze opisuju odnose između neurona i obično im se dodeljuje numerička težina.

\par
Postoji nekoliko različitih vrsta neuronskih mreža. Tipičan primer jesu neuronske mreže sa propagacijom unapred. Ime proističe iz činjenice da podaci teku od ulaza mreže do izlaza, bez postojanja ikakve povratne sprege. Neuronske mreže sa propagacijom unapred sastoje se iz slojeva neurona. Ukoliko se u ovaj model uvede neki tip povratne sprege, tada se govori o rekurentnim neuronskim mrežama. Pri radu sa slikama i raznim drugim vrstama signala, najčešće se koriste konvolutivne neuronske mreže, o kojima će biti reči kasnije. Ono što je zajedničko je da su neuronske mreže sposobne za izdvajanje određenih karakteristika u podacima koji se obrađuju. To znači da se vrši kreiranje novih atributa na osnovu već postojećih. Taj proces naziva se ekstrakcijom atributa i smatra se da je to jedan od najbitnijih razloga za delotvornost neuronskih mreža.
\par
Za uspeh neuronskih mreža zaslužna je njihova fleksibilnost ali su rezultati dobijeni najpre eksperimentisanjem. Naime, veliki deo zaključaka o ponašanju neuronskih mreža u raznim situacijama nije teorijski potkrepljen. Stoga, istraživački rad vezan za neuronske mreže zahteva dosta pokušaja da bi se došlo do uspeha.

\section{Neuronske mreže sa propagacijom unapred}

Neuronske mreže sa propagacijom unapred jedna su od najkorišćenijih vrsta neuronskih mreža. Gradivni elementi ovakvog modela, neuroni (koji se još nazivaju i jedinicama), organizuju se u slojeve koji se nadovezuju i time čine neuronsku mrežu. Organizacija neurona i slojeva, uključujući i veze između neurona, predstavlja arhitekturu mreže. Prvi sloj mreže naziva se ulaznim slojem dok se poslednji sloj naziva izlaznim slojem. Neuroni prvog sloja kao argumente primaju ulaze mreže dok neuroni svakog od preostalih slojeva kao svoje ulaze prihvataju izlaze prethodnog sloja. Svi slojevi koji svoje izlaze prosleđuju narednom sloju nazivaju se skrivenim slojevima. Mreže sa više od jednog skrivenog sloja nazivaju se dubokim neuronskim mrežama.  Broj slojeva mreže određuje njenu dubinu. Termin "duboko učenje" nastao je baš iz ove terminologije.


\par
Svaki neuron opisuje se pomoću vektora $w = (w_0, ..., w_n)$ koji se naziva vektorom težina. Ulazni parametrar $x = (x_1, ..., x_n)$ linearno se transformiše na sledeći način:
\begin{equation}
\label{eq:neuron}
		f_w(x) = w_0 + \sum_{i=1}^{n} x_nw_n 
\end{equation}
a zatim se primenjuje takozvana aktivaciona funkcija, $g$. Izlaz iz neurona je $g(f_w(x))$ i, uprkos linearnosti prve transformacije, izlaz ne mora biti linearna transformacija ulaza, tj. $g$ najčešće nije linearna funkcija. Za $g_i$ bira se nelinearna funkcija jer se u suprotnom kao celokupna transformacija koju neuron vrši dobija linearna funkcija; na ovaj način, mreža bi predstavljala linearnu funkciju i ne bi bilo moguće njom aproksimirati nelinearne funkcije dovoljno dobro. Vrednost $w_0$ naziva se slobodnim članom. Nekada se vektor $x$ transformiše tako da bude oblika $x = (1, x_1, ..., x_n)$ kako bi izraz \eqref{eq:neuron} imao kraći zapis $f_w(x) = w \cdot x$, gde $\cdot$ označava skalarni proizvod.\\

\begin{figure}
	\centering
	\resizebox{.5\linewidth}{!}{
\tikzset{basic/.style={draw,text width=1em,text badly centered}}
\tikzset{input/.style={basic,circle}}
\tikzset{weights/.style={text width=1em, text badly centered}}
\tikzset{functions/.style={basic,circle}}

    \begin{tikzpicture}
        \node[functions] (center) {};
        \node[below of=center,font=\scriptsize,text width=4em] {Aktivaciona funkcija};
        %\draw[thick] (0.5em,0.5em) -- (0,0.5em) -- (0,-0.5em) -- (-0.5em,-0.5em);
        %\draw (0em,0.75em) -- (0em,-0.75em);
        \draw (0em,0em) -- (-0.6em,0em);
        \draw (0em,0em) -- (0.5em, 0.5em) ;
        \node[right of=center] (right) {};
            \path[draw,->] (center) -- (right);
        \node[functions,left=3em of center] (left) {$\sum$};
            \path[draw,->] (left) -- (center);
        \node[weights,left=3em of left] (2) {} -- (2) node[input,left of=2] (l2) {$x_2$};
        \node[above=-0.1em of 2] (ww2) {$w_2$} ;
            %\path[draw,->] (l2) -- (2);
            \path[draw,->] (l2) -- (left);
        \node[below of=2] (dots) {} -- (dots) node[left of=dots] (ldots) {$\vdots$};
        \node[above=-0.7em of dots] (wn) {$w_n$} ;
        \node[weights,below of=dots] (n) {} -- (n) node[input,left of=n] (ln) {$x_n$};
            %\path[draw,->] (ln) -- (n);
            \path[draw,->] (ln) -- (left);
        \node[weights,above of=2] (1) {$w_1$} -- (1) node[input,left of=1] (l1) {$x_1$};
            %\path[draw,->] (l1) -- (1);
            \path[draw,->] (l1) -- (left);
        \node[weights,above of=1] (0) {} -- (0) node[input,left of=0] (l0) {$1$};
        \node[weights, below=-0.3em of 0] (w0) {$w_0$} ;
            %\path[draw,->] (l0) -- (0);
            \path[draw,->] (l0) -- (left);
        \node[below=0.5em of ln,font=\scriptsize] {Ulazi};
    \end{tikzpicture}}
	\caption{Neuron}
	\label{fig:neuron}
\end{figure}

Model, tj. neuronska mreža, formalno se definiše na sledeći način:
\begin{equation}
	\begin{gathered}
		h_0 = x  \\
	 	h_i = g_i(W_ih_{i-1} + w_{i0}) \text{,~za~} i=1, ..., L
	\end{gathered}
\end{equation}

gde je $x$ vektor ulaza u mrežu predstavljen kao kolona, $W_i$ je matrica čija $j$-ta vrsta predstavlja vektor težina $j$-tog neurona u sloju $i$ a $w_{i0}$ je kolona slobodnih članova svih jedinica u sloju $i$. Funkcije $g_i$ su nelinearna aktivaciona funkcija i za vektor $t=(t_1, ..., t_n)$, $g_i(t)$ predstavlja kolonu $(g_i(t_1), ..., g_i(t_n))^T$. Na ovaj način dobija se funkcija čiji su parametri $W_i$ i $w_{i0}$ za $i=1,...,L$. Ako se parametri označe sa $w$, tada se model zapisati kao $f_w$. Parametri $w$ mogu se pronaći matematičkom optimizacijom nekog kriterijuma kvaliteta modela. Taj proces opisan je u delu ~\ref{subsec:optimizacija}.

\subsection{Aktivacione funkcije}


Preteča neuronskih mreža, perceptron, je model koji se sastoji samo iz jednog neurona čija je aktivaciona funkcija data sledećim izrazom:
\begin{equation}
	g(x)=
	\begin{cases}
		1, 	& \text{~ako~} x \geq 0 \\
		0, 	& \text{~inače}
	\end{cases}
\end{equation} 

Definicija aktivacione funkcije perceptrona znači da njegova primena ima relativno jako ograničenje. S obzirom na to da će ulaz u funkciju $g$ biti linearna kombinacija ulaza i parametara, perceptron će moći da napravi podelu prostora određenu hiperravni\footnote{Hiperravan je uopštenje ravni u trodimenzionom prostoru; predstavlja potprostor dimenzije za $1$ manje od prostora u kom se nalazi.}. Ukoliko skup ulaznih podataka nije moguće podeliti linearnom funkcijom, ovakvo ponašanje nije zadovoljavajuće.
\par
Dakle, neophodno je naći druge funkcije koje služe kao aktivacione funkcije. Poželjna svojstva aktivacione funkcije su:
\begin{itemize}
	\item Nelinearnost: Kao što je objašnjeno ranije, kompozicija linearnih funkcija daje linearnu funkciju, što onemogućuje dovoljno preciznu aproksimaciju nelinearnih funkcija;
	\item Diferencijabilnost: Optimizacija se najčešće vrši metodima koji koriste gradijent funkcije;
	\item Monotonost: Ako aktivaciona funkcija nije monotona, povećavanje nekog od težinskih parametara neurona, umesto da poveća izlaz i time proizvede jači signal, može imati suprotan efekat;
	\item Ograničenost: Ukoliko vrednosti unutar neuronske mreže nisu ograničene, moguće je da dođe do pojavljivanja ogromnih vrednosti koje potencijalno dovode do prekoračenja. Ograničene aktivacione funkcije ovo znatno ublažuju. 
\end{itemize}

Dozvoljeno je da aktivaciona funkcija ne poseduje neko od navedenih svojstava ali ovime se može znatno smanjiti brzina konvergencije. 
\par
Najčešće korišćene aktivacione funkcije su:
\begin{itemize}
	\item Sigmoidna funkcija: $\sigma (x) = \cfrac{1}{1+e^{-x}}$
	\item Tangens hiperbolički: $tanh(x) = \cfrac{e^{2x}-1}{e^{2x}+1}$
	\item Ispravljena linearna jedinica: $ReLU(x) = max(0, x)$
\end{itemize}

Sigmoidna funkcija bila je najkorišćenija aktivaciona funkcija pri radu sa neuronskim mrežama. Ograničena je (sve slike nalaze se u intervalu $(-1, 1)$), monotona i diferencijabilna u svakoj tački skupa $\mathbb{R}$. Međutim, što se argument više udaljava od nule, to nagib funkcije postaje manji. To znači da će gradijent funkcije biti mali i da će učenje teći jako sporo. \\
Tangens hiperbolički srodan je sigmoidnoj funkciji ($tanh(x) = 2\sigma (x) - 1$) ali je imala veći uspeh od sigmoidne funkcije.  U okolini nule, ova funkcija slična je identičkoj, što olakšava učenje. Međutim, i pri korišćenju ove funkcije može se naići na problem sa malim gradijentima ukoliko se argument dovoljno udalji od nule.\\
Uprkos tome što za razliku od prethodne dve funkcije nije ni ograničena ni diferencijabilna u svim tačkama domena, danas je ispravljena linearna jedinica najpopularniji izbor za aktivacionu funkciju. Funkcija je jednaka identitetu desno od nule i stoga se gradijent ne menja. Takođe, verovatnoća da se traži gradijent u tački u kojoj funkcija nije diferencijabilna je mala. Ipak, ni ova funkcija nije bez mana; problem često pravi deo levo od nule, gde je funkcija konstantna. To znači da je gradijent nula i da se prilikom optimizacije težine neurona neće izmeniti. Zbog nedostatka promene, može se desiti da neki neuroni u mreži postanu pasivni, tj. da im izlaz postane $0$. Za ovaj problem postoje rešenja; jedno jeste da izlaz funkcije desno od nule ne bude konstanta $0$ već $\alpha x$, za neko malo $\alpha$. Ta modifikovana ReLU funkcija naziva se nakošena ispravljena linearna jedinica (eng. leaky rectified linear unit). 
\\
Iako sve ove funkcije imaju prednosti i mane u odnosu na preostale, ne postoji jedinstveni izbor nego je na osnovu problema neophodno zaključiti koju je aktivacionu funkciju najbolje koristiti.

\subsubsection{Izlazni sloj}

Neuronske mreže koriste se pri regresiji, određivanju funkcije koja opisuje vezu izmedju ulaza i izlaza, i klasifikaciji, svrstavanje ulaznih vektora u jednu od konačnog broja kategorija. Pri regresiji, u poslednjem sloju ne primenjuje se aktivaciona funkcija. Proces optimizacije svodi se na minimizaciju funkcije greške. Kod rešavanja problema klasifikacije (u $N$ kategorija), koristi se funkcija mekog maksimuma (eng. softmax):
\begin{equation}
	softmax(x) = \bigg( \cfrac{e^{x_1}}{\sum_{i=1}^{N} e^{x_i}}, ~...~, \cfrac{e^{x_N}}{\sum_{i=1}^{N} e^{x_i}} \bigg)
\end{equation}
Suma ovako dobijenog vektora je $1$ i stoga može predstavljati diskretnu raspodelu verovatnoća. Za vrednost aproksimacije uzima se kategorija kojoj odgovara najviša vrednost izlaznog vektora. Za optimizaciju pri radu sa probabilističkim problemima, kao što je problem klasifikacije, primenjuje se metod maksimalne verodostojnosti (eng. maximum likelihood estimate), odnosno traži se maksimum sledećeg izraza: 
\begin{equation}
	P(y_1, ..., y_N | x_1, ..., x_N)
\end{equation}.

\subsection{Optimizacija}
\label{subsec:optimizacija}

Ukoliko je neuronska mreža predstavljena kao funkcija $f_w$, gde su $w$ parametri mreže, neophodno je izvršiti minimizaciju\footnote{Naravno, optimizacioni metodi primenjuju se i na maksimizaciju ali je u slučaju mašinskog učenja najčešće neophodno minizovati funkciju greške.} funkcije koja predstavlja kriterijum kvaliteta aproksimacije. Problem optimizacije u slučaju neuronskih mreža težak je zbog nekonveksnosti. Ona čini neke metode teško primenljivim ili izuzetno sporim. Moguće je i završiti u lokalnom optimumu funkcije. Uobičajeno se koriste metodi zasnovani na gradijentu funkcije. Postoje metodi drugog reda, zasnovani na hesijanu \footnote{Hesijan je matrica parcijalnih izvoda drugog reda.} ali je njegovo računanje u slučaju većeg broja parametara preskupo. \par
Učenje funkcioniše na sledeći način za fiksirane ulaze $x$ posmatra se njima uparen izlaz $y$ i $f_w(x)$ a zatim i $L(y, f_w(x))$, odnosno funkcija greške između stvarne i očekivane vrednosti. Koristeći algoritam propagacije unazad (\ref{subsub:backprop}) uz neki od algoritama za optimizaciju, vrši se minimzacija funkcije $L$ u odnosu na parametre mreže, $w$. 
%Jedna od najčešće korišćenih funkcija greške pri regresiji je srednjekvadratna greška:
%\begin{equation}
%	L(y, \hat{y}) =\frac{1}{2}\norm{y-\hat{y}}_2^2 = \frac{1}{2}\sum_{i=1}^{2} (y_i - \hat{y}_i)^2
%\end{equation}
%Kako su pri minimizaciji $L$ za određeni par $x,y$ jedino težine nepoznate, u nastavku će %$L(w)$ označavati funkciju greške za neki par.

\subsubsection{Metod gradijentnog spusta i stohastičkog gradijentnog spusta}

Gradijent funkcije $f:\mathbb{R}^n \rightarrow \mathbb{R}$ u tački $x=(x_1, ..., x_n)$ označava se sa $\nabla f$ i predstavlja vektor parcijalnih izvoda u toj tački:
\begin{equation}
\nabla f(x) = \bigg( \cfrac{\partial f}{\partial x_1}(x),...,\cfrac{\partial f}{\partial x_n}(x) \bigg)
\end{equation}.

Gradijent funkcije u tački $x$ predstavlja pravac i smer najbržeg rasta funkcije pa $- \nabla f(x)$ predstavlja pravac smer najbržeg opadanja funkcije. Kako se najčešće minimizuje funkcija greške, u oznaci $L$, na dalje je korišćeno to ime za funkciju umesto $f$.
\par

Metod gradijentnog spusta jedan je od najstarijih metoda  optimizacije. Iterativnim pristupom minimizuje se konveksna diferencijabilna funkcija. Polazeći od nasumično odabrane tačke i prateći pravac i smer gradijenta u svakom koraku, dolazi se do minimuma funkcije. Iterativni korak definisan je na sledeći način:
\begin{equation}
	\label{eq:gradijentni_spust}
	w_{k+1} = w_k - \alpha_k \nabla L(w_k), k=0, 1, 2, ... ,
\end{equation}
gde je $w_0$ ta nasumično odabrana početna tačka a $\alpha_k$ je pozitivan realan broj koji se naziva veličinom koraka ili stopom učenja (eng. learning rate). 
Za funkciju greške u ovom slučaju uzima se srednjekvadratno odstupanje:
\begin{equation}
	\frac{1}{2N}\sum_{i=1}^{N} \norm{y_i - f_w(x_i)}_2^2
\end{equation}

Bitno je pažljivo odabrati veličinu koraka jer ova vrednost može uticati na konvergenciju. Jedan primer odabira veličine koraka jeste niz za koji važe Robins Monroovi uslovi\footnote{\url{https://en.wikipedia.org/wiki/Stochastic_approximation}} :
\begin{equation}
	 \sum_{k=0}^{\infty} \alpha_k = \infty \hspace{2cm} \sum_{k=0}^{\infty} \alpha_k^2 < \infty 
	\end{equation}
Jednostavniji pristup bio bi da se odabere mali pozitivan parametar $\alpha$ i da za svako $k$ važi $\alpha_k = \alpha$.
\par
Postavlja se i pitanje koliko koraka načiniti pre zaustavljanja. U praksi se koristi nekoliko kriterijuma kao što su zaustavljanje kada su dve uzastopne vrednosti $w_k$ i $w_{k+1}$ dovoljno bliske, kada su vrednosti funkcije za dve uzastopne vrednosti dovoljno bliske ali se može zaustaviti i nakon unapred određenog broja koraka. Postoji još kriterijuma i moguće ih je kombinovati.
\par
Iako jednostavan i široko primenljiv metod optimizacije, gradijentni spust nije najbolji izbor. Naime, pravac najbržeg uspona funkcije nije uvek i pravac koji osigurava najbrže približavanje optimumu funkcije. U praksi, gradijentni spust ume da proizvodi cik-cak kretanje koje dovodi do spore konvergencije. Takođe, za jedan iterativni korak neophodno je proći kroz sve parove ulaza i izlaza, što u slučaju velikog skupa podataka za obučavanje može biti jako velika količina podataka.
\par
Za obučavanje neuronskih mreža češće se koristi metod stohastičkog gradijentnog spusta. Pretpostavka je da je funkcija koja se optimizuje oblika:

\begin{equation}
		L(w) = \frac{1}{N}\sum_{i=1}^{N} L_i(w)	
\end{equation}
odnosno da se može predstaviti kao prosek nekih $N$ funkcija. Kako je neuronska mreža jedan od metoda mašinskog učenja, na raspolaganju je skup za obučavanje pa se funkcija greške na celom skupu može predstaviti kao prosek grešaka na pojedinačnim instancama skupa. Novi oblik 
jednakosti ~\eqref{eq:gradijentni_spust} je:
\begin{equation}
	w_{k+1} = w_k - \alpha_k \bigg( \frac{1}{N}\sum_{i=1}^{N} \nabla L_i(w_k) \bigg)\text{,~} k=0,1,2, ... 
\end{equation}.
Pri korišćenju stohastičkog gradijentnog spusta za minimizaciju funkcije greške, iterativni korak izgleda ovako:
\begin{equation}
	w_{k+1} = w_k - \alpha \nabla L_i(w_k) 
\end{equation}
Postoje razni načini za odabir $i$ u nekom koraku, kao što je $i=k (mod N) + 1$, gde je $N$ veličina skupa za obučavanje. Još jedan primer je nasumični odabir instance u svakom koraku. Kakav god način izbora bio, neophodno je iskoristiti sve greške. Moguće je proći greške iz skupa za obučavanje i nekoliko puta dok se ne postigne željeni nivo aproksimacije. 
\par
Kako ova aproksimacija može biti prilično neprecizna, pribegava se kompromisu: prilikom iterativnog koraka ne koriste se pojedinačne instance već prosek nekog podskupa skupa za obučavanje (eng. minibatch). Pri treniranju neuronskih mreža, ovo je uobičajeni pristup.
\par
Ovom aproksimacijom mogu se izbeći lokalni minimumi funkcije. Metod stohastičkog gradijentnog spusta manje je računski zahtevna od gradijentnog spusta ali je manje precizna i neophodan je veći broj iteracija kako bi se dostigao minimum.
\par
Postoje razni metodi optimizacije koji se koriste pri mašinskom učenju. Neki menjaju veličinu koraka u zavisnosti od prethodnih izračunatih koraka i gradijenata. Takvi metodi nazivaju se adaptivnim metodima optimizacije. Primer adaptivnih metoda optimizacije su Adam i RMSProp.

\subsubsection{RMSProp}

Algoritam RMSProp (eng. root mean square propagation) predložio je Džof Hinton na jednom od svojih predavanja na sajtu Kursera~\footnote{\url{http://www.cs.toronto.edu/~tijmen/csc321/slides/lecture_slides_lec6.pdf}}. Ovo je algoritam optimizacije korišćen prilikom razvijanja DQN algoritma. Glavna ideja je čuvanje dosadašnjeg otežanog proseka kvadrata gradijenta funkcije koji će biti obeležen sa $g_k$. Simbol $\odot$ obeležava pokoordinatno množenje dva vektora. Kako algoritam nije objavljen u radu, može se naći veliki broj implementacija. U nastavku je predstavljen algoritam u skladu sa implementacijom iz biblioteke Keras, koja je korišćena za implementaciju DQN algoritma u ovom radu.
\begin{equation}
	\begin{gathered}
		g_0 = 0 \\
		g_{k+1} = \gamma g_k + (1 - \gamma)\nabla L(w_k) \odot \nabla L(w_k) \\
		\alpha_0 = \alpha \\
		\alpha_{k+1} = \frac{\alpha_k}{1+d(k+1)}	
	\end{gathered}
\end{equation}
Tada se iterativni korak definiše: 
\begin{equation}
	w_{k+1} = w_k - \frac{\alpha_{k+1}}{\sqrt{g_{k+1}} + \varepsilon} \nabla L(w_k)
\end{equation}
Sve operacije vrše se pokoordinatno. Parametar $\gamma$ pripada poluotvorenom intervalu $\left[0, 1\right) $. U svom predavanju, Hinton predlaže da njegova vrednost bude $0.9$. Preporučena vrednost za veličinu koraka odnosno stopu učenja, u oznaci $\alpha$, je $0.001$ dok $d$ označava faktor opadanja za parametar $\alpha$. Parametar $\varepsilon$ služi da bi se izbeglo deljenje nulom i obično je reda veličine $10^{-8}$.

\subsubsection{Adam}

Adam (eng. adaptive moment estimation) jedan je od najčešćih algoritama za optimizaciju korišćen pri obučavanju neuronskih mreža. Algoritam Adam zasnovan je na korišćenju ocena prvog i drugog momenta gradijenta, datim sledećim formulama:

\begin{equation}
	\begin{gathered}
			m_0 = 0 \\
			v_0 = 0 \\ 
			m_{k+1} = \beta_1 m_k + (1-\beta_1) \nabla L(w_k) \\
			v_{k+1} = \beta_2 v_k + (1-\beta_2) \nabla L(w_k) \odot \nabla L(w_k) 
	\end{gathered}
\end{equation}

Ocena prvog momenta, $m_0$, predstavlja otežani prosek pravca kretanja dok ocena drugog momenta, $v_0$, predstavlja otežani prosek kvadrata norme gradijenata. Međutim, ove dve ocene su pristrasne ka početnoj vrednosti, u ovom slučaju $0$\footnote{Ovde se misli na $0$ vektor istih dimenzija kao $x_k$ u slučaju prvog momenta i skalar $0$ u slučaju drugog momenta}. Da bi se to ispravilo, vrši se sledeća korekcija:
\begin{equation}
	\begin{gathered}
		\hat{m}_{k+1} = \frac{m_{k+1} }{1 - \beta_1^{k+1}} \\
		\hat{v}_{k+1} = \frac{v_{k+1} }{1 - \beta_2^{k+1}}
	\end{gathered}
\end{equation}
Na kraju, iterativni korak dat je ispod. Dodavanje skalara $\varepsilon$ na vektor $\hat{v_{k+1}}$ predstavlja dodavanje tog skalara svakom članu datog vektora. Korenovanje, deljenje i oduzimanje vrše se pokoordinatno.
\begin{equation}
	w_{k+1} = w_k - \alpha \frac{\hat{m}_{k+1}}{\sqrt{\hat{v}_{k+1}} + \varepsilon}
\end{equation}

Parametar $\alpha$ naziva se veličina koraka ili stopa učenja. Vrednosti parametara $\beta_1$ i $\beta_2$ ograničene su na skup $\left[0, 1\right) $ i preporučene vrednosti su $0.9$ i $0.999$, redom, dok se za $\varepsilon$ preporučuje vrednost $10^{-8}$. Kao i u algoritmu RMSProp, svrha parametra $\varepsilon$ je izbegavanje deljenja sa nulom. Takođe nalik algoritmu RMSProp opisanom iznad, moguće je uvesti stopu opadanja parametra $\alpha$.
\par
Intuicija kojom se vodi algoritam Adam jeste da dužina svakog koraka zavisi od osobina funkcije u regionu u kom se trenutno vrši optimizacija. Ovaj algoritam pokazao se kao superioran u odnosu na ostale algoritme za optimizaciju, u opštem slučaju.



\subsubsection{Metod propagacije unazad}
\label{subsub:backprop}

Metod propagacije unazad jedan je od najznačajnijih algoritama pri radu sa neuronskim mrežama. Zasniva se na korišćenju parcijalnog izvoda složene funkcije kako bi se izračunao gradijent funkcije koju predstavlja neuronska mreža. Na primer, ukoliko su date funkcije $g: \mathbb{R}^n \rightarrow \mathbb{R}^m$ i $f: \mathbb{R}^m \rightarrow \mathbb{R}$, tada se parcijalni izvod funkcije $f \circ g $ odnosno $f(g)$ po $i$-toj promenljivoj za $i=1,...,n$ računa:
\begin{equation}
	\partial_i (f \circ g) = \sum_{j=1}^{m} (\partial_j f \circ g) \partial_i g_j 
\end{equation}
Svaka iteracija algoritma propagacije unazad sastoji se od tri koraka:
\begin{itemize}
	\item Proširivanja do sada izračunatog parcijalnog izvoda izvodom aktivacione funkcije za dati sloj po pravilima računanja parcijalnog izvoda složene funkcije;
	\item Računanja vrednosti gradijenta prema parametrima jedinica datog sloja. Kako se pre aktivacione funkcije računa linearna kombinacija, ovaj gradijent je vektor vrednosti koji taj sloj dobija i ovaj korak se vrši množenjem tim vektorom;
	\item Proširivanja do sada izračunatog parcijalnog izvoda izvodom te linearne kombinacije po ulazima prateći pravilo za računanje parcijalnog izvoda složene funkcije.
\end{itemize}

Jednostavan primer rada algoritma propagacije unazad za složenu funkciju $f \circ g \circ h$ izgleda:
\begin{equation}
	\begin{aligned}
		f(g(h(x)))' &= \\
		&= f'(g(h(x)))g(h(x))' \\ 
		&= f'(g(h(x)))g'(h(x))h(x)' \\ 
		&= f'(g(h(x)))g'(h(x))h'(x) \\
	\end{aligned}
\end{equation}

Sada su prikazani potpuni alati za optimizaciju neuronske mreže sa propagacijom unapred. Treba imati u vidu da je algoritam propagacije unazad računski skup i da pri radu sa velikim neuronskim mrežama proces učenja može biti jako skup. Takođe, sem težina same neuronske mreže, na proces učenja utiču razni drugi parametri kao što su arhitektura mreže, podela podataka na skupove za obučavanje i testiranje i parametri algoritama za optimizaciju. Ovi parametri nazivaju se metaparametrima i neretko je neophodno pokušavati razne njihove kombinacije dok ponašanje mreže ne dostigne željeni nivo. Često se umesto traženja metaparametara pribegava korišćenju unapred poznate kombinacije za koju je već pokazano da daje željene rezultate pri rešavanju nekog problema.

\section{Konvolutivne neuronske mreže}



%[STOHASTICKI SPUST] \\
%[BACKPROP] \\ 
%[ADAM/RMSPROP?] -- bar procitaj\\ 


[MOZDA O TEOREMI O UNIVERZALNOJ APROKSIMACIJI?]

[SLIKA NEKE MREZE]



[AKTIVACIONE FUNKCIJE]

Uvod  (prosiriti?) \\
Vrste NN \\
Aktivacione funkcije \\
Perceptron / neuron (da li?) \\
Gradijentni spust, Propagacija unazad \\
Primene i ograničenja \\
Nešto o normalizaciji??? Regularizaciji
\chapter{Učenje potkrepljivanjem}
\label{ch:rl}

Učenje potkrepljivanjem je vid mašinskog učenja koji podseća na učenje koje se može naći u prirodi: ljudi i životinje uče na osnovu interakcije sa svetom oko sebe. Razlikovanje povoljnog i nepovoljnog ponašanja nije unapred učinjeno već je neophodno da onaj koji uči donese taj zaključak. U učenju potkrepljivanjem kao grani mašinskog učenja, komunikacija sa okruženjem svodi se na preduzimanje akcija u nekoj situaciji i dobijanja odgovora u vidu numeričke nagrade i informacije o tome kako se situacija promenila. Entitet koji komunicira sa svetom naziva se (softverskim) agentom dok se svet naziva okruženjem. Kao i kod ostalih oblasti mašinskog učenja, učenje potkrepljivanjem podrazumeva skup problema koji imaju za cilj obučavanje agenta radi maksimizacije dugoročne numeričke nagrade. 

\section{Osnovni pojmovi}
[MOŽDA SKLONI SECTION] \\
Učenje potkrepljivanjem sastoji se iz četiri komponente: politike ponašanja (ili samo politike), nagrade, funkcije vrednosti i modela okruženja. Politika opisuje način na koji se agent ponaša. Nagrada predstavlja numerički signal koji agent dobija od okruženja. U toku učenja, cilj agenta je da maksimizuje ukupnu nagradu dobijenu od okruženja. Dakle, nagradom je implicitno objašnjeno šta je dobro a šta loše ponašanje. Prilikom rešavanja problema učenja potkrelpljivanjem, cilj je nalaženje optimalne politike, tj. politike čijim se praćenjem dobija najveća dugoročna nagrada. Funkcija vrednosti govori koliko je dobro naći se u nekom stanju okruženja. Za razliku od nagrade, funkcija vrednosti opisuje kvalitet nekog stanja na duže staze. Ova komponenta je neophodna jer je moguće da dolaskom u neko stanje agent dobije malu nagradu ali da dato stanje ima veliku vrednost, što znači da je dolaskom u to stanje moguće ostvariti veliku dugoročnu nagradu. Neki algoritmi učenja potkrepljivanjem koriste model okruženja kako bi planirali unapred. Moguće je koristit pristupe koji koriste model, pristupe koji ne koriste model već se uči iz iskustva, kao i pristupe koji koriste učenje iz iskustva radi učenja modela.
\par 
U učenju potkrepljivanjem javlja se potreba za uspostavljanjem balansa između istraživanja i iskorišćavanja već stečenog znanja (eng.~{\em exloration vs. exploitation}). Naime, na početku učenja, agent istražuje okruženje i time uči kako bi trebalo da se ponaša. Čak i kada se nauči neko ponašanje, često je neophodno nastaviti sa istraživanjem u nekoj meri. Kako već naučeno ponašanje možda nije najbolje, bez istraživanja je moguće završiti u nekom od lokalnih optimuma. Međutim, kako učenje teče i agent unapređuje svoje ponašanje, poželjno je pratiti politiku koja dovodi do velikih dugoročnih nagrada. Najčešće je stopa istraživanja velika na početku učenja i opada u toku ovog procesa. [OVO NEGDE DRUGDE?]

\section{Markovljevi procesi odlučivanja}
\label{sec:mdp}

Markovljevi procesi odlučivanja (eng. {\em Markov decision process}, skr. {\em MDP}) daju teorijski okvir u kome je moguće relativno jednostavno postaviti i rešiti problem učenja potkrepljivanjem. Markovljevi procesi odlučivanja opisuju okruženje s kojim je moguće komunicirati. Ta komunikacija sastoji se iz toga da se od okruženja može dobiti informacija o stanju i da se okruženju može poslati informacija o akciji koja se preduzima, na šta okruženje odgovara informacijama o novom stanju i o numeričkoj nagradi. Izuzetno je važno da se od okruženja ne može dobiti informacija o tome da li je preduzeta akcija prava ili ne već samo informacija o numeričkoj nagradi, koju na duže staze treba maksimizovati. Kako se pri učenju potkrepljivanjem pretpostavlja postojanje agenta, kao i postojanje okruženja, nadalje se podrazumeva da postoji neki agent koji komunicira sa okruženjem. Treba imati u vidu da Markovljevi procesi odlučivanja opisuju idealno okruženje. U praksi okruženja često nisu idealna i tada se pribegava metodima koji ne koriste MDP. Iako je veliki deo teorije u učenju potkrepljivanjem ograničen na MDP, iste ideje imaju širu primenu.
%Pretpostavlja se da postoji agent koji komunicira sa okruženjem. Agent u svakom trenutku ima informaciju o %stanju okruženja i u mogućnosti je da preduzme akciju na šta od okruženja dobija odgovor u vidu novog stanja %i numeričke nagrade. Jedna važna pretpostavka jeste da agent ne dobija informaciju o tome da li je preduzeta %akcija bila dobra ili ne. Okruženje se predstavlja pomoću Markovljevog procesa odlučivanja. U nastavku je dat %teorijski opis Markovljevih procesa odlučivanja.

\subsection{Osnovni pojmovi}
Pod pretpostavkom da se interakcija između agenta i okruženja izvršava u diskretnim trenucima, stanje okruženja u trenutku $t$ označava se sa $S_t$ dok se skup svih stanja označava sa $\mathcal{S}$.  Agent u stanju $S_t$ preduzima akciju $A_t \in \mathcal{A}(S_t)$ i prelazi u novo stanje, $S_{t+1}$ dobijajući nagradu $R_{t+1}$. $\mathcal{A}(s)$ označava skup dozvoljenih akcija u stanju $s$. Ukoliko se sa $\mathbb{A}$ označi skup svih akcija, $\mathcal{A}$ se može posmatrati kao funkcija:  $\mathcal{A}: \mathcal{S} \rightarrow  \mathbb{P}(\mathbb{A})$. Skup $\mathcal{R}$ je skup mogućih realnih nagrada. Sekvenca $S_0, A_0, R_1, S_1, A_1, R_2, S_2, ...$ naziva se putanjom i dobija se interakcijom agenta sa okruženjem. Putanja može biti ili konačna i beskonačna.
Neophodno je definisati i funkciju prelaska, $p$:
\begin{equation}
	p(s', r~|~s, a) \eqdef P(S_{t+1} = s', R_{t+1} = r ~|~ S_t = s, A_t = a)
\end{equation}
koja predstavlja verovatnoću prelaska u stanje $s'$ i dobijanje nagrade $r$ pod uslovom da je u stanju $s$ preduzeta akcija $a$. Još jedna bitna pretpostavka je da je $p$ raspodela verovatnoće, iz čega sledi da:
\begin{equation}
	\sum_{s', r}^{} p(s', r~|~s, a) = 1
\end{equation}
za sve $s$ i $a \in \mathcal{A}(s)$. Ova skraćena oznaka za dvostruku sumu biće korišćena nadalje i označava sumiranje po svim $r \in \mathcal{R}$ i svim $s' \in \mathcal{S}$.
\par 
Markovljevi procesi odlučivanja imaju takozvano Markovljevo svojstvo tj. osobinu da trenutno stanje i nagrada zavise isključivo od prethodnog stanja i u njemu preduzete akcije a ne od cele putanje koja je dovela do datog stanja. Formalno, ovo svojstvo zapisuje se sledećom jednakošću:
\begin{equation}
	P(S_t, R_t ~|~ S_0, A_0, R_1, ..., S_{t-1}, A_{t-1}) = P(S_t, R_t ~|~ S_{t-1}, A_{t-1})
\end{equation}

U odnosu na putanju od trenutka $t$, može se govoriti o dugoročnoj nagradi, koja se još naziva i dobitkom:
\begin{equation}
\label{eq:dug_suma}
	G_t = \sum_{i=0}^{\infty} \gamma^iR_{t+i+1}
\end{equation}
Metaparametar $\gamma$ naziva se umanjenjem i ukazuje na to koliko se značaja pridaje kasnije dobijenim nagradama u odnosu na neposrednu nagradu. Za $\gamma = 0$\footnote{$0^0$ definiše se kao $1$.}, buduće nagrade nisu bitne dok postavljanje vrednosti $\gamma$ na $1$ ukazuje na to da se smatra da su sve nagrade putanje jednako bitne. Iz ove sume uočava se i odnos sa dugoročnom nagradom od trenutka $t+1$:
\begin{equation}
	\begin{aligned}
		G_t &= \sum_{i=0}^{\infty} \gamma^iR_{t+i+1} \\
        	&= R_{t+1} + \gamma\sum_{i=0}^{\infty} \gamma^iR_{(t+1)+i+1} \\
        	&=R_{t+1} + \gamma G_{t+1}
	\end{aligned}
\end{equation}
\par 
Sada je moguće dati formalnu definiciju: Markovljev proces odlučivanja je uređena petorka $(\mathcal{S}, \mathcal{A}, \mathcal{R}, p, \gamma)$. Ova definicija je prilično jednostavna a ipak je dovoljno fleksibilna za formalne opise raznih modela. Ako su skupovi $\mathcal{S}$, $\mathcal{R}$, $\mathcal{A}(s)$, za svako $s$, konačni, tada se za Markovljev proces odlučivanja kaže da je konačan. Umesto zahteva da $\mathcal{A}(s)$ bude konačan za svako $s$, može se zahtevati da je skup $\mathbb{A}$ konačan.
\par 
Markovljevi procesi odlučivanja koriste se za modeliranje interakcije sa okruženjem i donošenje odluka. Ova primena pokazala se izuzetno pogodno za probleme učenja potrepljivanjem. Cilj agenta biće maksimizacija dugoročne nagrade prilikom interakcije sa okruženjem.


\subsubsection{Epizode}

U jednakosti \eqref{eq:dug_suma} pretpostavlja se da niz interakcija sa okruženjem, tj. putanja, traje beskonačno, što se vidi iz gornje granice u sumi. Međutim, često je prirodnije pretpostaviti da su putanje konačne i da se završavaju u nekom posebnom stanju iz kog nije moguće dalje preduzimati akcije. Ovakva stanja nazivaju se završnim stanjima. Takvih stanja može biti više ali, zbog načina na koji je definisana funkcija prelaska\footnote{Agent dobija numeričku nagradu za preduzimanje akcije u stanju a ne za dolazak u stanje.}, za time nema potrebe i bez gubitka opštosti se može pretpostaviti da je ovakvo stanje, ukoliko postoji, jedinstveno. Jedan niz interakcija agenta sa okruženjem koji se završava završnim stanjem naziva se epizodom. Epizode su međusobno nezavisne u smislu da ishod jedne epizode ni na koji način ne utiče na neku od narednih epizoda, što se tiče samog okruženja. Ukoliko agent treba da igra, na primer, šah, partije se mogu smatrati epizodama. 
\par 
Neki problemi ne mogu se razbiti na epizode. Ovi problemi predstavljaju dugoročne zadatke kao što su beskonačno balansiranje uspravnog štapa ili odbrana od nadolazećih talasa neprijatelja u slučaju nekih video igara. Kod ovakvih problema, jako je važno postaviti umanjenje na vrednost manju od $1$. Naime, ukoliko važi $\gamma < 1$ i niz nagrada $R_t$ je ograničen, tada će suma \eqref{eq:dug_suma} konvergirati. Ukoliko ta suma divergira, odnosno ako je njena vrednost $\infty$ ili neodređena, tada će maksimizacija postati trivijalna, odnosno nemoguća.
\par 
Sa stanovišta završnih stanja i dugoročne nagrade, ova dva slučaja mogu se objediniti bez izmene \eqref{eq:dug_suma}. Kod problema koji se ne mogu podeliti na epizode, suma ostaje ista uz zahtev da je umanjenje strogo manje od $1$. Kod problema koji se mogu podeliti na epizode, može se uvesti pretpostavka da će se iz završnog stanja sa verovatnoćom $1$ prelaziti u to isto stanje uz vrednost nagrade $0$. Ovo je neizvodljivo za implementaciju pa se u praksi koristi konačna suma oblika:
\begin{equation}
	G_t = \sum_{i=0}^{T} \gamma^iR_{t+i+1}
\end{equation}
gde je $t+T+1$ trenutak kraja epizode. 

\subsubsection{Politika; vrednosti stanja i akcije}

Kako je neophodno opisati pravila ponašanja agenta, uvodi se funkcija $\pi$, koja predstavlja verovatnoću da agent u stanju $s$ preduzme akciju $a$. Ova funkcija naziva se politikom i, ukoliko neki agent preduzima akciju $a$ u stanju $s$ sa verovatnoćom $\pi(a~|~s)$, za sva stanja $s$ i sve akcije $a$, kaže se da agent prati politiku $\pi$. Vrednosti $\pi(a~|~s)$, kad $a \notin \mathcal{A}(s)$, se ne razmatraju.

Ukoliko za neku politiku $\pi$ važi da za svako $s$ postoji stanje $a_s$ takvo da $\pi(a_s~|~s)=1$ i $\pi(a~|~s)=0$ za sve ostale akcije $a$, tada se kaže da je politika deterministička. Radi jednostavnosti, u tom slučaju može se napisati $\pi(s)=a_s$.
\par 
Ako je poznata politika $\pi$, može se definisati funkcija vrednosti stanja pri praćenju politike $\pi$:
\begin{equation}
	v_{\pi}(s) \eqdef \mathbb{E}_{\pi}[G_t~|~S_t=s]
\end{equation}
Na sličan način uvodi se i funkcija vrednosti akcije u stanju pri praćenju politike $\pi$:
\begin{equation}
	q_{\pi}(s, a) \eqdef \mathbb{E}_{\pi}[G_t~|~S_t=s, A_t=a]
\end{equation}
Simbol $\mathbb{E}_{\pi}$ označava matematičko očekivanje ako se podrazumeva da se pri preduzimanju akcija prati politika $\pi$. Radi jednostavnostu će u nastavku funkcija $v_{\pi}$, odnosno $q_{\pi}$, biti nazivana samo funkcija vrednosti stanja, odnosno funkcija vrednosti akcije u stanju, dok će politika biti zapisana u indeksu.
\par 
Mnogi algoritmi učenja potkrepljivanjem zasnivaju se na nalaženju optimalne politike, odnosno politike čijim se praćenjem dolazi do maksimalne dugoročne nagrade. Moguće je uvesti parcijalno uređenje politika definisano na sledeći način:
\begin{equation}
	\label{eq:nije_losija}
	\pi_1 \leq \pi_2 \eqivdef \big(\forall s \in \mathcal{S}\big) \big( v_{\pi_1}(s) \leq v_{\pi_2}(s)\big)
\end{equation}
i tada se kaže da politika $\pi_2$ nije lošija  od politike $\pi_1$. Ukoliko postoji neka politika $\pi_*$ koja nije lošija ni od jedne politike za dati Markovljev proces odlučivanja, tada se ona naziva optimalnom politikom. 

\subsubsection{Belmanove jednakosti; optimalna politika}

Funkcije $v_{\pi}$ i $q_{\pi}$ zadovoljavaju rekurentne relacije koje se zovu Belmanovim jednakostima. U daljim izvođenjima za funkciju $v_{\pi}$ podrazumeva se da jednakosti važe za sva stanja $s$ odnosno za sva stanja $s$ i dozvoljene akcije u tim stanjima, $a$, za $q_{\pi}$. Važi:
\begin{equation}
	\begin{aligned}
		v_{\pi}(s) &\eqdef \mathbb{E}_{\pi}[G_t~|~S_t=s] \\ 
		&= \mathbb{E}_{\pi}[R_{t+1} + \gamma G_{t+1}~|~S_t=s] \\
		&= \mathbb{E}_{\pi}[R_{t+1}~|~S_t=s] + \gamma \mathbb{E}_{\pi}[G_{t+1}~|~S_t=s]
	\end{aligned}
\end{equation}
Sada je neophodno izračunati oba sabirka. Prvi sabirak označava očekivanu neposrednu nagradu polazeći iz stanja $s$ i prateći politiku $\pi$:
\begin{equation}
	\begin{aligned}
		\mathbb{E}_{\pi}[R_{t+1}~|~S_t=s] = \sum_{a}^{} \pi(a~|~s)\sum_{s', r}^{}rp(s', r~|~s, a)
	\end{aligned}
\end{equation}
Iz stanja $s$ preduzima se akcija $a$ sa verovatnoćom $\pi(a~|~s)$, dok se za preduzetu akciju $a$ u stanju $s$ sa verovatnoćom $p(s', r~|~s, a)$ prelazi u stanje $s'$ uz dobijanje nagrade $r$. U drugoj sumi, indeksi su $s'$ i $r$, što označava dvostruku sumu, jednu po $s'$ i jednu po $r$.
\par 
Drugi sabirak, bez množioca $\gamma$, proširuje se na sledeći način:
\begin{equation}
	\begin{aligned}
		\mathbb{E}_{\pi}[G_{t+1}~|~S_t=s] &= \sum_{a}^{} \pi(a~|~s)\sum_{s',r}^{}p(s', r~|~s, a)\mathbb{E}_{\pi}[G_{t+1}~|~S_{t+1}=s'] \\
		&= \sum_{a}^{} \pi(a~|~s)\sum_{s',r}^{}p(s', r~|~s, a)v_{\pi}(s')
	\end{aligned}
\end{equation}
Sumiranje se vrši i po $r$ jer nije nužno da stanju $s'$ odgovara jedinstvena nagrada $r$.
\par 
Spajanjem dve jednakosti dobija se:
\begin{equation}
	\label{eq:bel_v}
	\begin{aligned}
		v_{\pi}(s)  &= \mathbb{E}_{\pi}[R_{t+1}~|~S_t=s] + \gamma \mathbb{E}_{\pi}[G_{t+1}~|~S_t=s] \\
		&= \sum_{a}^{} \pi(a~|~s)\sum_{s', r}^{}rp(s', r~|~s, a) + \gamma \sum_{a}^{} \pi(a~|~s)\sum_{s',r}^{}p(s', r~|~s, a)v_{\pi}(s') \\
		&= \sum_{a}^{} \pi(a~|~s)\sum_{s', r}^{}p(s', r~|~s, a) \big[ r+\gamma v_{\pi}(s') \big]
	\end{aligned}
\end{equation}
Analogno se izvodi rekurentna veza za funkciju vrednosti akcije u stanju, $q_{\pi}$:
\begin{equation}
	\label{eq:bel_q}
	\begin{aligned}
		q_{\pi}(s,a) = \sum_{a}^{} \pi(a~|~s)\sum_{s', r}^{}\big[ r+\gamma\sum_{a'}^{}\pi(a'~|~s') q_{\pi}(s', a') \big]
	\end{aligned}
\end{equation}
Jednakosti \eqref{eq:bel_v} i \eqref{eq:bel_q} nazivaju se Belmanovim jednakostima i ključne su za mnoge algoritme učenja potkrepljivanjem. 
\par 
Iz definicije optimalne politike, $\pi_*$, slede definicije optimalne funkcije vrednosti stanja i optimalne funkcije vrednosti akcije u stanju koje odgovaraju optimalnoj politici, u oznaci $v_*$ i $q_*$:
\begin{equation}
	\begin{aligned}
		v_*(s) &\eqdef \max_{\pi} v_{\pi} (s)	 \\
		q_*(s,a) &\eqdef \max_{\pi} q_{\pi}(s,a)
	\end{aligned}
\end{equation}
Ukoliko neki agent prati optimalnu politiku, ona će ga dovesti do maksimalne dugoročne nagrade. Ovo sledi iz činjenica da $v_{\pi}(s)$ predstavlja dugoročnu nagradu polazeći iz stanja $s$, i da $v_*$ za svako stanje $s$ predstavlja najveću vrednost stanja među svim politikama.
\par
Moguće je uspostaviti vezu između $v_*$ i $q_*$:
\begin{equation}
	\begin{aligned}
		v_*(s) &= \max_{a \in \mathcal{A}(s)} q_* (s,a)	 \\
		q_*(s,a) &= \mathbb{E}_{\pi}\big[R_{t+1} + \gamma v_*(s_{t+1})~|~S_t=s, A_t=a\big]
	\end{aligned}
\end{equation}
Relativno jednostavnim izvođenjem dolazi se do još jednog para jednakosti koje se nazivaju Belmanovim jednakostima optimalnosti.
	\begin{align}
		\label{eq:bel_v_opt} v_*(s) &= \max_{a}\sum_{s', r}^{} p(s', r~|~s,a)\big[r+\gamma v_*(s')\big]	 \\
		\label{eq:bel_q_opt} q_*(s,a) &= \sum_{s', r}^{} p(s', r~|~s,a)\big[\gamma \max_{a'}q_*(s',a')\big]
	\end{align}
\par 
U slučaju konačnih Markovljevih procesa odlučivanja, rekurentne relacije \eqref{eq:bel_v_opt} čine sistem od $n$ jednačina sa $n$ nepoznatih.\footnote{Kao što je pomenuto na početku dela o Belmanovim jednakostima, one važe za svako stanje $s$.} Isto važi i za relacije \eqref{eq:bel_q_opt}. Ovi sistemi sadrže funkciju $max$ pa stoga nisu linearni. Ukoliko je funkcija prelaska za dati MDP poznata, ovi sistemi se mogu rešiti; ta rešenja dobijena su bez prethodnog znanja o optimalnoj politici. 
\par 
Ukoliko je $v_*$ poznata, moguće je odrediti optimalnu politiku. Iz jednakosti \eqref{eq:bel_v_opt} jasno je da postoji jedna ili više akcija koje u stanju $s$ dovode do maksimalne sume. Bilo koja politika koja nekim od ovih akcija dodeljuje nenula vrednosti a svim ostalim dodeljuje vrednost $0$ je optimalna. Ukoliko je poznata funkcija $q_*$, nalaženje optimalne politike još je jednostavnije: u svakom stanju $s$ preduzima se akcija $a$ takva da se maksimizuje $q_*(s,a)$. Poznavanje $q_*$, dakle, omogućuje nalaženje optimalne politike bez ikakvog uvida u funkciju prelaska.  Međutim, ovaj direktni pristup nalaženju $v_*$ ili $q_*$ obično se ne koristi jer funkcija prelaska u praksi uglavnom nije poznata.

\section{Rešavanje Markovljevih procesa odlučivanja}

Uz pretpostavku da se radi o konačnom MDP, osnovni pristup njegovom rešavanju jeste koristeći principe dinamičkog programiranja. Kako je cilj učenja traženje optimalne politike, dve politike biće poređene u skladu sa definicijom \eqref{eq:nije_losija}, odnosno koristeći funkciju vrednosti politike. Stoga, prvo je neophodno naći metod kojim se računa funkcija vrednosti stanja koja odgovara nekoj politici $\pi$.  Način na koji se ovo radi je korišćenjem Belmanovih jednakosti, \eqref{eq:bel_v}. Počinje se od vektora $v_0$ koji sadrži proizvoljne vrednosti i dužine je koja odgovara broju stanja. Dalje se primenjuje iterativno ažuriranje po sledećem pravilu:
\begin{equation}
		v_{k+1}(s) = \sum_{a}^{} \pi(a~|~s)\sum_{s', r}^{}p(s', r~|~s, a) \big[ r+\gamma v_k(s') \big]
\end{equation}
za sve $s \in \mathcal{S}$. Pod uslovom da $\gamma < 1$, niz vektora $v_k$ konvergira ka $v_{\pi}$. Sada je i u praksi moguće uporediti dve politike.
\par 
Kada je data deterministička politika $\pi$, postavlja se pitanje da li je moguće unaprediti je, tj. da li u stanju $s$ odabrati akciju $\pi(s)$ ili neku drugu akciju? Tada je za sve akcije $a$ dozvoljene u $s$ neophodno ispitati vrednost $q_{\pi}(s,a)$ u odnosu na vrednost $v_{\pi}(s)$. Ako neko $s$ postoji akcija $a$ takva da je $q_{\pi}(s, a) \geq v_{\pi}(s)$, tada je moguće konstruisati politiku $\pi'$ takvu da važi $\pi' \geq \pi$. Ovo važi na osnovu teoreme o unapređenju politike (eng.~{\em policy improvement theorem}):
\begin{theorem*}
	Neka su $\pi$ i $\pi'$ dve determinističke politike takve da za sve $s \in \mathcal{S}$ važi $q_{\pi}(s, \pi'(s)) \geq v_{\pi}(s)$. Tada politika $\pi'$ nije lošija od politike $\pi$, odnosno $\pi' \geq \pi$.
\end{theorem*}
Dokaz teoreme relativno je jednostavan i može se naći u \cite{rli_tup}. Preostalo je da se za proizvoljnu politiku $\pi$ nađe politika $\pi'$ takva da za svako $s$ važi $q_{\pi}(s, \pi'(s)) \geq v_{\pi}(s)$. Ukoliko se u svakom stanju $s$ odabere akcija $a$ koja maksimizuje izraz $q_{\pi}(s, a)$, tada tako dobijena politika zadovoljava uslove teoreme. Dakle, od politike $\pi$ dobija se politika $\pi'$ tako što se za svako $s \in \mathcal{S}$ računa:
\begin{equation}
	\pi'(s) = \argmax_a q_{\pi}(s, a)
\end{equation}
Proces dobijanja $\pi'$ od $\pi$  naziva se unapređenjem politike (eng.~{\em policy improvement}) i omogućuje iterativni pristup traženju optimalne politike. Politika $\pi'$ naziva se pohlepnom politikom jer se njenim praćenjem uvek bira najbolja akcija u neposrednom smislu. Ova politika zadovoljava uslove teoreme jer
\begin{align}
	\max_a q_{\pi}(s,a) \geq&~ q_{\pi}(s,\pi(s)) \\=& v_{\pi}(s)
\end{align}
Druga jednakost sledi iz definicije funkcija $q_{\pi}$ i $v_{\pi}$.
\par 
Polazeći od nasumično kreirane determinističke politike i prateći ovaj postupak, naizmenično sa evaluacijom novodobijene politike, dobija se niz politika $\pi_0, \pi_1, ...$ koji u slučaju konačnih MDP konvergira optimalnoj politici, $\pi_*$. Ime ovog procesa je iterativno unapređenje politike (eng.~{\em policy iteration}). Nasumična deterministička politika $\pi$ kreira se tako što se za svaku vrednost $\pi(s)$ proizvoljno bira element iz $\mathcal{A}(s)$.
\par 
Iterativno unapređenje politike, iako relativno brzo konvergira u odnosu na broj stanja, zahteva evaluaciju politike u svakoj iteraciji. Pored ovog algoritma, postoji i algoritam iterativnog unapređenja funkcije vrednosti. Naime, umesto traženja optimalne politike, iterativno se traži optimalna funkcija vrednosti stanja, $v_*$. Na osnovu Belmanove jednakosti optimalnosti za funkciju vrednosti, \eqref{eq:bel_v_opt}, kreira se iterativno pravilo ažuriranja:
\begin{equation}
	\label{eq:val_iter} v_{k+1}(s) = \max_{a}\sum_{s', r}^{} p(s', r~|~s,a)[r+\gamma v_k(s')]
\end{equation}
gde se, kao i kod evaluacije politike, polazi od vektora $v_0$ čije su vrednosti nasumično odabrane. Proces se zaustavlja kada razlika dve uzastopne funkcije postane dovoljno mala, u skladu sa metrikom $\norm{\cdot}_{\infty}$ definisanom na sledeći način:
\begin{equation}
	\norm{v}_{\infty} = \max_x|v(x)|
\end{equation}
\par 
Iako su do sada opisani algoritmi prilično intuitivni, nije direktno jasno da moraju da konvergiraju. [OBRISATI]
Pravilo ažuriranja \eqref{eq:val_iter} može se predstaviti kao operator $B:R^{|S|} \rightarrow R^{|S|}$ koji se naziva Belmanovim operatorom:
\begin{equation}
	Bv(s) = \max_{a}\sum_{s', r}^{} p(s', r~|~s,a)\bigcap+\gamma v(s')\big]
\end{equation}
Ukoliko se pokaže da je operator $B$ kontrakcija, tada po teoremi o fiksnoj tački važi da postoji vektor $v_*$ takav da $Bv_*=v_*$. Opis i dokaz ove teoreme može se naći u \cite{num_met}. Uslov kontrakcije zapisuje se sledećim izrazom:
\begin{equation}
	\norm{Bv_1-Bv_2}_{\infty} \leq \alpha \norm{v_1-v_2}_{\infty}
\end{equation}
za neko $\alpha \in [0, 1)$. 
Dokaz da je Belmanov operator kontrakcija: [PROVERI TACNOST SVIH OVIH JEDNAKOSTI]
\begin{align}
	\norm{Bv_1-Bv_2}_{\infty} &= \max_s |Bv_1(s)-Bv_2(s)| \\
	&= \mabs*{\max_{a}\sum_{s', r}^{} p(s', r~|~s,a)[r+\gamma v_1(s')] - \max_{a}\sum_{s', r}^{} p(s', r~|~s,a)[r+\gamma v_2(s')]} \\
	&\leq \max_{a} \mabs*{\sum_{s', r}^{} p(s', r~|~s,a)[r+\gamma v_1(s')] - \sum_{s', r}^{} p(s', r~|~s,a)[r+\gamma v_2(s')]} \\
	&=\max_{a} \mabs*{ \sum_{s', r}^{} p(s', r~|~s,a) \big[(r+\gamma v_1(s'))-(r+\gamma v_2(s'))\big]} \\
	&=\gamma\max_{a} \mabs*{ \sum_{s', r}^{} p(s', r~|~s,a)\big[v_1(s')-v_2(s')\big]} \\
	&\leq \gamma \max_s \mabs*{v_1(s)-v_2(s)} \\
	&= \gamma \norm{v1-v2}_{\infty}
\end{align}
Ukoliko je $\gamma \in [0, 1)$, Belmanov operator je kontrakcija. Treći red sledi iz činjenice da 
\begin{equation}
	\mabs*{\max_{x} f(x) - \max_{x} g(x)} \leq \max_{x}\mabs*{f(x)-g(x)}
\end{equation}
dok pretposlednji red sledi iz činjenice da je $p$ raspodela verovatnoće za fiksirane parametre $s$ i $a$ [Da li da ovo pomenem u definiciji funkcije p?]. Pošto je dokazano da postoji jedinstvena fiksna tačka Belmanovog operatora i, kako po njegovoj konstrukciji važi $Bv \geq v$, vidi se da je niz dobijen njegovom uzastopnom primenom rastući. Stoga je fiksna tačka zaista i optimalna funkcija vrednosti.
\par 
Ranije je pokazano kako se, ukoliko je poznata funkcija $v_*$, može konstruisati optimalna politika. Iz postojanja $v_*$ sledi postojanje $\pi_*$ ali ovakvih politika može biti više zbog mogućnosti da iz jednog stanja više akcija vode ka stanjima sa istom, maksimalnom, vrednošću. 
\par 
Iako sporiji postupak, traženje $q_*$ funkcije pogodnije je od traženja $v_*$ u slučajevima kada funkcija prelaska nije poznata. Pristupi traženju ovih funkcija dinamičkim programiranjem podjednako su zastupljeni. 

\subsection{Ostali načini rešavanja MDP}

Funkciju vrednosti stanja i funkciju vrednosti akcije u stanju metodima linearnog programiranja. Jedan od primera je traženje $v_*$. Na osnovu jednakosti \eqref{eq:bel_v_opt} može se postaviti skup uslova:
\begin{equation}
	v(s) \geq \sum_{s', r}^{} p(s', r~|~s,a)\big[r+\gamma v(s')\big] \text{,~za~sve~} s \in \mathcal{S} \text{~i~sve~} a \in \mathcal{A}(s)
\end{equation}
Ovih uslova ima $\sum_{s \in \mathcal{S}}^{} |\mathcal{A}(s)| \leq |\mathcal{S}||\mathbb{A}|$. Pri ovim uslovima, treba minimizovati izraz $\sum_{s \in \mathcal{S}}^{} v(s)$. Rešavanje ovog problema daje $v_*$. Ovaj pristup rešavanju MDP-a pri određenim uslovima daje teorijski bržu konvergenciju od pristupa dinamičkim programiranjem. Uprkos tome, pri porastu broja stanja, pristup linearnim programiranjem brže postaje neizvodljiv od pristupa dinamičkim programiranjem.
\par 
Još jedan način rešavanja MDP-a predstavljaju takozvani Monte Karlo metodi, koji se oslanjaju na prikupljanje informacija o okruženju iz iskustva. Na primer, za približno određivanje optimalne funkcije vrednosti stanja čuva se prosek vrednosti dobijenih nakon dolaska u to stanje. Proširivanje ovog iskustva [u ovom slučaju] znači bolje određivanje optimalne funkcije vrednosti.

Međutim, nekada je broj stanja toliko veliki da ni jedan od ova dva pristupa nije izvodljiv.

% ostali metodi resavanja problema RL
% neka prica o opstem Q ucenju
% neki zakljucak -- kada je sta bolje koristiti









\chapter{DQN}
\label{ch:dqn}

\chapter{Detalji implementacije}
\label{ch:implementacija}

\chapter{Eskperimentisanje sa elementima algoritma DQN}
\label{ch:eksperimenti}

\addcontentsline{toc}{chapter}{Literatura}
\nocite{*}
\bibliography{tex/lit}{}

 
\end{document}

